
\documentclass[12pt,a4paper]{amsart}

\usepackage[margin=1in]{geometry}
\usepackage[T1]{fontenc}
\usepackage[utf8]{inputenc}
\usepackage{lmodern}
\usepackage{microtype}
\emergencystretch=3em
\usepackage{amsmath,amssymb,amsthm,mathtools}
\usepackage{enumitem}
\usepackage{aliascnt}
\usepackage[colorlinks=true,linkcolor=blue,citecolor=blue,urlcolor=blue]{hyperref}
\usepackage[capitalize,nameinlink]{cleveref}

\numberwithin{equation}{section}

\newtheorem{theorem}{Theorem}[section]
\newaliascnt{proposition}{theorem}
\newtheorem{proposition}[proposition]{Proposition}
\aliascntresetthe{proposition}
\newaliascnt{lemma}{theorem}
\newtheorem{lemma}[lemma]{Lemma}
\aliascntresetthe{lemma}
\newaliascnt{corollary}{theorem}
\newtheorem{corollary}[corollary]{Corollary}
\aliascntresetthe{corollary}
\newaliascnt{definition}{theorem}
\newtheorem{definition}[definition]{Definition}
\aliascntresetthe{definition}
\theoremstyle{remark}
\newaliascnt{remark}{theorem}
\newtheorem{remark}[remark]{Remark}
\aliascntresetthe{remark}
\newtheorem{hypothesis}[theorem]{Hypothesis}
\newtheorem{problem}[theorem]{Problem}

\crefname{theorem}{Theorem}{Theorems}
\Crefname{theorem}{Theorem}{Theorems}
\crefname{proposition}{Proposition}{Propositions}
\Crefname{proposition}{Proposition}{Propositions}
\crefname{lemma}{Lemma}{Lemmas}
\Crefname{lemma}{Lemma}{Lemmas}
\crefname{corollary}{Corollary}{Corollaries}
\Crefname{corollary}{Corollary}{Corollaries}
\crefname{definition}{Definition}{Definitions}
\Crefname{definition}{Definition}{Definitions}
\crefname{remark}{Remark}{Remarks}
\Crefname{remark}{Remark}{Remarks}
\crefname{hypothesis}{Hypothesis}{Hypotheses}
\Crefname{hypothesis}{Hypothesis}{Hypotheses}
\crefname{problem}{Problem}{Problems}
\Crefname{problem}{Problem}{Problems}

\DeclareMathOperator{\RS}{RS}
\DeclareMathOperator{\MCA}{MCA}
\DeclareMathOperator{\CA}{CA}
\DeclareMathOperator{\rank}{rank}
\DeclareMathOperator{\Span}{Span}
\DeclareMathOperator{\Spec}{Spec}
\DeclareMathOperator{\Hom}{Hom}
\DeclareMathOperator{\im}{im}
\DeclareMathOperator{\codim}{codim}
\DeclareMathOperator{\supp}{supp}
\DeclareMathOperator{\Fib}{Fib}
\DeclareMathOperator{\Sym}{Sym}
\DeclareMathOperator{\Gr}{Gr}
\DeclareMathOperator{\Ent}{H}
\DeclareMathOperator{\ord}{ord}
\DeclareMathOperator{\Tr}{Tr}
\DeclareMathOperator{\Gal}{Gal}
\DeclareMathOperator{\Res}{Res}
\DeclareMathOperator{\Aut}{Aut}
\DeclareMathOperator{\diag}{diag}
\DeclareMathOperator{\wt}{wt}

\newcommand{\F}{\mathbb F}
\newcommand{\B}{\mathbb B}
\newcommand{\K}{\mathbb K}
\newcommand{\Z}{\mathbb Z}
\newcommand{\A}{\mathbb A}
\newcommand{\Pj}{\mathbb P}
\newcommand{\C}{\mathbb C}
\newcommand{\Prob}{\mathbb P}
\newcommand{\E}{\mathbb E}
\newcommand{\eps}{\varepsilon}
\newcommand{\Omc}{\Omega^{\circ}}
\newcommand{\Ccal}{\mathcal C}
\newcommand{\Dcal}{\mathcal D}
\newcommand{\Ecal}{\mathcal E}
\newcommand{\Fcal}{\mathcal F}
\newcommand{\Gcal}{\mathcal G}
\newcommand{\Hcal}{\mathcal H}
\newcommand{\Ical}{\mathcal I}
\newcommand{\Jcal}{\mathcal J}
\newcommand{\Lcal}{\mathcal L}
\newcommand{\Mcal}{\mathcal M}
\newcommand{\Ncal}{\mathcal N}
\newcommand{\Pcal}{\mathcal P}
\newcommand{\Qcal}{\mathcal Q}
\newcommand{\Rcal}{\mathcal R}
\newcommand{\Scal}{\mathcal S}
\newcommand{\Tcal}{\mathcal T}
\newcommand{\Ucal}{\mathcal U}
\newcommand{\Vcal}{\mathcal V}
\newcommand{\Wcal}{\mathcal W}
\newcommand{\Xcal}{\mathcal X}
\newcommand{\Ycal}{\mathcal Y}
\newcommand{\Zcal}{\mathcal Z}
\newcommand{\Gord}{\Gamma^{\rm ord}}
\newcommand{\Gsid}{\Gamma^{\rm sid}}
\newcommand{\Gprim}{\Gamma^{\rm prim}}
\newcommand{\barN}{\overline N}
\newcommand{\Renyi}{R\'enyi}
\newcommand{\Hb}{H_2}
\newcommand{\UU}{\mathbb U}
\newcommand{\abs}[1]{\lvert #1\rvert}
\newcommand{\norm}[1]{\lVert #1\rVert}
\newcommand{\floor}[1]{\left\lfloor #1\right\rfloor}
\newcommand{\ceil}[1]{\left\lceil #1\right\rceil}
\newcommand{\defeq}{\coloneqq}
\newcommand{\one}{\mathbf 1}
\newcommand{\Fq}{\mathbb F_q}

\title{Asymptotic RS--MCA Frontiers}
\author{Przemek Chojecki}
\hypersetup{pdftitle={Asymptotic RS--MCA Frontiers},pdfauthor={Przemek Chojecki}}
\subjclass[2020]{94B35, 11T71, 05B40}
\keywords{Reed--Solomon codes, mutual correlated agreement, smooth domains,
circle codes, entropy frontier, additive energy}
\date{Revised July 2026}

\begin{document}

\begin{abstract}
We study mutual correlated agreement (MCA) for sequences of Reed--Solomon
codes: a slope is bad when its received-line word agrees with a low-degree
polynomial on many coordinates although the two line generators are not
both so explained on that same set.  Unconditionally, we determine the
high-agreement bad-slope count and translate exactly among shared leading
coefficients of polynomials vanishing on the agreement supports, lists of
nearby codewords, and intersections of a parity-check image line with
column spans; these identities give an exact deep-target frontier and an
unconditional shallow-prefix error exponent.  Over polynomial-size evaluation
fields in odd and characteristic two, exponentially large support families
with few repeated differences disprove a uniform bound based only on the
average fiber size of the leading-locator-coefficient map; the
odd-characteristic example also transfers to circle evaluation codes.  For
supports formed from complete polynomial-map
fibers and marked exceptional points, we prove an exact quotient/remainder
normal form and quantify the enlargement caused by a smaller quotient
field.  We also prove effective-span Fourier inversion, an exact
partial-occupancy census and add-back, and exact finite
occupancy, pair-incidence, and curve-degree compilers, removing hidden
normalization and asymptotic constants from the formal reductions.  The
same analysis proves that a degree-uniform pointwise Weil certificate and
automatic profile-scale ray multiplicity do not follow from smooth/circle
geometry; thus, outside the unconditional deep frontier and the verified
shallow Fourier range, the unrestricted positive-density frontier still
requires a witness-exhaustive atlas, image-scale MI and MA or a direct
Sidon payment, residual ray bounds, and comparison of the complete profile
envelope with the target and lower reserve.
\end{abstract}

\maketitle

\section{Introduction and main results}\label{sec:introduction}

By a finite Reed--Solomon \emph{row} we mean data
\((\F,D,k)\), where \(\F\) is a finite field,
\(D\subseteq\F\) is an \emph{evaluation domain} of size \(n\), and
\(1\le k<n\).  The associated code
\(C=\RS_\F(D,k)\subseteq\F^D\) consists of the evaluations on \(D\) of
polynomials of degree less than \(k\).  Thus a support is simply a subset
of the coordinate set \(D\).  A word \(u\in\F^D\) is
\emph{explained on a support} \(S\subseteq D\) if its restriction to
\(S\) agrees with one such degree-less-than-\(k\) polynomial.
When the evaluation points lie in a designated coefficient subfield, we
write \(D\subseteq\B\subseteq\F\); \(\B\) controls locator coefficients
and domain structure, while \(\F\) controls codeword coefficients and
slopes.
Throughout,
\(h(x)=-x\log x-(1-x)\log(1-x)\) denotes binary entropy in natural
units, and \(H_2(x)=h(x)/\log 2\) denotes binary entropy in bits, with the
usual convention \(0\log0=0\).

A received pair \(r=(r_0,r_1)\in(\F^D)^2\) determines the
\emph{received line} \(\{r_0+\gamma r_1:\gamma\in\F\}\).  The pair is
simultaneously explained on \(S\) if each of the two received words has an explanation
there, possibly by different polynomials.  A finite slope \(\gamma\) is
\emph{support-wise MCA-bad at agreement \(a\)} if some support
\(S\subseteq D\), \(\abs S\ge a\), explains
\(r_0+\gamma r_1\) but does not simultaneously explain the pair.  The
same support is used in both tests; this is the support-wise convention.
A tuple \((\gamma,S,h)\), where \(\deg h<k\) and
\(h|_S=(r_0+\gamma r_1)|_S\), is a \emph{witness} to that bad slope when
the pair is not simultaneously explained on \(S\), or equivalently is
\emph{support-wise nontrivial} there.  Finally,
\(B_C^{\rm MCA}(a)\) is the maximum, over all received pairs, of the
number of distinct bad slopes.  It counts slopes, not witnesses or
supports.

Here is the asymptotic formulation used throughout.  If
\(\varnothing\ne\Gamma\subseteq\F\) is the \emph{challenge set}, meaning
the set of slopes allowed in the test, let
\(B_{C,\Gamma}^{\rm MCA}(a)\) denote the same maximum with
\(\gamma\in\Gamma\), and at the radius grid point \(\delta=1-a/n\) put
\[
 e_{\rm MCA,\Gamma}(C,\delta)
   =\frac{B_{C,\Gamma}^{\rm MCA}(a)}{\abs\Gamma}.
 \tag{1.1}\label{eq:intro-normalized-mca}
\]
For the standard full-field challenge we abbreviate
\[
 e_{\rm MCA}(C,1-a/n)=\frac{B_C^{\rm MCA}(a)}{\abs\F},
 \qquad
 \delta_C^*(\eps)=\sup\{\delta<1-k/n:
                         e_{\rm MCA}(C,\delta)\le\eps\},
\]
with radii interpreted on the closed integer grid.  Taking
\(\Gamma=\F\) in \eqref{eq:intro-normalized-mca} gives exactly this
customary formulation.
By a \emph{row sequence} we mean a family
\(C_n=\RS_{\F_n}(D_n,k_n)\) indexed so that \(\abs{D_n}=n\); the fields,
domains, and dimensions may all vary with \(n\).  Every claimed upper
bound is uniform over received pairs, whereas a lower construction need
only exhibit the stated received pair.  For such a sequence,
challenge sets
\(\Gamma_n\), and targets \(0\le\eps_n\le1\), define
\[
 \begin{split}
 B_n^*&=\floor{\eps_n\abs{\Gamma_n}},\\
 a_n^*&=\min\{a\in\{k_n+1,\ldots,n\}:
              B_{C_n,\Gamma_n}^{\rm MCA}(a)\le B_n^*\},
 \qquad
 \delta_n^*=1-\frac{a_n^*}{n},
 \end{split}
 \tag{1.2}\label{eq:intro-asymptotic-threshold}
\]
whenever the set defining \(a_n^*\) is nonempty.  Thus the asymptotic
RS--MCA problem is to determine \(\delta_n^*\), or equivalently the first
\emph{safe agreement}: agreement \(a\) is safe when its bad-slope
numerator is at most \(B_n^*\), and unsafe otherwise.  Since the numerator
is nonincreasing in \(a\), the resulting safe/unsafe boundary is the
\emph{MCA frontier}.  This formulation retains the target and challenge
size; replacing them by a zero-exponent slogan loses real information when
the target is small.

\begin{theorem}[Unconditional asymptotic RS--MCA theorem in the deep regime]
\label{thm:intro-unconditional-asymptotic-deep}
Let \(C_n=\RS_{\F_n}(D_n,k_n)\) be any Reed--Solomon row sequence over
arbitrary characteristics, let
\(\varnothing\ne\Gamma_n\subseteq\F_n\), and put
\(B_n^*=\floor{\eps_n\abs{\Gamma_n}}\) as in
\eqref{eq:intro-asymptotic-threshold}.
\begin{enumerate}[label=\textup{(\roman*)},leftmargin=*]
\item If \(B_n^*=0\), no agreement is safe and the set defining
      \(a_n^*\) is empty.
\item If \(B_n^*=\abs{\Gamma_n}\), then \(a_n^*=k_n+1\).
\item Suppose eventually
\[
 1\le B_n^*<\abs{\Gamma_n},\qquad
 B_n^*\le n-k_n-1,\qquad
 3(B_n^*-1)\le n-k_n.                              \tag{1.3}
 \label{eq:intro-unconditional-asymptotic-deep}
\]
Then the threshold exists and is exactly
\[
 a_n^*=n-B_n^*+1,\qquad
 \delta_n^*=\frac{B_n^*-1}{n}.                    \tag{1.4}
 \label{eq:intro-unconditional-asymptotic-deep-answer}
\]
In particular, if \(k_n/n\to\rho\),
\(B_n^*/n\to b<(1-\rho)/3\), and
\(1\le B_n^*<\abs{\Gamma_n}\) eventually, then
\(\delta_n^*\to b\).  This conclusion applies, without a smoothness,
odd-characteristic, Fourier, atlas, or ray-compiler hypothesis, to every
multiplicative smooth RS row and to every circle code that has been proved
diagonally agreement-equivalent to an RS presentation.
\end{enumerate}
\end{theorem}

\begin{proof}
If \(B_n^*=0\), the universal tangent floor at agreement \(n\) is one,
so no grid point is safe.  If \(B_n^*=\abs{\Gamma_n}\), the numerator is
always at most the size of the challenge set, so the first grid point
\(k_n+1\) is safe.  In the remaining case put
\(a_0=n-B_n^*+1\).  Then \(n-a_0=B_n^*-1\), and the last inequality in
\eqref{eq:intro-unconditional-asymptotic-deep} places \(a_0\) in the
deep regime.  The exact deep numerator of
\cref{cor:exact-deep-numerator} is
\(\min\{\abs{\Gamma_n},B_n^*\}=B_n^*\), so \(a_0\) is safe.  The second
inequality in \eqref{eq:intro-unconditional-asymptotic-deep} puts
\(a_0-1\) on the admissible grid, and
\cref{prop:universal-tangent-floor} gives at least
\(\min\{\abs{\Gamma_n},B_n^*+1\}=B_n^*+1\) bad slopes there.  Hence
\(a_0-1\) is unsafe, and monotonicity proves
\eqref{eq:intro-unconditional-asymptotic-deep-answer}.  The asymptotic
specialization follows because its strict inequality implies
\eqref{eq:intro-unconditional-asymptotic-deep} for all sufficiently large
\(n\).
\end{proof}

\begin{theorem}[Unconditional shallow-prefix RS--MCA exponent]
\label{thm:unconditional-asymptotic-rs-mca}
Let \(C_n=\RS_{\F_n}(D_n,k_n)\), where
\(D_n\subseteq\B_n\subseteq\F_n\), \(\abs{D_n}=n\), and
\(k_n/n\to\rho\in(0,1)\).  Let \(a_n\ge k_n+1\), put
\(w_n=a_n-k_n-1\), and assume
\[
 \frac{a_n}{n}\to\alpha\in(0,1),\qquad
 w_n\log\abs{\B_n}=o(n),                           \tag{AS1}
\]
\(\abs{\F_n}>n\), and, for some finite \(\phi\ge0\),
\[
 \frac1n\log\abs{\F_n}\to\phi,
 \qquad
 \frac1n\log(\abs{\F_n}-n)\to\phi.               \tag{AS2}
\]
For the full-field challenge,
\[
 \frac1n\log B_{C_n}^{\rm MCA}(a_n)
   =\min\{h(\alpha),\phi\}+o(1),                   \tag{AS3}
\]
and therefore
\[
 \frac1n\log e_{\rm MCA}
   \left(C_n,1-\frac{a_n}{n}\right)
   =-\bigl(\phi-h(\alpha)\bigr)_++o(1).            \tag{AS4}
\]
Because \(\abs{\B_n}\ge n\), condition \textup{(AS1)} forces
\(\alpha=\rho\).  The result applies to every smooth multiplicative row
and to every standard circle code covered by the exact
agreement-preserving diagonal equivalence of
\cref{lem:circle-uniformization}; it assumes no Fourier payment, atlas, or
ray compiler.
\end{theorem}

\begin{proof}
Put \(M_n=\binom n{a_n}\) and
\(L_n=\ceil{M_n\abs{\B_n}^{-w_n}}\).  Exact-agreement reduction and
fixed-support slope uniqueness give the unconditional upper bound
\[
 B_{C_n}^{\rm MCA}(a_n)\le\min\{\abs{\F_n},M_n\}.  \tag{AS5}
\]
The exact locator-prefix list and collision-aware pole compiler give
\[
 B_{C_n}^{\rm MCA}(a_n)\ge
 \left\lceil
 \frac{L_n(\abs{\F_n}-n)}
 {\abs{\F_n}-n+k_n(L_n-1)}
 \right\rceil.                                     \tag{AS6}
\]
By Stirling's formula and \textup{(AS1)},
\(\log L_n=h(\alpha)n+o(n)\).  For positive \(x,y\),
\(xy/(x+k_ny)\ge\frac12\min\{y,x/k_n\}\); hence
\textup{(AS2)}, \(\log k_n=o(n)\), and the ceiling show that the
right side of \textup{(AS6)} has logarithm
\(n\min\{h(\alpha),\phi\}+o(n)\).  The upper bound
\textup{(AS5)} has the same exponent, proving \textup{(AS3)}; subtracting
\(\log\abs{\F_n}\) proves \textup{(AS4)}.  Finally
\(w_n\log n=o(n)\), so \((a_n-k_n)/n\to0\) and
\(\alpha=\rho\).
\end{proof}

\begin{corollary}[Unconditional large-challenge RS--MCA threshold]
\label{cor:unconditional-rs-mca-capacity}
Under the row hypotheses of
\cref{thm:unconditional-asymptotic-rs-mca}, let the full-field target
satisfy \(\eps_n=\exp(-\lambda n+o(n))\) for a finite
\(\lambda\ge0\).  If
\[
                         \lambda<\phi-h(\rho),      \tag{AS7}
\]
then, for every sufficiently large \(n\),
\[
 a_n^*=k_n+1,\qquad
 \delta_n^*=1-\frac{k_n+1}{n}=1-\rho+o(1).         \tag{AS8}
\]
Conversely, if \(\lambda>(\phi-h(\rho))_+\), every agreement sequence
satisfying \((a_n-k_n-1)\log\abs{\B_n}=o(n)\) is unsafe.  The critical
boundary \(\lambda=(\phi-h(\rho))_+\) is governed by the
exact finite constants and is not decided by exponent asymptotics.  This
is a large scalar-extension regime.  If \(\abs{\F_n}\) is polynomial in
\(n\), then \(\phi=0\) and \textup{(AS7)} cannot hold; the corollary does
not settle that Proximity-Prize parameter regime.
\end{corollary}

\begin{proof}
At \(a=k_n+1\), \textup{(AS5)} is
\(B_{C_n}^{\rm MCA}(a)\le\binom n{k_n+1}
=\exp(h(\rho)n+o(n))\).  Under \textup{(AS7)},
\(\eps_n\abs{\F_n}=\exp((\phi-\lambda)n+o(n))\) has strictly larger
exponent, and taking the integer floor does not change the comparison.
Thus the first permitted agreement is safe
and \textup{(AS8)} follows.  The converse follows from \textup{(AS4)}
and the strict reverse exponential inequality.
\end{proof}

\begin{remark}[Binary-field scope]
The deep-target theorem, the shallow-prefix exponent, the exact support
atlas, and the adjacent-row theorem are characteristic-free; they therefore
include binary extension fields whenever the evaluation domain exists.
The subsquare-root Fourier theorem
\cref{thm:unconditional-shallow-mi-ma} is deliberately odd-prime-field:
its proof uses \(m<p\) to control every Li--Wan cycle.  In characteristic
two, the exact cycle decomposition and the cubic obstruction remain the
correct interfaces; no binary Fourier-flatness claim is inferred from the
odd-characteristic theorem.
\end{remark}

The upper argument organizes witnesses as follows.  A \emph{compiler} is
a proved implication that converts one family of row certificates---for
example support-fiber, Fourier, or incidence bounds---into the next
distinct-slope or threshold bound.  A \emph{chart} is a
coordinate presentation of part of the witness incidence, including any
auxiliary quotient, planted, tangent, or rank parameters.  A \emph{cell}
is a specified subincidence in such a chart; after the received line and
auxiliary parameters are fixed, its \emph{realized cell} is an actual
finite witness subset.  A \emph{profile} records the cell type together
with those realized parameter values; a boundary profile additionally
records its \emph{full support slice} and \emph{boundary map}.  The slice
is the set of supports compatible with the fixed profile parameters before
first-match deletion; the map records the coefficient or syndrome data
whose equal fibers are counted.  We use ``full support slice'' and ``full
profile slice'' synonymously.

Here a quotient parameter records factorization through a finite folding
map, meaning a map with complete uniform fibers; a
planted parameter records a forced algebraic factor or support block, and
tangent or rank parameters record loci where a projection loses its
expected rank.  An ordered
\emph{first-match atlas} is witness-exhaustive if its realized cells cover
every witness; a slope occurring in several cells is assigned to the first
cell whose actual slope projection contains it.  A cell is \emph{paid}
only when a proved bound controls that distinct-slope projection at the
cell's natural scale---for a boundary profile, its full support count
divided by the size of its realized boundary image, and otherwise an
explicitly stated direct scale.

A \emph{ledger} is the atlas together with all such
payments, the treatment of its primitive residual profiles, and the final
\emph{add-back}, meaning the summation of all first-match profile budgets
into one row-level bound.  It is \emph{closed} when each required payment and ray bound is
a theorem for the row.  The \emph{profile envelope} is the row-level sum
of the natural scales of all realized profiles, maximized over received
lines; it is the budget to which a closed ledger compiles.

Here \emph{quotient descent} means factorization through a folding map,
whereas \emph{field descent} means that the data are defined over a proper
subfield.  \emph{Saturation} records the two actual projections
\((\gamma,S,h)\mapsto(\gamma,h)\mapsto\gamma\), together with their
fiber cardinalities; different explaining polynomials at one slope are not
identified with the slope until the second projection.  The setwise
stabilizer is taken in the folding-symmetry group acting on the evaluation
domain.

A residual profile is called \emph{primitive} only after the atlas has certified the
removal of its named quotient, field-descent, planted, rank, and ray-saturation
degeneracies; primitivity is a ledger certificate, not merely trivial
setwise stabilizer or a topological openness condition.  A \emph{leaf} is
one fixed realized profile remaining at the end of this first-match
process.

The domain families in the title carry additional folding structure.  A
\emph{smooth multiplicative domain} is a coset \(D=\theta H\) of a
subgroup \(H\le\B^\times\), equipped with those power maps
\(x\mapsto x^c\) for which every attained value has exactly \(c\)
preimages in \(D\).

Assume now that \(q\) is an odd prime power.  A \emph{norm-one torus} over \(\F_q\) has rational
points \(\{u\in\F_{q^2}^\times:u^{q+1}=1\}\).  A
\emph{Chebyshev twin-coset domain} is obtained from a disjoint inverse pair
\(gH\sqcup g^{-1}H\) in this group by the projection
\(\chi(u)=(u+u^{-1})/2\); the polynomial \(T_c\) satisfying
\(T_c(\chi(u))=\chi(u^c)\) supplies its complete-fiber folding maps.  We
call this projected set a \emph{circle \(x\)-domain}.

A \emph{standard circle code} instead evaluates bounded-total-degree
bivariate polynomials on points of \(x^2+y^2=1\), as in the circle-code
setting of \cite{HLP24}.  The two presentations
are \emph{diagonally equivalent} only after a coordinate bijection and
coordinatewise multiplication by a fixed nonzero vector have been proved;
that equivalence preserves all agreement sets.  The precise construction,
including the exceptional points at which \(u\) and \(u^{-1}\) coincide,
is \cref{def:circle-domain-code,def:circle-twin-domain}.

\paragraph{Proximity Prize and recent work.}
The Ethereum Foundation's preliminary Proximity Prize asks for the
largest radius at which smooth-domain Reed--Solomon codes of rates
\(1/2,1/4,1/8,1/16\) have MCA error at most a cryptographic target such
as \(2^{-128}\) \cite{ProximityPrize,ABF26}.  The foundational
Reed--Solomon proximity-gap theorem holds at radii below the Johnson bound
\(1-\sqrt{\rho}\) at code rate \(\rho\)
\cite{BCIKS20}.  Recent work sharpens the positive theory in that range
and proves substantial lower obstructions beyond it
\cite{BCHKS25,CS25,Habock25}; general reductions from list decoding and
polynomial-generator MCA supply further routes to proximity and
correlated agreement \cite{GaoCaiXuKan25,BordageChiesaGuanManzur25}.
Near capacity, Goyal and Guruswami prove MCA and proximity gaps for
folded Reed--Solomon, multiplicity, and subspace-design codes and transfer
them to random linear and random-evaluation Reed--Solomon ensembles
\cite{GoyalGuruswami25}.  The random-ensemble parameters were subsequently
improved through row-span-constrained local properties
\cite{GoyalGuruswamiSunWootters26}, while a direct syndrome-space approach
gives near-capacity results for random linear codes
\cite{YuanZhu26}.  Explicit constant-alphabet
Alon--Edmonds--Luby (AEL)-based subspace-design constructions inherit
broad random-linear-code local properties
\cite{GoyalGuruswamiHsieh26}; these are concatenated or additive code
families, not asymptotic ordinary Reed--Solomon codes over \(\F_2\).
The Reed--Solomon proximity
testing protocol WHIR is an important application and terminology source
for MCA \cite{WHIR}.

These results do not automatically cover ordinary Reed--Solomon codes on
a fixed multiplicative coset or a fixed circle domain: random evaluation
points, folding, multiplicity, and subspace-design hypotheses change the
code family.  The S-two whitepaper discusses related line-decoding
questions and conjectures \cite{CGHLL26}, while the explicit counterexamples
proved here apply directly to fixed smooth rows.  We therefore
do not claim the first MCA formalism, the first MCA counterexample, or an
unrestricted capacity theorem.

Our contribution is complementary.  We first give exact finite-row
compilers through syndromes and auxiliary simple poles: the former replaces
a word by its coset modulo the code, while the latter divides by
\(X-\alpha\), \(\alpha\notin D\), to place a list on one received line.
We then construct a target-aware profile envelope and an exact
quotient--remainder normal form separating whole folding fibers from the
\emph{support remainder}, the exceptional selected points outside those
complete fibers.  Odd- and characteristic-two
field-drop examples show how quotient coefficients in a proper subfield
enlarge every \emph{live prefix slot}, meaning a prefix coordinate not
forced by lacunarity (prescribed gaps among the highest locator
coefficients), planted data, or previously fixed quotient data and
therefore still ranging over its recorded coefficient field.  We close the
normalization and support-census interfaces by Fourier inversion on the
effective additive span and by an exact partial-occupancy disintegration.
Exact explanation occupancy, pair-incidence, and curve-degree formulas
then identify every finite ray loss, while an aperiodic pole-line example
shows that occupancy one is possible.  Thus the first-match formalism
isolates, rather than hides, the aggregate Fourier and residual
distinct-slope theorems still required for an unrestricted upper result.
Within explicit regimes we close those interfaces: the deep-target
frontier is exact for every row; shallow locator prefixes have an exact
MCA exponent; subexponential boundary groups give an exhaustive coarse
support-paid atlas and direct support-to-slope payment, while the stronger
active-size condition \(\log A_\lambda=o(N_\lambda)\) gives vacuous MI and
all-major MA on its live slices; and bounded kernels, one MDS circuit,
locator curves, pencils, and deep quotients have literal RC constants.
The first adjacent row is also computed exactly over a separating
scalar-extension field.
The finite ledgers and
reductions used by this formulation are developed in
\cite{Cho26CapV13,Cho26Grande}.

For an \(a\)-point support \(S\subseteq D\), its monic
\emph{locator} is
\(Q_S(X)=\prod_{t\in S}(X-t)=X^a+c_1(S)X^{a-1}+\cdots+c_a(S)\).
The \emph{depth-\(w\) locator prefix} is
\(\Phi_w(S)=(c_1(S),\ldots,c_w(S))\), and a \emph{prefix fiber} is a
set \(\Phi_w^{-1}(z)\) of supports with the same prescribed leading
coefficients.  The natural near-capacity lower construction takes
\(w=a-k-1\).  Its \emph{identity profile} is the unprofiled family of
all \(a\)-subsets of \(D\), with no quotient, planted, or field-descent
restriction.  Averaging its \(\binom na\) supports over the ambient
prefix space \(\B^w\) gives the ambient \emph{identity scale}
\[
        \barN_1(a)=\binom na\abs\B^{-w},
        \qquad D\subseteq\B\subseteq\F.
        \tag{1.5}\label{eq:intro-identity-scale}
\]
Here \(\B\) is the coefficient field containing the domain; the code and
challenge field \(\F\) may be a \emph{scalar extension}, meaning that the
same evaluation points are retained while polynomial coefficients and
slopes are allowed in a larger field \(\F\supseteq\B\).  It is tempting to
conjecture that \(\barN_1\), up to \(e^{o(n)}\), also bounds every
received line.  The \emph{identity entropy crossing} is the agreement at
which \(\log\barN_1(a)\) meets the logarithm of the target budget; for a
subexponential target this is the zero of the identity exponent
\(n^{-1}\log\barN_1(a)\).  The
countertheorem shows that the identity scale is not a universal bound without a
quotient envelope.  A large-depth identity prefix can contain a bounded-degree
\emph{quotient profile}, meaning supports assembled from complete fibers
of a folding map, whose live coefficients are defined over a much smaller
subfield.  At the identity entropy crossing that profile is exponentially
larger than \(\barN_1\).  The same construction survives on the circle
after the coordinate bijection and nonzero diagonal rescaling of
\cref{prop:stereographic-circle-equivalence}.

The second point is analytic.  We identify supports with their incidence
vectors in the torsion-free group \(\Z^D\).  For a support family \(F\),
its \emph{additive energy} is
\(E(F)=\abs{\{(x_1,x_2,x_3,x_4)\in F^4:x_1-x_2=x_3-x_4\}}\), normalized
as \(\Delta(F)=E(F)/\abs F^3\).  A large prefix fiber need not have high
energy.  We call a positive-rate contribution from fibers with
exponentially small \(\Delta\) the \emph{Sidon-heavy branch}; the name
records scarcity of repeated differences, not an assertion that the
fiber is a literal Sidon set.

More precisely, for a support map
\(\Phi:\Omega^0\to\Sigma\), write
\(F_s=\Phi^{-1}(s)\), \(f_s=\abs{F_s}\),
\(L=\abs{\Phi(\Omega^0)}\), and
\(\barN=\abs{\Omega^0}/L\).  Its normalized order-\(q\) moment is
\(L^{-1}\sum_s(f_s/\barN)^q\); the Sidon-heavy part restricts this sum to
\(\Delta(F_s)\le e^{-\sigma N}\), where \(N\) is the number of active
coordinates and \(\sigma>0\) is fixed.  A \emph{logarithmic moment order}
is an order \(q=q_N\to\infty\) growing at most polylogarithmically in
\(N\).  We split that moment mass into the low-energy branch and a
high-energy branch.  Here \emph{positive rate} means growth
\(e^{cN-o(N)}\) for a fiber, or \(e^{cNq-o(Nq)}\) for an order-\(q\)
moment, with some fixed \(c>0\).

If the Sidon branch is paid by an explicit Fourier
estimate, the Balog--Szemeredi--Gowers (BSG) theorem and the Boolean-cube
difference-growth theorem, namely
\(\abs{U-U}\ge\abs U^{3/2}\) for \(U\subseteq\{0,1\}^d\), rule out a
positive-rate high-energy fiber.  This is a genuine theorem, but it does
not pay the Sidon branch at the deployed \emph{deep prefix}---the large
locator-prefix depth used at the entropy crossing, unrelated to the
high-agreement syndrome range called the deep regime below---merely by
naming that branch a cell.

Fourier inversion must first be performed on the \emph{effective additive
span}, the group generated by differences of attained column values;
ambient characters annihilating this span cancel exactly against the
ambient normalization.  On the resulting effective dual, the exact
minor-character requirement is the aggregate estimate \textup{(MI)}.
The Li--Wan distinct-coordinate sieve combined with Weil square-root
cancellation \cite{LiWanSieve,LiWanSubset,Weil} supplies \textup{(MI)} only under its displayed
degree-uniform pointwise certificate \textup{(PF)}.  The \emph{rational
phase} is the rational function inside the corresponding
additive-character sum; it is \emph{nondegenerate} when it is nonconstant
and not Artin--Schreier.  The remaining designated \emph{major
characters}, including nonconstant modes requiring quotient or recursive
treatment, require the separate aggregate estimate \textup{(MA)}.  A
Frobenius coboundary \(G^p-G\) plus a constant has constant trace phase;
when it is defined on the whole active set, effective-span normalization
removes it as an annihilator mode.  Neither smoothness nor the name of an
algebraic cell proves the remaining \textup{(MI)} or \textup{(MA)} sum.

The third point concerns the MCA numerator.  The \emph{Q statistic} is the
largest support fiber of a profile boundary map, and \emph{SP} is its
second moment, equivalently the number of ordered support pairs with the
same boundary value.  These quantities count supports or support pairs,
whereas MCA counts distinct slopes.  A raw witness
\((\gamma,S,h)\) first projects to its \emph{explanation state}
\((\gamma,h)\) and only then to the finite slope \(\gamma\); different
explanation states may have the same slope.  Equivalently, a finite slope
may be represented by the projective line parameter \([1:\gamma]\), but
it is not the same object as the explaining polynomial or codeword.  A
\emph{ray compiler} is a proved bound for this final slope projection,
obtained directly or from a proved lower fiber multiplicity.  The
one-parameter moving-root
theorem supplies such a count after a chart has genuinely been reduced to
one projective locator pencil \(Q_0+\lambda Q_1\).  It does not exhaust the high-dimensional
\emph{balanced-core} kernels, meaning residual locator families with a
factored common core (a shared locator factor), equal-degree monic moving locators with a common
prefix, and more than one live projective parameter, occurring
at the frontier.  The exact direct-bound-or-multiplicity hypothesis used by
the compiler is denoted \textup{(RC)}.

For the finite theorem below, a \emph{syndrome} is the parity-check image
of a received word.  A syndrome line meets the span \(V_E\) of the
parity-check columns indexed by a possible error set \(E\).  The meeting is
\emph{transverse} when the whole syndrome line is not contained in
\(V_E\), so that this support supplies at most one slope.  Thus a
``transverse syndrome-line incidence'' is a distinct line parameter
\(\gamma\) for which that meeting occurs.

\begin{theorem}[Unconditional finite-row theorem package]
\label{thm:main-unconditional}
The following assertions hold for every finite field row in their displayed
ranges.
\begin{enumerate}[label=\textup{(\roman*)},leftmargin=*]
\item In the \emph{deep range} \(3(n-a)\le n-k\), where three permitted
      error supports are still too small to support a nonzero RS codeword,
      the full-field MCA numerator is
      exactly \(n-a+1\); for a challenge subset \(\Gamma\subseteq\F\) it is
      \(\min\{\abs\Gamma,n-a+1\}\).
\item The bad-slope numerator is exactly the transverse syndrome-line
      incidence \eqref{eq:exact-secant-numerator}; in particular one fixed
      \emph{noncommon support}---a support on which the two received words
      are not both explained---carries at most one finite bad slope.
\item Writing
      \(L=\lceil\binom na\abs\B^{-(a-k-1)}\rceil\), locator prefixes give
      an exact list of size at least \(L\) in the dimension-\(k+1\)
      code; when \(\abs\F>n\), the collision-aware pole transform gives at least
      \[
      \left\lceil\frac{L(\abs\F-n)}{\abs\F-n+k(L-1)}\right\rceil
      \]
      distinct MCA-bad slopes in the dimension-\(k\) code.
\item Every complete depth-\(w\) prefix fiber
      \(\Fcal_z=\Phi_w^{-1}(z)\) of cardinality \(N\), with
      \(0\le w\le m-2\), is realized after any scalar extension satisfying
      \(\abs\F-n>(m-w-1)\binom N2\) by one received line for
      \(\RS_\F(D,m-w-1)\) with exactly \(N\) bad slopes; the fiber, slope
      set, and exact-\(m\) witness incidence are in bijection.  Under the
      corresponding condition for an \(N\)-element subfamily, that
      subfamily furnishes at least \(N\) distinct bad slopes, but the line
      may also contain slopes from the rest of the complete fiber.
\item For every complete degree-\(c\) folding map, a depth-\(w\)
      complete-fiber-plus-remainder prefix with remainder degree \(r<c\)
      has the exact finite-row normal form of
      \cref{thm:exact-quotient-remainder-normal-form}.  Along every
      sequence satisfying \cref{prop:identity-quotient-comparison}, fixed
      \(c,r\) give the field-sensitive identity-crossing exponent
      \eqref{eq:qr-comparison-crossing}.
\item Q and SP support statistics do not determine the syndrome-line
      incidence.  The precise multiplicity condition under which a pair
      count compiles to distinct rays is
      \eqref{eq:pair-ray-multiplicity}.
\item For every challenge set,
      \(B_{C,\Gamma}^{\rm MCA}(a)\le
      \min\{\abs\Gamma,\binom na\}\) exactly.  At the first adjacent
      agreement \(a=k+1\), if
      \(\abs\F>\max\{\binom n{k+1},
      \binom{\binom n{k+1}}2\}\), the full-field numerator is
      exactly \(\binom n{k+1}\), and the adjacent threshold alternatives
      are the literal comparisons in \textup{(AD2)}.
\end{enumerate}
\end{theorem}

The assertions are proved in
\cref{sec:syndrome-exact,prop:exact-support-upper,thm:exact-first-adjacent-row,sec:exact-prefix-pole,sec:quotient-obstruction,thm:prefix-to-line-hardness,cor:exact-prefix-ray-realization,sec:compiler}.

\begin{theorem}[Exact closure formulas and barriers]
\label{thm:intro-closure-barrier-package}
The following identities and barriers require no ledger-admissibility
hypothesis.  Where an item assumes an incidence-degree or curve-cover
condition, the unconditional assertion is the resulting exact implication;
it does not assert that the condition holds for every row.  In particular,
outside the explicit regimes stated below, this package does not discharge
\textup{(MI)}, \textup{(MA)}, or \textup{(RC)} for unrestricted rows.
\begin{enumerate}[label=\textup{(\roman*)},leftmargin=*]
\item Fourier inversion for a prefix chart has the exact effective-span
      form \textup{(EF2)}--\textup{(EF4)}.  Ambient annihilator modes are
      removed with no loss, and the finite minor and major losses are their
      literal normalized \(\ell^1\)-Fourier masses.  On the unweighted
      prime-field coset and twin-coset power-sum charts satisfying all
      hypotheses of \cref{thm:unconditional-shallow-mi-ma},
      \cref{thm:unconditional-shallow-mi-ma} verifies effective
      \textup{(MI)} and \textup{(MA)} unconditionally.
\item Every complete uniform-fiber folding map has the exact canonical
      partial-occupancy disintegration \textup{(PO1)}; all such slices add
      back exactly to \(\binom na\) by \textup{(PO2)}, and their remaining
      Fourier obligation is exactly \textup{(PO3)}--\textup{(PO4)}.  The
      canonical support atlas is unconditionally exhaustive.  Condition
      \textup{(PO8)} closes only the coarse support-paid envelope
      \textup{(PO9)}; the automatic empty-minor/all-major verification of
      slice MI and MA uses \textup{(PO10)}, and a target frontier
      additionally requires the literal reserve comparisons in
      \textup{(SB2)}.
\item Witnesses project through explanation states to slopes with the
      exact occupancy identities \textup{(RC$_{\rm occ}$)}--\textup{(RC2)}.
      Under their displayed incidence-degree and curve-cover hypotheses,
      pair incidences satisfy \textup{(RC$_{\rm pair}$)} and locator
      curves satisfy \textup{(RC$_{\rm curve}$)}.  The finite
      first-match compiler is the integer formula \textup{(FC1)}--
      \textup{(FC2)}, and the adjacent identity ratio is exactly
      \textup{(AR1)}.  Bounded residual kernels, one fixed MDS circuit,
      exact-support curve covers and pencils, and deep invariant quotient
      cells have the unconditional direct payments
      \textup{(RC$_{\rm ker}$)}, \textup{(RC$_{\rm circ}$)},
      \textup{(RC$_{\rm loc}$)}, and \textup{(RC$_{\rm quot}$)}.
\item The degree-uniform Weil certificate \textup{(PF)} fails numerically
      at the usual deep polynomial-field entropy prefix, even on
      \(\F_p^\times\).  Before effective-image routing, the square-field
      construction forces an ambient PF/MA dichotomy.  In characteristic
      two, an exponentially large aperiodic prefix fiber can be realized
      bijectively as exact-agreement bad slopes with retained occupancy one.
\end{enumerate}
Consequently the effective-normalization identities and support census are
exact, and the finite compiler displays every loss as a row-dependent
finite parameter.  An unrestricted smooth/circle profile-envelope theorem
still requires a witness-exhaustive first-match atlas, image-scale
effective Fourier payment or a direct Sidon payment, residual ray bounds
for the remaining higher-dimensional balanced cores and uncontrolled
unions of MDS-circuit slope projections, complete envelope domination, and
the target and lower reserve comparisons.
\end{theorem}

\begin{proof}
Part~\textup{(i)} is
\cref{lem:effective-span-fourier,def:effective-fourier-payment,cor:exact-finite-fourier-constants,thm:unconditional-shallow-mi-ma};
part~\textup{(ii)} is
\cref{thm:exact-partial-occupancy,prop:partial-occupancy-fourier,thm:canonical-partial-occupancy-atlas};
part~\textup{(iii)} is
\cref{lem:exact-occupancy-compiler,cor:finite-pair-ray-compiler,prop:curve-degree-ray-compiler,thm:bounded-residual-kernel-ray,thm:single-mds-circuit-ray,cor:automatic-locator-curve-ray,prop:deep-invariant-quotient-ray,thm:exact-finite-profile-compiler,lem:exact-adjacent-identity-ratio};
and part~\textup{(iv)} is
\cref{prop:pf-deep-prefix-barrier,prop:pf-ma-square-dichotomy,thm:aperiodic-ray-obstruction}.
\end{proof}

\begin{theorem}[Countertheorem to unrestricted identity-only bounds]
\label{thm:intro-countertheorem}
Fix \(0<\alpha<1/2\).
\begin{enumerate}[label=\textup{(\alph*)},leftmargin=*]
\item There are odd-characteristic smooth multiplicative rows with
      \(\log|\B_n|=O(\log n)\), \(a_n/n\to\alpha\), and
      \(\barN_1(a_n)=e^{o(n)}\), but a square-quotient prefix fiber---that
      is, a profile of complete fibers of \(x\mapsto x^2\)---has size
      \(\exp((h(\alpha)/4+o(1))n)\), exponentially small normalized
      additive energy, and, after scalar extension, that many distinct
      MCA-bad slopes.  The same obstruction transfers to standard
      bivariate circle codes.
\item There are characteristic-two smooth multiplicative rows with the
      same ambient identity-scale behavior but a cubic-quotient fiber of
      complete \(x\mapsto x^3\) fibers of size
      \(\exp((h(\alpha)/6+o(1))n)\), again giving exponentially many
      distinct MCA-bad slopes after scalar extension.
\end{enumerate}
Consequently smoothness and \(\log|\B_n|=o(n)\) do not imply a uniform
ambient identity-normalized Sidon payment or an identity-scale MCA
numerator bound.  They also invalidate any unrestricted identity-frontier
formula derived sharply from those assertions; threshold counterexamples result
for target budgets below the constructed distinct-slope count.
\end{theorem}

\begin{proof}
Part~\textup{(a)} is
\cref{thm:smooth-quotient-obstruction,cor:circle-quotient-obstruction};
part~\textup{(b)} is
\cref{thm:binary-cubic-obstruction}.  The MCA statements use the exact
prefix-to-line compiler of \cref{thm:prefix-to-line-hardness}.
\end{proof}

\begin{remark}[Threshold and field-size scope]
\label{rem:intro-countertheorem-scope}
Here \(w_n/n=\Theta(1/\log n)=o(1)\).  The countertheorem therefore
refutes ambient identity-normalized Sidon payment, identity-only numerator
bounds, and \emph{row-sharp} threshold derivations (derivations retaining
the row's actual sublinear displacement rather than absorbing it into an
unspecified \(o(1)\) radius error), but it does not by itself
refute a coarse radius identity with an \(o(1)\) error.  The domain field
\(\B_n\) is polynomial-size, whereas the scalar extension used to separate
all exponentially many slopes has \(\log\abs{\F_n}=\Theta(n)\).  A claim
for a prescribed Proximity-Prize target must still use the exact
target-aware bracket and verify its field and reserve hypotheses.
\end{remark}

For the upper theorem, a realized first-match profile \(\lambda\) has a
\emph{full support slice} \(\Omega^0_\lambda\), namely all supports
compatible with its fixed profile parameters before first-match deletion,
a coefficient field \(\B_\lambda\), and \(R_\lambda\) effective prefix
equations.  It comes with a proved \emph{boundary map}
\(\Phi_\lambda:\Omega^0_\lambda\to\B_\lambda^{R_\lambda}\).
For a received line \((r_0,r_1)\) and agreement \(a\), let
\(\Lambda(r_0,r_1;a)\) be the finite set of its realized profiles whose
first-match cells are nonempty at that agreement.  Put
\[
 A_\lambda=\abs{\B_\lambda}^{R_\lambda},\qquad
 L_\lambda=\abs{\Phi_\lambda(\Omega^0_\lambda)},\qquad
 \barN_\lambda=\frac{\abs{\Omega^0_\lambda}}{L_\lambda},
\]
and define the image-normalized profile envelope
\[
 \mathfrak E_n(a)=1+(n-a+1)+
 \sup_{(r_0,r_1)}
 \sum_{\lambda\in\Lambda(r_0,r_1;a)}(1+\barN_\lambda).
 \tag{1.6}\label{eq:profile-envelope}
\]
The formal codomain has \(A_\lambda\) points, the realized image has
\(L_\lambda\) points, and \(\barN_\lambda\) is therefore the average
full-slice fiber size at the realized image scale.  For each received line
the sum includes the identity profile and all
quotient, Chebyshev, planted, and remainder profiles that meet that line;
the supremum makes the envelope a row-level quantity.  With subexponentially
many profiles, the sum and maximum have the same exponential scale.  The
ambient scale
\(\barN_\lambda^{\rm amb}=|\Omega^0_\lambda|/A_\lambda\) may replace
\(\barN_\lambda\) only after a full-image certificate
\[
 L_\lambda\ge e^{-o(n)}A_\lambda                       \tag{FI}
\]
has been proved uniformly.  Otherwise effective-image collapse is
\emph{routed}---assigned by the first-match rule to an earlier profile so
that its slopes and all witnesses above them are removed from the later
residual---as an effective-image rank-collapse profile, and the image scale
is retained.

Here
\emph{effective-image collapse} means that the realized image is
exponentially smaller than the formal codomain.  This is distinct from a
ray-saturation profile, which records many raw witness representatives
producing the same explanation or slope ray.

The countertheorem is exactly a row for which a quotient profile in
\(\mathfrak E_n\) is exponentially larger than the ambient identity term.
A \emph{shift-pair chart} records two residual supports with the same
boundary value; a balanced-core chart records the higher-dimensional
locator kernel that may remain after their common core is factored.

\begin{definition}[Ledger-admissible smooth/circle sequence]
\label{def:admissible-sequence}
A sequence \((C_n,a_n)\), with \(C_n=\RS_{\F_n}(D_n,k_n)\), is
ledger-admissible if the following conditions hold uniformly in every
received line.
\begin{enumerate}[label=\textup{(A\arabic*)},leftmargin=*]
\item The domains are among the structured classes of
      \cref{def:structured-folding,def:circle-twin-domain}, or are circle
      codes diagonally equivalent to such RS rows, and the folding maps
      have complete uniform fibers.
\item A first-match atlas covers every bad-slope witness and has
      \(e^{o(n)}\) profiles.  The total \emph{distinct-slope} contribution
      of its algebraic cells is at most \(e^{o(n)}\mathfrak E_n(a_n)\).
      Partial occupancies of every uniform folding map use the canonical
      slices of \cref{thm:exact-partial-occupancy}; any comparison of
      their image-normalized budgets records the image-coverage and overlap
      factors of \cref{lem:exact-profile-addback}.
\item Every primitive profile uses the image normalization above.  Its
      leaf has an \emph{active coordinate set} of size \(N\), namely the
      evaluation positions still free to enter or leave the support after
      the profile parameters have been fixed, and has fixed support size
      \(m_N\), and density
      \(m_N/N\in[\alpha,1-\alpha]\) for one row-uniform constant
      \(\alpha>0\).  Writing \(L_N\) for the size of its realized boundary
      image, there is a logarithmic order \(q_N\to\infty\),
      \(q_N\le(\log N)^C\), satisfying
      \(\log L_N/q_N=o(N)\), unless it carries a separately proved
      primitive-Q bound.  Any use of the ambient scale is accompanied by
      \textup{(FI)}.
\item Fourier inversion on every primitive prefix leaf is normalized on
      the effective span of \cref{lem:effective-span-fourier}, and the
      leaf carries the certified effective partition of
      \cref{def:effective-major-minor}.  The aggregate minor payment \textup{(MI)} of
      \cref{def:aggregate-minor-payment} and the effective major payment
      \textup{(MA)} hold there, or the leaf has a separately proved
      image-normalized Sidon/Fourier moment payment of the strength in
      \cref{def:sidon-paid-cell}.  The pointwise condition \textup{(PF)}
      may be used only as a sufficient certificate for \textup{(MI)} in
      the range where its numerical inequality and all certified-lift
      hypotheses are verified.
\item Primitive columns are genuine weighted Vandermonde columns, meaning
      columns of the form
      \(\rho(t)(1,t,\ldots,t^{R_N-1})^{\mathsf T}\) at distinct
      \(t\)'s with \(\rho(t)\ne0\).  Every
      use of power-sum coordinates satisfies
      \(R_N<\operatorname{char}\B_N\); small-characteristic leaves instead
      retain elementary coordinates or carry a direct Sidon/Q theorem.
\item Every residual shift-pair and balanced-core chart satisfies the ray
      compiler \textup{(RC)} of \cref{hyp:ray-compiler}, or has a direct
      distinct-slope bound at most its contribution to
      \(e^{o(n)}\mathfrak E_n(a_n)\).
\item The profile envelope in \eqref{eq:profile-envelope}, including all
      quotient subfields \(\B_\phi\) with
      \(\phi(D)\subseteq\eta\B_\phi\) for some scalar \(\eta\)
      (the \emph{scaled quotient coefficient fields}) and all profile fields of
      definition, is used in
      the final budget; it is not replaced by the identity term unless that
      comparison is proved for the row.
\end{enumerate}
\end{definition}

\begin{theorem}[Conditional profile-envelope compiler]
\label{thm:main-smooth-circle}
Every ledger-admissible smooth/circle sequence satisfies
\[
        B_{C_n}^{\rm MCA}(a_n)
        \le e^{o(n)}\mathfrak E_n(a_n).
        \tag{1.7}\label{eq:conditional-numerator}
\]
\end{theorem}

\begin{proof}
Condition~\textup{(A2)} supplies a witness-exhaustive first-match atlas
and pays its algebraic slope projections.  On every remaining primitive
profile, condition~\textup{(A4)} supplies the certified Sidon/Fourier
payment, either directly or through
\cref{prop:effective-mi-ma-flatness}, while conditions~\textup{(A3)}--\textup{(A5)} and
\cref{thm:primitive-q} give image-normalized Q.  Condition~\textup{(A6)}
and \cref{prop:q-implies-sp} convert those bounds to distinct slopes, and
\textup{(A7)} places every profile at its term in
\(\mathfrak E_n(a_n)\).  Summing the \(e^{o(n)}\) first-match payments and
applying \cref{lem:first-match-bound} proves
\eqref{eq:conditional-numerator}.
\end{proof}

For the lower side, a \emph{certified profile-list construction} is an
explicit list inside one identity, quotient, Chebyshev, or remainder
profile, together with the pole and challenge-intersection estimates that
prove its bad-slope count exceeds the target.  A profile-envelope budget is
\emph{safe} only when the closed ledger proves its total distinct-slope
count is at most the target.  The next theorem brackets the frontier
between one certified unsafe agreement and one certified safe agreement.

\begin{theorem}[Target-aware threshold bracket]
\label{thm:intro-asymptotic-rs-mca}
Let \(a_{-,n}<a_{+,n}\) be agreements for which a certified profile-list
construction is unsafe at \(a_{-,n}\) and the closed-ledger
profile-envelope budget is safe at \(a_{+,n}\).  Then the threshold
in \eqref{eq:intro-asymptotic-threshold} exists and
\[
 a_{-,n}<a_n^*\le a_{+,n},\qquad
 1-\frac{a_{+,n}}n\le\delta_n^*<1-\frac{a_{-,n}}n.
 \tag{1.8}\label{eq:intro-threshold-bracket}
\]
In particular, if \(a_{+,n}-a_{-,n}=o(n)\), \(k_n/n\to\rho\), and the
common crossing is \(a_n=k_n+1+g_nn+o(n)\), then
\(\delta_n^*=1-\rho-g_n+o(1)\).
\end{theorem}

We call a range of consecutive agreement values, or an \emph{agreement
window}, \emph{identity-dominant} when the entire profile envelope
is, up to \(e^{o(n)}\), bounded by the identity scale together with the
universal constant and deep terms.  Only in such a window can the profile
compiler reduce to the familiar one-variable identity frontier.

\begin{corollary}[Identity-frontier corollary]
\label{cor:intro-identity-frontier}
Assume the full-field challenge \(\Gamma_n=\F_n\), and suppose in a
window that
\[
 \mathfrak E_n(a)\le e^{o(n)}
 \bigl(1+n-a+\barN_1(a)\bigr).
\]
Put \(\rho_n=k_n/n\), \(\beta_n=\log_2|\B_n|\),
\(\tau_n=n^{-1}\log_2(1+B_n^*)\), and
define
\[
 F_n(g)=H_2(\rho_n+g)-\beta_ng,\qquad
 g_{T,n}=\sup\bigl(\{g\in[0,1-\rho_n]:F_n(g)\ge\tau_n\}\cup\{0\}\bigr).
 \tag{1.9}\label{eq:intro-target-crossing}
\]
The \emph{right crossing} is this rightmost point of the superlevel set
\(F_n(g)\ge\tau_n\).  It is \emph{interior} when
\(0<g_{T,n}<1-\rho_n\) and the function passes from the superlevel side
to the sublevel side there.  By continuity, at an interior right crossing,
\[
 H_2(\rho_n+g_{T,n})-\beta_ng_{T,n}=\tau_n.
\]
If the exact safe and unsafe tests furnish agreements
\(a_{\pm,n}=k_n+1+g_{T,n}n+o(n)\), then
\[
 \delta_n^*=1-\rho_n-g_{T,n}+o(1).
\]
If additionally \(\log(1+B_n^*)=o(n)\), put
\[
 g^*(r,b)=\sup\{g\in[0,1-r]:H_2(r+g)\ge bg\}.
\]
Then \(g_{T,n}=g^*(\rho_n,\beta_n)+o(1)\) and
\[
 \delta_n^*
 =1-\rho_n-g^*(\rho_n,\beta_n)+o(1).              \tag{1.10}
\]
For a proper challenge subset the exact bracket remains valid, but this
corollary additionally requires a certified challenge-intersection lower
bound; \(|\Gamma_n|\) alone does not determine it.
\end{corollary}

The proof is not a change of notation: the upper implication uses the
distinct-ray ledger, the lower implication uses an actual profile list and
the collision-aware pole line, and monotonicity then gives
\eqref{eq:intro-threshold-bracket}.  The complete argument, including the
Stirling error and the \emph{target reserve}---the quantitative margin by
which the safe and unsafe integer bounds clear \(B_n^*\)---is given in
\cref{sec:frontier}.

The profile-envelope compiler is conditional only on the explicit
enumerative statements in \textup{(A1)}--\textup{(A7)}, not on an unnamed
closed ledger.  The analytic implication
\[
 \begin{aligned}
 &\text{Sidon payment}+\text{positive-rate moment excess}
   \Longrightarrow\text{high energy},\\
 &\text{high energy}\Longrightarrow\text{BSG small doubling}
   \Longrightarrow\text{contradiction}
 \end{aligned}
\]
is proved in full.

The polynomial-field construction below proves that a
uniform ambient identity-normalized Sidon payment on the unrestricted deep
prefix is false, even after
exact stabilizers have been destroyed by a bounded planted remainder.

The exact bounded-remainder comparison identifies a missing exponential
profile term rather than removing it.  Arbitrary positive-density
\emph{partial-fiber occupancy}---supports selecting neither zero nor all
points in a positive-density family of folding fibers---is not exhausted
by that normal form, but its support atlas and add-back are closed by
\cref{thm:canonical-partial-occupancy-atlas}.  Outside the deep threshold
regime, the shallow exponent and large-challenge threshold regime, and the
low-boundary upper-compiler regime proved above, what
remains genuinely open is a Sidon payment on the residual obtained
after every quotient, subfield, planted, and remainder profile has been
removed, together with uniform profile-envelope domination of the exact
partial-occupancy terms and a higher-dimensional transverse-secant bound
beyond the bounded-kernel, single-circuit, curve, pencil, and deep-quotient
cases proved here.  These are logically coupled conditions; they cannot
be replaced by the false unrestricted identity numerator theorem.
The low-boundary compiler by itself is an upper theorem and does not
locate a target frontier; the latter additionally needs the literal safe
and unsafe reserves in \textup{(SB2)} or the exact deep comparison.

\begin{definition}[Closed asymptotic ledger]\label{def:closed-asymptotic-ledger}
A ledger-admissible sequence has a \emph{closed asymptotic ledger} when,
uniformly in every received line, all of the following are supplied as
theorems rather than left as hypotheses:
\begin{enumerate}[label=\textup{(L\arabic*)},leftmargin=*]
\item its \emph{exact-agreement witness incidence}, meaning the witnesses
      restricted to supports \(S\) with \(\abs S=a_n\), has a witness-exhaustive
      first-match atlas in the sense of \cref{def:first-match}, with
      \(e^{o(n)}\) realized profiles;
\item every algebraic or directly counted profile has a certified bound
      for its actual first-match slope set \(Z_i^\circ\), meaning the
      slopes assigned to profile \(i\) but to no earlier profile, and the sum of
      those bounds is at most its stated contribution to
      \(e^{o(n)}\mathfrak E_n(a_n)\);
\item every primitive first-match residual of
      \cref{def:primitive-first-match-residual} has its image-normalized
      scale, Sidon/Fourier or direct primitive-Q estimate, and distinct-ray
      compiler proved at its stated profile scale; and
\item effective-image collapse is either routed to an earlier profile or
      the full-image certificate \textup{(FI)} is proved before an ambient
      scale is used, and all remaining profile budgets add back to at most
      \(e^{o(n)}\mathfrak E_n(a_n)\).
\end{enumerate}
Thus closure is an enumerative statement about distinct slope projections.
Constructibility, a raw support count, or a support-pair moment alone does
not close a ledger.
\end{definition}

\begin{theorem}[Compiler for closed ledgers]\label{thm:main-ledger}
A closed asymptotic ledger satisfies \eqref{eq:conditional-numerator} and
therefore supplies the safe side of the target problem.  Whenever a
certified profile-list construction supplies the unsafe side, the threshold
bracket of \cref{thm:intro-asymptotic-rs-mca} follows.  Under the additional
identity-dominance, target-reserve, and exact lower-side hypotheses, this
specializes to \cref{cor:intro-identity-frontier}.
\end{theorem}

\begin{proof}
The primitive analytic theorem is proved in
\cref{sec:sidon-split,sec:fourier,sec:primitive-q-proof}.  The ray step under
\textup{(RC)} is \cref{prop:q-implies-sp}, first-match summation is
\cref{lem:first-match-bound}, and the target-aware entropy comparison is in
\cref{sec:frontier}.  Adding the algebraic, primitive, and direct-ray
budgets gives \eqref{eq:conditional-numerator}.
\end{proof}

\begin{theorem}[Primitive Sidon-heavy implication]
\label{thm:intro-sidon-heavy-repair}
For each \(N\), let \(\Omega_N^0\subseteq\{0,1\}^{T_N}\) be a
full profile slice with \(\abs{T_N}=N\) and fixed support weight \(m_N\),
where \(m_N/N\in[\alpha,1-\alpha]\) for some constant \(\alpha>0\).  Let
\(\Phi_N:\Omega_N^0\to\B_N^{R_N}\), and put
\[
 \Scal_N=\Phi_N(\Omega_N^0),\qquad L_N=\abs{\Scal_N},
 \qquad \barN_N=\abs{\Omega_N^0}/L_N.
\]
Let \(\Omc_N\subseteq\Omega_N^0\) be the primitive first-match residual
after the algebraic profiles have been removed.  For \(s\in\Scal_N\), put
\(F_{N,s}=\Omc_N\cap\Phi_N^{-1}(s)\) and
\(f_{N,s}=\abs{F_{N,s}}\).  Assume there is
a logarithmic order \(q_N\to\infty\), \(q_N\le(\log N)^C\), such that
\[
             \frac{\log L_N}{q_N}=o(N).
 \tag{1.11}\label{eq:intro-moment-accessibility}
\]
If the image-normalized Sidon-heavy moment at this order is
\(e^{o(Nq_N)}\) for every fixed energy cutoff
\(\Delta(F_{N,s})\le e^{-\sigma N}\), \(\sigma>0\), then the leaf satisfies
primitive Q:
\[
        \max_{s\in\Scal_N}\abs{\Omc_N\cap\Phi_N^{-1}(s)}
        \le e^{o(N)}\barN_N .
\]
The same conclusion follows from an ambient Fourier payment only when that
payment is proved directly; such a bound also certifies
\(\abs{\Scal_N}\ge e^{-o(N)}\abs{\B_N}^{R_N}\) by
\cref{rem:flatness-certifies-image}.
Consequently the Sidon-heavy branch is not an additional inverse theorem: once
its Fourier/effective-image payment is verified, any remaining positive-rate
\Renyi{} excess---exponential excess of the normalized high moment
\(L_N^{-1}\sum_s(f_{N,s}/\barN_N)^{q_N}\)---is high-energy and is impossible
in a Boolean primitive leaf.
\end{theorem}

\begin{proof}
This is \cref{thm:primitive-q}.  For completeness, if primitive Q failed
at positive exponential rate, the elementary lower moment inequality
\(L_N^{-1}\sum_sR_s^{q_N}\ge L_N^{-1}(\max_sR_s)^{q_N}\), with
\(R_s=f_{N,s}/\barN_N\), and
\eqref{eq:intro-moment-accessibility} would make the ordinary normalized
moment at the stipulated order have positive-rate excess.
The Sidon payment removes the low-energy contribution.  The
ordinary-moment split of \cref{prop:ordinary-moment-split}, which
separates low- and high-energy fiber contributions, then produces a fiber
of size \(e^{\Omega(N)-o(N)}\barN_N\), hence
\(e^{\Omega(N)-o(N)}\), with normalized energy \(e^{-o(N)}\).  By
Balog--Szemeredi--Gowers and Boolean-cube difference growth,
\cref{prop:high-energy-impossible}, such a fiber has only
\(e^{o(N)}\) points, a contradiction.
\end{proof}

\begin{remark}[Scope of the primitive implication]
\label{rem:intro-sidon-heavy-scope}
The theorem is the precise primitive implication ``pay structure, flatten
Sidon, inverse high energy.''  It does not prove the source estimate
\textup{(MI)} or \textup{(MA)} (with \textup{(PF)} only a sufficient
minor certificate), window-uniformity in the agreement
(uniform constants for every integer agreement in the stated interval), the
effective-image certificate, profile add-back, the lower-side reserve, or
the conversion from support pairs to distinct MCA rays.  Those remain the
separate normalization, profile, Fourier, and ray interfaces recorded in
\cref{def:admissible-sequence}.
\end{remark}

\paragraph{Terminology map.}
The following summary fixes the meaning of the recurring paper-specific
terms; the later definitions supply their scheme-theoretic and finite-row
forms.
\begin{description}[leftmargin=0pt,labelindent=0pt]
\item[Row and domain.] A row is \((\F,D,k)\), with coordinate/evaluation
domain \(D\); a row sequence allows all three data to vary with
\(\abs D=n\).
\item[Explanation, bad slope, and witness.] Explanation means agreement
with a degree-less-than-\(k\) polynomial on the named support.  A bad slope
has a witness \((\gamma,S,h)\) for which the line point is explained on
\(S\) but the received pair is not simultaneously explained on that same
\(S\).
\item[Locator, prefix, and identity profile.] The locator \(Q_S\) is the
monic polynomial with root set \(S\); a prefix fiber fixes its leading
nontrivial coefficients.  The identity profile is the full unprofiled
support family, and \(\barN_1\) is its ambient average scale.
\item[Folding, quotient, and circle domain.] A folding map has complete
uniform fibers; a quotient support is a union of those fibers.  A smooth
multiplicative domain uses power foldings on a multiplicative coset, while
a Chebyshev twin-coset domain projects two inverse torus cosets and uses the
descended Chebyshev foldings.
\item[Chart, cell, profile, and atlas.] A chart gives witness coordinates,
a cell specifies a subincidence, and its realization is the corresponding
finite witness subset after all parameters are fixed.  A realized profile
records the cell type and parameter values; a boundary profile also records
its full support slice and boundary map.  An
ordered atlas is witness-exhaustive and uses actual first-match slope
projections.
\item[Payment and ledger.] A payment is a proved distinct-slope bound at
the relevant profile scale.  A ledger is the atlas, all its payments, the
primitive residual analysis, and their add-back; it is closed only when
every required bound is verified for the row.
\item[Image scale and envelope.] For a full profile slice, the image scale
is support count divided by the size of the realized boundary image.  The
ambient codomain scale may replace it only after \textup{(FI)}.  The
profile envelope sums these realized scales, together with the universal
finite terms, uniformly over received lines.
\item[Q, SP, and rays.] Q is a profile max-fiber bound, SP is the
same-boundary support-pair count, and a ray bound counts distinct slopes.
The passage from Q/SP to rays is precisely the separate \textup{(RC)}
interface.
\item[Syndrome, energy, and Fourier payment.] A syndrome is a parity-check
image, and the secant compiler counts where its received line meets error-
support spans.  Additive energy counts repeated support differences.  The
Sidon-heavy moment is the low-energy part of a normalized fiber moment;
\textup{(MI)} is the required aggregate minor payment, \textup{(PF)} is
a degree-uniform sufficient certificate for it, and \textup{(MA)} pays
the exceptional effective characters in aggregate.
\item[Frontier.] The target-aware frontier is the first safe agreement, or
equivalently its radius \(1-a/n\); the identity frontier is its conditional
specialization when the identity profile dominates the full envelope.  A
target reserve is the explicit margin by which a certified numerator lies
on the required side of the target budget.
\end{description}

\section{MCA witnesses and first-match ledgers}\label{sec:mca-ledger}

Throughout this section \(\F\) is a finite field, \(D\subseteq\F\) has
size \(n\), and \(C=\RS_\F(D,k)\).  Write
\(\F[X]_{<k}=\{f\in\F[X]:\deg f<k\}\), and let
\(\operatorname{ev}_S:\F[X]_{<k}\to\F^S\) be restriction of the
evaluation map.

The terminology in this section turns the informal statement ``a point on
the received affine line looks like a codeword on many coordinates'' into a
finite incidence problem.  A \emph{witness} records all data certifying that
statement: the line parameter \(\gamma\), the agreement support \(S\), and
the explaining polynomial \(h\).  A \emph{profile} records a specified
structural reason why witnesses may occur in abundance.  An \emph{atlas} is
an ordered collection of such profiles covering all witnesses, and a
\emph{ledger} attaches a proved distinct-slope budget to every entry.  These
words have no additional geometric meaning beyond the exact definitions
below.

\begin{definition}[Explanation and support-wise nontriviality]
\label{def:explanation}
For \(u\in\F^D\) and \(S\subseteq D\), an \emph{explanation of \(u\) on
\(S\)} is a polynomial \(f\in\F[X]_{<k}\) with
\(\operatorname{ev}_S(f)=u|_S\).  We write
\(\operatorname{Expl}_C(u;S)\) when such an \(f\) exists.  A pair
\(r=(r_0,r_1)\) is \emph{simultaneously explained on \(S\)} when
\(\operatorname{Expl}_C(r_0;S)\) and
\(\operatorname{Expl}_C(r_1;S)\) both hold, with no requirement that the
two explaining polynomials coincide.  It is \emph{support-wise
nontrivial on \(S\)} when it is not simultaneously explained there.
We denote this last condition by \(\operatorname{NT}_C(r;S)\).
\end{definition}

Thus support-wise nontriviality is a condition on the same support that
explains the affine combination.  It does not mean that both coordinates
are individually unexplained; failure of either explanation is enough.

\begin{definition}[MCA-bad slope]\label{def:mca-bad}
Let \(r=(r_0,r_1)\in(\F^D)^2\) and let \(a\) be an agreement threshold.
A finite slope \(\gamma\in\F\) is support-wise MCA-bad at agreement \(a\)
if there are \(S\subseteq D\), \(\abs S\ge a\), and
\(h\in\F[X]_{<k}\) such that
\(\operatorname{ev}_S(h)=(r_0+\gamma r_1)|_S\) and
\(\operatorname{NT}_C(r;S)\) holds.  Define
\[
        B_C^{\rm MCA}(a)
        =\max_{r\in(\F^D)^2}\abs{\{\gamma\in\F:
          \gamma\text{ is MCA-bad for }r\text{ at }a\}}.
\]
\end{definition}

The finite-slope restriction costs at most the point at infinity, which may
also be treated by exchanging the two coordinates.  The support-wise
condition is essential: explanation and nontriviality are tested on the
same set.

\begin{definition}[Exact-agreement witness incidence]
\label{def:exact-witness-incidence}
Assume \(k+1\le a\le n\).  For a received line \(r=(r_0,r_1)\), its
exact-\(a\) witness incidence is
\[
 \Wcal_a(r)=\left\{(\gamma,S,h):
 \begin{array}{l}
   \gamma\in\F,\quad S\in\binom Da,\quad h\in\F[X]_{<k},\\
   \operatorname{ev}_S(h)=(r_0+\gamma r_1)|_S,\quad
   \operatorname{NT}_C(r;S)
 \end{array}\right\}.
 \tag{2.1}\label{eq:exact-witness-incidence}
\]
Its slope projection is
\(\pi_\gamma(\gamma,S,h)=\gamma\), and its exact bad-slope set is
\(Z_a(r)=\pi_\gamma(\Wcal_a(r))\subseteq\F\).
\end{definition}

\begin{lemma}[Reduction to exact agreement]
\label{lem:exact-agreement-reduction}
For \(a\ge k+1\), the set \(Z_a(r)\) is precisely the set of slopes that
are MCA-bad at threshold \(a\) in \cref{def:mca-bad}.
\end{lemma}

\begin{proof}
An exact-\(a\) witness is a threshold-\(a\) witness.  Conversely, let
\((\gamma,S,h)\) witness badness with \(\abs S\ge a\).  At least one of
\(r_0|_S,r_1|_S\) is outside
\(\operatorname{ev}_S(\F[X]_{<k})\); call it \(u|_S\).  Interpolate \(u\)
on any fixed \(k\)-subset \(K\subseteq S\).  If the resulting polynomial
agreed with \(u\) at every point of \(S\setminus K\), it would explain
\(u\) on \(S\), a contradiction.  Hence some \(x\in S\setminus K\)
makes \(u\) unexplained on \(K\cup\{x\}\).  Extend this set to an
\(a\)-subset \(S'\subseteq S\).  Non-explanation persists on supersets,
while the same \(h\) explains \(r_0+\gamma r_1\) on \(S'\).
\end{proof}

\begin{proposition}[Exact support-atlas upper bound]
\label{prop:exact-support-upper}
For every Reed--Solomon row, every nonempty challenge set
\(\Gamma\subseteq\F\), and every \(k+1\le a\le n\),
\[
 B_{C,\Gamma}^{\rm MCA}(a)
 \le\min\left\{\abs\Gamma,\binom na\right\}.        \tag{2.2}
\]
This is a literal finite-row bound: it has no polynomial or
subexponential loss.
\end{proposition}

\begin{proof}
By \cref{lem:exact-agreement-reduction}, choose one noncommon exact-\(a\)
support for every bad slope.  One noncommon support carries at most one
slope, since two distinct explained slopes would separately explain the
two received words on that support.  The chosen-support map is therefore
injective.  Intersect with the challenge set.
\end{proof}

\begin{definition}[Universal witness incidence]
\label{def:universal-witness-incidence}
Let \(\mathfrak S_a(D)=\coprod_{S\in\binom Da}\Spec\F\) be the finite
reduced scheme of exact-\(a\) supports.  The \emph{universal witness
incidence} is the row-wise locally closed set
\[
 \mathfrak W_a\subseteq
 (\A^n)^2\times\A^1_\gamma\times
 \mathfrak S_a(D)\times\A^k_h
 \]
whose fiber over \(r\in(\F^D)^2\) is \(\Wcal_a(r)\).  On the component
indexed by \(S\), it is cut out by
\(\operatorname{ev}_S(h)=(r_0+\gamma r_1)|_S\) and by the open condition
\((r_0|_S,r_1|_S)\notin
 \operatorname{im}(\operatorname{ev}_S)^2\).
\end{definition}

Here \(\mathfrak S_a(D)\) is simply the disjoint union of one point for each
\(a\)-element support.  The scheme notation packages all supports into one
algebraic object; readers may equivalently read \(\mathfrak W_a\) as the
finite family of polynomial equation-and-inequation systems indexed by
those supports.  ``Locally closed'' means an intersection of a Zariski-open
and a Zariski-closed set, and ``constructible'' means a finite union of such
sets.

The evaluation equations are polynomial and the simultaneous-explanation
locus is a closed linear subspace.  Thus \(\mathfrak W_a\) is locally
closed on each support component, and hence locally constructible.  It is
generally not a closed subvariety: support-wise nontriviality is an open
condition.  Eliminating \(h\), or adjoining locators, quotient data, and
\emph{rank pivots}, produces constructible images by Chevalley's theorem.
A rank pivot is a specified matrix minor required to be nonzero, thereby
selecting a chart on which a chosen rank is visible.  We use Chevalley's
theorem only in its standard form that the algebraic projection of a
constructible set is constructible.

\begin{definition}[Profile cell and realized profile]
\label{def:profile-cell}
A \emph{profile type} \(\lambda\) consists of a parameter space
\(P_\lambda\) and a locally constructible incidence
\(\mathfrak C_\lambda\subseteq\mathfrak W_a\times P_\lambda\).
For a received line \(r\) and a parameter
\(\xi\in P_\lambda(\F)\), the corresponding \emph{realized profile cell}
is the actual finite witness set
\[
 \Ccal_{\lambda,r,\xi}
   =\{w\in\Wcal_a(r):(w,\xi)\in
       \mathfrak C_\lambda(\F)\}.
\]
It is a realized profile when this set is nonempty.  A boundary profile
also specifies a finite full support slice \(\Omega^0_{\lambda,r,\xi}\)
and a proved boundary map
\(\Phi_{\lambda,r,\xi}:\Omega^0_{\lambda,r,\xi}\to\Sigma_{\lambda,r,\xi}\).
Its image-normalized scale is
\[
 L_{\lambda,r,\xi}
   =\abs{\Phi_{\lambda,r,\xi}(\Omega^0_{\lambda,r,\xi})},
 \qquad
 \barN_{\lambda,r,\xi}
   =\frac{\abs{\Omega^0_{\lambda,r,\xi}}}
          {L_{\lambda,r,\xi}}.
 \tag{2.2}\label{eq:realized-profile-scale}
\]
The ambient codomain size is not substituted for \(L_{\lambda,r,\xi}\)
without the full-image certificate \textup{(FI)}.
\end{definition}

The parameters \(\xi\) are part of the profile, not hidden existential
choices.  Consequently an assertion that there are \(e^{o(n)}\) profiles
means that the number of nonempty pairs \((\lambda,\xi)\) occurring for a
received line is \(e^{o(n)}\), uniformly in that line.

\begin{definition}[Witness-exhaustive first-match atlas]
\label{def:first-match}
Fix \(r\) and order its realized profile cells as
\(\Ccal_1,\ldots,\Ccal_s\subseteq\Wcal_a(r)\).  Define their \emph{actual}
slope projections and first-match parts by
\[
 Z_i=\pi_\gamma(\Ccal_i),\qquad
 Z_i^\circ=Z_i\setminus\bigcup_{j<i}Z_j,\qquad
 \Ccal_i^\circ=\Ccal_i\cap\pi_\gamma^{-1}(Z_i^\circ).
 \tag{2.3}\label{eq:first-match-projections}
\]
The ordered family is a \emph{witness-exhaustive first-match atlas} if
\(\Wcal_a(r)=\bigcup_{i=1}^s\Ccal_i\).  A certified upper ledger for this
atlas consists of supersets \(E_i\supseteq Z_i^\circ\) and numbers
\(U_i\) satisfying \(\abs{E_i}\le U_i\).
\end{definition}

Witness exhaustivity is stronger than merely covering the set of slopes,
and is the property needed to assert that no unclassified witness remains.
The sets \(Z_i^\circ\), rather than arbitrary slope supersets, define the
first-match assignment.  In particular,
\(Z_a(r)=\coprod_i Z_i^\circ\).

\begin{definition}[Profile payment]
\label{def:profile-payment}
For a row sequence \(C_n=\RS_{\F_n}(D_n,k_n)\), with
\(\abs{D_n}=n\to\infty\), a \emph{scaled realized profile cell} is a
realized cell equipped with either its image-normalized scale \(\barN_i\)
or a declared direct scale \(M_i\).  Such a cell is \emph{paid} in a
given first-match atlas if, uniformly in the received line and its realized
parameter, its certified budget satisfies
\[
 \abs{Z_i^\circ}\le U_i
   \le \exp(o(n))\bigl(1+\barN_i\bigr).
 \tag{2.4}\label{eq:profile-payment}
\]
More generally a direct profile scale \(M_i\) is paid when
\(U_i\le e^{o(n)}M_i\).  An atlas is paid at an envelope \(E_n\) when
\(\sum_iU_i\le e^{o(n)}E_n\).  All \(o(n)\) terms are uniform over the
profiles and received lines under consideration.
\end{definition}

At image scale one always has \(\barN_i\ge1\), so the additive \(1\) is
harmless.  It keeps the same payment notation valid for separately
certified ambient or direct profile scales that can lie below one.
Payment is an enumerative assertion about an actual slope projection;
constructibility alone is not payment.

\begin{definition}[Primitive first-match residual]
\label{def:primitive-first-match-residual}
After the first \(j\) cells of an ordered atlas have been routed, define
the remaining witness incidence by
\[
 \Wcal_{>j}(r)=
 \Wcal_a(r)\cap
 \pi_\gamma^{-1}\!\left(\F\setminus\bigcup_{i\le j}Z_i\right).
 \tag{2.5}\label{eq:first-match-residual}
\]
A later realized boundary profile is a \emph{primitive first-match
residual} when its witnesses lie in \(\Wcal_{>j}(r)\), every algebraic
degeneracy named by the atlas has been routed among the first \(j\) cells,
and the remaining boundary map carries the row-specific primitivity
certificate used by the analytic and ray arguments.  Its residual support
set \(\Omc_\lambda\) is the subset of
\(\Omega^0_\lambda\) induced by \eqref{eq:first-match-residual}.
\end{definition}

The residual in \cref{def:primitive-first-match-residual} is a
first-match excision, not a topological assertion.  It is constructible in
the row-wise finite incidence after taking the displayed projections, but
need not be Zariski open.  Nor is ``primitive'' synonymous with having
trivial stabilizer: it records the specific quotient, field-descent, rank,
planted, and ray-saturation exclusions certified by the atlas.

\begin{lemma}[First-match bound]\label{lem:first-match-bound}
If every received line admits a witness-exhaustive first-match atlas and a
certified upper ledger with \(\sum_iU_i\le U\), then
\(B_C^{\rm MCA}(a)\le U\).
\end{lemma}

\begin{proof}
For a fixed received line, exact-agreement reduction and witness
exhaustivity give the disjoint union
\(Z_a(r)=\coprod_iZ_i^\circ\).  Since
\(Z_i^\circ\subseteq E_i\), its size is at most \(U_i\).  Summing and then
maximizing over received lines proves the claim.
\end{proof}

The first-match rule is the combinatorial form of excision in the witness
incidence.  Profile cells need not be disjoint.  Once a slope is charged to
the first profile in whose actual projection it occurs, all of its later
witnesses are removed from the residual by
\eqref{eq:first-match-residual}.

\section{Exact syndrome geometry and the deep regime}
\label{sec:syndrome-exact}

The \emph{syndrome} of a word is the result of applying a parity-check
matrix; it vanishes exactly on codewords.  Thus syndrome space forgets the
part of a received word already lying in the code and retains only its error
class.  This description separates possible error supports from distinct
line parameters.  Put \(R=n-k\), the redundancy (and syndrome-space
dimension), and choose the generalized Vandermonde parity-check matrix
\(H\) with columns
\[
 h_x=\lambda_x(1,x,\ldots,x^{R-1})^{\mathsf T},\qquad
 \lambda_x=\left(\prod_{y\in D\setminus\{x\}}(x-y)\right)^{-1}.
 \tag{3.1}\label{eq:parity-columns}
\]
Lagrange interpolation gives
\(\sum_{x\in D}\lambda_xx^j=0\) for \(0\le j\le n-2\), so this matrix
annihilates every degree-less-than-\(k\) evaluation.  Every set of at most
\(R\) columns is independent.  For \(E\subseteq D\), write
\(V_E=\Span\{h_x:x\in E\}\); it is exactly the set of syndromes of words
supported on \(E\).  We write \(s(u)=Hu^{\mathsf T}\) for the syndrome of
\(u\).

\begin{proposition}[Syndrome-line normal form]
\label{prop:syndrome-line-normal-form}
Let \(S=D\setminus E\).  A word \(u\) is explained on \(S\) if and only if
\(s(u)\in V_E\).  Consequently \(\gamma\) is MCA-bad on \(S\) for a pair
\((u_0,u_1)\) if and only if
\[
 s(u_0)+\gamma s(u_1)\in V_E,
 \qquad
 \{s(u_0),s(u_1)\}\not\subseteq V_E.
 \tag{3.2}\label{eq:syndrome-line}
\]
For fixed \(E\), there is at most one such finite slope.
\end{proposition}

\begin{proof}
If \(u=c+e\), where \(c\in C\) and \(e\) is supported on \(E\), then
\(s(u)=s(e)\in V_E\).  Conversely, if
\(s(u)=\sum_{x\in E}e_xh_x\), subtract the word supported on \(E\) with
values \(e_x\); the result has zero syndrome and belongs to \(C\).  Apply
this equivalence to the line point and to both coordinates.  Passing
\eqref{eq:syndrome-line} to \(\F^R/V_E\) shows that two distinct solutions
would put both projected syndromes at zero, contradicting badness.
\end{proof}

Since \(\lambda_x\ne0\), one has
\([h_x]=[1:x:\cdots:x^{R-1}]\); these are the points of the rational
normal curve in projective syndrome space, while the weights
\(\lambda_x\) merely choose vector representatives.  A \emph{secant span} is the linear span
\(V_E\) of a selected set of those points.  We call the meeting of a
syndrome line with \(V_E\) \emph{transverse} when the entire line is not
contained in \(V_E\), equivalently when its two generators are not both in
\(V_E\).  The next result is a \emph{compiler} in the sense that it
translates the bad-slope problem, exactly and in both directions, into this
line--secant incidence problem.  A \emph{ray} means a one-dimensional
direction in syndrome or coefficient space; different raw witnesses may
represent the same ray and hence the same slope.

\begin{theorem}[Exact syndrome--secant compiler]
\label{thm:syndrome-secant-exact}
Put \(t=n-a\), and for a syndrome line
\(\ell_{y_0,y_1}=\{y_0+\gamma y_1:\gamma\in\F\}\) define
\[
 \Theta_t(y_0,y_1)=
 \abs{\{\gamma\in\F:\exists E\subseteq D,\ \abs E\le t,
 y_0+\gamma y_1\in V_E,
 \ \{y_0,y_1\}\not\subseteq V_E\}}.
 \tag{3.3}\label{eq:transverse-secant-count}
\]
Then
\[
 B_C^{\rm MCA}(a)=\max_{y_0,y_1\in\F^{n-k}}
                    \Theta_t(y_0,y_1).
 \tag{3.4}\label{eq:exact-secant-numerator}
\]
For each fixed \(E\), the affine line has at most one transverse
intersection parameter with \(V_E\).  For a Reed--Solomon code the
spaces \(V_E\) are the spans of at most \(t\) points of the weighted
rational normal curve in \eqref{eq:parity-columns}.  Thus the
higher-dimensional ray problem is exactly a transverse line--secant
incidence problem, with repeated intersection parameters deduplicated.
\end{theorem}

\begin{proof}
Apply \cref{prop:syndrome-line-normal-form} with \(S=D\setminus E\), and
take the union over \(\abs E\le t\).  The syndrome map is surjective, so
maximizing over received pairs is the same as maximizing over
\((y_0,y_1)\).  In the quotient \(\F^{n-k}/V_E\), a transverse
intersection solves
\(\overline y_0+\gamma\overline y_1=0\) with the two projected vectors
not both zero, and therefore has at most one finite solution.  The last
assertion is the Vandermonde description of the parity columns.
\end{proof}

A \emph{parity-column circuit} is a minimally linearly dependent set of
columns of \(H\).  The next lemma says that many distinct transverse rays
cannot all be supported inside an independent union of columns.

\begin{lemma}[Independent-union ray bound]
\label{lem:independent-union-rays}
Let every set of at most \(R\) parity columns be independent, put
\(t<R\), and fix a syndrome line.  For each slope in a transverse set
\(Z\), choose an error set \(E_\gamma\), \(\abs{E_\gamma}\le t\),
witnessing \eqref{eq:transverse-secant-count}.  If
\(U=\bigcup_{\gamma\in Z}E_\gamma\) has size at most \(R\), then
\(\abs Z\le t+1\).  Consequently, more than \(t+1\) transverse slopes
force the witness union to contain a parity-column circuit.
\end{lemma}

\begin{proof}
If \(\abs Z\le1\), there is nothing to prove; assume that \(Z\) contains
two distinct slopes.  Two distinct line points lie in \(V_U\), hence
\(y_0,y_1\in V_U\).
Column independence gives unique coefficient lifts
\(e_0,e_1\in\F^U\), and every chosen lift is
\(e_\gamma=e_0+\gamma e_1\).  Let \(U'\) be the set of \(s\) coordinates
where the pair \((e_0,e_1)\) is nonzero.  If \(s\le t\) and
\(E_\gamma\) contained all of \(U'\), then
\(y_0,y_1\in V_{U'}\subseteq V_{E_\gamma}\), contradicting the
transversality of this particular witness.  Hence some
\(x\in U'\setminus E_\gamma\) exists, and uniqueness of the lift gives
\(e_0(x)+\gamma e_1(x)=0\).  Each nonzero affine coordinate function has
at most one root, so \(\abs Z\le s\le t\).  If \(s>t\), every chosen lift of
weight at most \(t\) has at least \(s-t\) zero coordinates.  Double
counting the pairs \((\gamma,x)\) with
\(e_0(x)+\gamma e_1(x)=0\) gives
\(\abs Z(s-t)\le s\), whence
\(\abs Z\le\floor{s/(s-t)}\le t+1\).  If \(\abs U\le R\), independence
holds; the contrapositive proves the circuit assertion.
\end{proof}

\begin{theorem}[Bounded residual-kernel ray compiler]
\label{thm:bounded-residual-kernel-ray}
Let \(H\) be a parity-check matrix over \(\F\), let \(U\) be a set of
columns, and put
\[
 \kappa(U)=\dim_\F\ker(H_U:\F^U\longrightarrow\F^{n-k}).
\]
Fix a syndrome line \(y_0+\gamma y_1\), an integer \(t\ge0\), and a set
\(Z\) of finite transverse slopes.  Suppose that for every \(\gamma\in Z\)
there is \(E_\gamma\subseteq U\), \(\abs{E_\gamma}\le t\), such that
\[
 y_0+\gamma y_1\in V_{E_\gamma},\qquad
 \{y_0,y_1\}\not\subseteq V_{E_\gamma}.
\]
Then
\[
                  \abs Z\le(t+1)\abs{\F}^{\kappa(U)}. \tag{RC$_{\rm ker}$}
\]
For a Reed--Solomon parity check of redundancy \(R=n-k\),
\(\kappa(U)=\max\{0,\abs U-R\}\).  Hence a certified residual chart with
\((\abs U-R)_+\log\abs\F=o(n)\) has a direct subexponential ray payment
with the displayed exact finite loss.
\end{theorem}

\begin{proof}
The assertion is immediate when \(\abs Z\le1\).  Otherwise two line
points belong to \(V_U\), so \(y_0,y_1\in V_U\).  Choose lifts
\(b_0,b_1\in\F^U\) with \(H_Ub_i=y_i\).  For every \(\gamma\), choose a
lift \(c_\gamma\) supported in \(E_\gamma\) and put
\(z_\gamma=c_\gamma-b_0-\gamma b_1\in\ker H_U\).

Fix \(z\in\ker H_U\), and consider the slopes for which
\(z_\gamma=z\).  Put
\[
 U_z=\{x\in U:(b_0(x)+z(x),b_1(x))\ne(0,0)\},
 \qquad s=\abs{U_z}.
\]
Every active coordinate is a nonzero affine function of \(\gamma\), and
therefore vanishes for at most one slope.  Hence the total number of active
zero incidences in this kernel class is at most \(s\).  If \(s>t\), the
weight bound on \(c_\gamma\) forces at least \(s-t\) active zeros, so
\[
 \abs{Z_z}(s-t)\le s,\qquad
 \abs{Z_z}\le\floor{s/(s-t)}\le t+1.
\]
If \(s\le t\), transversality forces
at least one active zero: without one,
\(U_z=\supp c_\gamma\subseteq E_\gamma\), while both \(b_0+z\) and
\(b_1\) are supported in \(U_z\), putting \(y_0,y_1\) in
\(V_{E_\gamma}\).  Thus this kernel class contains at most \(s\le t\)
slopes.  Summing over the \(\abs{\F}^{\kappa(U)}\) kernel elements proves
the bound.  The MDS column property gives the Reed--Solomon formula for
\(\kappa(U)\).
\end{proof}

\begin{theorem}[Single MDS-circuit ray compiler]
\label{thm:single-mds-circuit-ray}
Let \(H\) be a Reed--Solomon parity check of redundancy \(R\), let
\(U\subseteq D\) have size \(R+1\), put \(t<R\), and fix a syndrome line
\(y_0+\gamma y_1\).  If every slope in a transverse set \(Z\) has a witnessing error set
\(E_\gamma\subseteq U\) of size at most \(t\), then
\[
                         \abs Z\le\binom{R+1}{2}.   \tag{RC$_{\rm circ}$}
\]
Thus one fixed MDS-circuit profile has a field-independent polynomial ray
payment.
\end{theorem}

\begin{proof}
There is nothing to prove when \(\abs Z\le1\).  Otherwise two distinct
line points lie in \(V_U\); subtracting them and solving for the constant
term gives \(y_0,y_1\in V_U\).  Choose lifts
\(b_0,b_1\in\F^U\).  The circuit kernel is generated by a vector
\(z\) whose coordinates are all nonzero, since a zero coordinate would
give a proper dependent subset of the minimal circuit \(U\).  For each slope choose
\(c_\gamma\in\F\) so that
\(e_\gamma=b_0+\gamma b_1+c_\gamma z\) is supported in
\(E_\gamma\).  For \(x\in U\), set
\[
             f_x(\gamma)=-(b_0(x)+\gamma b_1(x))/z(x).
\]
A zero coordinate is an equality \(c_\gamma=f_x(\gamma)\), and at least
\(R+1-t\ge2\) such equalities hold.

Every transverse slope makes two nonidentical functions \(f_x,f_{x'}\)
equal.  Indeed, otherwise all vanishing coordinates belong to one class
\(C\) of identical affine functions, say
\(f_x=\alpha+\beta\gamma\) on \(C\), and no function outside \(C\) has
that value at \(\gamma\).  Then
\(\supp e_\gamma=U\setminus C\subseteq E_\gamma\), while
\(b_0+\alpha z\) and \(b_1+\beta z\) are also supported in
\(U\setminus C\).  This puts both \(y_0,y_1\) in \(V_{E_\gamma}\), a
contradiction.  Two nonidentical affine functions agree at at most one
finite slope.  Charging each slope to one unordered pair proves the
displayed bound.
\end{proof}

For a fixed received line, the exact bad-slope set is therefore the set of
unique cancelling line parameters contributed by the transverse quotient
equations.  For each \(E\), \(\abs E\le n-a\), project the two
syndromes to \(\F^R/V_E\).  If the second projection is nonzero and the
first is a scalar multiple of it, record the unique cancelling scalar;
otherwise record nothing.  Deduplicating these scalars gives the numerator.
In particular, the number of supports, the number of pairs of supports, and
the number of distinct slopes are three different quantities.

We call \(3(n-a)\le d-1\) the \emph{deep regime}: the allowed error
radius \(n-a\) is so small relative to the minimum codeword weight \(d\)
that the union of three permitted error supports still cannot contain a
nonzero codeword support.

\begin{theorem}[Deep-regime upper bound]\label{thm:deep-regime-upper}
Let \(C\subseteq\F^n\) be linear with minimum nonzero weight \(d\), let
\(0\le a\le n\), put \(r=n-a\), and assume \(3r\le d-1\).  Then
\[
 B_C^{\rm MCA}(a)\le r+1.
 \tag{3.5}\label{eq:deep-upper}
\]
\end{theorem}

\begin{proof}
Fix \((u_0,u_1)\), and let \(Z\) be the slopes for which \(u_0+\gamma u_1\)
is within distance \(r\) of \(C\).  If \(\abs Z\le r+1\), there is nothing
to prove.  Otherwise choose distinct \(\gamma_1,\gamma_2\in Z\), codewords
\(p_1,p_2\), and error sets \(E_1,E_2\), each of size at most \(r\), such
that \(u_0+\gamma_i u_1=p_i\) off \(E_i\).  Solving the two affine
equations gives \(c_0,c_1\in C\) for which the error pair
\(e_i=u_i-c_i\) is supported on \(T=E_1\cup E_2\), \(\abs T\le2r\).

For any \(\gamma\in Z\), let \(p\) be a closest explaining codeword.  The
codeword \(p-(c_0+\gamma c_1)\) is supported on at most \(3r\le d-1\)
positions and is zero.  Thus \(\wt(e_0+\gamma e_1)\le r\).  Let \(T'\) be
the support of the error pair.  If \(\abs{T'}\ge r+1\), every
\(\gamma\in Z\) makes at least \(\abs{T'}-r\) nonzero affine coordinate
functions \(e_{0,x}+\gamma e_{1,x}\) vanish.  Each coordinate has at most
one root, so \(\abs Z(\abs{T'}-r)\le\abs{T'}\) and \(\abs Z\le r+1\).

It remains that \(\abs{T'}\le r\).  If \(\gamma\) is MCA-bad on a set
\(S\), \(\abs S\ge n-r\), its explaining codeword equals
\(c_0+\gamma c_1\), since their difference would otherwise have weight at
most \(2r\le d-1\).  The set \(S\) must meet \(T'\), and at a point of
\(S\cap T'\) the nonzero pair satisfies
\(e_{0,x}+\gamma e_{1,x}=0\).  Such a point determines \(\gamma\), giving
at most \(\abs{T'}\le r\) bad slopes.
\end{proof}

Recall that for a nonempty challenge set \(\Gamma\subseteq\F\),
\(B_{C,\Gamma}^{\rm MCA}(a)\) denotes the same maximum with slopes restricted to
\(\Gamma\).

The following ``tangent floor'' is an elementary family of neighboring
coordinate-span intersections.  The name does not assert a hidden
differentiability or multiplicity condition.

\begin{proposition}[Universal tangent floor]\label{prop:universal-tangent-floor}
For every Reed--Solomon code, every \(a\ge k+1\), and every nonempty
\(\Gamma\subseteq\F\),
\[
 B_{C,\Gamma}^{\rm MCA}(a)
 \ge\min\{\abs\Gamma,n-a+1\}.
 \tag{3.6}\label{eq:tangent-floor}
\]
\end{proposition}

\begin{proof}
Put \(j=n-a\), \(t=\min\{\abs\Gamma,j+1\}\), and choose distinct
\(\gamma_1,\ldots,\gamma_t\in\Gamma\) and independent parity columns
\(h_{x_1},\ldots,h_{x_t}\).  Define
\(s_1=\sum_i h_{x_i}\) and \(s_0=-\sum_i\gamma_i h_{x_i}\), and lift them
to received words.  For \(E_i=\{x_\ell:\ell\ne i\}\),
\(s_0+\gamma_i s_1\in V_{E_i}\), whereas \(s_1\notin V_{E_i}\).
\Cref{prop:syndrome-line-normal-form} makes every \(\gamma_i\) bad on
\(D\setminus E_i\), which has size at least \(a\).
\end{proof}

\begin{corollary}[Exact deep numerator]\label{cor:exact-deep-numerator}
If \(C=\RS_\F(D,k)\), \(1\le k<n\), \(k+1\le a\le n\),
\(\varnothing\ne\Gamma\subseteq\F\), and \(3(n-a)\le n-k\), then
\[
 B_{C,\Gamma}^{\rm MCA}(a)
 =\min\{\abs\Gamma,n-a+1\}.
\]
\end{corollary}

\begin{proof}
The upper bound is \cref{thm:deep-regime-upper}, intersected with
\(\Gamma\).  The hypotheses imply \(a\ge k+1\), so
\cref{prop:universal-tangent-floor} gives equality.
\end{proof}

\begin{theorem}[Exact first adjacent row over a separating field]
\label{thm:exact-first-adjacent-row}
Let \(C=\RS_\F(D,k)\), put \(R=n-k\) and
\(M=\binom n{R-1}=\binom n{k+1}\), and define
\[
                    Q_{\rm sep}(M)=\max\left\{M,\binom M2\right\}.
\]
If \(\abs\F>Q_{\rm sep}(M)\), then for the full-field challenge
\[
                         B_C^{\rm MCA}(k+1)=M.      \tag{AD1}
\]
If \(R\ge2\) and \(b=\floor{\eps\abs\F}\), the exact target threshold
therefore satisfies
\[
 \begin{array}{ll}
 b\ge M &\Longrightarrow\quad a^*=k+1,\\[2mm]
 \displaystyle\binom n{R-2}\le b<M
   &\Longrightarrow\quad a^*=k+2.
 \end{array}                                       \tag{AD2}
\]
Thus the first finite adjacent comparison has literal binomial constants;
no \(e^{o(n)}\) or unspecified polynomial factor occurs.
\end{theorem}

\begin{proof}
The upper bound in \textup{(AD1)} is
\cref{prop:exact-support-upper}.  For the reverse inequality use the
Reed--Solomon parity check of redundancy \(R\).  For every
\((R-1)\)-set \(E\subseteq D\), its column span \(V_E\) is a hyperplane
in \(\F^R\), and distinct sets give distinct hyperplanes: otherwise the
union of two distinct \((R-1)\)-sets would contain at least \(R\) columns
in one \((R-1)\)-space, contradicting the MDS column property.

Every proper linear hyperplane in \(\F^R\) has
\(\abs{\F}^{R-1}\) points.  Since \(M<\abs\F\), the union of the \(M\)
spaces \(V_E\) has fewer than \(\abs{\F}^{R}\) points; choose \(y_1\)
outside it.  For each \(E\),
choose a nonzero linear functional \(\ell_E\) with kernel \(V_E\).
The line \(y_0+\gamma y_1\) meets \(V_E\) at
\[
                 \gamma_E=-\ell_E(y_0)/\ell_E(y_1).
\]
For \(E\ne E'\), equality \(\gamma_E=\gamma_{E'}\) cuts out a proper
hyperplane in the variable \(y_0\), because the normalized functionals
\(\ell_E/\ell_E(y_1)\) and
\(\ell_{E'}/\ell_{E'}(y_1)\) are distinct.  There are at most
\(\binom M2<\abs\F\) collision hyperplanes; their union also has fewer
than \(\abs{\F}^{R}\) points, so choose \(y_0\) outside it.  Since
\(y_1\notin V_E\), the line is not contained in \(V_E\), and every
intersection is transverse.  The \(M\) parameters \(\gamma_E\) are then
finite and pairwise distinct.  Syndrome surjectivity and
\cref{thm:syndrome-secant-exact} realize them as \(M\) bad slopes at
agreement \(k+1\), proving \textup{(AD1)}.

If \(b\ge M\), the support upper bound makes the first allowed agreement
safe.  If \(\binom n{R-2}\le b<M\), equality \textup{(AD1)} makes
\(k+1\) unsafe, whereas
\cref{prop:exact-support-upper} makes \(k+2\) safe.  Monotonicity proves
\textup{(AD2)}.
\end{proof}

\begin{theorem}[Subfield confinement]\label{thm:subfield-confinement-full}
Let \(\B\subseteq\F\), \(D\subseteq\B\), and let \(u_0,u_1\in\B^D\).  If
\(\gamma\in\F\setminus\B\) and \(u_0+\gamma u_1\) is explained on \(S\),
then the pair is explained on the same \(S\).  Hence every MCA-bad slope of
a \(\B\)-valued received line lies in \(\B\).
\end{theorem}

\begin{proof}
Choose a \(\B\)-basis \(1,\beta_1,\ldots\) of \(\F\).  Decompose the
coefficients of the explaining polynomial and of \(\gamma\) in this basis.
Because \(D\subseteq\B\), comparison of basis components on \(S\) gives a
codeword explaining \(u_1\) from any nonzero nonbase component of
\(\gamma\), and then a codeword explaining \(u_0\) from the base component.
\end{proof}

\section{Exact locator prefixes and collision-aware pole lines}
\label{sec:exact-prefix-pole}

Assume \(D\subseteq\B\subseteq\F\).  For an \(m\)-set \(S\subseteq D\),
write
\[
 Q_S(X)=\prod_{x\in S}(X-x)
       =X^m+c_1(S)X^{m-1}+\cdots+c_m(S),
 \qquad
 \Phi_w(S)=(c_1(S),\ldots,c_w(S)).
\]

The monic polynomial \(Q_S\), whose simple roots are precisely the points
of \(S\), is the \emph{support locator}.  Its \emph{depth-\(w\) locator
prefix} is the vector of the first \(w\) coefficients below the leading
coefficient.  For \(z\in\B^w\), the inverse image
\(\Phi_w^{-1}(z)\) is the \emph{prefix fiber over \(z\)}: the family of
all \(m\)-element supports having that same prefix.  Thus ``fiber'' in this
paper always refers to the inverse image of a displayed map.

\begin{proposition}[Exact prefix-list bijection]\label{prop:exact-prefix-list}
Let \(1\le K\le m\le n\), put \(w=m-K\), and set
\(U_z(X)=X^m+\sum_{i=1}^wz_iX^{m-i}\) for \(z\in\B^w\).  The codewords of
\(\RS_\F(D,K)\) agreeing with \(U_z|_D\) on at least \(m\) points are
exactly
\[
        (U_z-Q_S)|_D,\qquad S\in\Phi_w^{-1}(z).
 \tag{4.1}\label{eq:prefix-bijection}
\]
No such codeword agrees on more than \(m\) points.  In particular, some
received word has list size at least
\(\lceil\binom nm\abs\B^{-w}\rceil\).
\end{proposition}

\begin{proof}
For \(S\in\Phi_w^{-1}(z)\), the leading term and next \(w\) coefficients
cancel, so \(\deg(U_z-Q_S)<m-w=K\).  Conversely, if \(P\), \(\deg P<K\),
agrees with \(U_z\) at \(m\) points, then \(U_z-P\) is monic of degree
\(m\) with those \(m\) roots and equals their locator.  Coefficient
comparison gives the prefix.  A nonzero degree-\(m\) polynomial has no
additional root.  Pigeonhole the \(\binom nm\) supports over
\(\abs\B^w\) prefixes.
\end{proof}

The \emph{simple-pole conversion} chooses an auxiliary point
\(\alpha\notin D\) and divides coordinatewise by \(X-\alpha\).  This
turns one list of nearby polynomials into points on a received affine line
for the dimension-\(k\) code.  A \emph{collision} is an equality
\(P_i(\alpha)=P_j(\alpha)\) for two distinct listed polynomials; the next
formula pays such equalities rather than assuming all slopes are distinct.

\begin{theorem}[Collision-aware simple-pole conversion]
\label{thm:collision-aware-pole}
Let \(C=\RS_\F(D,k)\), \(C^+=\RS_\F(D,k+1)\), \(1\le k<n\), and
\(q=\abs\F>n\).  If a word has \(L\) distinct codewords of \(C^+\)
agreeing on at least \(m\ge k+1\) points, then one received line for \(C\)
has at least
\[
 M(L)=\left\lceil
 \frac{L(q-n)}{q-n+k(L-1)}
 \right\rceil
 \tag{4.2}\label{eq:collision-aware-pole}
\]
distinct MCA-bad slopes at agreement \(m\).
\end{theorem}

\begin{proof}
Let the received word be \(U\in\F^D\), and let the listed polynomials be
\(P_1,\ldots,P_L\), each of degree at most \(k\).  For
\(\alpha\in\F\setminus D\), define the word coordinatewise by
\(f_\alpha(x)=U(x)/(x-\alpha)\) and set
\(g_\alpha(x)=-1/(x-\alpha)\).  On an agreement set for \(P_i\), the line
point of slope \(P_i(\alpha)\) equals
\((P_i(x)-P_i(\alpha))/(x-\alpha)\), a degree-less-than-\(k\) evaluation.
The pair is not simultaneously explained on \(k+1\) points: an explanation
of \(g_\alpha\) by \(G\), \(\deg G<k\), would make
\((X-\alpha)G(X)+1\) vanish at \(k+1\) points although its degree is at
most \(k\) and its value at \(\alpha\) is \(1\).

For \(i\ne j\), the equality \(P_i(\alpha)=P_j(\alpha)\) holds at at most
\(k\) poles.  Some pole has at most \(k\binom L2/(q-n)\) colliding pairs.
If its value multiplicities are \(m_1,\ldots,m_M\), then
\(\sum m_i^2\le L+kL(L-1)/(q-n)\).  Cauchy--Schwarz gives
\(L^2\le M\sum m_i^2\), which yields \eqref{eq:collision-aware-pole}.
\end{proof}

\begin{corollary}[Identity-prefix MCA floor]\label{cor:identity-prefix-floor}
For \(m\ge k+1\), put \(w=m-k-1\) and
\(L_m=\lceil\binom nm\abs\B^{-w}\rceil\).  If \(q>n\), then
\[
 B_C^{\rm MCA}(m)\ge M(L_m).
 \tag{4.3}\label{eq:identity-prefix-floor}
\]
\end{corollary}

For equal-size sets, the \emph{Johnson distance} is
\(d_J(S,T)=\abs{S\setminus T}=\abs{T\setminus S}\).  It counts the
number of point replacements required to turn one support into the other.

\begin{proposition}[Prefix rigidity and its packing limit]
\label{prop:prefix-rigidity-full}
Two distinct \(m\)-sets in one depth-\(w\) prefix fiber have Johnson
distance at least \(w+1\).  Consequently, with
\(t=\floor{w/2}\), every fiber has size at most
\[
 \frac{\binom nm}{\sum_{i=0}^t\binom mi\binom{n-m}{i}}.
 \tag{4.4}\label{eq:packing-fiber-cap}
\]
At \(w\asymp n/\log\abs\B\), this packing loss is only
\(e^{o(n)}\) and does not supply the random factor \(\abs\B^{-w}\).
\end{proposition}

Prefix rigidity says that sharing many leading locator coefficients forces
distinct supports to be far apart in this metric.  Here the
\emph{polynomial-field regime} means \(\abs\B=n^{O(1)}\), equivalently
\(\log\abs\B=O(\log n)\).

\begin{proof}
If \(S,T\) have Johnson distance \(j\le w\), factor the locator of their
intersection, of degree \(m-j\), from \(Q_S-Q_T\).  The residual factor is
nonzero, so the difference has degree at least \(m-j\ge m-w\), whereas the
common prefix makes its degree at most \(m-w-1\), a contradiction.  Radius
\(t\) Johnson balls about the fiber members are disjoint and each has the
denominator in \eqref{eq:packing-fiber-cap}.  For
\(w=n/\log\abs\B\), the logarithm of that ball volume is
\(O(w\log(n/w))=o(n)\) in the polynomial-field regime.
\end{proof}

\begin{theorem}[Prefix fibers embed in received lines]
\label{thm:prefix-to-line-hardness}
Let \(\mathcal G\subseteq\Phi_w^{-1}(z)\) be a fiber of \(N\) many
\(m\)-sets, with \(0\le w\le m-2\), and put \(k=m-w-1\).  Every finite
extension \(\F/\B\) satisfying
\[
        \abs\F-n>k\binom N2
        \tag{4.5}\label{eq:extension-separation}
\]
admits one received line for \(\RS_\F(D,k)\) with at least \(N\) distinct
MCA-bad slopes furnished by \(\mathcal G\).
\end{theorem}

\begin{proof}
For distinct \(S,T\in\mathcal G\), the polynomial \(Q_S-Q_T\) is nonzero
of degree at most \(k\).  Condition \eqref{eq:extension-separation} leaves a
pole \(\alpha\in\F\setminus D\) avoiding the roots of all these
differences.  Thus the values \(Q_S(\alpha)\) are pairwise distinct.  By
\cref{prop:exact-prefix-list}, the polynomials \(U_z-Q_S\) form one list for
the dimension-\(k+1\) code; the pole line makes their values, and hence the
resulting MCA slopes, pairwise distinct.
\end{proof}

\begin{theorem}[Exact list--line bijection at a separating pole]
\label{thm:exact-list-line-bijection}
Let \(D\subseteq\B\subseteq\F\), let \(1\le k<n\), let \(m\ge k+1\),
and let \(U\in\B^D\).  Write
\[
 \Lcal_m(U)=\{P\in\B[X]:\deg P\le k,
       \ \abs{\{x\in D:U(x)=P(x)\}}\ge m\}.
\]
Suppose \(\alpha\in\F\setminus D\) and the evaluation map
\(P\mapsto P(\alpha)\) is injective on \(\Lcal_m(U)\).  Define
\[
 f_\alpha(x)=\frac{U(x)}{x-\alpha},\qquad
 g_\alpha(x)=-\frac1{x-\alpha}\qquad(x\in D).
\]
Then \(P\mapsto P(\alpha)\) is a bijection from \(\Lcal_m(U)\) onto
the set of finite MCA-bad slopes of the received line
\((f_\alpha,g_\alpha)\) for \(\RS_\F(D,k)\) at agreement \(m\).  The
corresponding explaining polynomial is
\[
       h_P(X)=\frac{P(X)-P(\alpha)}{X-\alpha},
\]
and its agreement set with
\(f_\alpha+P(\alpha)g_\alpha\) is exactly the agreement set of \(P\)
with \(U\).

If \(L=\abs{\Lcal_m(U)}\), such a pole exists in every finite extension
satisfying
\[
       \abs\F>n+k\binom L2.                         \tag{4.6}
\]
\end{theorem}

\begin{proof}
For \(P\in\Lcal_m(U)\), the polynomial \(h_P\) has degree less than
\(k\), and on every point where \(U=P\) one has
\[
 f_\alpha+P(\alpha)g_\alpha
   =\frac{P-P(\alpha)}{X-\alpha}=h_P.
\]
The word \(g_\alpha\) cannot be explained by a polynomial \(G\) of degree
less than \(k\) on \(k+1\) points, since then
\((X-\alpha)G+1\) would have degree at most \(k\), at least \(k+1\)
roots, and value one at \(\alpha\).  Thus every displayed slope is
MCA-bad.

Conversely, suppose \(f_\alpha+\gamma g_\alpha\) is explained by
\(h\in\F[X]_{<k}\) on at least \(m\) points.  Then
\(P(X)=(X-\alpha)h(X)+\gamma\) has degree at most \(k\) and agrees with
\(U\) there.  Interpolation on any \(k+1\) of those points uses a
Vandermonde system over \(\B\) with right-hand side in \(\B\), so
\(P\in\B[X]\).  Hence \(P\in\Lcal_m(U)\) and
\(\gamma=P(\alpha)\).  Injectivity proves the bijection, and the same
calculation proves equality of the agreement sets.

For distinct \(P,Q\in\Lcal_m(U)\), the nonzero polynomial \(P-Q\) has
degree at most \(k\), and therefore at most \(k\) roots.  The union of
\(D\) and the root sets of all \(\binom L2\) differences has size at most
\(n+k\binom L2\); a larger extension contains an admissible pole.
\end{proof}

\begin{corollary}[Exact prefix-fiber ray realization]
\label{cor:exact-prefix-ray-realization}
Let \(0\le w\le m-2\), put \(k=m-w-1\), and let
\(\Fcal_z=\Phi_w^{-1}(z)\) be a complete depth-\(w\) fiber of
\(m\)-subsets of \(D\).  If \(\alpha\) separates the values
\(Q_S(\alpha)\), \(S\in\Fcal_z\), then the pole line associated with
\(U_z=X^m+\sum_{i=1}^wz_iX^{m-i}\) has exactly
\(\abs{\Fcal_z}\) MCA-bad slopes at agreement \(m\).  The maps
\[
 S\longmapsto U_z(\alpha)-Q_S(\alpha)
 \longmapsto
 \left(U_z(\alpha)-Q_S(\alpha),S,
 \frac{U_z-Q_S-(U_z(\alpha)-Q_S(\alpha))}{X-\alpha}\right)
\]
are bijections between the prefix fiber, the bad-slope set, and the
exact-\(m\) witness incidence.  In particular every retained-support
occupancy in \cref{lem:exact-occupancy-compiler} equals one.
\end{corollary}

\begin{proof}
The exact prefix-list bijection identifies \(\Fcal_z\) with the complete
dimension-\((k+1)\) list for \(U_z\).  Apply
\cref{thm:exact-list-line-bijection}.  Since
\(U_z-(U_z-Q_S)=Q_S\) has precisely the roots in \(S\), every agreement
is exactly \(m\) and supplies exactly that support.
\end{proof}

\begin{theorem}[Aperiodic ray obstruction below the profile frontier]
\label{thm:aperiodic-ray-obstruction}
There is a sequence of smooth multiplicative rows
\[
 D_n=\B_n^\times\subseteq\F_n^\times,
 \qquad \log\abs{\B_n}=O(\log n)=o(n),
\]
in characteristic two, with rates tending to \(1/2\), and agreements
\(m_n=k_n+O(n/\log n)\), for which one received line has
\(2^{\Omega(n)}\) distinct finite MCA-bad slopes.  Every threshold-
\(m_n\) witness support of that line has trivial multiplicative
stabilizer, and after enlarging \(\F_n\) with only
\(\log\abs{\F_n}=O(n)\), the bad slopes and exact-agreement supports are
in bijection.  Thus smoothness, quotient-aperiodicity, Q/SP support data,
and exact agreement do not force retained-support occupancy greater than
one.  This is a sharpness result for occupancy-based compilers, not a
counterexample to a direct distinct-slope theorem at the correctly
normalized profile scale.
\end{theorem}

\begin{proof}
Let \(q_r=2^r\), put \(\B_r=\F_{q_r}\),
\(D_r=\B_r^\times\), and \(n_r=q_r-1\).  Set
\[
 m_r=2^{r-1}-1=\frac{n_r-1}{2}.
\]
Fix \(0<\eta<\log2\), take
\(w_r=\floor{\eta n_r/\log q_r}\), and put
\(k_r=m_r-w_r-1\).  Then \(k_r/n_r\to1/2\) and
\(m_r-k_r=w_r+1=O(n_r/\log n_r)\).

Pigeonholing the \(m_r\)-set locators over their \(q_r^{w_r}\) prefixes
gives a fiber \(\Fcal_r\) satisfying
\[
 \log\abs{\Fcal_r}
 \ge \log\binom{n_r}{m_r}-w_r\log q_r
 \ge(\log2-\eta)n_r-o(n_r).
\]
Moreover
\[
 \gcd(m_r,n_r)=
 \gcd(2^{r-1}-1,2^r-1)=2^{\gcd(r-1,r)}-1=1.
\]
If an \(m_r\)-support were invariant under a nontrivial subgroup of
\(D_r\) of order \(c>1\), it would be a union of \(c\)-element orbits,
forcing \(c\mid m_r\) and \(c\mid n_r\), a contradiction.  Thus every
member of \(\Fcal_r\) is aperiodic.

Choose a finite extension \(\F_r/\B_r\) with
\[
 \abs{\F_r}>n_r+k_r\binom{\abs{\Fcal_r}}2.
\]
Since \(\abs{\Fcal_r}\le2^{n_r}\), the least extension degree meeting
this inequality has \(\log\abs{\F_r}=O(n_r)\).  By
\cref{cor:exact-prefix-ray-realization}, a separating pole gives exactly
\(\abs{\Fcal_r}=\exp(\Omega(n_r))\) distinct bad slopes, one for each
aperiodic support, and every occupancy equals one.
\end{proof}

\begin{remark}[Scope of the aperiodic obstruction]
The construction lies below the entropy--profile crossing: its natural
prefix scale is itself exponential.  It therefore does not contradict a
profile-envelope upper bound proved at the correct scale.  It does rule
out a field-independent polynomial ray bound, division of SP by the full
support slice without a retained multiplicity theorem, and the assertion
that quotient-aperiodicity alone forces ray collisions.
\end{remark}

\begin{definition}[Profiled asymptotic row datum]
\label{def:asymptotic-row}
A \emph{profiled asymptotic row datum} consists of a row sequence
\(C_n=\RS_{\F_n}(D_n,k_n)\), \(\abs{D_n}=n\), with coefficient fields
\(D_n\subseteq\B_n\subseteq\F_n\), together with agreements \(a_n\)
and, for every received line, a witness-exhaustive first-match atlas of
realized profiles with their fields, full slices, boundary maps, and actual
slope projections recorded.  The identity term is
\(\barN_1=\binom n{a_n}\abs{\B_n}^{-(a_n-k_n-1)}\); the correct global
scale is the profile envelope \(\mathfrak E_n(a_n)\) of
\eqref{eq:profile-envelope}.
\end{definition}

The identity term is the random size of one unprofiled prefix fiber.  It is
an exact lower scale by \cref{prop:exact-prefix-list}, but it is an upper
scale only after every quotient/remainder term has been compared with it and
the primitive max-fiber and ray conditions have been proved.

\section{Cells as moduli strata}\label{sec:cells}

The word \emph{cell} does not mean a topological cell decomposition.  The
ambient object is the locally constructible universal witness incidence
\(\mathfrak W_a\) of \cref{def:universal-witness-incidence}; a profile
cell is the parameterized incidence of \cref{def:profile-cell}.  The
qualification ``algebraic'' below records how such a profile is presented,
not whether its slope projection has been paid.

A \emph{moduli stratum} here simply means a parameterized family of
witnesses having the same declared structural features.  ``Description
complexity \(e^{o(n)}\)'' means that the number and degrees of equations,
inequations, auxiliary parameters, and Boolean pieces together contribute
only a subexponential factor to the final count.  It is a bookkeeping
condition, not a claim that every family has bounded dimension.

The \emph{structure maps} of a profile are its displayed projections,
boundary maps, and auxiliary-coordinate maps.  A \emph{planted divisor}
is a fixed factor common to the support locators in the profile; a
\emph{tangent datum} records a rank-defective slope projection; and a
\emph{locator pencil} is a projective polynomial family
\(Q_0+\lambda Q_1\).  The \emph{field of definition} is the smallest
recorded subfield over which these parameter spaces, equations, and maps
are defined.

\begin{definition}[Algebraic profile cell]\label{def:cell}
An \emph{algebraic profile cell} is a profile type
\(\mathfrak C_\lambda\subseteq\mathfrak W_a\times P_\lambda\) for which
\(P_\lambda\), the incidence, and all structure maps are locally
constructible of uniform subexponential description complexity.  A
description may use polynomial equalities, inequations, rank conditions,
finite Boolean operations, and projections.  Quotient maps, planted
divisors, tangent points, fields of definition, and locator-pencil
parameters are recorded in \(P_\lambda\); they are not suppressed when
counting realized profiles.
\end{definition}

\begin{definition}[Paid cell]\label{def:paid-cell}
Fix a witness-exhaustive ordered atlas.  A scaled realized cell \(\Ccal_i\)
is \emph{paid at its profile scale} when its actual first-match projection
\(Z_i^\circ\) has a certified bound
\[
       \abs{Z_i^\circ}\le U_i
          \le e^{o(n)}(1+\barN_i),
 \tag{5.1}\label{eq:paid-cell-scale}
\]
uniformly in the received line and the realized profile parameter.  A
directly counted cell is paid at a stated scale \(M_i\) when
\(U_i\le e^{o(n)}M_i\).  An ordered family is paid at an envelope when
the sum of its certified \(U_i\)'s is at most that envelope up to
\(e^{o(n)}\).  This is the specialization of
\cref{def:profile-payment} to the moduli cells used below.
\end{definition}

Payment depends on the chosen order because it bounds \(Z_i^\circ\), not a
formal cell name.  It is an enumerative certificate, not a topological
property.  A cell may be Zariski dense in an ambient chart but paid because
its actual slope projection is small after first-match saturation.
Conversely, a proper subvariety is not automatically paid: its degree,
components, parameter census, and projection fibers must fit the profile
budget.  The term \(1+\barN_i\) in \eqref{eq:paid-cell-scale} also prevents
a nonempty cell from being assigned a budget below one.

In the next lemma an \emph{extension locus} records data that factor only
after enlarging the coefficient field; a \emph{differential-locator
locus} is a rank-loss locus of the locator-prefix derivative; a
\emph{saturation locus} identifies raw witnesses with the same projective
explanation direction; and a \emph{fixed-dimensional pencil locus} is a
bounded-dimensional projective family of locators.  The catalogue below
then explains how these loci are used.

\begin{lemma}[Constructibility of the algebraic cells]\label{lem:constructible-cells}
The quotient, Chebyshev, planted, tangent, extension, differential-locator,
saturation, and fixed-dimensional pencil loci in
\cref{sec:cell-catalogue} are locally constructible profile cells over
\(\mathfrak W_a\).
The Sidon moment component is analytic and is not asserted to be a
bounded-complexity constructible cell.
\end{lemma}

\begin{proof}
Each listed algebraic condition is expressed by polynomial equations, rank
conditions, and nonvanishing conditions.  Common roots and planted blocks
use divisibility or resultant conditions, where
\(\operatorname{Res}(f,g)=0\) detects a common root; quotient pullbacks use
coefficient identities after composition; tangency and differential
defects use minors; field descent uses Frobenius-fixed coefficient equations
\(c^{\abs{\B'}}=c\), which force \(c\) into the recorded subfield
\(\B'\);
and saturation or fixed-dimensional pencils are projective incidence
conditions.  Chevalley's theorem and stability under finite Boolean
operations and projections prove constructibility.  The exceptional
Fourier characters may themselves form a constructible subset of the
character parameter space, but constructibility alone does not imply the
aggregate estimate \textup{(MA)}.
\end{proof}

\subsection{Cell catalogue}\label{sec:cell-catalogue}

The catalogue below is a language for row-specific proofs, not a theorem
that every displayed locus is automatically paid.  In a concrete family,
each item must be accompanied by a distinct-slope projection bound at its
natural profile scale.  The conditional compiler uses only those bounds,
not constructibility itself.

\paragraph{Quotient and periodic cells.}
A quotient cell consists of pullbacks along a nontrivial map
\(\pi:D\to D'\): support membership is constant on the relevant fibers
of \(\pi\), so the support descends to the smaller domain \(D'\).  Such a
support is called \emph{periodic} with respect to \(\pi\).  Algebraically,
the locator or boundary map factors
through the quotient coordinate.  The smaller support count is offset by
the descended prefix depth and possibly a smaller coefficient field, so
the cell is budgeted at its natural quotient term in \(\mathfrak E_n\);
no identity-scale payment is automatic.

\paragraph{Dihedral and Chebyshev cells.}
These are quotient cells with an additional involution, usually inversion
\(u\mapsto u^{-1}\), or with the dihedral symmetry generated by inversion
and multiplication by a root of unity.  The invariant coordinate is
\(x=(u+u^{-1})/2\), and the Chebyshev polynomial \(T_d\) is characterized
by \(T_d(x)=(u^d+u^{-d})/2\).  On polynomial coordinates these cells are
detected by composition with Dickson or Chebyshev maps; in the unnormalized
coordinate the Dickson polynomial is characterized by
\(D_d(u+u^{-1})=u^d+u^{-d}\).  Equivalently they are detected by inversion
invariance after a multiplicative change of variables.  They must be
separated because a naive quotient count may miss the two-to-one or
dihedral fiber structure.  The precise circle domains are defined in
\cref{sec:smooth-circle-domains}.

\paragraph{Planted-block cells.}
A \emph{planted block} is a predetermined group of support positions
forced by an
algebraically controlled common divisor \(P\).  Here \(P\) is a polynomial
factor common to every support locator \(Q_S\), and its roots in \(D\) are
the planted positions.  Divisibility makes the
locus constructible, but payment additionally requires a subexponential
census of allowed \(P\), the residual prefix estimate, and its slope
projection; arbitrary planted subsets are not one profile.

\paragraph{Tangent, deep, and common-line cells.}
For a support \(S\), the \emph{Reed--Solomon evaluation variety} means the
linear image \(\operatorname{ev}_S(\A^k)\subseteq\A^S\).  These cells
record rank-defective contact between the received line and the
corresponding support-restricted evaluation incidence.  The deep regime is controlled exactly by
\cref{thm:deep-regime-upper}.  General tangent charts require the
independent-minor and slope-projection criterion of
\cref{prop:tangent-payment}; rank drop alone is not payment.
A \emph{common-line cell} is a locus on which several nominal witness
charts project to the same received line, while a \emph{tangent cell} is a
determinantal locus where the differential of the slope projection has
smaller than its expected rank.

\paragraph{Extension and field-descent cells.}
A witness lies in an extension or field-descent cell if its data is defined over
a proper subfield (field descent) or factors only after enlarging the coefficient
field (extension).  Frobenius invariance is
constructible.  Its natural subfield profile can be larger than the
identity profile, so a direct field-sensitive slope count is required.

\paragraph{Differential-locator cells.}
The derivative of the locator-prefix map may be expressed by a Vandermonde
or Hasse-derivative Jacobian.  The Hasse derivatives are defined by
\(f(X+T)=\sum_{j\ge0}D^{(j)}f(X)T^j\); unlike division by \(j!\), this
convention remains valid in small characteristic.  A
\emph{differential-locator cell} is the
locus where that matrix loses its expected rank.  It is determinantal, but
its rank loss
need not equal its codimension.  Payment requires independent minors and a
slope-projection estimate as in \cref{prop:rank-pivot-payment}.

\paragraph{Saturation and effective-image-collapse cells.}
The exact projections are those of
\cref{def:explanation-occupancy}: a raw witness first maps to its
explanation state \((\gamma,h)\) and then to its slope \(\gamma\).
Removing common codeword, quotient, and planted factors may identify many
raw witnesses or many explanation states, but distinct explaining
polynomials can still occur at one slope.  Thus \emph{saturation} means
passing through these displayed images with their actual fiber
cardinalities; it is not an assertion that explanation rays and slopes are
the same set.  MCA counts the final slope image, and the occupancy or
image--fiber bound must replace an uncorrected support count.  The cell is
constructible in the projective locator and explanation incidence, but its
projection degree remains an enumerative input.
\emph{Effective-image collapse} is the related event that a boundary map
reaches exponentially fewer boundary values than its ambient codomain contains.

\paragraph{Balanced-core and split-pencil cells.}
A \emph{common core} is a locator factor shared by every member of a
family.  A \emph{balanced core} is a pair of equal-degree monic residual
locators with a common depth-\(w\) prefix after common and planted factors
have been removed.  A \emph{balanced-core chart} is a parameter family of
such cores whose projective dimension exceeds one.  In projective
dimension one, the residual family is a \emph{split pencil}
\(Q_0+\lambda Q_1\) whose members split into roots in the evaluation
domain.  The \emph{moving-root incidence} of the pencil is
defined on the projective parameter line
\(\Pj^1(\F)=\F\cup\{\infty\}\) by
\[
 \{(\lambda,t)\in\Pj^1\times D:
        Q_0(t)+\lambda Q_1(t)=0,
        (Q_0(t),Q_1(t))\ne(0,0)\}.
\]
It is then bounded by
\cref{prop:split-pencil-payment}.  Higher-dimensional balanced-core charts
require a proved decomposition or a direct ray estimate.

\paragraph{Fourier/Sidon-heavy analytic component.}
A primitive prefix fiber may be too large while having exponentially small
additive energy (few additive quadruples relative to its size), so
approximate-group inverse theory cannot handle it.  Here that phrase means
inverse results that extract a large subset with small sumset or difference
set from a set with large additive energy; exponentially small energy gives
no such input.
The required input is the analytic moment bound of
\cref{def:sidon-heavy,def:sidon-paid-cell}.  Li--Wan and Weil control its
minor characters in aggregate when \textup{(MI)} is proved; the
degree-uniform condition \textup{(PF)} is one sufficient certificate.
The remaining effective major characters require \textup{(MA)} or a
direct max-fiber theorem.  These are additive Fourier characters of the
effective boundary target, not automatically of its ambient formal
codomain.  This analytic
component is deliberately kept separate from the constructible atlas and
from the distinct-slope ray compiler.

\begin{remark}
The catalogue is intentionally geometric.  A row proof should not merely state that an exceptional family is bounded; it should identify the moduli stratum, the projection being counted, and the entropy or degree reason why the projection is small.  This is what makes the ledger compatible with limits.
\end{remark}


\section{Counting principles for paid cells}\label{sec:payment-principles}

The closed-ledger hypothesis is verified in practice by a small collection of counting principles.  We record them to make the role of payment precise.  The statements are intentionally asymptotic and profile-uniform; finite rows require the same arguments with constants.

Recall that \(h\) is the natural-log binary entropy fixed in the
Introduction.  A map has \emph{generic fiber size} \(d\) when all points
outside a lower-dimensional exceptional set have exactly \(d\) preimages.
``Transverse'' equations are locally independent equations, so each removes
the expected finite-field dimension.  A locus has \emph{codimension
\(h\)} when its local dimension is \(h\) below that of the ambient witness
chart.  The \emph{support divisor} of \(S\) is
\(\sum_{x\in S}[x]\), equivalently the root data encoded by its locator
\(Q_S\).

\begin{lemma}[Entropy loss from a fixed planted block]\label{lem:planted-entropy}
Let \(P\subseteq D\) have size \(b\), and consider supports \(S\subseteq D\) of size \(a\) with \(P\subseteq S\).  If \(b=\theta n+o(n)\), then their number is
\[
        \binom{n-b}{a-b}
        =\exp\left(n\bigl((1-\theta)h((a/n-\theta)/(1-\theta))+o(1)\bigr)\right).
\]
For fixed density \(a/n\), this is exponentially smaller than \(\binom na\) whenever \(\theta>0\).  If \(b=o(n)\), the loss is absorbed into the profile factor \(e^{o(n)}\).
\end{lemma}

\begin{proof}
This is Stirling's formula applied to \(\binom{n-b}{a-b}\).  Strict entropy loss for \(\theta>0\) follows from the strict concavity of the binary entropy function and from the fact that requiring a positive-density block removes independent support choices.
\end{proof}

\begin{lemma}[Quotient entropy principle]\label{lem:quotient-entropy}
Suppose a quotient cell factors the moving support through a map \(\pi:D\to D'\) of generic fiber size \(d\ge2\), and that the support is constant on a positive-density family of quotient fibers.  Then the number of such supports is exponentially smaller than \(\binom na\), unless the number of constrained quotient fibers is \(o(n)\).  In the latter case the quotient data contributes only \(e^{o(n)}\) profiles.
\end{lemma}

\begin{proof}
On each constrained fiber, a support choice is replaced by a bounded list of invariant choices rather than arbitrary subsets of the fiber.  For a positive-density number of fibers this reduces the entropy per fiber by a positive amount, uniformly away from the endpoints.  Summing over quotient fibers gives an exponential entropy drop.  If only \(o(n)\) fibers are constrained, the number of choices for them is at most \(C^{o(n)}=e^{o(n)}\) for a bounded constant \(C\).
\end{proof}

\begin{lemma}[Rank-defect payment criterion]\label{lem:rank-defect-payment}
Let \(M(x)\) be a matrix of polynomial functions on a witness chart.  The
locus \(\rank M\le r-h\) is cut out by minors.  It has codimension at least
\(h\) only if one exhibits \(h\) independent local equations (or proves
the corresponding determinantal transversality), and a numerator payment
also requires a bound for its projection to slope space.  Under those two
additional hypotheses a codimension-\(h\) finite-field estimate gives the
claimed factor \(\abs\B^{-h+o(n)}\); generic rank alone gives no such
conclusion.
\end{lemma}

\begin{proof}
The rank locus is defined by the stated minors.  If \(h\) independent
equations cut it transversely, bounded-degree point counting gives at most
\(\abs\B^{\dim X-h+o(n)}\) chart points.  The assumed slope-projection
bound converts this to a ray estimate.  Without independence the loss can
be only one even when the rank drops by \(h\) (for example
\(M(x)=xI_h\)), and without projection control a point count need not be a
distinct-slope bound.
\end{proof}

\begin{lemma}[One-pencil moving-root payment]\label{lem:pencil-payment}
Let \(L_{[s:t]}=sQ_0+tQ_1\) be a projective locator pencil, factor its
common \(D\)-root divisor, and let \(g\) be the number of common roots.  If
every counted parameter has at least \(h\ge1\) remaining roots in \(D\),
then the number of counted parameters, and hence the number of slopes when
the slope-to-pencil map is injective, is at most
\[
             \floor{(n-g)/h}.
\]
\end{lemma}

\begin{proof}
For a remaining root \(x\), the pair
\((Q_0(x),Q_1(x))\ne(0,0)\), so
\(sQ_0(x)+tQ_1(x)=0\) has exactly one solution in
\(\Pj^1(\F)\), including the point at infinity when appropriate.  Count
parameter--root incidences.
\end{proof}

\begin{lemma}[Saturation principle]\label{lem:saturation-principle}
Let \(\Ccal\subseteq\Wcal_a(r)\).  If every explanation state in
\(\Rcal(\Ccal)\) has at least \(H\) retained support witnesses, then
\[
 \abs{Z(\Ccal)}\le\abs{\Rcal(\Ccal)}
 \le\floor{\abs{\Ccal}/H}.
\]
If no such lower occupancy is proved, the cell requires a direct bound on
its explanation or slope image.  A lower fiber bound for the first
projection bounds explanation states; the second projection can only
decrease cardinality.
\end{lemma}

\begin{proof}
This is \cref{lem:exact-occupancy-compiler}.
\end{proof}

These principles explain the asymptotic add-back condition in \cref{def:closed-asymptotic-ledger}.  Each paid cell either has a genuine exponential loss, or it is a bounded-complexity profile with only subexponentially many choices.  The primitive leaf is precisely the part where no such loss or profile explanation remains.


\section{Smooth and circle evaluation domains}\label{sec:smooth-circle-domains}

This section records the domain geometry used to close the ledger.  The
multiplicative and circle cases are treated simultaneously through
uniform-fiber folding maps.  The point is not merely that such maps exist,
but that their fibers are complete inside the evaluation domain and their
locators have a \emph{lacunary} top part, meaning that prescribed gaps occur
among the highest-degree coefficients.  These two facts make quotient
cells visible in the witness moduli and keep the bounded-prefix Fourier
calculation stable under planted cores.

\begin{definition}[Complete-fiber folding map and smooth multiplicative coset]
\label{def:structured-folding}
Let \(D\subseteq\B\) be finite.  A polynomial map
\(\varphi:D\to Q:=\varphi(D)\) is a \emph{\(c\)-fold
complete-fiber folding map} if every \(y\in Q\) has exactly \(c\)
preimages in \(D\).  A support is \emph{complete at this folding scale}
if it is a union of whole fibers, equivalently if it is
\(\varphi^{-1}(E)\) for some \(E\subseteq Q\).  Passing from \(D\) to
\(Q\), and from a complete support to \(E\), is the corresponding
\emph{quotient} or \emph{folding}.

A \emph{multiplicative coset} is a set \(D=\theta H\subseteq\B^\times\),
where \(H\le\B^\times\) is a subgroup and \(\theta\in\B^\times\).
For every divisor \(c\mid\abs H\), the restriction
\[
             \pi_c:D\longrightarrow D^{(c)},\qquad
             \pi_c(x)=x^c,\qquad
             D^{(c)}=\{x^c:x\in D\},
\]
is a \(c\)-fold complete-fiber folding map.  We call \(D\)
\emph{smooth at scale \(c\)} when this power folding is among the
quotients retained by the row, and \emph{smooth} for a family of scales
when it is smooth at every scale in that family.  Thus the adjective
\emph{smooth} here is an explicit uniform-fiber property, not a claim
about nonsingularity of an algebraic variety.  The locator of a complete
support is \(V_E(X^c)\); its nonleading terms occur only in degree gaps
divisible by \(c\).  This controlled pattern of missing top coefficients
is what \emph{lacunary} means in this section.
\end{definition}

\begin{definition}[Circle domain and circle code]
\label{def:circle-domain-code}
Let \(\B_0\) have odd characteristic and let
\(\mathcal C/\B_0\) be the affine conic
\(\mathcal C:x^2+y^2=1\).  A \emph{standard circle domain} is a finite
set \(E\subseteq\mathcal C(\B_0)\).  If \(\F\supseteq\B_0\), the
degree-\(w\) \emph{circle code} on \(E\) is
\[
 \mathcal C_w(\F,E)
 =\{(f(P))_{P\in E}:
      f\in\operatorname{im}(\F[x,y]_{\le w}
          \to\F[x,y]/(x^2+y^2-1))\}.
\]
The degree is the total-degree filtration before passage to the quotient
ring.  A \emph{circle \(x\)-domain} is instead a set of first
coordinates obtained from the torus construction below; an RS code on
that \(x\)-domain and a bivariate circle code are different presentations
until a diagonal equivalence such as
\cref{lem:circle-uniformization} has been exhibited.
\end{definition}

\begin{definition}[Norm-one torus, twin coset, and Chebyshev domain]
\label{def:circle-twin-domain}
Let \(p\equiv3\pmod4\), choose \(i\in\F_{p^2}\) with \(i^2=-1\), and
write
\[
 \UU=\ker\!\left(N_{\F_{p^2}/\F_p}:\F_{p^2}^{\times}
                      \to\F_p^\times\right)
     =\{u\in\F_{p^2}^{\times}:u^{p+1}=1\}.
\]
The algebraic kernel of the norm morphism is the
\emph{norm-one torus}; \(\UU\) is its cyclic group of
\(\F_p\)-rational points.  Put
\(\chi(u)=(u+u^{-1})/2\in\F_p\).  If \(H\le\UU\) and
\(g\in\UU\) satisfies \(g^2\notin H\), then
\[
 \Dcal(g,H)=gH\sqcup g^{-1}H
\]
is a disjoint union exchanged by \(u\mapsto u^{-1}\); it is called a
\emph{twin coset}.  Its projected set
\[
             D(g,H)=\chi(\Dcal(g,H))\subseteq\F_p
\]
has cardinality \(\abs H\) and is the \emph{Chebyshev twin-coset
\(x\)-domain}.

For \(c\ge1\), let \(T_c\) be the Chebyshev polynomial normalized by
\[
       T_c\!\left(\frac{u+u^{-1}}2\right)
       =\frac{u^c+u^{-c}}2 .
\]
Thus \(T_c\circ\chi=\chi\circ(u\mapsto u^c)\).  Whenever
\(c\mid\abs H\) and the two image cosets remain disjoint, the restriction
of \(T_c\) to \(D(g,H)\) is the \emph{Chebyshev folding map}; a
Chebyshev twin-coset domain equipped with these complete-fiber maps is
the circle analogue of a smooth multiplicative coset.
\end{definition}

\begin{remark}[Characteristic two and binary-extension rows]
\label{rem:characteristic-two-rows}
The entropy function \(h\), the locator-prefix bijection, the
syndrome--secant identity, the pole-line compiler, and the
quotient/remainder normal form are characteristic-free.  The use of
power sums \(\sum_{t\in S}t^j\) is not: Newton's identities, which triangularly relate elementary
locator coefficients to power sums, give a change from the first \(R\)
elementary locator coefficients to the first \(R\) power sums is
invertible only when \(R<\operatorname{char}\B\).  In characteristic two
the elementary locator prefix therefore remains the canonical exact
coordinate system, while any power-sum statement must be proved directly.

A literal Reed--Solomon code over \(\F_2\) has length at most two.  Thus an
asymptotic ``binary-field'' RS row means a tower
\(\F_{2^{s_n}}\) (or a scalar extension of such a field), with
\(s_n\ge\log_2 n\); the regime in this paper permits
\(s_n=o(n)\).  This is distinct from taking a binary subfield subcode
\(C\cap\F_2^D\) or a trace code obtained by applying the field trace
coordinatewise, operations that do not preserve the RS agreement problem
used here.  The notation \(H_2\) denotes base-two entropy and has no
connection with the coefficient-field characteristic.  Interleaving an
extension-field row or expanding its symbols into binary coordinates also
changes the code family and requires a separate comparison theorem.

The standard affine circle does not survive unchanged in characteristic
two: \(x^2+y^2-1=(x+y+1)^2\), so it is a doubled line rather than the
smooth norm-one conic used in
\cref{lem:circle-uniformization,prop:stereographic-circle-equivalence}.
Smooth characteristic-two norm
conics can instead be written in a form such as
\(x^2+xy+\delta y^2=1\), with the quadratic polynomial chosen
irreducible, but transferring the circle ledger to that model requires a
separate torus uniformization.  Accordingly
\cref{thm:binary-cubic-obstruction} below is a theorem for smooth
multiplicative rows and makes no claim for the standard characteristic-two
circle.
\end{remark}

\begin{definition}[Map-smooth domain]\label{def:map-smooth}
Let \(\B\subseteq\F\) be finite fields, let \(\varphi\in\B[X]\) have
degree \(c\ge1\), and let \(D\subseteq\B\) have size \(n\).  The set
\(D\) is \((\varphi,c)\)-smooth if every nonempty fiber
\(\varphi^{-1}(y)\cap D\), \(y\in\varphi(D)\), has exactly \(c\)
elements.  We write \(Q=\varphi(D)\) and \(N=n/c\).
\end{definition}

A multiplicative coset \(D=\theta H\subseteq\B^\times\) is \((X^a,a)\)-smooth for every divisor \(a\mid\abs H\).  The circle case uses Chebyshev maps rather than power maps; see \cref{lem:cheb-smooth}.  The following lemma is the common quotient-support construction.  It is included because it is the lower-side source of the entropy frontier and because the same lacunarity is used by the upper ledger to identify quotient cells.

\begin{lemma}[Map-smooth fiber construction]\label{lem:map-smooth-fiber}
Let \(D\) be \((\varphi,a)\)-smooth, \(\abs D=n\), and \(C^+=\RS_\F(D,k+1)\).  Put \(\ell=\lfloor k/a\rfloor+2\) and \(A=a\ell\).  If \(\ell\le N-1\), then there is \(z\in\B\) and a word \(u_z\in\B^D\) such that \(u_z\) has at least \(\binom N\ell/\abs\B\) distinct codewords of \(C^+\) agreeing with it on at least \(A\) positions.  Moreover \(k+a+1\le A\le k+2a\), with equality \(A=k+2a\) if \(a\mid k\).
\end{lemma}

\begin{proof}
For each \(\ell\)-subset \(E\subseteq Q\), let \(P_E(Y)=\prod_{b\in E}(Y-b)=Y^\ell-e_1(E)Y^{\ell-1}+R_E(Y)\), with \(\deg R_E\le\ell-2\).  The composed locator \(P_E(\varphi(X))\) vanishes on the union of the \(a\)-point fibers above \(E\), hence on \(A=a\ell\) points of \(D\).  Set \(z_E=-e_1(E)\) and \(c_E=-R_E(\varphi)\).  Since \(\deg c_E\le a(\ell-2)\le k\), the word \(c_E\) lies in \(C^+\), and it agrees on that fiber union with \(u_{z_E}(x)=\varphi(x)^\ell+z_E\varphi(x)^{\ell-1}\).  Among the \(\binom N\ell\) subsets \(E\), some value of \(z_E\in\B\) occurs at least \(\binom N\ell/\abs\B\) times.  If two such subsets with the same \(z\) gave the same codeword, then \(R_E(\varphi)=R_{E'}(\varphi)\), hence \(R_E=R_{E'}\) because composition with a nonconstant polynomial is injective, and then \(P_E=P_{E'}\).  Thus the codewords are distinct.  The bounds on \(A\) follow from \(\ell=\lfloor k/a\rfloor+2\).
\end{proof}

\begin{proposition}[Collision-aware lower cap for map-smooth domains]\label{prop:map-smooth-cap}
Let \(D\) be \((\varphi,a)\)-smooth, \(n<\abs\F=q\), and let
\(C=\RS_\F(D,k)\).  With \(\ell,A\) as in
\cref{lem:map-smooth-fiber}, assume \(\ell\le N-1\), and put
\[
 L=\left\lceil\binom N\ell\abs\B^{-1}\right\rceil.
\]
Then for every integer threshold \(k+1\le m\le A\),
\[
 B_C^{\rm MCA}(m)\ge
 \left\lceil\frac{L(q-n)}{q-n+k(L-1)}\right\rceil .
\]
\end{proposition}

\begin{proof}
The list in \cref{lem:map-smooth-fiber} has size at least \(L\) at
agreement \(A\).  Apply the collision-aware pole conversion
\cref{thm:collision-aware-pole}.  Its line pair is not simultaneously
explained on any \(k+1\) points, so an agreement set of size \(A\) also
witnesses badness for every threshold \(m\) in the displayed range.
\end{proof}

We now verify the complete-fiber property for the construction in
\cref{def:circle-twin-domain}.  Thus
\(\Dcal=gH\sqcup g^{-1}H\subseteq\UU\),
\(D=\chi(\Dcal)\), and \(T_a(\chi(u))=\chi(u^a)\).

\begin{lemma}[Chebyshev smoothness]\label{lem:cheb-smooth}
Let \(\Dcal=gH\cup g^{-1}H\subseteq\UU\) be a twin coset, let \(D=\chi(\Dcal)\), and let \(a\mid\abs H\).  If \(g^{2a}\notin H^{(a)}=\{h^a:h\in H\}\), then \(D\) is \((T_a,a)\)-smooth.  If \(\ord(g)\nmid2\abs H\), this holds for all \(a\mid\abs H\).
\end{lemma}

\begin{proof}
The map \(u\mapsto u^a\) has kernel of size \(a\) on \(\UU\), and that kernel lies in \(H\).  Hence it maps \(gH\) and \(g^{-1}H\) onto \(g^aH^{(a)}\) and \(g^{-a}H^{(a)}\), respectively, with fibers of size \(a\).  The hypothesis says these image cosets are disjoint, so the image is again a twin coset.  Since \(\chi\) is exactly two-to-one on both twin cosets and identifies only inverse points, the equation \(T_a(x)=y\) on \(D\) has exactly \(a\) solutions for each \(y\in T_a(D)\).  Taking \(a=\abs H\) shows that the all-scale condition is equivalent to \(g^{2\abs H}\ne1\), and the implication to every smaller divisor follows by raising to \(\abs H/a\).
\end{proof}

\begin{lemma}[Diagonal uniformization of circle codes]\label{lem:circle-uniformization}
Assume \(\F\supseteq\F_{p^2}\).  Let
\(E\subseteq\{(x,y)\in\F_p^2:x^2+y^2=1\}\), define
\(\iota(x,y)=x+iy\), put \(E'=\iota(E)\subseteq\UU\), and let
\(\Ccal_w(\F,E)\) be the circle code of
\cref{def:circle-domain-code}.  Under the coordinate bijection \(\iota\)
and the diagonal multiplier \(t_u=u^{-w}\), the code
\(\Ccal_w(\F,E)\) is \(t\circ\RS_\F(E',2w+1)\).
Consequently list sizes, ordinary correlated-agreement (CA) numerators,
and mutual-correlated-agreement (MCA) numerators are preserved.
\end{lemma}

\begin{proof}
The map \(u=x+iy\) identifies \(\F[x,y]/(x^2+y^2-1)\) with \(\F[u,u^{-1}]\).  The subspace of total degree at most \(w\) maps to \(u^{-w}\F[u]_{\le2w}\), and both spaces have dimension \(2w+1\).  Multiplication by the nonzero diagonal vector \((u^{-w})_{u\in E'}\) preserves coordinatewise equality and hence every agreement notion used here.  Therefore it preserves list, CA, and MCA counts exactly.
\end{proof}

\begin{definition}[Artin--Schreier phase]
\label{def:artin-schreier-phase}
Let \(\B=\F_{p^s}\), let \(X/\B\) be the affine line or the norm-one
torus, and write a nontrivial additive character as
\(\psi_a(t)=\exp(2\pi i\,\operatorname{Tr}_{\B/\F_p}(at)/p)\) with
\(a\in\B^\times\).  A rational phase \(R\in\B(X)\) is
\emph{Artin--Schreier degenerate for \(\psi_a\)} if
\[
                 aR=G^p-G+c
\]
for some \(G\in\B(X)\) and \(c\in\B\).  Away from the poles,
\(\psi_a(R)\) then has constant trace phase, so the square-root
cancellation asserted by the ordinary Weil estimate is unavailable.
A phase not admitting such a representation is called
\emph{non-Artin--Schreier}.  This qualification is always relative to
the chosen additive character.
\end{definition}

The Fourier/Sidon payment uses bounded-prefix flatness as a sufficient minor-arc estimate, but the ledger also records low-energy heavy fibers as a separate first-match cell.  We state here the smooth and circle flatness forms and then prove the Sidon-resolved payment.  The proof and the general weighted statement are given in \cref{sec:fourier,sec:li-wan-details}; the Sidon-resolved payment is \cref{thm:sidon-resolved-payment}.  For a full profile slice write
\[
        \Omega^0=\Omega_{T,m}=\{x\in\{0,1\}^{T}:\sum_{t\in T}x_t=m\},
        \qquad
        \barN_0=\abs\B^{-R}\binom{\abs T}{m}.
\]
Here \(\barN_0\) is the ambient profile scale.  The theorems below prove
an ambient max-fiber bound, which is stronger than image-normalized Q and,
by \cref{rem:flatness-certifies-image}, simultaneously certifies that the
realized image has full exponential size.
The prefix functions in a primitive chart are weighted rational functions
\(g_1,\ldots,g_R\) on \(T\), and the full profile map is
\(\Psi(x)=\sum_{t\in T}x_t(g_1(t),\ldots,g_R(t))\).  For
\(\alpha\in\B^R\), put
\(R_\alpha(t)=\alpha\cdot(g_1(t),\ldots,g_R(t))\).  The algebraic
conditions listed in \cref{def:major-arc} define an ambient exceptional
locus of phase parameters.  They do not yet define a major/minor partition
of the effective dual, because restriction to the effective span has
nonunique ambient lifts.  The effective partition and, for every character
treated by a Weil estimate, a certified ambient lift whose phase satisfies
the stated nondegeneracy hypotheses are part of the row-specific chart
data.  Naming the corresponding primal stratum neither removes its
effective character from Fourier inversion nor proves its aggregate
payment.  Finally,
\[
 e_m(a_t:t\in T)=\sum_{S\in\binom Tm}\prod_{t\in S}a_t
\]
is the \(m\)-th elementary symmetric sum.  Naming the
corresponding primal stratum does not remove these characters from Fourier
inversion on the full slice, so their aggregate is a separate input.

There is, however, an exact normalization that must be performed before
declaring any character major.  Assume \(1\le m\le\abs T-1\), as on every
fixed-density primitive leaf, put \(p=\operatorname{char}\B\), and fix
\(S_0\in\binom Tm\) and \(t_0\in T\).
Then
\[
 \Span_{\F_p}\{\Psi(S)-\Psi(S_0):S\in\tbinom Tm\}
 =\Span_{\F_p}\{g(t)-g(t_0):t\in T\}.               \tag{EF0}
\]
Define the \emph{effective Fourier span}
\[
 V_g=\Span_{\F_p}\{g(t)-g(t_0):t\in T\}\subseteq\B^R,
 \qquad A_{\rm eff}=\abs{V_g}.
 \tag{EF1}\label{eq:effective-fourier-span}
\]
The fixed-weight sum map takes values in the affine space
\(m g(t_0)+V_g\).  Thus Fourier inversion belongs on the additive group
\(V_g\), not automatically on the ambient group \(\B^R\).

\begin{lemma}[Effective-span Fourier inversion]
\label{lem:effective-span-fourier}
For \(z\in V_g\), let
\[
 N_g(z)=\abs{\left\{S\in\binom Tm:
       \sum_{t\in S}g(t)=m g(t_0)+z\right\}}.
\]
Then
\[
 N_g(z)=\frac1{A_{\rm eff}}
 \sum_{\chi\in\widehat{V_g}}\overline{\chi(z)}
 e_m\bigl(\chi(g(t)-g(t_0)):t\in T\bigr).           \tag{EF2}
\]
In particular,
\[
 \max_zN_g(z)\le\frac1{A_{\rm eff}}
 \left(\binom{\abs T}m+
 \sum_{1\ne\chi\in\widehat{V_g}}
 \left|e_m\bigl(\chi(g(t)-g(t_0)):t\in T\bigr)\right|\right). \tag{EF3}
\]
Every character of \(V_g\) extends to a trace character of \(\B^R\),
and a nontrivial effective character is not constant on all the generators
\(g(t)-g(t_0)\).  Ambient characters that are trivial on \(V_g\) occur
with multiplicity \(\abs{\B}^R/A_{\rm eff}\), which cancels exactly
against the ambient Fourier normalization.
\end{lemma}

\begin{proof}
The inclusion from left to right in \textup{(EF0)} follows by expanding a
difference of two support sums.  Conversely, for distinct \(t,u\in T\),
choose an \((m-1)\)-set \(A\subseteq T\setminus\{t,u\}\); such a set
exists because \(1\le m\le\abs T-1\).  Then
\(\Psi(A\cup\{t\})-\Psi(A\cup\{u\})=g(t)-g(u)\), so the two spans agree.
The first formula is character orthogonality on the finite additive group
\(V_g\), applied after translating by \(m g(t_0)\); the second follows by
the triangle inequality.  The trace pairing on the \(\F_p\)-space
\(\B^R\) is nondegenerate, so every character of its subspace \(V_g\)
extends to the ambient space.  A character trivial on each displayed
generator is trivial on their span.  Finally, the annihilator
\(V_g^\perp\) has size \(\abs{\B}^R/A_{\rm eff}\); grouping ambient
characters by their restriction to \(V_g\) proves the cancellation claim.
\end{proof}

\begin{definition}[Certified effective major/minor partition]
\label{def:effective-major-minor}
Let
\(\operatorname{res}:\widehat{\B^R}\twoheadrightarrow\widehat{V_g}\)
be restriction.  A \emph{certified effective Fourier partition} is a
disjoint union
\[
 \widehat{V_g}\setminus\{1\}
 =\mathfrak m_{\rm eff}\sqcup\mathfrak M_{\rm eff}
\]
such that every \(\chi\in\mathfrak m_{\rm eff}\) is supplied with an
ambient lift whose rational phase satisfies the character-sum hypotheses
used to pay it.  Characters without such a certified lift may be placed in
\(\mathfrak M_{\rm eff}\), but their aggregate must then be paid by
\textup{(MA)}.  This designation is data of the chart; it is not a
property of an unspecified ambient lift.
\end{definition}

\begin{definition}[Effective Fourier payment]
\label{def:effective-fourier-payment}
The full slice has an \emph{effective Fourier payment with constant
\(\kappa\ge1\)} if
\[
 \sum_{1\ne\chi\in\widehat{V_g}}
 \left|e_m\bigl(\chi(g(t)-g(t_0)):t\in T\bigr)\right|
 \le (\kappa-1)\binom{\abs T}m.                    \tag{EFP}
\]
By \cref{lem:effective-span-fourier}, this implies the exact max-fiber
bound
\[
               \max_zN_g(z)\le
               \kappa\binom{\abs T}m/A_{\rm eff}. \tag{EF4}
\]
It also implies that the realized image contains at least
\(A_{\rm eff}/\kappa\) points.  The ambient instances of \textup{(PF)}
and \textup{(MA)} are one sufficient way to prove \textup{(EFP)}, but
\textup{(EFP)} is phase-dependent and does not pay an ambient annihilator
twice.
\end{definition}

\begin{proof}[Proof of the image-size assertion]
The fibers sum to \(\binom{\abs T}m\), and each nonempty fiber is at most
\(\kappa\binom{\abs T}m/A_{\rm eff}\) by \textup{(EF4)}.
Hence at least \(A_{\rm eff}/\kappa\) fibers are nonempty.
\end{proof}

\begin{lemma}[Constant and Artin--Schreier modes are annihilator modes]
\label{lem:constant-modes-annihilator}
Suppose an ambient trace character has phase constant on \(T\).  Then its
restriction to \(V_g\) is trivial.  In particular, if an ambient rational
phase is Artin--Schreier of the form \(G^p-G+c\) and \(G\) is defined at
the \(\B\)-rational points of \(T\), then its character is an annihilator
mode and disappears from the nontrivial sum in \textup{(EF2)}.
\end{lemma}

\begin{proof}
The quotient of the character values at \(g(t)\) and \(g(t_0)\) is one
for every \(t\), so the character is trivial on the generators of
\(V_g\).  For an Artin--Schreier phase,
\(\Tr_{\B/\F_p}(G(t)^p-G(t))=0\), hence its additive-character value is
constant on all points where it is defined.
\end{proof}

\begin{remark}[Effective normalization does not prove the remaining aggregate]
The effective-span reduction removes exactly the constant annihilator
multiplicity.  A nonconstant character may still factor through a quotient
map or have a phase-dependent estimate too weak for \textup{(EFP)}.  Such
modes require a recursive quotient estimate, the aggregate
\textup{(MA)}, or a direct full-slice bound.  Thus effective normalization
repairs the Fourier denominator but is not a proof of unrestricted deep-prefix
flatness.
\end{remark}

\begin{definition}[Major-character aggregate]
\label{def:major-arc-aggregate}
Let \(A_*\) be the additive target on which Fourier inversion is being
performed, let \(h:T\to A_*\), and let
\(\mathfrak M_*\subseteq\widehat{A_*}\setminus\{1\}\) be the designated
major characters.  The chart satisfies \textup{(MA)} on \(A_*\) if
\[
 \tag{MA}\label{eq:major-arc-aggregate}
 \sum_{\chi\in\mathfrak M_*}
 \left|e_m\bigl(\chi(h(t)):t\in T\bigr)\right|
 \le e^{o(\abs T)}\binom{\abs T}m .
\]
Unless the ambient target is explicitly named, \(A_*=V_g\) and
\(h(t)=g(t)-g(t_0)\).  For the ambient target \(A_*=\B^R\), the same
condition is equivalently
\[
 \abs\B^{-R}\sum_{\alpha\in\mathfrak M_{\rm amb}}
 \left|e_m\bigl(\psi(\alpha\cdot g(t)):t\in T\bigr)\right|
 \le e^{o(\abs T)}\barN_0.
\]
The condition is vacuous when the designated major set is empty.  It can
be proved by a recursive quotient census or replaced by a direct
full-slice max-fiber bound; first-match terminology alone does not imply
it.
\end{definition}

\begin{lemma}[Sparse effective major set]
\label{lem:sparse-major-payment}
On the effective group \(V_g\), let \(\mathfrak M_{\rm eff}\) be any set
of nontrivial characters designated major.  Its exact normalized Fourier
loss is at most \(\abs{\mathfrak M_{\rm eff}}\).  Consequently
\(\abs{\mathfrak M_{\rm eff}}=e^{o(\abs T)}\) proves the effective form
of \textup{(MA)}; for a finite row, the literal cardinality is a valid
major-character constant.
\end{lemma}

\begin{proof}
Every fixed-weight Fourier coefficient is a sum of \(\binom{\abs T}m\)
complex numbers of modulus one, and hence has absolute value at most that
binomial coefficient.  Sum this bound over the designated set.
\end{proof}

\begin{theorem}[Exact low-boundary-entropy closure]
\label{thm:small-effective-dual-closure}
Let \(\Omega=\binom Tm\), let \(W\) be a finite additive
\(\F_p\)-space, and let \(\Psi:\Omega\to W\) be a fixed-weight additive
boundary map.  Fix \(S_0\in\Omega\), put \(z_0=\Psi(S_0)\), and define
the effective difference span
\[
 V=\Span_{\F_p}\{\Psi(S)-\Psi(S_0):S\in\Omega\}\subseteq W.
\]
Then \(\Psi(\Omega)\subseteq z_0+V\).  Put
\[
       A=\abs V,\qquad L=\abs{\Psi(\Omega)},
       \qquad M=\binom{\abs T}m,\qquad \bar N=M/L.
\]
Designate every nontrivial character of \(V\) as major and no character
as minor.  Then the literal finite-row constants are
\[
 C_{\rm min}=0,\qquad C_{\rm maj}\le A-1,\qquad
 \max_z\abs{\Psi^{-1}(z)}\le L\bar N=M.            \tag{SE1}
\]
For every exact-agreement witness cell whose support projection is
contained in \(\Omega\), its first-match slope set satisfies
\[
                         \abs Z\le L\bar N=M.       \tag{SE2}
\]
Consequently, if \(\log A=o(\abs T)\), effective \textup{(MI)},
effective \textup{(MA)}, image-normalized Q, and the direct alternative
in \textup{(RC)} all hold with subexponential loss.  Taking the
support-restricted incidence over \(\Omega\) as one ordered cell is
exhaustive for that restricted incidence; it is exhaustive for the whole
exact-agreement incidence when \(\Omega=\binom Da\), or after an
exhaustive partition into such slices.  For a locator prefix with
\(\abs T=\Theta(n)\), the concrete condition
\((a-k-1)\log\abs\B=o(n)\) implies this hypothesis because
\(A\le\abs\B^{a-k-1}\).
\end{theorem}

\begin{proof}
Fourier inversion is performed on the effective group \(V\).  The minor
sum is empty.  Every fixed-weight Fourier coefficient has
modulus at most \(M\), so the normalized major loss is at most \(A-1\).
The max-fiber estimate in \textup{(SE1)} is the trivial bound by the whole
support slice, written as the exact image-normalized multiplier \(L\).
By exact-agreement reduction, every bad slope in the cell has a noncommon
\(m\)-support, and a fixed noncommon support carries at most one slope.
Choosing one support for each slope is therefore injective and proves
\textup{(SE2)}.  Since \(L\le A\), all three finite multipliers
\(A-1,L,L\) are \(e^{o(\abs T)}\) when \(\log A=o(\abs T)\).
The one-cell atlas is exhaustive for its support-restricted incidence by
construction; the stated whole-incidence alternatives give global
exhaustivity.  If it is refined by a partial-occupancy partition,
\textup{(PO1)}--\textup{(PO2)} add the same support set back exactly and
do not change the conclusion.
\end{proof}

\begin{definition}[Aggregate minor payment]
\label{def:aggregate-minor-payment}
After effective-span normalization, let \(\mathfrak m_{\rm eff}\) be the
nontrivial characters not designated major.  The chart satisfies
\textup{(MI)} if
\[
 \sum_{\chi\in\mathfrak m_{\rm eff}}
 \left|e_m\bigl(\chi(g(t)-g(t_0)):t\in T\bigr)\right|
 \le e^{o(\abs T)}\binom{\abs T}m.                 \tag{MI}
\]
For a finite row the corresponding exact loss is this sum divided by
\(\binom{\abs T}m\).  The pointwise condition \textup{(PF)} is only one
sufficient way to prove \textup{(MI)}; phase-dependent estimates may be
strictly stronger.
\end{definition}

\begin{proposition}[Effective MI plus MA]
\label{prop:effective-mi-ma-flatness}
Partition the nontrivial effective characters as
\(\widehat{V_g}\setminus\{1\}
=\mathfrak m_{\rm eff}\sqcup\mathfrak M_{\rm eff}\).  Suppose
\textup{(MI)} and the effective form of \textup{(MA)} in
\cref{def:major-arc-aggregate} hold.
Then the full slice and every first-match residual satisfy
\[
 \max_z\abs{\{S:\Psi(S)=z\}}
 \le e^{o(\abs T)}\binom{\abs T}m/A_{\rm eff}.     \tag{EF5}
\]
The realized image has size at least
\(e^{-o(\abs T)}A_{\rm eff}\), so the same assertion holds at image
normalization.  For a finite row the exact multiplier is
\[
 1+\binom{\abs T}m^{-1}
 \sum_{1\ne\chi\in\widehat{V_g}}
 \left|e_m\bigl(\chi(g(t)-g(t_0)):t\in T\bigr)\right|.
 \tag{EF6}
\]
\end{proposition}

\begin{proof}
Add the minor and major sums and apply
\cref{lem:effective-span-fourier,def:effective-fourier-payment}.  A
first-match residual only deletes supports from a full-slice fiber.
\end{proof}

\begin{corollary}[Exact finite minor and major constants]
\label{cor:exact-finite-fourier-constants}
Put \(M=\binom{\abs T}m\) and, for an arbitrary partition of the
nontrivial effective dual, define
\[
 C_{\rm min}=M^{-1}\sum_{\chi\in\mathfrak m_{\rm eff}}\abs{E_m(\chi)},
 \qquad
 C_{\rm maj}=M^{-1}\sum_{\chi\in\mathfrak M_{\rm eff}}\abs{E_m(\chi)},
\]
where
\(E_m(\chi)=e_m(\chi(g(t)-g(t_0)):t\in T)\).  Then
\[
              \max_zN_g(z)\le
              \frac{M}{A_{\rm eff}}
              (1+C_{\rm min}+C_{\rm maj}).          \tag{EF7}
\]
If every minor character has power sums bounded by \(\Lambda\) through
order \(m\), then
\[
 C_{\rm min}\le
 \frac{\abs{\mathfrak m_{\rm eff}}}{M}
 \binom{\ceil\Lambda+m-1}{m}.                      \tag{EF8}
\]
Thus \textup{(EF7)}, with the two actual finite sums, records every finite
Fourier loss in an adjacent-row certificate exactly.
\end{corollary}

\begin{proof}
Formula \textup{(EF7)} is \textup{(EF3)} after splitting the nontrivial
characters.  The cycle-index coefficient bound gives
\(\abs{E_m(\chi)}\le\binom{\ceil\Lambda+m-1}{m}\), proving
\textup{(EF8)}.
\end{proof}

\begin{definition}[Pointwise minor-arc range]\label{def:prefix-flat-range}
A primitive weighted prefix chart satisfies the pointwise minor-arc range if the characteristic is larger than \(m\), every certified minor lift has a rational phase of total pole-degree at most \(\Delta_{\rm pole}\) and satisfies the nondegeneracy hypotheses of \cref{prop:weighted-weil-minor-arcs}, and, with
\[
        \Lambda=C_0(\Delta_{\rm pole}+1)\sqrt{\abs\B}+\abs P
\]
for the absolute Weil constant \(C_0\) and planted exceptional set \(P\), one has
\[
        \tag{PF}
        R\log\abs\B+
        \log\binom{\Lambda+m-1}{m}
        -\log\binom{\abs T}{m}\le o(\abs T).
\]
For circle twin-cosets the same definition is used with \(2C_0(\Delta_{\rm pole}+1)\sqrt{\abs\B}\) in place of \(C_0(\Delta_{\rm pole}+1)\sqrt{\abs\B}\), accounting for the two inverse-coset lifts on the torus before the involution quotient.
\end{definition}

\begin{proposition}[The degree-uniform PF certificate fails at deep polynomial-field prefixes]
\label{prop:pf-deep-prefix-barrier}
Let \(N=\abs T\), let \(m/N\to\theta\in(0,1)\), and suppose the
effective target has size \(A_N\) with
\[
             \log A_N=\log\binom Nm+o(N).           \tag{PFb1}
\]
Here the effective version replaces \(R\log\abs\B\) by \(\log A_N\);
because \(A_N\le\abs\B^R\), failure of this smaller effective expression
also implies failure of the displayed ambient certificate.
If the uniform cycle-index majorant used for the nontrivial characters has
parameter \(\Lambda_N\ge cN\) for some fixed \(c>0\), then the
degree-uniform sufficient condition \textup{(PF)} fails: its left side is
at least
\[
 N(c+\theta)h\!\left(\frac{\theta}{c+\theta}\right)+o(N),
 \tag{PFb2}
\]
and hence is positive of linear order.

More intrinsically, fix \(\theta\in(0,1)\), put
\(T_p=\F_p^\times\), \(N=p-1\), and choose \(m_p/N\to\theta\).  Let
\(x_p=\log_p\binom N{m_p}\).  By the prime number theorem one may choose
a prime
\[
 R_p+1\le x_p,\qquad R_p=(1-o(1))x_p,\qquad R_p\nmid p-1;
\]
indeed the two largest primes below \(x_p-1\) are both
\((1-o(1))x_p\), while their product eventually exceeds \(p-1\), so at
most one can divide \(p-1\).  Use the depth-\((R_p+1)\) power-sum chart,
whose coordinates include \(X^{R_p}\).
Then
\[
 0\le\log\binom N{m_p}-(R_p+1)\log p=o(N),
\]
so its ambient mean is \(e^{o(N)}\).  The phase \(X^{R_p}\) is
nonconstant, has no proper power-map stabilizer because
\(\gcd(R_p,p-1)=1\), is not Artin--Schreier because \(R_p<p\), and has
no field-descent alternative over the prime field.  Thus it is a minor
phase under the algebraic classification, and every displayed
degree-uniform Weil certificate has \(\Delta_{\rm pole}\ge R_p\).
Consequently its uncapped parameter satisfies
\(\Lambda_p\gg R_p\sqrt p\gg N\), and
\[
 (R_p+1)\log p+\log\binom{\Lambda_p+m_p-1}{m_p}
 -\log\binom N{m_p}=\Omega(N\log N).
\]
Capping the power-sum bound by the trivial value \(N\) still leaves a
positive \(\Omega(N)\) error, so the degree-uniform \textup{(PF)}
certificate fails.  The analogous two-branch circle majorant is no
smaller.
\end{proposition}

\begin{proof}
Since \(A_N\le\abs\B^R\), \textup{(PFb1)} shows that the left side of
\textup{(PF)} is at least
\(\log\binom{\Lambda_N+m-1}{m}+o(N)\).  This expression is increasing
in \(\Lambda_N\).  Stirling's formula at \(\Lambda_N=cN+O(1)\) gives
\textup{(PFb2)}, whose coefficient is strictly positive.  The remaining
assertions follow from \(R_p\asymp p/\log p\), the degree of
\(X^{R_p}\) on the projective line, and the elementary fact that a prime
\(R_p\nmid p-1\) is coprime to \(p-1\).
\end{proof}

\begin{remark}[The barrier concerns the certificate, not every phase]
The proposition does not prove that a deep prefix is nonflat.  It proves
that substituting one worst-degree Weil number into every cycle of every
minor character cannot establish flatness there.  Indeed exponentiation
by the prime \(R_p\) in its construction permutes \(\F_p^\times\), and
for every \(1\le j\le m_p<p\)
\[
        \sum_{t\in\F_p^\times}\psi(jt^{R_p})=-1,
\]
although the degree-uniform Weil majorant is much larger.  Any proof at
the deep frontier must therefore retain phase-dependent power sums or
prove their aggregate directly.  Smoothness alone supplies neither
estimate.
\end{remark}

\begin{proposition}[Ambient PF/MA dichotomy on the square-field row]
\label{prop:pf-ma-square-dichotomy}
In the odd-characteristic rows of
\cref{thm:smooth-quotient-obstruction}, use the power-sum presentation of
the depth-\(w\) prefix, with \(w\) even.  There is an \(\F_p\)-subspace
\(U\subseteq\B^w\) of cardinality \(p^{w/2}\) such that every
\(\xi\in U\) has constant additive-character value one on the domain.
Consequently, no partition of the nonzero \emph{ambient} characters into
minor and major sets can satisfy both \textup{(PF)} and \textup{(MA)} at
the identity entropy crossing.  Classifying a nonzero member of \(U\) as
minor makes the pointwise cycle parameter at least \(n\); classifying all
of \(U\setminus\{0\}\) as major gives a normalized major loss at least
\[
       p^{w/2}-1
       =\exp\bigl((h(\alpha)/4+o(1))n\bigr).         \tag{PFMA}
\]
The entire subspace \(U\) is an annihilator of the effective Fourier span,
so \cref{lem:effective-span-fourier} removes it exactly.  This proves the
necessity of effective-image and quotient routing, but it does not bound
the remaining nonconstant effective characters.
\end{proposition}

\begin{proof}
Write \(\B=\F_{p^2}\), let
\(K=\ker\Tr_{\F_{p^2}/\F_p}\), and use the construction in which
\(D^2=\theta^2\F_p^\times\).  Put
\[
 U=\left\{\xi\in\B^w:
 \xi_{2j}=b_j\theta^{-2j},\ b_j\in K,\quad
 \xi_{2j-1}=0\right\}.
\]
There are \(w/2\) free \(\F_p\)-coordinates, so \(\abs U=p^{w/2}\).
If \(t^2=\theta^2u\), with \(u\in\F_p^\times\), then
\(R_\xi(t)=\sum_jb_ju^j\in K\).  Hence the canonical trace character
equals one on \(R_\xi(t)\), and its \(r\)-th powers do also for
\(1\le r<p\).  Therefore the fixed-weight Fourier coefficient is
\(\binom na\) for every \(\xi\in U\).  A minor classification cannot
have a uniform power-sum bound below \(n\); at the identity crossing the
PF left side is therefore
\(\log\binom{n+a-1}{a}+o(n)=\Omega(n)\).  A major classification
contributes \(\abs U-1\) times the zero-character mass.  Finally the
identity crossing gives
\((w/2)\log p=(h(\alpha)/4+o(1))n\).  Constancy on all generators says
exactly that \(U\) lies in the annihilator of the effective span.
\end{proof}

\begin{theorem}[Bounded-prefix flatness on smooth cosets]\label{thm:bounded-prefix-equidistribution}
Let \(T\) be a multiplicative coset or a coset with \(o(\abs T)\) planted
deletions over \(\B\).  Let \(g_1,\ldots,g_R\) be the weighted rational
prefix functions of a primitive smooth-domain chart.  If the chart
satisfies the ambient instances of the pointwise minor-arc condition
\textup{(PF)} and the aggregate condition \textup{(MA)}, then for every
target \(z\in\B^R\),
\[
        \abs{\{x\in\Omega_{T,m}:\Psi(x)=z\}}
        \le e^{o(\abs T)}\barN_0.
\]
The same bound holds for every first-match residual
\(\Omc\subseteq\Omega_{T,m}\).
\end{theorem}

\begin{proof}
For a minor character \(\alpha\), the subgroup indicator expands as an
average of multiplicative characters.  Weil's bound for a nonconstant
rational phase gives
\(\abs{\sum_{t\in T}\psi(jR_\alpha(t))}\le
C_0(\Delta_{\rm pole}+1)\sqrt{\abs\B}+\abs P=\Lambda\) for \(1\le j\le m\).  The
cycle-index estimate---the generating-function count of distinct-coordinate
tuples grouped by permutation cycle lengths---used in
\cref{thm:prefix-flatness-power-sum} bounds
the complete minor-character contribution by
\(\binom{\Lambda+m-1}{m}\).  The zero character contributes \(\barN_0\),
and \textup{(MA)} controls the omitted major characters.  Condition
\textup{(PF)} makes the minor contribution at most
\(e^{o(\abs T)}\barN_0\).  Passing to a first-match residual can only decrease
fibers.
\end{proof}

\begin{theorem}[Bounded-prefix flatness on circle domains]\label{thm:circle-prefix-equidistribution}
Let \(T=\chi(gH\cup g^{-1}H)\) be a circle twin-coset \(x\)-domain, with
the branch points \(u=u^{-1}\) of \(\chi\) and any repeated projected
coordinates already isolated as planted exceptions.  For the primitive weighted
prefix functions on a circle or diagonally uniformized circle-code chart,
assume the ambient circle version of the pointwise minor-arc range and
the ambient instance of \textup{(MA)}.  Then every fiber
\(\Psi^{-1}(z)\) of the full profile slice
and every first-match residual fiber
is bounded by \(e^{o(\abs T)}\barN_0\).
\end{theorem}

\begin{proof}
Pull the phase back by \(x=\chi(u)=(u+u^{-1})/2\) to the norm-one torus.
A minor arc becomes a nonconstant rational function on the genus-zero
torus, not of Artin--Schreier form.  Weil's bound on the two torus branches
gives the same estimate as above with
\(2C_0(\Delta_{\rm pole}+1)\sqrt{\abs\B}\).  The rest is identical to
the proof of \cref{thm:bounded-prefix-equidistribution}.  A character
whose pullback factors through a power or dihedral quotient is designated
major unless a certified minor lift still satisfies the displayed Weil
hypotheses; it must then be paid in aggregate by \textup{(MA)} or by a
separately proved quotient recursion.  Such dual factorization alone does
not put the primal witness support in a quotient cell.
\end{proof}

\begin{corollary}[Sidon moment bound from full Fourier control]\label{cor:fourier-sidon-paid-smooth-circle}
On every primitive smooth or circle chart satisfying the ambient
instances of the pointwise minor-arc estimate \textup{(PF)} and the
aggregate estimate \textup{(MA)}, the
Fourier/Sidon-heavy cell has image-normalized logarithmic moment
\(e^{o(nq)}\) for every logarithmic \(q\le(\log n)^A\).
\end{corollary}

\begin{proof}
The full-profile max-fiber theorem gives the ambient estimate
\(f_s\le e^{o(N)}\barN_0\) for every residual fiber.  It therefore
certifies \textup{(FI)}, and
\cref{lem:image-ambient-moment-conversion} gives the image-normalized
Sidon moment.  The estimate \textup{(MA)}, not the name of a first-match
cell, controls the major-character contribution.
\end{proof}

\section{Exact quotient--remainder profiles and the field-drop obstruction}
\label{sec:quotient-obstruction}

The quotient and remainder variables can be separated exactly at the
locator-prefix level.  This removes a spurious factor coming from summing
over marked remainder sets, but it also exposes a genuine scaled quotient-
coefficient-field term which prevents an unrestricted comparison with the
identity profile.

\begin{definition}[Quotient/remainder profile and field drop]
\label{def:quotient-remainder-profile}
Let \(\phi:D\to Q\) be a \(c\)-fold complete-fiber folding map in the
sense of \cref{def:structured-folding}.  A
\emph{complete-fiber-plus-remainder support} is
\[
                  S=\phi^{-1}(E)\sqcup R,
\]
where \(E\subseteq Q\) and
\(R\subseteq D\setminus\phi^{-1}(E)\).  The set \(R\) is the
\emph{support remainder}: it records the selected points in fibers that
are not declared complete.  A \emph{quotient/remainder profile} fixes
\(\phi\), its degree \(c\), the size \(r=\abs R\), the relevant
coefficient field, and any planted remainder data, while \(E\) and the
unfixed part of \(R\) range subject to the displayed disjointness.
A fixed-\(R\) profile is the subprofile in which the remainder support
itself is part of the label.  The adjective \emph{Euclidean} used below
refers to separating the locator factor of degree \(r<c\) from the
composed quotient locator; it does not mean that an arbitrary support
remainder has size less than \(c\).

We call \(\B_\phi\subseteq\B\) a \emph{scaled quotient coefficient
field} when, for some \(\eta\in\B^\times\),
\[
                         Q=\phi(D)\subseteq\eta\B_\phi .
\]
When this field is recorded in a profile, it is taken minimal among the
subfields with this property.
Then the \(j\)-th elementary quotient coefficient lies in
\(\eta^j\B_\phi\).  A \emph{field drop} is the use of a proper such
subfield.  It is a \emph{positive-rate field drop} along a sequence when
\[
       1-\frac{\log\abs{\B_\phi}}{\log\abs\B}
\]
stays bounded away from zero and a positive-density number on the
entropy scale of quotient-prefix slots is live.  The payment saved in
each live slot is the factor \(\abs\B/\abs{\B_\phi}\); this is the term
quantified in \cref{prop:identity-quotient-comparison}.
\end{definition}

\begin{definition}[Balanced quotient core]
\label{def:balanced-quotient-core}
A \emph{split locator} is a monic squarefree polynomial all of whose roots
lie in the evaluation domain \(D\).
After common locator factors, fixed quotient data, and planted remainder
points have been removed, let \(A\) and \(B\) be two monic split
locators of the same degree \(e\).  They form a
\emph{depth-\(w\) balanced core} if they have the same first \(w\)
nonleading coefficients, equivalently
\[
                         \deg(A-B)\le e-w-1.
\]
It is a \emph{balanced quotient core} when the remaining locator factors
are pullbacks through the same folding map \(\phi\).  The word
\emph{balanced} records equality of degree and leading normalization;
the \emph{core} consists of the roots still allowed to move after common
and planted factors are removed.  If the resulting projective locator
family is one-dimensional it is a split pencil
\(Q_0+\lambda Q_1\); in higher dimension this definition supplies no
distinct-ray bound, and the separate ray compiler is required.
\end{definition}

For a monic polynomial
\[
 A(X)=X^d+a_1X^{d-1}+\cdots+a_d
\]
and an integer $0\le u\le d$, write
\(\operatorname{pref}_u(A)=(a_1,\ldots,a_u)\).  Equivalently, if
\(A^\vee(Z):=Z^dA(Z^{-1})\), then
\(\operatorname{pref}_u(A)\) is the coefficient vector of
\(A^\vee(Z)\bmod Z^{u+1}\), with its constant coefficient omitted.

\begin{theorem}[Exact quotient--remainder prefix normal form]
\label{thm:exact-quotient-remainder-normal-form}
Let $\B\subseteq\F$ be finite fields, let $D\subseteq\B$ have $n$
distinct elements, and let $\phi\in\B[X]$ be monic of degree $c\ge2$.
Assume that the restriction
\[
             \phi:D\longrightarrow Q:=\phi(D)
\]
has complete fibers of size $c$; in particular, $N:=|Q|=n/c$.
Fix integers $m,r$ with $0\le r<c$ and $cm+r\le n$.  For
\[
 E\in\binom Qm,
 \qquad
 R\in\binom{D\setminus\phi^{-1}(E)}r,
\]
put
\[
 \begin{split}
 V_E(Y)&:=\prod_{y\in E}(Y-y)
       =Y^m+\sum_{j=1}^m v_j(E)Y^{m-j},\\
 P_R(X)&:=\prod_{x\in R}(X-x)
       =X^r+\sum_{i=1}^r p_i(R)X^{r-i},\\
 L_{E,R}(X)&:=P_R(X)V_E(\phi(X)).
 \end{split}                                      \tag{QR1}\label{eq:qr-locators}
\]
Then $L_{E,R}$ is the monic locator of the support
\(\phi^{-1}(E)\sqcup R\), of degree $A=cm+r$.

Let $0\le w\le A$ and put $d=\lfloor w/c\rfloor$.

\begin{enumerate}[label=\textup{(\roman*)},leftmargin=*]
\item If $w\ge r$, every nonempty fiber of
      \((E,R)\mapsto\operatorname{pref}_w(L_{E,R})\) determines $R$
      uniquely and, after $R$ is recovered, is exactly one quotient-prefix
      fiber.  More precisely, for every prefix value $z$ there are unique
      data $R_z$ and $\eta_z\in\B^d$, whenever the fiber is nonempty, such
      that
      \[
      \begin{split}
       &\{(E,R):\operatorname{pref}_w(L_{E,R})=z\}\\
       &\quad=
       \{(E,R_z): E\in\tbinom{Q\setminus\phi(R_z)}m,
           \ \operatorname{pref}_d(V_E)=\eta_z\}.
      \end{split}                                  \tag{QR2}\label{eq:qr-fiber-exact}
      \]
      Thus no factor counting possible remainder sets occurs in a fixed
      global prefix fiber.

\item If $w<r$, then necessarily $w<c$, and the quotient set $E$ is
      invisible to the first $w$ coefficients.  There is a fixed
      unitriangular bijection $T_{\phi,m,w}:\B^w\to\B^w$ such that
      \[
          \operatorname{pref}_w(L_{E,R})
          =T_{\phi,m,w}(\operatorname{pref}_w(P_R))                 \tag{QR3}
      \]
      for every admissible $(E,R)$.  Consequently, with the convention
      \(\binom uv=0\) for $v>u$,
      \[
      \begin{split}
       &\#\{(E,R):\operatorname{pref}_w(L_{E,R})=z\}\\
       &\quad=
       \sum_{\substack{R\in\binom Dr\\
          \operatorname{pref}_w(P_R)=T_{\phi,m,w}^{-1}(z)}}
          \binom{N-|\phi(R)|}{m}.
      \end{split}                                  \tag{QR4}\label{eq:qr-shallow-remainder}
      \]
      Thus this case is a remainder-prefix problem, multiplied by the
      number of quotient sets still available after $R$ is fixed.
\end{enumerate}
\end{theorem}

\begin{proof}
For a monic degree-$e$ polynomial $A$, set
\(A^\vee(Z)=Z^eA(Z^{-1})\).  Also set
\[
       \Phi(Z):=Z^c\phi(Z^{-1})
       =1+\phi_1Z+\cdots+\phi_cZ^c.
\]
Taking reciprocal polynomials in \eqref{eq:qr-locators} gives the exact
formal-power-series identity
\[
 L_{E,R}^\vee(Z)=P_R^\vee(Z)
 \left(
   \Phi(Z)^m+
   \sum_{j=1}^m v_j(E)Z^{cj}\Phi(Z)^{m-j}
 \right).                                          \tag{QR5}\label{eq:qr-infinity}
\]
All factors on the right have constant coefficient one except for the
displayed scalar coefficients $v_j(E)$.

Suppose first that $w\ge r$.  Since $r<c$, no term containing a $v_j$
occurs below degree $c$.  For $1\le i\le r$, the coefficient of $Z^i$
on the right of \eqref{eq:qr-infinity} is
\(p_i(R)\) plus an expression depending only on
\(p_1(R),\ldots,p_{i-1}(R)\), on $m$, and on $\Phi$.  The first $r$
coefficients therefore recover $p_1(R),\ldots,p_r(R)$ successively.
They recover the monic polynomial $P_R$, hence its root set $R$, uniquely.

We may now divide \(L_{E,R}^\vee\) by the unit
\(P_R^\vee\) modulo $Z^{w+1}$.  In the remaining factor of
\eqref{eq:qr-infinity}, the coefficient of $Z^{cj}$ is
\(v_j(E)\) plus an expression depending only on
\(v_1(E),\ldots,v_{j-1}(E)\), on $m$, and on $\Phi$.  Hence the
coefficients at degrees $c,2c,\ldots,dc$ recover
\(v_1(E),\ldots,v_d(E)\) successively.  Conversely, $P_R$ and these $d$
coefficients determine the right side of \eqref{eq:qr-infinity} modulo
$Z^{w+1}$, because every term containing $v_j$ with $j>d$ starts in degree
$cj>w$.  This proves \eqref{eq:qr-fiber-exact}; admissibility of $E$ is
exactly the condition $E\cap\phi(R)=\varnothing$.

If $w<r<c$, every term involving a $v_j$ starts in degree at least
$c>w$, so modulo $Z^{w+1}$ the bracket in \eqref{eq:qr-infinity} is the
fixed unit $\Phi(Z)^m$.  Multiplication by this unit sends the first $w$
coefficients of $P_R^\vee$ to the first $w$ coefficients of
$L_{E,R}^\vee$ by a unitriangular transformation.  It is independent of
$E$ and invertible.  For a fixed $R$, the admissible $E$ are precisely the
$m$-subsets of $Q\setminus\phi(R)$, of which there are
\(\binom{N-|\phi(R)|}{m}\).  Summing over the indicated remainder-prefix
fiber proves \eqref{eq:qr-shallow-remainder}.
\end{proof}

\begin{proposition}[Complete support factorization and its limitation]
\label{prop:complete-support-factorization}
Under the complete-fiber hypotheses of
\cref{thm:exact-quotient-remainder-normal-form}, every support
\(S\subseteq D\) has the unique decomposition
\[
 E(S)=\{y\in Q:\phi^{-1}(y)\subseteq S\},\qquad
 R(S)=S\setminus\phi^{-1}(E(S)),
 \tag{QR5a}\label{eq:complete-support-factorization}
\]
and
\(Q_S(X)=P_{R(S)}(X)V_{E(S)}(\phi(X))\).  The remainder meets every
nonfull \(\phi\)-fiber in at most \(c-1\) points, but its total size may
be \(\Theta(n)\).  The reciprocal identity \eqref{eq:qr-infinity} remains
valid for this arbitrary remainder.  When \(\abs{R(S)}\ge c\), however,
the quotient coefficients and the remainder coefficients begin to
interlace at degree \(c\), so the triangular recovery theorem above does
not give a comparison without an additional partial-occupancy atlas.
\end{proposition}

\begin{proof}
A fiber is put into \(E(S)\) exactly when all of its points lie in \(S\);
all other selected points form \(R(S)\).  This proves existence,
disjointness, the occupancy bound, and uniqueness.  Multiplying the
corresponding linear factors gives the locator identity.  The reciprocal
calculation used for \eqref{eq:qr-infinity} is purely multiplicative and
does not require \(\deg P_R<c\).  If \(\deg P_R\ge c\), both
\(p_c(R)\) and \(v_1(E)\) occur in degree \(c\), which proves the stated
limitation of the triangular separation.
\end{proof}

\begin{theorem}[Exact partial-occupancy disintegration]
\label{thm:exact-partial-occupancy}
Let \(D=D_0\sqcup X\), where \(\abs X=b\), and let
\(\phi:D_0\twoheadrightarrow Q\) have exactly \(N=\abs Q\) fibers, each
of cardinality \(c\).  For \(S\subseteq D\), put
\[
 X_S=S\cap X,\qquad
 E(S)=\{y\in Q:\phi^{-1}(y)\subseteq S\},\qquad
 R(S)=(S\cap D_0)\setminus\phi^{-1}(E(S)).
\]
Write
\[
 t=\abs{X_S},\qquad m=\abs{E(S)},\qquad
 p=\abs{\phi(R(S))},\qquad r=\abs{R(S)}.
\]
Then this datum is unique, every fiber meeting \(R(S)\) contains between
one and \(c-1\) points of \(R(S)\), and \(\abs S=t+cm+r\).  If
\(\Omega_{t,m,p,r}\) denotes the family of supports with these four
parameters, then
\[
 \abs{\Omega_{t,m,p,r}}
 =\binom bt\binom Np\binom{N-p}m
   [x^r]\bigl((1+x)^c-1-x^c\bigr)^p.                 \tag{PO1}
\]
Consequently the nonempty families \(\Omega_{t,m,p,r}\) partition
\(\binom Da\), and their exact add-back is
\[
 \sum_{\substack{t,m,p,r\\t+cm+r=a}}\abs{\Omega_{t,m,p,r}}
 =\binom{cN+b}{a}.                                   \tag{PO2}
\]
If \(\phi\) is the restriction of a monic degree-\(c\) polynomial and
every fiber \(\phi^{-1}(y)\) is the complete simple root set in \(D_0\)
of \(\phi(X)-y\), then their locators satisfy
\[
             Q_S=Q_{X_S}P_{R(S)}V_{E(S)}(\phi).
\]
\end{theorem}

\begin{proof}
Full fibers, partial fibers, and empty fibers are mutually exclusive, so
the displayed decomposition is canonical.  Choose the \(t\) exceptional
points, the \(p\) partial fibers, and the \(m\) full fibers in
\(\binom bt\binom Np\binom{N-p}m\) ways.  On one partial fiber the
weight enumerator of the allowed nonempty proper subsets is
\((1+x)^c-1-x^c\); hence the coefficient in \textup{(PO1)} counts the
choices having total partial weight \(r\).  Summing \textup{(PO1)} and
extracting the coefficient of \(x^a\) gives
\[
 [x^a](1+x)^b
 \left(1+x^c+\bigl((1+x)^c-1-x^c\bigr)\right)^N
 =[x^a](1+x)^{b+cN}=\binom{b+cN}{a}.
\]
Multiplying the linear factors belonging to the three disjoint parts of
the support proves the locator identity.
\end{proof}

\begin{proposition}[Partial-occupancy Fourier factorization]
\label{prop:partial-occupancy-fourier}
In the setting of \cref{thm:exact-partial-occupancy}, let \(G\) be a
finite abelian group, let \(g:D\to G\), and put
\(\Phi(S)=\sum_{x\in S}g(x)\).  For a character
\(\chi\in\widehat G\) and a regular fiber \(F_y=\phi^{-1}(y)\), define
\[
 A_y(\chi)=\prod_{x\in F_y}\chi(g(x)),\qquad
 P_y(\chi;z)=
 \prod_{x\in F_y}\bigl(1+z\chi(g(x))\bigr)-1-z^cA_y(\chi).
\]
Then
\[
 \sum_{S\in\Omega_{t,m,p,r}}\chi(\Phi(S))
 =[u^tv^m\zeta^pz^r]
 \prod_{x\in X}\bigl(1+u\chi(g(x))\bigr)
 \prod_{y\in Q}
 \bigl(1+vA_y(\chi)+\zeta P_y(\chi;z)\bigr).          \tag{PO3}
\]
Consequently, for every \(s\in G\),
\[
 \abs{\{S\in\Omega_{t,m,p,r}:\Phi(S)=s\}}
 =\frac1{\abs G}\sum_{\chi\in\widehat G}
   \overline{\chi(s)}
   \sum_{S\in\Omega_{t,m,p,r}}\chi(\Phi(S)).         \tag{PO4}
\]
At the trivial character, \textup{(PO3)} is exactly the support count
\textup{(PO1)}.  Thus support add-back is unconditional, while a finite
flatness theorem is exactly a bound for the nontrivial-character sum in
\textup{(PO4)}.
\end{proposition}

\begin{proof}
For a fixed regular fiber the three terms in the second product encode,
respectively, an empty fiber, a full fiber, and a nonempty proper subset.
The variables \(v,\zeta,z\) record the numbers \(m,p,r\), while \(u\)
records the exceptional weight \(t\).  Multiplicativity of \(\chi\) turns
the character of the support sum into the product of its selected-point
characters, proving \textup{(PO3)} by coefficient extraction.
Formula \textup{(PO4)} is character orthogonality on \(G\).
\end{proof}

\begin{remark}[Effective normalization of a partial-occupancy slice]
Fix a nonempty slice \(\Omega_\lambda=\Omega_{t,m,p,r}\), choose
\(S_0\in\Omega_\lambda\), and let
\[
 G_\lambda=
 \left\langle\Phi(S)-\Phi(S_0):S\in\Omega_\lambda\right\rangle
 \le G.
\]
The attained boundary values lie in the affine coset
\(\Phi(S_0)+G_\lambda\), and the effective inversion is
\[
 \abs{\{S\in\Omega_\lambda:\Phi(S)=s\}}
 =\frac1{\abs{G_\lambda}}
   \sum_{\chi\in\widehat{G_\lambda}}
   \overline{\chi(s-\Phi(S_0))}
   \sum_{S\in\Omega_\lambda}
      \chi(\Phi(S)-\Phi(S_0)).                       \tag{PO5}
\]
Every character of \(G_\lambda\) extends to the ambient finite abelian
group, so the inner sum may be evaluated from \textup{(PO3)} after
multiplication by the fixed factor
\(\overline{\chi(\Phi(S_0))}\).  Different extensions agree on support
differences.  Equivalently, ambient characters group by their restriction
to \(G_\lambda\), and the annihilator multiplicity cancels exactly against
the ambient denominator.  This identifies the correct Fourier denominator;
it does not assert that the realized image fills the affine group.
\end{remark}

\begin{remark}[What partial-occupancy add-back does and does not prove]
The exponentially many local proper-subset choices are already contained
inside the coarse slices \(\Omega_{t,m,p,r}\); they incur no separate
profile-count loss in \textup{(PO2)}.  The theorem does not bound the
nontrivial Fourier terms in \textup{(PO4)}, the size of a realized
boundary image, or the projection of a witness cell to distinct slopes.
Those are separate analytic and ray payments.
\end{remark}

\begin{theorem}[Canonical partial-occupancy atlas and bounded-boundary closure]
\label{thm:canonical-partial-occupancy-atlas}
Fix a received line and exact agreement \(a\).  For every nonempty
canonical slice \(\Omega_\lambda=\Omega_{t,m,p,r}\) of
\cref{thm:exact-partial-occupancy}, let \(\Ccal_\lambda\) be the witnesses
whose support lies in \(\Omega_\lambda\), and order these cells
arbitrarily.  Then they form a witness-exhaustive first-match atlas and,
with \(Z_\lambda^\circ\) its first-match slope projections,
\[
 \abs{Z_\lambda^\circ}\le\abs{\Omega_\lambda},
 \qquad
 \sum_\lambda\abs{Z_\lambda^\circ}
 \le\sum_\lambda\abs{\Omega_\lambda}=\binom na.    \tag{PO6}
\]

More generally, give each slice a boundary map
\(\Phi_\lambda:\Omega_\lambda\to G_\lambda^{\rm amb}\) into a finite
abelian group of order \(A_\lambda\), let \(N_\lambda\) be its
active-coordinate count, let \(L_\lambda\) be its realized image size,
and put
\(\bar N_\lambda=\abs{\Omega_\lambda}/L_\lambda\).  The exact finite ray
multiplier and challenge-set bound are
\[
             \abs{Z_\lambda^\circ}
             \le L_\lambda\bar N_\lambda.          \tag{PO7}
\]
\[
 B_{C,\Gamma}^{\rm MCA}(a)
 \le\min\left\{\abs\Gamma,\
       \sup_r\sum_{\lambda\in\Lambda(r;a)}
              L_\lambda\bar N_\lambda\right\}
 \le\min\left\{\abs\Gamma,\binom na\right\}.        \tag{PO7a}
\]
Define the coarse partial-occupancy envelope
\[
 \mathfrak E_n^{\rm PO}(a)
 =1+\sup_r\sum_{\lambda\in\Lambda(r;a)}(1+\bar N_\lambda).
\]
If, uniformly over received lines and the polynomially many nonempty
canonical slices,
\[
                         \log A_\lambda=o(n),       \tag{PO8}
\]
then this atlas has a closed asymptotic profile-envelope compiler:
\[
 B_{C}^{\rm MCA}(a)\le
 e^{o(n)}\mathfrak E_n^{\rm PO}(a).                \tag{PO9}
\]
Condition \textup{(PO8)} discharges atlas exhaustivity, support add-back,
and the direct alternative in \textup{(RC)}, whose exact multiplier is
\(L_\lambda\le A_\lambda\).  If, in addition,
\[
                         \log A_\lambda=o(N_\lambda) \tag{PO10}
\]
on every slice to which effective Fourier terminology is applied, then
the effective inversion \textup{(PO5)} may designate every nontrivial
effective character major: its minor sum is empty and its exact major
multiplier is at most \(A_\lambda-1\).  Thus \textup{(PO10)} verifies the
slice analogues of MI and MA rather than assuming them.  A concrete
specialization is a polynomial-size coefficient field with uniformly
bounded boundary depth and \(N_\lambda=\Theta(n)\) on every live canonical
folding profile.  This is a coarse support-paid envelope, not an
identity-scale or refined quotient-envelope comparison.  It gives a safe
target whenever the literal right side of \textup{(PO7a)} is at most
\(B^*\); a sharp frontier additionally uses the unsafe inequality in
\textup{(SB2)}.  Condition \textup{(PO8)} alone does not locate a target
frontier.
\end{theorem}

\begin{proof}
The support slices partition \(\binom Da\) by \textup{(PO2)}, hence the
corresponding witness cells partition the exact witness incidence.  For
each first-match slope choose one witness support in its cell.  A fixed
noncommon support carries at most one slope, so this choice is injective
and proves the first inequality in \textup{(PO6)}.  Exact add-back gives
the equality there, and \textup{(PO7)} is the same inequality rewritten
using \(\abs{\Omega_\lambda}=L_\lambda\bar N_\lambda\).

The exact challenge bound follows by summing \textup{(PO7)}, taking the
maximum over received lines, and using \textup{(PO2)}.  Since
\(L_\lambda\le A_\lambda\), condition \textup{(PO8)} turns
\textup{(PO7)} into
\(e^{o(n)}\bar N_\lambda\).  There are only polynomially many quadruples
\((t,m,p,r)\), so their sum proves \textup{(PO9)}.  If
\(A_\lambda\le\abs\B^{R_\lambda}\), a polynomial-size \(\B\) and
bounded \(R_\lambda\) imply \textup{(PO8)}.  Finally, the effective
difference group in \textup{(PO5)} has at most \(A_\lambda\) characters.
Under \textup{(PO10)}, designating all of its nontrivial characters major
makes the minor sum empty and bounds the normalized major aggregate by
\(A_\lambda-1=e^{o(N_\lambda)}\).
\end{proof}

\begin{remark}[Exact dimension rounding]
\label{rem:qr-rounding}
If the global locator degree is $A=cm+r$, the list-code dimension is $K$,
and $w=A-K$, then the quotient depth in
\cref{thm:exact-quotient-remainder-normal-form} is
\[
 d=\left\lfloor\frac{A-K}{c}\right\rfloor
   =m-\left\lceil\frac{K-r}{c}\right\rceil.         \tag{QR5b}\label{eq:qr-rounding}
\]
Thus the descended list dimension is exactly
\(K_c=\lceil(K-r)/c\rceil\).  More explicitly, if \(w=cu+s\) with
\(0\le s<c\), then
\[
 \abs{\B_\phi}^{-\floor{w/c}}
 =\abs{\B_\phi}^{s/c}\abs{\B_\phi}^{-w/c}.          \tag{QR5c}
\]
Thus the exact rounding multiplier is \(\abs{\B_\phi}^{s/c}\), lying
between one and \(\abs{\B_\phi}^{(c-1)/c}\).  Retaining the integer
\(\floor{w/c}\) avoids even this comparison; no polynomial or
\(e^{o(n)}\) placeholder is needed in a finite certificate.
\end{remark}

\begin{remark}[Nonmonic maps and Chebyshev maps]
\label{rem:qr-chebyshev}
If the leading coefficient of $\phi$ is $b\ne0$, replace $\phi$ by
$b^{-1}\phi$ and $Q$ by $b^{-1}Q$.  This changes quotient coefficients by
fixed nonzero scalar factors and leaves every support and every fiber
cardinality unchanged.  The theorem therefore applies to the Chebyshev
folding map $T_c$ after this normalization.  On a circle twin-coset domain
on which $T_c$ has complete fibers of size $c$, it gives the same statement
verbatim.  The fixed or ramified endpoints form an exceptional set \(X\)
of cardinality \(b\) and are included exactly by the factor \(\binom bt\)
in \textup{(PO1)}; the crude total exceptional-point multiplier is at most
\(2^b\), with \(b\le2\) for involution-fixed endpoints.  Coincident twin
cosets are not a bounded planted set: they form the separate dihedral
collapse profile of \cref{lem:circle-edge-cases-repaired}.
\end{remark}

For a declared complete-fiber-plus-Euclidean-remainder profile, the exact
normal-form theorem separates algebraic bookkeeping from prefix
equidistribution only when the remainder degree is less than \(c\).  The
exact disintegration of \cref{thm:exact-partial-occupancy} exhausts
arbitrary partial-occupancy supports, but when \(\deg P_R\ge c\) the
quotient and remainder coefficients interlace.  Thus that disintegration
does not by itself prove prefix flatness on the resulting coarse slices.
Let
$\B_\phi\subseteq\B$ be a subfield and suppose
\(Q\subseteq\eta\B_\phi\) for some $\eta\in\B^\times$.  Then
\(v_j(E)\in\eta^j\B_\phi\).  Hence, in case
\textup{(i)}, the natural average scale of the quotient-prefix fiber with
fixed $R$ is
\[
 \overline N_{\phi,R}(w)
 :=\binom{N-|\phi(R)|}{m}
       |\B_\phi|^{-\lfloor w/c\rfloor}.             \tag{QR6}\label{eq:qr-natural-scale}
\]
Pigeonholing proves a fiber of size at least
\(\lceil\overline N_{\phi,R}(w)\rceil\).  An upper bound of the same order
is not a consequence of the normal form: it is precisely the uniform
prefix-flatness, or Sidon/Fourier, input on the quotient row.

\begin{proposition}[Identity--quotient exponent comparison]
\label{prop:identity-quotient-comparison}
Consider a sequence to which
\cref{thm:exact-quotient-remainder-normal-form} applies with fixed
$c\ge2$ and fixed $0\le r<c$.  Suppose
\[
 \frac an\longrightarrow\alpha\in(0,1),\qquad
 \log|\B|=o(n),\qquad
 \log\binom na-w\log|\B|=o(n),                     \tag{QR7}\label{eq:identity-crossing-qr}
\]
where $a=cm+r$, and suppose
\(Q\subseteq\eta\B_c\) for a subfield $\B_c\subseteq\B$.
Write \(\overline N_{c,r}(w)=\overline N_{\phi,R}(w)\) for the natural
scale in \eqref{eq:qr-natural-scale}, with this fixed remainder profile.
Then $w\to\infty$, so $w\ge r$ eventually, and
\[
 \log\overline N_{c,r}(w)
 =\frac1c\log\!\left(\binom na|\B|^{-w}\right)
  +\frac wc\log\frac{|\B|}{|\B_c|}+o(n),           \tag{QR8}\label{eq:qr-comparison-general}
\]
where replacing $N-|\phi(R)|$ by $N$ changes only the $o(n)$ term.
Equivalently, if
\(\lambda_c=\log|\B_c|/\log|\B|\) has a limit, then at the identity
crossing
\[
       \frac1n\log\overline N_{c,r}(w)
       =\frac{h(\alpha)}c(1-\lambda_c)+o(1),         \tag{QR9}\label{eq:qr-comparison-crossing}
\]
using the entropy convention fixed in the Introduction.
Thus the natural average scale of this bounded-degree quotient profile is
comparable with the identity scale at exponential accuracy exactly when
\(\log|\B_c|/\log|\B|=1-o(1)\).  A proper field drop of fixed relative
degree produces a positive exponential quotient term.
\end{proposition}

\begin{proof}
Stirling's formula, $a/n\to\alpha$, and fixed $c,r$ give
\[
 \log\binom{n/c+O(1)}{(a-r)/c}
 =\frac1c\log\binom na+o(n).
\]
Also
\(\lfloor w/c\rfloor\log|\B_c|
=(w/c)\log|\B_c|+o(n)\), because the rounding error is at most
\(\log|\B_c|\le\log|\B|=o(n)\).  This proves
\eqref{eq:qr-comparison-general}.  Since
\(\log\binom na=nh(\alpha)+o(n)\),
\eqref{eq:identity-crossing-qr} gives
\(w\log|\B|=nh(\alpha)+o(n)\).  In particular $w\to\infty$ because
\(\log|\B|=o(n)\).  Substitution in
\eqref{eq:qr-comparison-general} proves
\eqref{eq:qr-comparison-crossing}.
\end{proof}

\begin{remark}[Unbounded quotient degrees]
If $c=c_n\to\infty$, the complete-fiber choice contributes at most
\(2^{n/c}=e^{o(n)}\).  When $w\ge r$, the remainder is recovered exactly
and no further remainder multiplicity occurs.  When $w<r<c$, the exact
formula \eqref{eq:qr-shallow-remainder} reduces the profile to an ordinary
prefix problem for the remainder support, multiplied by $e^{o(n)}$
quotient choices.  Thus only bounded quotient degrees can create a new
positive-rate field-drop term; large degrees either cost $e^{o(n)}$ or
return to an identity/remainder prefix problem.  This does not prove
flatness of the remaining prefix problem.
\end{remark}


We now show that the field-drop term is a real Sidon-heavy obstruction in
the polynomial-field regime defined in \cref{sec:exact-prefix-pole}, rather
than an artifact of the comparison.  A
\emph{complete square-fiber support} is a pullback
\(\{x\in D:x^2\in E\}\) under the two-fold power map; the analogous
complete cube-fiber support is defined using \(x\mapsto x^3\).
We identify a support with its incidence vector in \(\{0,1\}^D\).
Recall that for a support family \(\mathcal A\), its additive energy and
normalized energy are
\[
 E(\mathcal A)
 =\abs{\{(A_1,A_2,A_3,A_4)\in\mathcal A^4:
                 A_1-A_2=A_3-A_4\}},
 \qquad
 \Delta(\mathcal A)=\frac{E(\mathcal A)}{\abs{\mathcal A}^3}.
\]
At the subexponential identity scales used below, an exponentially large
prefix fiber with \(\Delta\le e^{-\sigma n}\) is called a
\emph{Sidon-heavy obstructing fiber}: it contributes to the previously
defined Sidon-heavy branch by carrying large ordinary moment mass while
having exponentially little additive closure.

\begin{theorem}[Polynomial-field Sidon and MCA obstruction]
\label{thm:smooth-quotient-obstruction}
Fix \(0<\alpha<1/2\).  Let \(p\) tend to infinity through odd primes and
put
\[
 n=2(p-1),\qquad \B_0=\F_p,\qquad \B=\F_{p^2}.
 \tag{6.1}\label{eq:polynomial-field-tower}
\]
Let \(H\le\B^\times\) be the subgroup of order \(n\), choose a generator
\(\theta\) of \(\B^\times\), and put \(D=\theta H\).  Choose even integers
\(a=a_p\) with \(a/n\to\alpha\), let \(w\) be the largest even integer not
exceeding \(\log_{\abs\B}\binom na\), and put \(k=a-w-1\).  Then
\[
 \log\abs\B=O(\log n)=o(n),\qquad
 \operatorname{char}\B=p>a,\qquad
 w=\Theta(n/\log n),
 \tag{6.2}\label{eq:polynomial-field-parameters}
\]
and the identity scale satisfies
\[
 1\le \barN_1=\binom na\abs\B^{-w}<\abs\B^2=e^{o(n)}.
 \tag{6.3}\label{eq:polynomial-identity-scale}
\]
For \(E\in\binom{D^2}{a/2}\), put
\(S_E=\{x\in D:x^2\in E\}\).  There is a depth-\(w\) prefix \(z\) for
which the complete-square-fiber subfamily
\[
 \Qcal_z=\{S_E:E\in\tbinom{D^2}{a/2},\ \Phi_w(S_E)=z\}
 \subseteq \Fcal_z:=\Phi_w^{-1}(z),
\]
and the full prefix fiber obeys
\[
 \abs{\Fcal_z}\ge\abs{\Qcal_z}\ge
 \binom{n/2}{a/2}p^{-w/2}
 =\exp\left(\left(\frac14h(\alpha)+o(1)\right)n\right),
 \tag{6.4}\label{eq:polynomial-large-fiber}
\]
with \(h\) as fixed in the Introduction.
For the complete-square profile
\(\Omega_{\rm sq}=\{S_E:E\in\binom{D^2}{a/2}\}\), one also has
\[
 L_{\rm sq}:=\abs{\Phi_w(\Omega_{\rm sq})}\le p^{w/2},
 \qquad
 \frac{\abs{\Omega_{\rm sq}}}{L_{\rm sq}}
 \ge\exp\left(\left(\frac14h(\alpha)+o(1)\right)n\right).
 \tag{6.4$'$}\label{eq:polynomial-image-profile}
\]
Thus the square quotient term is exponentially large in image
normalization as well.

Moreover there is a constant \(\sigma=\sigma(\alpha)>0\) such that
\[
 \max\{\Delta(\Fcal_z),\Delta(\Qcal_z)\}\le e^{-\sigma n}.
 \tag{6.5}\label{eq:polynomial-sidon-energy}
\]
For the full \(a\)-subset slice define the ambient identity-normalized
moment
\[
 \Gamma^{\rm sid,amb}_{r,\sigma}
 =\abs\B^{-w}
   \sum_{\substack{z'\in\B^w\\
          \Delta(\Fcal_{z'})\le e^{-\sigma n}}}
   \left(\frac{\abs{\Fcal_{z'}}}{\barN_1}\right)^r.
\]
For every logarithmic moment order \(r_n\to\infty\),
\[
 \liminf_{n\to\infty}
 \frac{1}{nr_n}\log\Gamma^{\rm sid,amb}_{r_n,\sigma}
 \ge\frac14h(\alpha)>0.
 \tag{6.6}\label{eq:polynomial-sidon-moment}
\]
Thus no ambient identity-normalized Sidon payment can hold on an unrestricted
deep prefix before quotient, subfield, planted, and remainder profiles
have been routed, even when \(\log\abs\B=o(n)\).

Finally, every scalar extension \(\F/\B\) satisfying
\[
 \abs\F-n>k\binom{\abs{\Fcal_z}}2
 \tag{6.7}\label{eq:polynomial-extension-separation}
\]
has a received line for \(C=\RS_\F(D,k)\) with at least
\[
 B_C^{\rm MCA}(a)\ge\abs{\Fcal_z}
 \ge\exp\left(\left(\frac14h(\alpha)+o(1)\right)n\right)
 \tag{6.8}\label{eq:polynomial-mca-obstruction}
\]
distinct bad slopes.  Such extensions exist with degree
\([\F:\B]=O(n/\log n)\).
\end{theorem}

\begin{proof}
The cyclic group \(\B^\times\) has order \((p-1)(p+1)\), divisible by
\(n=2(p-1)\).  Squaring \(H\) gives its subgroup of order \(p-1\), which
is the unique such subgroup and hence equals \(\B_0^\times\).  Therefore
\[
 D^2=\theta^2\B_0^\times,
 \tag{6.9}\label{eq:square-quotient-domain}
\]
and the square map \(D\to D^2\) is exactly two-to-one.  The two roots of
\(1+X^2\) have order four and lie in \(H\), while \(D=\theta H\ne H\);
thus \(D\) avoids the stereographic poles.  Since
\(a/n\to\alpha<1/2\), one has \(a<p\) eventually.  Stirling's formula and
the definition of \(w\) give
\(w=nh(\alpha)/(2\log p)+o(n/\log n)\); the even rounding leaves a
remainder in \([0,2)\), proving
\eqref{eq:polynomial-field-parameters} and
\eqref{eq:polynomial-identity-scale}.

For \(E\in\binom{D^2}{a/2}\), let \(S_E\) be its full square preimage.
Then
\[
 Q_{S_E}(X)=\prod_{y\in E}(X^2-y)
 =X^a-e_1(E)X^{a-2}+e_2(E)X^{a-4}-\cdots .
 \tag{6.10}\label{eq:polynomial-even-locator}
\]
Among the first \(w\) nonleading coefficients, the odd ones vanish and
the coefficient in gap \(2j\) lies in
\(\theta^{2j}\B_0\).  Thus at most \(p^{w/2}\) prefixes occur.  Pigeonholing
the \(\binom{n/2}{a/2}\) supports gives a value \(z\) with
\(\abs{\Qcal_z}\ge\binom{n/2}{a/2}p^{-w/2}\).  Since
\(\Qcal_z\subseteq\Fcal_z\), this proves
\eqref{eq:polynomial-large-fiber}.  Also
\(w\log p=\tfrac12nh(\alpha)+o(n)\), so Stirling gives the stated
exponent.

By \cref{prop:high-energy-impossible}, there is an absolute \(A>0\) such
that
\[
 E(F)\ge\abs F^3/K,\quad F\subseteq\{0,1\}^D
 \quad\Longrightarrow\quad \abs F\le K^A.
 \tag{6.11}\label{eq:quantitative-cube-energy}
\]
Choose \(0<\sigma<h(\alpha)/(8A)\).  If either
\(\Fcal_z\) or \(\Qcal_z\) had normalized energy greater than
\(e^{-\sigma n}\), apply \eqref{eq:quantitative-cube-energy} to that
fiber with \(K=e^{\sigma n}\).  It gives
\(\abs F\le e^{A\sigma n}\), contradicting
\eqref{eq:polynomial-large-fiber} for large \(n\).  This proves
\eqref{eq:polynomial-sidon-energy}.

The formal prefix target space \(\B^w\) has size \(\abs\B^w\).  Retaining only the
summand of \(\Fcal_z\) in the Sidon-heavy moment and using
\(\log\barN_1=o(n)\) gives
\[
 \log\Gamma^{\rm sid,amb}_{r_n,\sigma}
 \ge-w\log\abs\B+
 r_n\bigl(\log\abs{\Fcal_z}-\log\barN_1\bigr)
 =-nh(\alpha)+o(n)
  +r_n\left(\frac14h(\alpha)n+o(n)\right).
\]
Division by \(nr_n\) proves \eqref{eq:polynomial-sidon-moment}.
Here the moment really is the full unrestricted prefix moment: because
\(w<a<p\), Newton's identities give an invertible triangular change of
coordinates between the first \(w\) locator coefficients and the first
\(w\) power sums.  Hence their fibers, and therefore their normalized
Sidon moments, coincide.

Finally \(\cref{thm:prefix-to-line-hardness}\), with
\(m=a\), applies directly to \(\Fcal_z\) and gives
\eqref{eq:polynomial-mca-obstruction}.  Since
\(\log\abs{\Fcal_z}=O(n)\) and \(\log\abs\B=\Theta(\log n)\), an extension
satisfying \eqref{eq:polynomial-extension-separation} can be chosen of
degree \(O(n/\log n)\).
\end{proof}

\begin{remark}[Normalization of the square obstruction]
\label{rem:odd-ambient-image}
Equation \eqref{eq:polynomial-sidon-moment} is an ambient
identity-normalized statement and does not use an unproved reverse
ambient-to-image transfer.  The corresponding image-normalized
profile-envelope obstruction is the independent estimate
\eqref{eq:polynomial-image-profile}; the locator-prefix fiber itself
supplies the MCA lower construction.
\end{remark}

\begin{theorem}[Characteristic-two cubic quotient obstruction]
\label{thm:binary-cubic-obstruction}
Fix \(0<\alpha<1/2\).  Let \(s\to\infty\) through odd integers, put
\[
 q=2^s,\qquad \B_0=\F_q,\qquad \B=\F_{q^2},\qquad n=3(q-1),
 \tag{6.11a}\label{eq:binary-cubic-tower}
\]
and let \(H\le\B^\times\) be the subgroup of order \(n\).  Choose a
generator \(\theta\) of \(\B^\times\) and put \(D=\theta H\).  Choose
multiples of three \(a=a_s\) with \(a/n\to\alpha\), let \(w\) be the
largest multiple of three not exceeding
\(\log_{\abs\B}\binom na\), and put \(k=a-w-1\).  Then
\[
 \log\abs\B=O(\log n)=o(n),\qquad
 w=\Theta(n/\log n),\qquad
 1\le\bar N_1:=\binom na\abs\B^{-w}<\abs\B^3=e^{o(n)}.
 \tag{6.11b}\label{eq:binary-identity-scale}
\]

For an \(a\)-set \(S\subseteq D\), let
\[
 \Phi_w^{\rm el}(S)=(c_1(S),\ldots,c_w(S)),\qquad
 \Phi_w^{\rm ps}(S)=(p_1(S),\ldots,p_w(S)),
\]
where \(Q_S=X^a+c_1(S)X^{a-1}+\cdots\) and
\(p_j(S)=\sum_{x\in S}x^j\).  Put
\[
 S_E:=\{x\in D:x^3\in E\},\qquad
 \Omega_{\rm cube}=\{S_E:E\in\tbinom{D^3}{a/3}\}.
\]
For each
\(\star\in\{\mathrm{el},\mathrm{ps}\}\) there is a prefix
\(z_\star\in\B^w\) such that
\[
 \begin{split}
 \mathcal F^\star_{z_\star}
   &:=\{S\in\tbinom Da:\Phi_w^\star(S)=z_\star\},\\
 \mathcal Q^\star_{z_\star}
   &:=\mathcal F^\star_{z_\star}\cap\Omega_{\rm cube},\\
 \abs{\mathcal F^\star_{z_\star}}
   &\ge\abs{\mathcal Q^\star_{z_\star}}
    \ge \binom{n/3}{a/3}q^{-w/3}
    =\exp\left(\left(\frac16h(\alpha)+o(1)\right)n\right).
 \end{split}
 \tag{6.11c}\label{eq:binary-cubic-large-fibers}
\]
These are two separately proved assertions; no Newton-coordinate
equivalence between the elementary and power-sum fibers is asserted.
For either coordinate system
\[
 L_{\rm cube}^\star
   :=\abs{\Phi_w^\star(\Omega_{\rm cube})}\le q^{w/3},
 \qquad
 \frac{\abs{\Omega_{\rm cube}}}{L_{\rm cube}^\star}
 \ge\exp\left(\left(\frac16h(\alpha)+o(1)\right)n\right).
 \tag{6.11c$'$}\label{eq:binary-cubic-image-profile}
\]
Thus the image-normalized cubic quotient profile is exponentially larger
than the ambient identity term; this conclusion does not depend on
ambient-to-image moment transfer.

There is \(\sigma=\sigma(\alpha)>0\) for which
\[
 \max\{\Delta(\mathcal F^\star_{z_\star}),
        \Delta(\mathcal Q^\star_{z_\star})\}\le e^{-\sigma n}
 \quad(\star\in\{\mathrm{el},\mathrm{ps}\}).
 \tag{6.11d}\label{eq:binary-cubic-sidon-energy}
\]
More explicitly, define the ambient identity-normalized Sidon moment by
\[
 \mathcal G^\star_{r,\sigma}
 =\abs\B^{-w}
   \sum_{\substack{z\in\B^w\\
          \Delta(\mathcal F^\star_z)\le e^{-\sigma n}}}
   \left(\frac{\abs{\mathcal F^\star_z}}{\bar N_1}\right)^r .
 \tag{6.11e}\label{eq:binary-cubic-sidon-moment-definition}
\]
For every logarithmic order \(r=r_n\to\infty\),
\[
 \liminf_{n\to\infty}\frac1{nr}
       \log\mathcal G^\star_{r,\sigma}
 \ge\frac16h(\alpha)
 \quad(\star\in\{\mathrm{el},\mathrm{ps}\}).
 \tag{6.11f}\label{eq:binary-cubic-sidon-moment}
\]

Finally, for the elementary-prefix fiber, every finite extension
\(\F/\B\) satisfying
\[
 \abs\F-n>k\binom{\abs{\mathcal F^{\rm el}_{z_{\rm el}}}}2
 \tag{6.11g}\label{eq:binary-extension-separation}
\]
has a received line for \(C=\RS_\F(D,k)\) with
\[
 B_C^{\rm MCA}(a)\ge
 \abs{\mathcal F^{\rm el}_{z_{\rm el}}}
 \ge\exp\left(\left(\frac16h(\alpha)+o(1)\right)n\right).
 \tag{6.11h}\label{eq:binary-cubic-mca-obstruction}
\]
Such extensions exist with degree \([\F:\B]=O(n/\log n)\).  This last
claim uses only the elementary locator prefix and does not transfer to
the degenerate standard characteristic-two circle.
\end{theorem}

\begin{proof}
Since \(s\) is odd, \(q\equiv-1\pmod3\), and hence
\(3(q-1)\mid(q^2-1)\).  The group \(\B^\times\) is cyclic, so \(H\)
exists.  Cubing \(H\) gives the unique subgroup of order \(q-1\), namely
\(\B_0^\times\), and the kernel of cubing on \(H\) has order three.
Consequently
\[
 D^3=\theta^3\B_0^\times,
 \qquad x\longmapsto x^3:D\longrightarrow D^3
 \quad\text{is exactly three-to-one}.
 \tag{6.11i}\label{eq:binary-cube-domain}
\]
Stirling's formula gives
\(\log_{\abs\B}\binom na=nh(\alpha)/(2\log q)+o(n/\log n)\).
The rounding error in \(w\) is less than three, which proves
\eqref{eq:binary-identity-scale}; it also gives \(w<a-1\) for all large
\(n\).

For \(E\in\binom{D^3}{a/3}\), let \(S_E\) be its complete inverse image
under cubing.  Its locator is
\[
 Q_{S_E}(X)=\prod_{y\in E}(X^3-y)
 =X^a+c_3(E)X^{a-3}+c_6(E)X^{a-6}+\cdots .
 \tag{6.11j}\label{eq:binary-cubic-locator}
\]
The coefficient in gap \(3j\) belongs to the scalar copy
\(\theta^{3j}\B_0\), and all other coefficients vanish.  Hence the
complete-cube supports assume at most \(q^{w/3}\) elementary prefixes.

The same count for power sums is direct rather than an application of
Newton's identities.  If \(\mu_3\le H\) is the group of cube roots of
unity, then, in characteristic two,
\[
 p_j(S_E)=\sum_{y\in E}\sum_{x^3=y}x^j
 =\begin{cases}
 0,&3\nmid j,\\
 \displaystyle\sum_{y\in E}y^{j/3},&3\mid j,
 \end{cases}
 \tag{6.11k}\label{eq:binary-cubic-power-sums}
\]
because \(\sum_{\zeta\in\mu_3}\zeta^j\) is zero unless \(3\mid j\),
and is \(3=1\) in \(\B\) otherwise.  The coordinate in position \(3j\)
lies in \(\theta^{3j}\B_0\), so there are at most \(q^{w/3}\)
power-sum prefixes as well.  Pigeonholing
\(\binom{n/3}{a/3}\) complete-cube supports proves the displayed lower
bound for \(\mathcal Q^\star_{z_\star}\), and hence for
\(\mathcal F^\star_{z_\star}\), in each coordinate system.
Moreover
\[
 \log\binom{n/3}{a/3}-\frac w3\log q
 =\frac n3h(\alpha)-\frac n6h(\alpha)+o(n),
\]
which proves its exponent.

By \cref{prop:high-energy-impossible}, an absolute constant \(A>0\)
satisfies
\[
 E(F)\ge\abs F^3/K,\quad F\subseteq\{0,1\}^D
 \quad\Longrightarrow\quad \abs F\le K^A.
\]
Choose \(0<\sigma<h(\alpha)/(12A)\).  If any of the full or quotient
fibers in \eqref{eq:binary-cubic-large-fibers} had
\(\Delta>e^{-\sigma n}\), the preceding implication with
\(K=e^{\sigma n}\) would contradict its exponential lower bound.
This proves \eqref{eq:binary-cubic-sidon-energy}.  Retaining its one
summand in \eqref{eq:binary-cubic-sidon-moment-definition} gives
\[
 \log\mathcal G^\star_{r,\sigma}
 \ge-w\log\abs\B+
 r\bigl(\log\abs{\mathcal F^\star_{z_\star}}-\log\bar N_1\bigr).
\]
Here \(w\log\abs\B=nh(\alpha)+o(n)\) and
\(\log\bar N_1=o(n)\).  Divide by \(nr\) and use \(r\to\infty\) to obtain
\eqref{eq:binary-cubic-sidon-moment}.

Finally \cref{thm:prefix-to-line-hardness}, applied to the elementary
fiber with \(m=a\), gives \eqref{eq:binary-cubic-mca-obstruction}.
Because its cardinality has logarithm \(O(n)\), while
\(\log\abs\B=\Theta(\log n)\), an extension satisfying
\eqref{eq:binary-extension-separation} may be chosen with degree
\(O(n/\log n)\).  Neither this step nor
\cref{prop:stereographic-circle-equivalence} identifies the doubled
characteristic-two affine circle with the present multiplicative row.
\end{proof}

\begin{remark}[Ambient moment versus cubic profile scale]
\label{rem:binary-ambient-image}
The moments in \eqref{eq:binary-cubic-sidon-moment-definition} are
ambient identity-normalized statements on the full \(a\)-subset slice.
They do not, without \textup{(FI)}, assert failure of an
image-normalized primitive payment.  In particular,
\(p_{2j}=p_j^2\) in characteristic two, so the power-sum map has a
Frobenius image collapse.  The robust profile-envelope obstruction is
\eqref{eq:binary-cubic-image-profile}, and the exact MCA construction uses
the elementary locator prefix only.
\end{remark}

\begin{definition}[Aperiodic support]
\label{def:aperiodic-support}
For \(S\subseteq D=\theta H\), put
\(\operatorname{Stab}_H(S)=\{h\in H:hS=S\}\).  The support is
\emph{aperiodic} if \(\operatorname{Stab}_H(S)=\{1\}\).  This removes
exact multiplicative periodicity; it does not by itself remove a planted
quotient/remainder explanation.
\end{definition}

\begin{proposition}[Aperiodic planted Sidon version]
\label{prop:aperiodic-planted-obstruction}
The obstruction in \cref{thm:smooth-quotient-obstruction} persists after
requiring every support to have trivial stabilizer under multiplication by
\(H\).  More precisely, fix \(x_0\in D\), omit its square fiber, and use
\[
 S_E=\{x_0\}\sqcup\{x\in D:x^2\in E\},
 \qquad
 E\in\binom{D^2\setminus\{x_0^2\}}{a/2}.
 \tag{6.12}\label{eq:aperiodic-planted-family}
\]
At total support size \(a+1\) and depth \(w\), some prefix fiber in this
aperiodic residual still has size
\(\exp((h(\alpha)/4+o(1))n)\), has normalized energy at most
\(e^{-\sigma n}\) after reducing \(\sigma\) if necessary, and has identity
scale \(e^{o(n)}\).  Every one of its supports has trivial \(H\)-stabilizer.
\end{proposition}

\begin{proof}
The locator is
\((X-x_0)\prod_{y\in E}(X^2-y)\).  Its first \(w\) coefficients are
determined by at most \(w/2\) elements of scalar copies of \(\B_0\);
therefore some prefix \(z\) contains an exponential constructed subfamily
\(\Pcal_z\) of the form \eqref{eq:aperiodic-planted-family}, with only
\(e^{o(n)}\) changes caused by deleting one square fiber and increasing
the support size by one.  Let
\[
 \Fcal_z^{\rm ap}=\{S\in\Phi_w^{-1}(z):
                    \operatorname{Stab}_H(S)=\{1\}\}.
\]
Every member of \(\Pcal_z\) has exactly one square fiber on which the involution
\(x\mapsto-x\) is incomplete, namely \(\{x_0,-x_0\}\).  Multiplication by
\(h\in H\) commutes with this involution.  If \(hS_E=S_E\), it must
preserve that unique incomplete fiber, so \(h=\pm1\); the value \(-1\)
does not preserve \(S_E\), because \(-x_0\notin S_E\).  Hence \(h=1\).
Thus \(\Pcal_z\subseteq\Fcal_z^{\rm ap}\), so the latter is exponential.
Applying \cref{prop:high-energy-impossible} to
\(\Fcal_z^{\rm ap}\subseteq\{0,1\}^D\) gives its stated low energy.
The identity scale changes by only a polynomial factor.
The same one-summand calculation as in
\eqref{eq:polynomial-sidon-moment} therefore gives
\(\exp(\Omega(nr_n))\) Sidon moment for the stabilizer-deleted residual
at every logarithmic order \(r_n\to\infty\).
\end{proof}

The theorem disproves two proposed unrestricted steps at once.  The
identity-scale quotient comparison fails by
\eqref{eq:qr-comparison-crossing}, and the large exceptional fiber is
Sidon-heavy rather than an approximate group.  Exact stabilizer deletion
does not restore the claim: \cref{prop:aperiodic-planted-obstruction} is a
planted quotient/remainder profile and must be charged at its own scale.
The remaining possible Sidon theorem is therefore a statement only about
the primitive residual after all such profiles have been routed.

\begin{proposition}[Stereographic circle equivalence]
\label{prop:stereographic-circle-equivalence}
Let \(\F\) have odd characteristic, let
\(\mathcal C:x^2+y^2=1\), and let \(E\subseteq\mathcal C(\F)\) avoid
\((-1,0)\).  Put \(t=y/(1+x)\), \(T=t(E)\), and let
\(\mathcal C_u(E)\) be the evaluation code of the degree-at-most-\(u\)
filtered residue-class space in
\(\F[x,y]/(x^2+y^2-1)\).  Then \(\mathcal C_u(E)\) is diagonally
equivalent to \(\RS_\F(T,2u+1)\).  The equivalence preserves every list,
CA, and support-wise MCA bad-slope set.
\end{proposition}

\begin{proof}
The inverse parametrization is
\(x=(1-t^2)/(1+t^2)\), \(y=2t/(1+t^2)\).  On \(E\),
\(1+t^2=2/(1+x)\ne0\).  Modulo the circle equation, the
filtered degree-\(u\) space is
\[
 \{f_0(x)+yf_1(x):\deg f_0\le u,\ \deg f_1\le u-1\},
\]
with the second summand absent at \(u=0\).  It has dimension \(2u+1\).
After substitution and multiplication by \((1+t^2)^u\), every element
becomes a polynomial of degree at most \(2u\).  The induced map of filtered
function spaces is injective before evaluation, because stereographic
substitution is an isomorphism of the localized coordinate rings.  The
dimensions agree, so it is an isomorphism.  Coordinatewise multiplication
by the nonzero values \((1+t^2)^u\) preserves equality on each support and
commutes with affine combinations.
\end{proof}

\begin{corollary}[Circle quotient obstruction]
\label{cor:circle-quotient-obstruction}
In \cref{thm:smooth-quotient-obstruction}, put
\(u=(k-1)/2=(a-w-2)/2\) and let \(E\subseteq\mathcal C(\F)\) be the inverse
stereographic image of \(D\).  Then the standard bivariate circle code
\(\mathcal C_u(E)\) has the same bad-slope sets and satisfies
\eqref{eq:polynomial-mca-obstruction}.  In particular the polynomial-field
quotient obstruction transfers without loss to the circle.
\end{corollary}

\begin{proof}
The integers \(a,w\) are even, so \(k=a-w-1=2u+1\).  The coset \(D\) avoids the
zeros of \(1+t^2\), and
\cref{prop:stereographic-circle-equivalence} transfers the received line
coordinate by coordinate.
\end{proof}

The example is the first nontrivial member of a general quotient envelope.

\begin{proposition}[Necessary power-quotient envelope]
\label{prop:necessary-quotient-envelope}
Let \(D=\theta H\), \(\abs H=n\), write
\(D^{(c)}=\{x^c:x\in D\}\), let \(c\) divide both \(n\) and \(m\), and let
\(0\le w\le m\), and suppose the
power quotient \(D^{(c)}\) is contained in a scalar copy of a subfield
\(\B_c\).  For global prefix depth \(w\), some prefix fiber contains at
least
\[
 \left\lceil
 \binom{n/c}{m/c}\abs{\B_c}^{-\floor{w/c}}
 \right\rceil
 \tag{6.13}\label{eq:necessary-quotient-envelope}
\]
complete-fiber supports.
\end{proposition}

\begin{proof}
The locator of the preimage of \(E\subseteq D^{(c)}\) is
\(\prod_{y\in E}(X^c-y)\).  Among its first \(w\) nonleading
coefficients, only the gaps \(c,2c,\ldots,\floor{w/c}c\) can be nonzero;
they are the first \(\floor{w/c}\) elementary symmetric functions of
\(E\), up to fixed scalar factors.  Pigeonhole the
\(\binom{n/c}{m/c}\) quotient supports over at most
\(\abs{\B_c}^{\floor{w/c}}\) prefixes.
\end{proof}

When \(m=cM\) and \(K=m-w\), the descended prefix depth is
\[
 M-\ceil{K/c}=\floor{w/c}.
 \tag{6.14}\label{eq:quotient-rounding}
\]
Thus an \(O(1)\) rounding difference costs
\(O(\log\abs{\B_c})\) bits and is subexponential when every relevant
coefficient field has logarithmic size \(o(n)\).  The positive-rate effect
in \cref{thm:smooth-quotient-obstruction} instead comes from the field
drop \(\abs\B/\abs{\B_c}\) repeated across \(\floor{w/c}\) live slots,
as quantified in \eqref{eq:qr-comparison-general}.  Remainder supports and
Chebyshev quotients contribute analogous profile terms and must be retained
in \(\mathfrak E_n\).


\section{Algebraic profile criteria and balanced-core scope}
\label{sec:ledger-closure}

Constructibility identifies where a witness lives; payment is a separate
enumerative statement about the image in slope space.  This section records
the algebraic facts that are unconditional and then states the precise
criteria required by the upper compiler.  In particular, quotient
classification does not imply payment at the identity scale, and a
codimension-one support locus does not by itself give the required number of
independent prefix equations.

Here a \emph{balanced-core chart} means a parameter family of equal-degree
monic residual locator pairs with a common prefix, obtained after all
common factors, quotient pullbacks, planted blocks, and already certified
rank degeneracies have been removed, whose remaining moving coefficient
space has projective dimension greater than one.  Thus
``balanced core'' is a description of what remains to be counted, not a
smallness assertion.  A \emph{ray bound} counts the distinct projective
line directions---equivalently, the distinct finite slopes in the chosen
affine chart---in the slope projection of that residual incidence.  The
upper \emph{compiler} is the collection of proved implications that turns
profile-wise support, fiber, and ray bounds into a bound for the MCA
numerator.  A \emph{saturation profile} records the projection that
identifies raw support/explanation witnesses having the same residual
projective direction; paying it requires either a lower bound on the fibers
of that projection or a direct bound on its image.

\begin{theorem}[Conditional algebraic-ledger criterion]
\label{thm:smooth-circle-ledger-closure}
Suppose the quotient profiles are budgeted at their natural terms in
\(\mathfrak E_n\); the planted, tangent, extension, rank-pivot, and
saturation profiles satisfy the direct projection criteria below; and every
remaining balanced-core chart is either covered by the one-pencil theorem
or satisfies the ray compiler \textup{(RC)}.  If the number of profiles is
\(e^{o(n)}\), then the algebraic part of the first-match ledger contributes
at most \(e^{o(n)}\mathfrak E_n\).
\end{theorem}

\begin{proof}
Each criterion below supplies a bound for a set of distinct slopes, not only
for its support incidence.  The profile envelope contains every quotient
and coefficient-field term.  The first-match sets are disjoint after
ordering, and there are \(e^{o(n)}\) of them, so their certified bounds sum
to \(e^{o(n)}\mathfrak E_n\).
\end{proof}

\begin{definition}[Periodicity scale]\label{def:periodicity-scale}
Let a finite group \(G\) act on \(D\).  For \(S\subseteq D\), the periodicity scale \(c(S)\) is the order of the stabilizer of \(S\) in the quotient tower used by the row.  In a multiplicative coset \(D=\theta H\), this is the largest \(c\mid\abs H\) such that \(S\) is a union of cosets of the order-\(c\) subgroup.  In a circle \(x\)-domain it is the largest Chebyshev scale for which \(S\) is a union of complete \(T_c\)-fibers.
\end{definition}

\begin{proposition}[Stabilizer exhaustion and quotient budget]
\label{prop:stabilizer-payment}
In a smooth multiplicative or circle row, every support with \(c(S)>1\)
lies in a power or Chebyshev quotient cell.  A scale-\(c\) complete-fiber
cell is budgeted at its natural descended prefix term, including the field
of definition and the rounding depth \(\floor{w/c}\).  Summing these terms
over the quotient profiles gives the quotient part of
\(\mathfrak E_n\).
\end{proposition}

\begin{proof}
For multiplicative cosets, the subgroup lattice of \(H\) partitions
supports by stabilizer.  If \(c(S)=c\), then \(S\) is the pullback of a
support \(\bar S\) on \(D/H_c\).  If the received and explaining data are
constant on the same fibers, invariant quotient descent applies.  Otherwise the
noninvariant residues require a separate first-match chart and direct
certificate.

The number of divisors of \(\abs H\) is subexponential.  The locator of a
complete-fiber support is lacunary, so its live depth and coefficient field
are those in \cref{prop:necessary-quotient-envelope}.  That proposition is
a lower bound and shows why the quotient term cannot be replaced by the
identity term.  The corresponding upper budget is therefore recorded as a
separate profile term.  The Chebyshev case is identical after lifting to the
torus and replacing \(X^c\) by \(T_c\).  Summation gives precisely the
quotient portion of \(\mathfrak E_n\), without asserting an identity-scale
comparison.
\end{proof}

\begin{proposition}[Exact invariant quotient descent]\label{prop:quotient-descent}
Let \(D=\theta H\), let \(c\mid\abs H\), and suppose
\(h\in\F[X]\), \(\deg h<n\), satisfies \(h(\zeta x)=h(x)\) for every
\(x\in D\) and every \(\zeta\in\mu_c\).  Then
\(h(X)=\bar h(X^c)\) with
\(\deg\bar h=\floor{\deg h/c}\).  The analogous statement on a circle
twin coset holds for an invariant symmetric Laurent polynomial and the
Chebyshev quotient.
\end{proposition}

\begin{proof}
Write \(h=\sum_{j=0}^{c-1}X^jh_j(X^c)\).  For \(\zeta\in\mu_c\), the
difference \(h(\zeta X)-h(X)\) has degree less than \(n\) and vanishes at
all \(n\) points of \(D\), so it is zero.  Orthogonality of the characters
\(\zeta\mapsto\zeta^j\) forces \(h_j=0\) for \(j>0\).  On the torus the
same character decomposition is compatible with the involution
\(z\mapsto z^{-1}\), and invariant Laurent monomials descend to Chebyshev
polynomials.  Constancy only on a selected witness support is weaker and
requires a separate quotient-cell certificate; it is not asserted here.
\end{proof}

\begin{proposition}[Deep RC on invariant quotient cells]
\label{prop:deep-invariant-quotient-ray}
Let \(\pi:D\to\bar D\) be a complete uniform power or Chebyshev folding,
and suppose a realized witness cell consists entirely of pullbacks of
witnesses for a quotient Reed--Solomon row
\(\bar C=\RS_\F(\bar D,\bar k)\): the support, both received words, and
the explaining polynomial descend through \(\pi\).  Let the quotient
agreement be \(\bar a\), put \(\bar n=\abs{\bar D}\) and
\(\bar t=\bar n-\bar a\), and assume
\(3\bar t\le\bar n-\bar k\).  Then the actual slope image of the cell,
restricted to \(\Gamma\subseteq\F\), has size at most
\[
                         \min\{\abs\Gamma,\bar t+1\}. \tag{RC$_{\rm quot}$}
\]
If the cell contains every pullback of the quotient tangent construction,
this bound is attained in the row maximum.
\end{proposition}

\begin{proof}
By the pullback hypothesis, every slope in the realized cell maps to a
noncommon quotient witness with the same slope and quotient agreement
\(\bar a\).  Thus the cell's slope image injects into---and need not
equal---the quotient bad-slope set \(Z_{\bar a}(\bar r)\cap\Gamma\).
Invariant quotient descent divides the complete-fiber support size by the
folding degree.  Since
\(3(\bar n-\bar a)\le\bar n-\bar k\),
\cref{cor:exact-deep-numerator} gives
\[
 \abs{Z_{\bar a}(\bar r)\cap\Gamma}
 \le\min\{\abs\Gamma,\bar t+1\}.
\]
If, for some quotient received line, the cell contains every pullback used
by the quotient tangent construction, those slopes remain noncommon
pullback witnesses, so equality holds for that cell and its cellwise row
maximum.
\end{proof}

\begin{proposition}[Planted-block payment criterion]\label{prop:planted-payment}
Let the allowed planted divisors of size \(\ell\) form a family of
\(e^{o(n)}\) candidates, and suppose the residual prefix fiber for each
candidate has size at most its profile term times \(e^{o(n)}\).  Then the
planted cells are paid at the sum of those terms.
\end{proposition}

\begin{proof}
After fixing \(P\), the residual support is an \((a-\ell)\)-subset of
\(D\setminus P\), and multiplication by \(Q_P\) gives a triangular shift
of the prefix target.  The asserted residual bound is therefore a direct
fiber estimate.  Summing it over the stipulated \(e^{o(n)}\) algebraic
candidates proves the criterion.  Arbitrary choices among
\(\binom n\ell\) subsets are not planted profiles.
\end{proof}

\begin{proposition}[Tangent and determinantal payment criterion]
\label{prop:tangent-payment}
The \emph{deep regime} is the high-agreement range
\(3(n-a)\le n-k\); in that range the deep cell has the exact bound of
\cref{thm:deep-regime-upper}.  A
tangent, common-line, or rank-drop chart is paid if a nonzero collection of
independent pivot minors cuts the \emph{slope image}---the set of slopes
actually attained by that chart---by the number of
prefix equations charged to that chart, with degree and chart count
\(e^{o(n)}\).
\end{proposition}

\begin{proof}
The deep assertion was proved exactly in \cref{sec:syndrome-exact}.  For a
determinantal chart, the nonzero independent minors give the displayed
codimension and a finite-field point bound of the expected order.  The
criterion deliberately requires independence and controls the projection
to slopes.  A matrix such as \(X I_r\) can lose rank \(r\) on a
codimension-one locus, so rank loss \(r\) alone does not certify codimension
\(r\).
\end{proof}

Here \emph{restriction of scalars} means choosing a basis of a bounded
extension field over its recorded coefficient subfield and replacing every
extension-valued variable by its tuple of coordinates over that smaller
field, which we call the base field in this proposition.

\begin{proposition}[Extension-cell criterion]\label{prop:extension-payment}
Restriction of scalars converts a bounded-degree extension cell to a
bounded number of base-field incidence charts.  Such a cell is paid only
after the resulting charts satisfy a direct slope-image bound at their
coefficient-field profile scales.
\end{proposition}

\begin{proof}
Choose a base-field basis.  Coefficient variables become bounded tuples and
the witness equations become bounded systems over the base field; the
operation does not itself estimate their slope projection.  Once the
displayed direct estimate is provided, bounded extension degree and
subexponential profile count preserve it under summation.
\end{proof}

\begin{proposition}[Rank-pivot criterion]\label{prop:rank-pivot-payment}
A differential-locator rank-pivot cell is paid when its selected
minor is nonzero on the \emph{ambient primitive chart}, meaning the full
primitive chart before the rank conditions are imposed, the required independent
minor conditions have the charged codimension, and the remaining fibers of
the slope projection have the displayed profile bound.
\end{proposition}

\begin{proof}
Under the hypotheses, standard determinantal incidence and finite-field
point counting give the charged factor for the ambient chart.  The stated
fiber hypothesis transfers it to distinct slopes.  Neither nonvanishing of
small Vandermonde minors nor constructibility alone proves these conditions
for the remaining high-dimensional balanced-core chart.
\end{proof}

\begin{proposition}[Saturation criterion]\label{prop:saturation-payment}
A saturation cell is paid if the projection from raw witnesses to the
explanation image and then the slope image of
\cref{def:explanation-occupancy} has a proved lower fiber bound, or if the
final image has a direct distinct-slope estimate at the profile scale.
\end{proposition}

\begin{proof}
This is \cref{lem:exact-occupancy-compiler}.  A lower retained-support
occupancy bounds explanation states, and projection to slopes can only
decrease cardinality.  If no such lower bound is known, only a direct image
estimate can certify payment.  An upper bound on the number of profile lifts
goes in the wrong direction.
\end{proof}

\begin{proposition}[One-pencil moving-root bound]
\label{prop:split-pencil-payment}
Let \(A,B\in\F[X]\), let
\(G=\gcd(A,B,\prod_{x\in D}(X-x))\), and put \(g=\deg G\).  If every
parameter \([s:t]\) in a set \(\mathcal Z\subseteq\mathbb P^1(\F)\)
makes \((sA+tB)/G\) have at least \(h\) distinct roots in
\(D\setminus Z(G)\), then
\[
        \abs{\mathcal Z}\le\floor{(n-g)/h}.
        \tag{7.1}\label{eq:moving-root-bound}
\]
\end{proposition}

\begin{proof}
Count incidences \(([s:t],x)\) with \(x\in D\setminus Z(G)\) and
\(sA(x)+tB(x)=0\).  Each parameter contributes at least \(h\) incidences.
For a fixed moving root, \((A(x),B(x))\ne(0,0)\), so exactly one projective
parameter works.  There are \(n-g\) moving points, proving
\eqref{eq:moving-root-bound}.
\end{proof}

\begin{remark}[Balanced-core exhaustion remains separate]
\label{rem:balanced-core-exhaustion}
The preceding proposition applies only after a chart is genuinely a
one-dimensional locator pencil and the slope injects into its projective
parameter.  It neither covers a higher-dimensional coefficient family nor
proves that such a family decomposes into subexponentially many pencils.
That structural statement is part of \textup{(RC)} or must be established
by a direct ray count.
\end{remark}

\begin{lemma}[Subexponential profile atlas]\label{lem:profile-atlas}
For a ledger-admissible sequence, the number of first-match profiles is
\(e^{o(n)}\).
\end{lemma}

\begin{proof}
This is condition (A2) in \cref{def:admissible-sequence}.  For cyclic
quotient towers the divisor count is subexponential, and for dyadic
Chebyshev towers it is logarithmic.  The statement does not include
arbitrary planted subsets or an unproved decomposition of a
higher-dimensional pencil; including either could create exponentially many
profiles.
\end{proof}

The criteria above are modular.  A concrete row must supply their missing
projection and exhaustion estimates; replacing them by the word
``constructible'' would not be a proof.
\section{The primitive prefix leaf}\label{sec:primitive-leaf}

After first-match removal of paid cells, a residual prefix-boundary chart
has a fixed set of active coordinates \(T\), \(\abs T=N\), and a fixed
support size \(m\).  The full profile slice is
\[
        \Omega_{T,m}\defeq\{x\in\{0,1\}^{T}:\sum_{t\in T}x_t=m\},
\]
and the primitive residual is a subset \(\Omc\subseteq\Omega_{T,m}\).
The superscript circle records a first-match residual; no Zariski openness
is asserted.  The normalization is taken from the full profile slice, not
from the possibly much smaller residual mass.  It is nevertheless the
\emph{realized image}, rather than the formal codomain, that supplies the
default number of target fibers.

The three words in \emph{primitive first-match residual} have separate
content.  ``First match'' means that a slope is removed as soon as one of
its witnesses occurs in an earlier ordered cell; ``residual'' means that
only witnesses surviving those removals remain; and ``primitive'' means
that the row-specific quotient, planted, field-descent, rank, and ray-saturation
certificates named by the atlas have all been applied.  Primitivity does
not mean merely trivial setwise stabilizer, and it does not assume the
Fourier, Sidon, max-fiber, or ray estimates that will be imposed below.

The prefix syndrome map is a linear map \(\Phi:\Z^T\to\B^R\), restricted
first to the full slice and then to \(\Omc\).  In the unprofiled locator
chart, when \(R<\operatorname{char}\B\), it is equivalent through Newton's
triangular change of coordinates to
\[
        \Phi(x)=\sum_{t\in T}x_t(t,t^2,\ldots,t^R).
        \tag{9.1}\label{eq:exact-power-sum-map}
\]
Equivalently, the augmented columns
\(\widetilde v_t=(1,t,\ldots,t^R)\) have constant first coordinate
\(\sum_tx_t=m\), and the other \(R\) coordinates are \(\Phi\).  A residual
quotient or rational chart is allowed to replace these columns by
\(v_t=\rho(t)(1,t,\ldots,t^{R-1})\) only when that genuine weighted
Vandermonde representation, with \(\rho(t)\ne0\), has been proved for the
chart.  It is not a consequence merely of constructibility or of deleting
a Jacobian locus.  The number \(R\) is the number of independent prefix
equations; in the unprofiled MCA boundary leaf \(R=w=a-K\).

Put \(M=\abs{\Omega_{T,m}}\), and write
\[
 \Scal=\Phi(\Omega_{T,m}),\qquad L=\abs\Scal,
 \qquad A=\abs\B^R,
 \qquad \barN^{\rm img}=M/L,
 \qquad \barN^{\rm amb}=M/A,
 \tag{9.2}\label{eq:image-ambient-scales}
\]
and for \(s\in\Scal\) put
\(F_s=\Omc\cap\Phi^{-1}(s)\) and \(f_s=\abs{F_s}\).  Unless an
ambient full-image certificate has been proved, we set
\(\barN=\barN^{\rm img}\).  The certificate is
\[
        L\ge e^{-o(N)}A.                       \tag{FI}\label{eq:full-image-certificate}
\]
The set \(\Scal\) is the \emph{effective image}: it consists of the
prefix targets actually attained by the full profile slice.  The formal
ambient target space \(\B^R\) can be much larger.  Condition \textup{(FI)},
the \emph{full-image certificate}, says that these two target counts differ
by at most a subexponential factor.  An \emph{effective-image collapse} is
the failure of this comparison; it must be retained as an effective-image
rank-collapse profile, or handled throughout with image normalization.
Under \eqref{eq:full-image-certificate}, the two scales differ by only
\(e^{o(N)}\), so either may be used at exponential accuracy.  If the
certificate fails, effective-image collapse must be retained as an
effective-image rank-collapse profile, or all estimates on this leaf must stay
at image scale.  Deleting earlier cells can only decrease every fiber.

The image-scale frontier condition is
\(\log M-\log L=o(N)\).  The familiar ambient equation
\(\log\binom Nm-R\log\abs\B=o(N)\) is interchangeable with it only under
\eqref{eq:full-image-certificate}.  This distinction is immaterial for a
surjective prefix map but essential for a collapsed primitive image.

\begin{lemma}[Exact image--ambient moment conversion]
\label{lem:image-ambient-moment-conversion}
For \(q\ge1\), extend \(f_s\) by zero from \(\Scal\) to \(\B^R\) and
define
\[
 \Gamma_q^{\rm img}=\frac1L\sum_{s\in\Scal}
       \left(\frac{f_s}{\barN^{\rm img}}\right)^q,
 \qquad
 \Gamma_q^{\rm amb}=\frac1A\sum_{s\in\B^R}
       \left(\frac{f_s}{\barN^{\rm amb}}\right)^q.
\]
Then
\[
 \Gamma_q^{\rm amb}
   =\left(\frac{A}{L}\right)^{q-1}\Gamma_q^{\rm img}.
 \tag{9.3}\label{eq:image-ambient-moment-identity}
\]
Thus an ambient-normalized moment bound implies the corresponding
image-normalized bound.  The converse is valid with subexponential loss at
logarithmic order only under \eqref{eq:full-image-certificate}.
\end{lemma}

\begin{proof}
Since \(\barN^{\rm amb}=M/A\) and \(\barN^{\rm img}=M/L\), zero extension
and one rearrangement give
\[
 \Gamma_q^{\rm amb}
 =\frac1A\sum_{s\in\Scal}\left(\frac{Af_s}{M}\right)^q
 =\left(\frac AL\right)^{q-1}
   \frac1L\sum_{s\in\Scal}\left(\frac{Lf_s}{M}\right)^q.
\]
Because \(L\le A\), the asserted one-way implication follows.  Under
\eqref{eq:full-image-certificate},
\((q-1)\log(A/L)=o(Nq)\), which proves the reverse transfer at the scale
used here.
\end{proof}

\begin{remark}[Full-slice flatness certifies the image size]
\label{rem:flatness-certifies-image}
If a Fourier argument proves
\(\max_s\abs{\Omega_{T,m}\cap\Phi^{-1}(s)}
 \le e^{o(N)}\barN^{\rm amb}\), then averaging over the nonempty image
fibers gives \(A/L\le e^{o(N)}\).  Hence such an ambient flatness theorem
itself proves \eqref{eq:full-image-certificate}; image collapse cannot be
silently assumed away before that estimate has been established.
\end{remark}

\begin{definition}[Primitive Q]\label{def:primitive-q}
The \emph{Q branch} is the primitive max-fiber problem: it asks whether one
effective prefix target contains substantially more residual supports than
the image-normalized average.  A primitive prefix leaf satisfies Q with
loss \(e^{o(N)}\) if
\[
        \max_{s\in\Scal}f_s\le e^{o(N)}\barN^{\rm img}.
\]
A sequence of primitive leaves satisfies logarithmic-moment Q if, at every
logarithmic order \(q=q_N\to\infty\), \(q\le(\log N)^C\), used by the
argument,
\[
 \Gord_q\defeq\Gamma_q^{\rm img}
 =L^{-1}\sum_{s\in\Scal}
       \left(\frac{f_s}{\barN^{\rm img}}\right)^q
 \le e^{o(Nq)}.
\]
Here a \emph{logarithmic moment order} means
\(q=q_N\to\infty\) with \(q\le(\log N)^C\) for a fixed \(C\).  The
displayed normalized \(q\)-th moment is a R\'enyi-type probe of the fiber
distribution: letting \(q\) grow slowly makes it approximate the largest
normalized fiber without introducing an exponential moment-order cost.
An ambient formulation is permitted only when it is proved directly or
when \eqref{eq:full-image-certificate} has already been certified.
\end{definition}

The next lemma is the primitive logarithmic moment equivalence used in the earlier RS--MCA reductions.  It is included because it is the exact interface between moment estimates and max-fiber Q.

\begin{lemma}[Logarithmic moments and primitive Q]\label{lem:logmoment-q}
Assume \(q=q_N\to\infty\) and \(\log L/q=o(N)\).  Then
\(\Gord_q\le e^{o(Nq)}\) implies primitive Q.  Conversely, primitive Q
implies \(\Gord_q\le e^{o(Nq)}\).
\end{lemma}

\begin{proof}
Let \(M_*=\max_s f_s/\barN^{\rm img}\).  Since
\(M_*^q\le\sum_s(f_s/\barN^{\rm img})^q=L\Gord_q\), the moment bound gives
\(\log M_*\le(\log L+\log\Gord_q)/q=o(N)\).  Conversely,
\(\sum_sf_s\le\abs{\Omega_{T,m}}=L\barN^{\rm img}\), so
\(\Gord_q\le M_*^{q-1}\).  Primitive Q gives \(M_*=e^{o(N)}\), hence
\(\Gord_q=e^{o(Nq)}\).
\end{proof}

The genuine Vandermonde condition removes the ordinary rank-defect alternative.

\begin{lemma}[Augmented Vandermonde independence]\label{lem:vandermonde}
Let \(t_1,\ldots,t_r\in\B\) be distinct and \(r\le R+1\).  Then the
columns \(\widetilde v_{t_i}=(1,t_i,\ldots,t_i^R)\) are linearly
independent.  The same holds for \(\rho(t_i)(1,t_i,\ldots,t_i^{R-1})\)
when \(r\le R\) and every \(\rho(t_i)\ne0\).
\end{lemma}

\begin{proof}
In either case the leading \(r\times r\) minor, after removing the nonzero
column weights when present, is \((t_i^j)_{0\le j<r,1\le i\le r}\).  Its
determinant is \(\prod_{i<j}(t_j-t_i)\ne0\).
\end{proof}

In the high-energy argument below the Vandermonde lemma is not used:
Boolean-cube growth gives the contradiction directly.  The lemma certifies
only the exact locator chart, or a residual chart for which the displayed
weighted form has separately been verified; it does not classify arbitrary
rank defects as paid cells.


\section{Prefix coordinates and the Vandermonde map}\label{sec:prefix-coordinates}

This section explains why the primitive residual has the linear Vandermonde form used in \cref{sec:primitive-leaf}.  Let \(S\subseteq D\) have size \(m\), and write its monic locator as
\[
        Q_S(X)=\prod_{t\in S}(X-t)=X^m+q_1(S)X^{m-1}+\cdots+q_m(S).
\]
The first \(R\) boundary coordinates may be taken either as elementary symmetric coefficients \(q_1,\ldots,q_R\) or, when \(R<\operatorname{char}\B\), as power sums \(p_j(S)=\sum_{t\in S}t^j\), \(1\le j\le R\).  Newton identities give triangular polynomial changes of variables between these two coordinate systems; the diagonal entries are \(1,2,\ldots,R\), so the change is invertible in characteristic greater than \(R\).

For the unprofiled locator chart, fixing the support size makes the zeroth
power sum constant and Newton's identities identify the first \(R\)
locator coefficients with the power sums
\(\sum_tx_tt^j\), \(1\le j\le R\).  Thus
\eqref{eq:exact-power-sum-map} is exact.  For a general residual chart one
first fixes the planted roots, quotient data, common factors, and all
coordinates charged to earlier cells.  The fixed part contributes a
constant syndrome \(s_0\).  To use the primitive theorem one must then
exhibit the remaining map explicitly, either as the same power-sum map, as
a genuine common-weight Vandermonde map, or as a map for which the required
Sidon estimate has been proved directly.  Coordinate-dependent rational
weights do not become a Vandermonde system merely because they are
nonzero.  Here ``common-weight'' means that one scalar \(\rho(t)\)
multiplies every monomial coordinate of the column at \(t\), rather than
using a different weight for each coordinate.  This verification is condition~\textup{(A5)} of
\cref{def:admissible-sequence}.

The fixed-density condition is also intrinsic.  In a prefix leaf the total
support size is fixed by the agreement threshold, and all paid planted
blocks have fixed size.  Thus the number of remaining choices is fixed.  If
a profile leaves several colors of coordinates, one obtains a product of
fixed-density slices; the proof applies to each color after encoding the
colors in a bounded alphabet.  For the main statement we present the
Boolean case because it is the MCA support case and because the
Boolean-cube difference bound applies directly.

\begin{lemma}[Newton-coordinate fiber equivalence]\label{lem:newton-equivalence}
Assume \(R<\operatorname{char}\B\).  For supports of fixed size, the fibers of the first \(R\) elementary locator coefficients are exactly the fibers of the first \(R\) power sums.  Consequently the Q max-fiber problem may be stated using either coordinate system.
\end{lemma}

\begin{proof}
Newton identities give, for \(1\le j\le R\),
\[
        p_j+q_1p_{j-1}+\cdots+q_{j-1}p_1+jq_j=0.
\]
Since \(j\) is invertible in \(\B\), \(q_j\) is determined by \(p_1,\ldots,p_j\), and conversely \(p_j\) is determined by \(q_1,\ldots,q_j\).  The transformations are triangular inverse polynomial maps on the first \(R\) coordinates.
\end{proof}

The \emph{description entropy} of a profile family is the logarithm of the
number of allowed auxiliary/profile descriptions.  It measures the census
of cells and is distinct from the algebraic description complexity of one
cell, which measures the number and degrees of its defining equations.

\begin{lemma}[Primitive profiles have subexponential multiplicity]\label{lem:profile-multiplicity}
If the cell ledger fixes only bounded-complexity quotient, planted, tangent,
extension, and split-pencil profiles, and the total description entropy of
all data outside the moving support is \(o(n)\) bits, then the number of
primitive profiles is \(e^{o(n)}\).
\end{lemma}

\begin{proof}
There are at most \(2^{o(n)}=e^{o(n)}\) descriptions by hypothesis.  In
particular, a planted or common block of size \(o(n)\) has
\(\sum_{j=o(n)}\binom nj=e^{o(n)}\) possible supports, but recording
\(s=o(n)\) freely chosen field values costs \(s\log_2\abs\B\) bits and is
subexponential only when that product is \(o(n)\).  Positive-density data
requires its own profile census and entropy payment; it cannot be discarded
as a bounded-complexity label.
\end{proof}

\begin{remark}
The subexponential profile lemma is the asymptotic counterpart of the
finite ledger census.  It is the reason a row proof must distinguish
between a genuine cell payment and a mere change of variables.  A change
of variables with exponentially many charts would consume the whole
frontier budget.
\end{remark}

\section{The Sidon split}\label{sec:sidon-split}

The Q moment can be large for two opposite reasons.  A fiber may contain
many repeated additive differences, which is the high-energy situation
seen by inverse additive combinatorics, or it may be large while its
differences almost never repeat, which is the Sidon-like situation.  The
split below separates these mechanisms before either one is estimated.

Every support is identified with its incidence vector in the torsion-free
group \(\Z^T\).  In particular, all sums and differences in this section
are integer operations, even when the Reed--Solomon coefficient field has
characteristic two.  For a finite set \(F\subseteq\{0,1\}^T\), define its
additive energy
\[
        E(F)=\abs{\{(a,b,c,d)\in F^4:a-b=c-d\}}.
\]
Let \(X,X'\) be independent uniform points of \(F\).  Then \(E(F)=\abs F^4\sum_z\Prob(X-X'=z)^2\).  It is convenient to normalize by
\[
        \Delta(F)=\frac{E(F)}{\abs F^3}.
\]
The quantity \(\Delta(F)\) is an additive closure probability:
\[
        \Delta(F)=\Prob_{a,b,c\in F}(a-b+c\in F).
\]
Indeed, once \(a,b,c\) are chosen, the fourth point in an energy quadruple is forced to be \(d=a-b+c\).

Thus additive energy counts repeated differences, and \(\Delta(F)\) is
their normalized density.  Values \(\Delta(F)=e^{-o(N)}\) indicate
substantial additive closure at exponential scale, whereas
\(\Delta(F)\le e^{-\sigma N}\) is the low-energy, Sidon-like regime.
The logarithmic moments below are the slowly growing normalized moments
from \cref{def:primitive-q}; they measure how much of the Q obstruction is
carried by fibers of either type.
Throughout these sections, \emph{positive rate} means a fixed positive
exponential rate along a subsequence: for a fiber it means an excess of
the form \(e^{cN-o(N)}\), and for a logarithmic moment it means an excess
of the form \(e^{cNq-o(Nq)}\), for some constant \(c>0\).

For the fiber \(F_s\), write \(\Delta_s=E(F_s)/f_s^3\), with \(\Delta_s=0\) if \(f_s=0\).

\begin{definition}[Sidon-heavy logarithmic moment]\label{def:sidon-heavy}
For \(\sigma>0\), the Sidon-heavy part of the normalized moment is
\[
        \Gsid_{q,\sigma}=L^{-1}\sum_{\Delta_s\le e^{-\sigma N}}
        \left(\frac{f_s}{\barN}\right)^q.
\]
A sequence has no positive-rate Sidon-heavy obstruction if \(\Gsid_{q,\sigma}\le e^{o(Nq)}\) for every fixed \(\sigma>0\) and every logarithmic \(q\).
Here ``Sidon'' refers to the exponentially small normalized energy
\(\Delta_s\), while ``heavy'' refers to the positive-exponential
contribution of those fibers to the normalized logarithmic moment; it does
not assert that every low-energy fiber is individually large.
\end{definition}

The following dichotomy is the correct replacement for the naive assertion that collision excess directly creates small doubling.

\begin{proposition}[Ordinary moment split]\label{prop:ordinary-moment-split}
Fix \(\eta>0\) and \(\sigma>0\).  If \(\Gord_q\ge e^{\eta Nq}\) and \(\Gsid_{q,\sigma}\le \frac12 e^{\eta Nq}\), then there is a syndrome \(s\) such that
\[
        f_s\ge e^{\eta N-o(N)}\barN,
        \qquad
        \Delta_s>e^{-\sigma N}.
\]
If no positive-rate Sidon-heavy obstruction is present, then any positive-rate moment failure gives, after taking \(\sigma\to0\) slowly, a fiber with \(f_s\ge e^{\eta N-o(N)}\barN\) and \(E(F_s)\ge f_s^3e^{-o(N)}\).
\end{proposition}

\begin{proof}
The high-energy part of the moment is at least \(\frac12e^{\eta Nq}\).  Since there are \(L\) syndromes, some high-energy syndrome satisfies
\[
        \left(\frac{f_s}{\barN}\right)^q
        \ge \frac{L}{2}\,e^{\eta Nq}\cdot L^{-1}
        =\frac12 e^{\eta Nq}.
\]
Thus \(f_s\ge e^{\eta N-o(N)}\barN\), and by construction \(\Delta_s>e^{-\sigma N}\).  If \(\Gsid_{q,\sigma}\le e^{o(Nq)}\) for every fixed \(\sigma\), apply the first part along a sequence \(\sigma_j\downarrow0\) and diagonalize; the normalization \(\Delta_s=E(F_s)/f_s^3\) gives the energy conclusion.
\end{proof}

The proof also explains the no-go point.  Suppose a fiber is heavy but \(\Delta_s\le e^{-\sigma N}\).  A closure-rich diagram using tests \(a-b+c\in F_s\) pays a factor \(\Delta_s\) for each nonredundant closure.  Such diagrams cannot account for the ordinary contribution \(f_s^q\) of \(q\) independent choices in \(F_s\).  Therefore Sidon-heavy fibers must be removed by a separate cell; they cannot be forced into the high-energy inverse branch.

\begin{definition}[Sidon moment payment]\label{def:sidon-paid-cell}
A primitive prefix leaf has a Sidon moment payment if
\(\Gsid_{q,\sigma}=e^{o(Nq)}\) for every fixed \(\sigma>0\) and every
logarithmic moment order used by the primitive-Q argument, uniformly over
profiles.  This is an analytic support-fiber statement.  It is not
equivalent to a first-match distinct-slope payment; the latter still
requires \textup{(RC)} or a direct ray bound.
\end{definition}

The equivalence in the definition follows from \cref{lem:logmoment-q} restricted to the low-energy set of syndromes.  Operationally, one proves payment either directly by bounding these fibers or indirectly by proving Fourier flatness on the primitive complement of algebraic major arcs.


\section{Why the Sidon cell is necessary}\label{sec:sidon-necessity}

It is tempting to try to prove primitive Q by expanding the logarithmic moment into diagrams and then showing that every primitive diagram contains many additive closure tests.  That strategy is correct only for a closure-rich submoment.  It cannot capture the ordinary moment on a heavy low-energy fiber.  The obstruction is simple enough to record as a lemma.

Let \(F=F_s\) be a fiber of size \(f\).  A closure test is an equation of the form \(a-b+c\in F\), with \(a,b,c\in F\) already chosen.  Its success probability is \(\Delta(F)=E(F)/f^3\).  A closure-rich diagram with \(\ell\) nonredundant closure tests and \(r\) free choices contributes, up to a subexponential diagram factor,
\[
        f^r\Delta(F)^\ell .
\]
After the free coordinates have been chosen, in the image normalization of a primitive prefix leaf, the remaining \(q-r\) coordinates are supplied only at the random scale \(\barN\).  Thus its contribution is at most
\[
        e^{o(Nq)}\barN^{q-r}f^r\Delta(F)^\ell .
\]

\begin{lemma}[No-go for ordinary moment capture]\label{lem:no-go}
Assume that \(f\ge e^{\eta N}\barN\), \(\Delta(F)\le e^{-\sigma N}\), and that a family of diagrams has \(r=o(q)\) free fiber choices and \(\ell\ge c q\) nonredundant closure tests.  Then the total contribution of these diagrams is smaller than the ordinary fiber contribution \(f^q\) by a factor
\[
        \exp\bigl(-(\eta+c\sigma)Nq+o(Nq)\bigr)
\]
up to changing \(\eta\) by \(o(1)\).
\end{lemma}

\begin{proof}
For one diagram the ratio to \(f^q\) is at most
\[
        e^{o(Nq)}\left(\frac{\barN}{f}\right)^{q-r}\Delta(F)^\ell
        \le e^{o(Nq)}\exp(-\eta N(q-r)-c\sigma Nq).
\]
Since \(r=o(q)\), this is \(\exp(-(\eta+c\sigma)Nq+o(Nq))\).  The number of bounded-complexity logarithmic diagrams is \(e^{O(q\log q)}=e^{o(Nq)}\), so summing over them does not change the exponential rate.
\end{proof}

The lemma has an important conceptual consequence.  A large Sidon-like fiber is not a hidden approximate group.  It is a different phenomenon: many independent choices land in the same syndrome, but their pairwise differences almost never repeat.  Additive-combinatorics inverse tools are designed for the opposite situation, where repeated differences are abundant.  Therefore the first-match ledger must contain a Sidon/Fourier cell before the primitive inverse argument begins.

The valid implication is exactly \cref{prop:ordinary-moment-split}.  Ordinary
moment mass splits into a low-energy part, which needs a separately proved
Sidon moment estimate, and a high-energy part, which is dispatched by
\cref{prop:high-energy-impossible}.

This mechanism does not replace a prefix-flatness or Sidon-payment theorem;
it identifies exactly where one is needed.  The first-match order is:
algebraic major arcs first, then a separately certified Sidon/Fourier cell,
and only then the high-energy primitive inverse step.  Naming a large
low-energy fiber a cell does not pay it.

\section{Fourier flatness as a Sidon payment}\label{sec:fourier}\label{sec:sidon-closure}

This section records the general weighted form of the Fourier/Sidon payment.  It is phrased for an arbitrary finite field \(\B\), because the primitive residual charts that come from the moduli ledger need not have unweighted power-sum coordinates.  Let \(T\) be a finite set, \(\abs T=N\), let \(g:T\to\B^R\) be any function, and set
\[
        \Omega_{T,m}=\{x\in\{0,1\}^T: \sum_{t\in T}x_t=m\},
        \qquad
        \Psi(x)=\sum_{t\in T}x_tg(t).
\]
In applications \(g(t)=(\rho_0(t),\rho_1(t)t,\ldots,\rho_{R-1}(t)t^{R-1})\) with nonzero rational weights \(\rho_i\) on the primitive chart.

Fourier inversion has three pieces.  The zero character produces the
average-fiber main term.  On the effective dual, the chart-specific
partition of \cref{def:effective-major-minor} designates certified minor
characters and the complementary major characters.  The ambient algebraic
exceptional locus of \cref{def:major-arc} helps construct that partition
but does not canonically determine it, because an effective character has
many ambient lifts.  The pointwise condition \textup{(PF)} controls every
certified minor lift, while the aggregate condition \textup{(MA)} controls
the sum of all designated major characters.  Neither estimate substitutes
for the other.

The condition \textup{(PF')} below is the raw power-sum flatness
inequality when one has a uniform bound for every nonzero character.
Condition \textup{(PF)} is its row-specific Weil specialization for minor
characters after substituting the pole-degree bound; the omitted major
characters are then handled by \textup{(MA)}.

\begin{theorem}[Prefix flatness from power-sum control]\label{thm:prefix-flatness-power-sum}
Fix a nontrivial additive character \(\psi\) of \(\B\).  For \(\alpha\in\B^R\setminus\{0\}\) and \(1\le j\le m\), put
\[
        P_j(\alpha)=\sum_{t\in T}\psi\bigl(j\alpha\cdot g(t)\bigr).
\]
Assume \(\abs{P_j(\alpha)}\le \Lambda\) for every \(\alpha\ne0\) and every \(1\le j\le m\).  Then, for every \(z\in\B^R\),
\[
        \abs{\Psi^{-1}(z)}
        \le \abs\B^{-R}\binom Nm+\binom{\Lambda+m-1}{m}.
\]
Consequently, if
\[
        \tag{PF'}
        R\log\abs\B+\log\binom{\Lambda+m-1}{m}-\log\binom Nm\le o(N),
\]
then \(\max_z\abs{\Psi^{-1}(z)}\le e^{o(N)}\abs\B^{-R}\binom Nm\).
\end{theorem}

\begin{proof}
By orthogonality,
\[
        \abs{\Psi^{-1}(z)}=\abs\B^{-R}
        \sum_{\alpha\in\B^R}\psi(-\alpha\cdot z)
        e_m\bigl(\psi(\alpha\cdot g(t)):t\in T\bigr),
\]
where \(e_m\) denotes the \(m\)-th elementary symmetric coefficient.  The term \(\alpha=0\) is the main term \(\abs\B^{-R}\binom Nm\).  For \(\alpha\ne0\), write \(a_t=\psi(\alpha\cdot g(t))\).  The generating identity
\(\sum_{r\ge0}e_r(a_t)u^r=\exp(\sum_{j\ge1}(-1)^{j-1}(\sum_ta_t^j)u^j/j)\) and the hypothesis \(\abs{\sum_ta_t^j}\le\Lambda\) for \(j\le m\) give
\(\abs{e_m(a_t:t\in T)}\le[u^m](1-u)^{-\Lambda}=\binom{\Lambda+m-1}{m}\).  Summing the nontrivial characters and using the outside factor \(\abs\B^{-R}\) gives the first bound.  The displayed condition \textup{(PF')} makes the error term subexponential relative to the random-fiber main term.
\end{proof}

\begin{remark}[Small characteristic cycles]\label{rem:small-characteristic-cycles}
If some integers \(j\le m\) vanish in the characteristic, the formal
Li--Wan cycle index separates the cycles divisible by
\(p=\operatorname{char}\B\); the corresponding coefficient majorant is
\([u^m](1-u)^{-\Lambda}(1-u^p)^{-(N-\Lambda)/p}\).  This bookkeeping
identity does not supply the missing character-sum bounds for those
cycles, exclude Artin--Schreier phases, or prove a uniform Sidon payment.
The pointwise theorem above is therefore deliberately stated in the
large-characteristic range.  A small-characteristic leaf needs a direct
cycle-by-cycle certificate or a separate image-normalized Q/Sidon theorem.
\end{remark}

\begin{proposition}[Weighted Weil minor arcs]\label{prop:weighted-weil-minor-arcs}
Let \(T\) be either a multiplicative coset in \(\B^\times\) or the \(x\)-coordinate of a torus twin coset.  Suppose that every certified minor lift \(\alpha\in\B^R\) has phase \(R_\alpha(t)=\alpha\cdot g(t)\) which is a nonconstant rational function on the corresponding genus-zero curve, is not of Artin--Schreier form, and has total pole-degree at most \(\Delta_{\rm pole}\).  If \(m<\operatorname{char}\B\), then the stated power-sum bound holds for every certified minor lift, with
\[
        \Lambda=C_0(\Delta_{\rm pole}+1)\sqrt{\abs\B}+\abs P
\]
for multiplicative cosets, and with \(2C_0(\Delta_{\rm pole}+1)\sqrt{\abs\B}+\abs P\) for circle twin-cosets.  Here \(P\) is the planted exceptional set and \(C_0\) is an absolute Weil constant.
\end{proposition}

\begin{proof}
For a multiplicative coset \(\theta H\subseteq\B^\times\), write the coset indicator as the average of the multiplicative characters of \(\B^\times/H\).  Each nontrivial power sum becomes an average of mixed multiplicative-additive character sums \(\sum_x\chi(x)\psi(jR_\alpha(\theta x))\).  Since \(1\le j<m<\operatorname{char}\B\), multiplication by \(j\) does not make the rational function constant or Artin--Schreier.  Weil's bound for one-variable mixed character sums gives \(C_0(\Delta_{\rm pole}+1)\sqrt{\abs\B}\).  Removing planted exceptional points changes the sum by at most \(\abs P\).  The circle case is the same after pulling back by \(x=(u+u^{-1})/2\) to the norm-one torus; the two torus branches give the factor two.
\end{proof}

\begin{theorem}[Unconditional shallow-prefix MI and MA]
\label{thm:unconditional-shallow-mi-ma}
Let \(p\to\infty\) through odd primes.  For each \(p\), let \(T_p\) be
either a multiplicative coset in \(\F_p^\times\) or a Chebyshev
twin-coset \(x\)-domain over \(\F_p\), in the latter case with
\(p\equiv3\pmod4\), and put \(N_p=\abs{T_p}\).
Fix \(0<\alpha<1/2\), choose integers \(m_p,R_p\) with
\[
 1\le R_p,\qquad
 \alpha\le \frac{m_p}{N_p}\le1-\alpha,
 \qquad g_p(t)=(t,t^2,\ldots,t^{R_p}),
 \qquad R_p\sqrt p=o(N_p),                         \tag{SFM1}
\]
and let \(\Psi_p(S)=\sum_{t\in S}g_p(t)\) on
\(\binom{T_p}{m_p}\).  Then \(V_{g_p}=\F_p^{R_p}\), and the certified effective partition may take
every nontrivial effective character to be minor and no character to be
major.  For some \(c_\alpha>0\),
\[
 \sum_{1\ne\chi\in\widehat{V_{g_p}}}
 \left|e_{m_p}\bigl(\chi(g_p(t)-g_p(t_0)):t\in T_p\bigr)\right|
 \le e^{-c_\alpha N_p}\binom{N_p}{m_p}.            \tag{SFM2}
\]
Consequently effective \textup{(MI)} holds, \textup{(MA)} is vacuous,
and
\[
 \max_z\abs{\Psi_p^{-1}(z)}
 \le\bigl(1+e^{-c_\alpha N_p}\bigr)
       p^{-R_p}\binom{N_p}{m_p}.                   \tag{SFM3}
\]
Indeed every \(z\in\F_p^{R_p}\) is attained for all sufficiently large
rows, and uniformly in \(z\),
\[
 \abs{\Psi_p^{-1}(z)}
 =p^{-R_p}\binom{N_p}{m_p}
  \left(1+O\bigl(e^{-c_\alpha N_p}\bigr)\right).   \tag{SFM4}
\]
Thus these smooth and circle prefix charts have unconditional effective
Fourier payment and image coverage in the shallow range \textup{(SFM1)}.
\end{theorem}

\begin{proof}
Condition \textup{(SFM1)} implies \(R_p<p\), \(R_p<N_p\), and
\(R_p\log p=o(N_p)\).  Moreover \(N_p\le p+1\), so the fixed upper
density bound gives \(m_p<p\) for every sufficiently large \(p\).
For a nonzero ambient parameter, let
\(d=\max\{j:a_j\ne0\}\).  Then \(1\le d\le R_p<p\), and
\(P(X)=\sum_{j=1}^{R_p}a_jX^j\) has a pole of order \(d\) at infinity.
Every pole of a nonconstant Artin--Schreier function \(G^p-G+c\) has
order divisible by \(p\), so \(P\) is non-Artin--Schreier.  In the circle
case \(P((u+u^{-1})/2)\) has poles of the same order \(d<p\) at the two
ends of the torus compactification and is likewise non-Artin--Schreier.

For a multiplicative coset, the character expansion of its indicator and
the mixed Weil bound give, uniformly for \(1\le j\le m_p<p\),
\[
 \left|\sum_{t\in T_p}\psi(jP(t))\right|=O(R_p\sqrt p)=o(N_p).
\]
For a circle domain, pullback by
\(x=(u+u^{-1})/2\) gives the same conclusion with an inessential factor
two.  Put \(\Lambda_p=O(R_p\sqrt p)=o(N_p)\).  The cycle-index estimate
therefore bounds every nontrivial fixed-weight Fourier coefficient by
\[
 \binom{\ceil{\Lambda_p}+m_p-1}{m_p}=e^{o(N_p)}.
\]
If a nonzero ambient character annihilated \(V_{g_p}\), its character
value would be constant on \(T_p\), making the power sum have modulus
\(N_p\), contrary to the preceding \(o(N_p)\) bound.  Thus the
annihilator is zero and \(V_{g_p}=\F_p^{R_p}\).  There are
\(p^{R_p}=e^{o(N_p)}\) effective characters, whereas
Stirling's formula and the density interval give
\(\binom{N_p}{m_p}\ge
e^{(h(\alpha)+o(1))N_p}\).  Summing the coefficient bounds proves
\textup{(SFM2)}, after decreasing the positive constant
\(c_\alpha\).  All nontrivial effective characters have just received a
certified minor lift, so the major set is empty.  Finally
Fourier inversion shows that the absolute nonzero-character error in each
fiber is \(e^{o(N_p)}\), whereas its main term is
\(p^{-R_p}\binom{N_p}{m_p}=e^{(h(m_p/N_p)+o(1))N_p}\).
This proves \textup{(SFM3)}--\textup{(SFM4)} and eventual surjectivity.
\end{proof}

\begin{definition}[Ambient algebraic exceptional locus]\label{def:major-arc}
An ambient phase parameter is \emph{algebraically exceptional} if its
weighted phase is constant on a routed quotient fiber, is Artin--Schreier
degenerate, has a pole or zero supported on a positive-density moving
divisor, or factors through a proper quotient, Chebyshev map,
field-descent locus, tangent rank drop, differential-locator defect, or
ray-saturation collapse.  On an effective chart this locus is used only
to construct the certified partition of
\cref{def:effective-major-minor}; it is not itself a numerical estimate.
In particular, an Artin--Schreier phase defined on the whole active set
restricts to the trivial effective character by
\cref{lem:constant-modes-annihilator}.
\end{definition}

\begin{proposition}[Major-arc algebraic interface]\label{prop:major-arcs-are-cells}
Each condition in \cref{def:major-arc} is a constructible condition on the
dual phase parameters and points to a quotient, Chebyshev, field-descent,
planted-divisor, differential, or rank-loss profile.  This classification
does not imply that the character vanishes after primal first-match
deletion and does not prove \textup{(MA)}.  A flatness argument must also
bound \eqref{eq:major-arc-aggregate}, recursively or directly.
\end{proposition}

\begin{proof}
Constancy on a quotient fiber, Chebyshev factorization,
Artin--Schreier degeneracy, a positive-density pole divisor, and a rank
drop are respectively coefficient, Frobenius, resultant, or determinantal
conditions.  They therefore identify constructible dual loci on which a
recursive estimate can be attempted.  Fourier inversion sums characters,
not primal witnesses, so no support-level payment follows from this fact.
\end{proof}

\begin{corollary}[Sidon payment from prefix flatness]\label{cor:sidon-payment-from-prefix-flatness}
If a primitive smooth or circle chart satisfies the pointwise minor-arc
range of \cref{def:prefix-flat-range} and the major-arc aggregate
in the ambient instance of \textup{(MA)}, then its Fourier/Sidon-heavy logarithmic moment is
\(e^{o(Nq)}\), in image normalization, for every logarithmic
\(q\le(\log N)^C\), for a fixed \(C\).
\end{corollary}

\begin{proof}
Write \(\barN^{\rm amb}=\abs\B^{-R}\binom Nm\).  The zero character gives
this ambient random-fiber term.  By
\cref{thm:prefix-flatness-power-sum,prop:weighted-weil-minor-arcs}, the
minor characters contribute at most \(e^{o(N)}\barN^{\rm amb}\), and
\textup{(MA)} gives the same bound for the major characters.  Hence every
full-slice, and therefore every residual, fiber is at most
\(e^{o(N)}\barN^{\rm amb}\).  By
\cref{rem:flatness-certifies-image}, this also proves the full-image
certificate; the exact identity
\eqref{eq:image-ambient-moment-identity} then gives the stated
image-normalized moment bound.
\end{proof}

\begin{proposition}[Frontier Weil separation]\label{prop:frontier-weil-separation}
Let \(N=\abs T\), let \(m=\theta N+o(N)\) with \(0<\theta<1\), and suppose a primitive smooth or circle chart has minor-arc power-sum parameter \(\Lambda=o(N)\).  If the ambient random-fiber scale \(\barN=\abs\B^{-R}\binom Nm\) satisfies \(\log\barN=o(N)\), or more generally \(\log\barN\ge -o(N)\), then the prefix-flat inequality \textup{(PF')} holds.  In particular, minor arcs alone cannot create a positive-rate max-fiber or Sidon-heavy moment at this near-zero ambient entropy scale.
\end{proposition}

\begin{proof}
The elementary estimate \(\binom{\Lambda+m-1}{m}=\binom{\Lambda+m-1}{\Lambda-1}\le (e(m+\Lambda)/\Lambda)^{\Lambda}\) gives \(\log\binom{\Lambda+m-1}{m}=o(N)\), since \(\Lambda=o(N)\).  The relative minor-arc error exponent is
\[
        \log\binom{\Lambda+m-1}{m}-\log\barN
        =R\log\abs\B+\log\binom{\Lambda+m-1}{m}-\log\binom Nm.
\]
The assumed frontier inequality makes this at most \(o(N)\); it may be
negative by a linear amount, which is only stronger.  Thus the minor-arc
contribution to every full-profile fiber is at most
\(e^{o(N)}\barN\).  This statement concerns only minor characters; a
full flatness conclusion also requires \textup{(MA)}.
\end{proof}

\begin{theorem}[Sidon payment in the pointwise flat range]\label{thm:sidon-resolved-payment}\label{thm:no-sidon-heavy-fourier}\label{thm:fourier-sidon-closure}
Let \(\Omc\subseteq\Omega_{T,m}\) be a primitive first-match residual of a
smooth or circle chart, with weighted prefix map \(\Psi\) and ambient
scale \(\barN^{\rm amb}=\abs\B^{-R}\binom Nm\).  Put
\(A_{\rm amb}=\abs\B^R\) and
\(L=\abs{\Psi(\Omega_{T,m})}\).  Suppose the pointwise minor-arc
condition \textup{(PF)} and the ambient instance of the aggregate
condition \textup{(MA)} hold.
Then, for every fixed \(\sigma>0\) and every
logarithmic \(q\le(\log N)^C\), for a fixed \(C\),
\[
        \Gsid_{q,\sigma}(\Omc)\le e^{o(Nq)},
\]
where \(\Gsid\) uses the image normalization of
\cref{def:sidon-heavy}.
Equivalently, no such pointwise-flat residual leaf supports positive-rate
low-energy heavy fibers.
\end{theorem}

\begin{proof}
It is enough to bound every residual fiber, and a residual is contained in
the full slice.  Fourier inversion on that slice splits into the zero,
minor, and major characters.  The zero character contributes
\(\barN^{\rm amb}\).
The cycle-index estimate and \cref{prop:weighted-weil-minor-arcs}, together
with \textup{(PF)}, bound the complete minor contribution by
\(e^{o(N)}\barN^{\rm amb}\).  Condition \textup{(MA)} bounds the complete major
contribution by the same quantity.  Thus every full-slice, and hence every
residual, fiber is \(e^{o(N)}\barN^{\rm amb}\).  This proves the
full-image certificate by \cref{rem:flatness-certifies-image}, and
\cref{lem:image-ambient-moment-conversion} transfers the ambient moment
bound to image normalization.
\end{proof}

\begin{remark}[Role of the Sidon-heavy cell]
The theorem implements ``pay structure, flatten Sidon, inverse high energy''
only in the range where both \textup{(PF)} and \textup{(MA)} are verified.
At entropy depth a phase
can have degree up to \(R\), and the raw Weil parameter is of order
\(R\sqrt{\abs\B}\); this is not automatically \(o(N)\) when
\(R\asymp N/\log\abs\B\).  Outside the pointwise-flat range a different
Sidon payment remains necessary.
\end{remark}

\section{External additive combinatorics}\label{sec:additive-tools}

Only two external additive-combinatorics inputs are used in the primitive high-energy proof.

\begin{theorem}[Balog--Szemeredi--Gowers, energy form]\label{thm:bsg}
There is an absolute constant \(C\) such that if \(A\) is a finite subset of an abelian group and \(E(A)\ge \abs A^3/K\), then there is \(A'\subseteq A\) with \(\abs{A'}\ge K^{-C}\abs A\) and \(\abs{A'-A'}\le K^C\abs{A'}\).
\end{theorem}

This is the usual energy form of the Balog--Szemeredi--Gowers theorem; polynomial dependence on \(K\) follows from Gowers's refinement of the Balog--Szemeredi theorem \cite{BalogSzemeredi,GowersSzemeredi}.  The difference-set form follows from the sumset form by replacing \(A\) by \(-A\) and using standard Ruzsa calculus \cite{TaoVu}, or directly by applying the graph form to differences.

\begin{theorem}[Boolean-cube difference growth]\label{thm:quasicube}
If \(U\subseteq\{0,1\}^d\), then
\[
                  \abs{U-U}\ge \abs U^{3/2}.
\]
\end{theorem}

This is the Boolean-cube specialization of a broader sumset inequality
of Green--Matolcsi--Ruzsa--Shakan--Zhelezov
\cite{GreenMatolcsiRuzsaShakanZhelezov}.  It also follows from the
induced-doubling identity of Matolcsi--Ruzsa--Shakan--Zhelezov
\cite{MatolcsiRuzsaShakanZhelezov}.  Only the displayed Boolean statement
is used below.

\begin{remark}
Entropy inverse theorems of Tao and the entropy-doubling theory of Green--Manners--Tao are compatible with this proof and are often useful after one has produced small-doubling trades \cite{TaoEntropy,GreenMannersTao}.  They are not needed as hypotheses here.  The Boolean-cube growth theorem gives a direct contradiction to a large small-doubling subset of a primitive Boolean support fiber.
\end{remark}


\section{Entropy inverse theory and its role}\label{sec:entropy-role}

The terminology ``Tao--Gowers method'' can refer to two related but
distinct steps.  The first is BSG, used directly in
\cref{prop:high-energy-impossible}.  The second is Freiman-type structural
theory, meaning inverse theorems that model a set with small sumset or
difference set by a generalized arithmetic progression or a coset
progression.  It is not needed for the Boolean contradiction, but it can suggest
candidate geometric strata in more general bounded-alphabet residuals;
those strata still require direct payment proofs.

A \emph{bounded integer trade} is a vector \(Y\in\Z^T\) whose
coordinates are bounded independently of \(N\), with \(\Phi(Y)=0\) and,
on a fixed-weight slice, \(\sum_tY_t=0\).  If \(Y\) is random, \(Y'\) is
an independent copy, and
\(\Ent(Y-Y')\le \Ent(Y)+o(N)\), then entropy BSG and entropy Freiman
theory produce, after conditioning and losing \(o(N)\) entropy, a set
model with doubling \(e^{o(N)}\).  Tao's entropy inverse theorem gives
this in \emph{coset-progression} language---a subgroup coset plus a
generalized arithmetic progression---for Shannon entropy \cite{TaoEntropy};
Green--Manners--Tao recast the theory in terms of entropic doubling and
Ruzsa distance \cite{GreenMannersTao}; Goh formulates an entropic analogue
of additive energy and the entropic BSG connection \cite{GohEnergy}.  In a
torsion model one may also use polynomial Freiman--Ruzsa technology after
reducing bounded integer trades modulo a large prime with no wraparound.

In the present paper the Boolean-cube difference theorem makes this
structural compression unnecessary for the primitive Q step.  Once BSG
finds \(A'\subseteq F_s\subseteq\{0,1\}^T\) with
\(\abs{A'-A'}\le e^{o(N)}\abs{A'}\), Boolean-cube difference growth
already forces \(\abs{A'}=e^{o(N)}\).  This contradicts the exponential
size supplied by moment excess.  Therefore there is no separate
entropy-Freiman hypothesis in the proof.

For non-Boolean residual trades, entropy inverse theory can suggest
coset-progression incidence strata or rank-defect tests.  This is a useful
heuristic for constructing a row-specific atlas, not an exhaustive
dichotomy: each suggested stratum still needs the profile census and
slope-projection payment in \textup{(A2)}, and arbitrary weighted charts
are not covered by \cref{lem:vandermonde}.  No such heuristic is used in
the proof of the Boolean primitive theorem.

\section{Primitive Q after the Sidon cut}\label{sec:primitive-q-proof}

The key primitive theorem is now short and self-contained modulo \cref{thm:bsg,thm:quasicube}.  It also addresses the main obstruction: ordinary \Renyi{} mass can be Sidon-heavy, but after that cell is removed it must be high-energy, and high-energy is impossible in a large Boolean fiber.

\begin{proposition}[High-energy Boolean fibers are small]\label{prop:high-energy-impossible}
Let \(F\subseteq\{0,1\}^T\), \(\abs T=N\).  If
\(E(F)\ge \abs F^3/K\), then \(\abs F\le K^{O(1)}\).  Consequently an
exponential set \(\abs F\ge e^{cN-o(N)}\) cannot have
\(E(F)\ge\abs F^3/e^{o(N)}\).
\end{proposition}

\begin{proof}
Apply \cref{thm:bsg}.  There is \(F'\subseteq F\) with \(\abs{F'}\ge K^{-C}\abs F\) and \(\abs{F'-F'}\le K^C\abs{F'}\).  Since \(F'\subseteq\{0,1\}^T\), \cref{thm:quasicube} gives \(\abs{F'-F'}\ge \abs{F'}^{3/2}\).  Hence \(\abs{F'}^{1/2}\le K^C\), so \(\abs F\le K^{3C}\) after increasing \(C\).  Taking \(K=e^{o(N)}\) rules out \(\abs F\ge e^{cN-o(N)}\).
\end{proof}

\begin{theorem}[Primitive Q after a Sidon moment payment]\label{thm:primitive-q}
Let \(\Omega_N^0\subseteq\{0,1\}^{T_N}\) be fixed-density full profile
slices, with \(\abs{T_N}=N\) and
\(m_N/N\in[\alpha,1-\alpha]\), and let
\(\Omc_N\subseteq\Omega_N^0\) be primitive first-match residuals.  Suppose
there is a logarithmic moment
order \(q_N\to\infty\) with
\(\log L_N/q_N=o(N)\), where
\(L_N=\abs{\Phi_N(\Omega_N^0)}\) and
\(\barN_N=\abs{\Omega_N^0}/L_N\), for which
\(\Gsid_{q_N,\sigma}\le e^{o(Nq_N)}\) for every fixed \(\sigma>0\).
Then the leaves satisfy primitive Q:
\[
        \max_s\abs{\Omc_N\cap\Phi_N^{-1}(s)}
        \le e^{o(N)}\barN_N .
\]
\end{theorem}

\begin{proof}
Assume primitive Q fails at positive exponential rate.  The lower half of
the moment sandwich gives
\(\Gord_{q_N}\ge e^{\eta Nq_N-o(Nq_N)}\) along a subsequence.  The Sidon
moment hypothesis and \cref{prop:ordinary-moment-split} then give a fiber
\(F_s\) with \(f_s\ge e^{\eta N-o(N)}\barN_N\) and
\(E(F_s)\ge f_s^3e^{-o(N)}\).  Since
\(\barN_N=\abs{\Omega_N^0}/L_N\ge1\), this fiber has size
\(e^{\eta N-o(N)}\), contradicting
\cref{prop:high-energy-impossible} with \(K=e^{o(N)}\).  Thus no
positive-rate max-fiber failure occurs.
\end{proof}

\section{Q, SP, and the MCA compiler}\label{sec:compiler}

The \emph{Q branch} is the maximum-prefix-fiber problem.  The \emph{SP
branch} (support-pair, or shift-pair, branch) is its second-moment problem:
it counts ordered pairs of supports in one prefix fiber.  These names label
two counting interfaces and carry no assertion that either bound has
already been proved.  Neither count is automatically the MCA numerator: a
bad slope begins with one support, while converting support pairs to
projective slope directions requires a proved multiplicity or incidence
bound.  We first record the exact moment identities and then state that
separate ray compiler.

\begin{definition}[Q and SP ledgers]\label{def:q-sp}
For a first-match residual profile \(\lambda\), let \(\Omega^0_\lambda\)
be its full profile slice, let
\(\Omc_\lambda\subseteq\Omega^0_\lambda\) be the residual, and put
\[
 \Scal_\lambda=\Phi_\lambda(\Omega^0_\lambda),\qquad
 L_\lambda=\abs{\Scal_\lambda},\qquad
 \barN_\lambda=\abs{\Omega^0_\lambda}/L_\lambda,
 \qquad
 Q_\lambda(s)=\abs{\Omc_\lambda\cap\Phi_\lambda^{-1}(s)}.
\]
The Q ledger is discharged if
\(\max_sQ_\lambda(s)\le e^{o(n)}\barN_\lambda\) for every primitive
profile.  The support-pair SP statistic is
\(\sum_sQ_\lambda(s)^2\).  A distinct-ray ledger is discharged only by a
direct slope bound or by \textup{(RC)} below.
If the ambient target count \(\abs{\B_\lambda}^{R_\lambda}\) is used in
place of \(L_\lambda\), the ledger must also record the full-image
certificate \eqref{eq:full-image-certificate} or a direct ambient theorem.
\end{definition}


The compiler uses two elementary inequalities repeatedly.  They are included to make explicit how the logarithmic moment, max-fiber Q, and shift-pair second moment interact.

\begin{lemma}[Moment sandwich]\label{lem:moment-sandwich-full}
Let \(N(s)\ge0\) be fiber sizes over a finite syndrome set of size \(L\), let \(\barN=L^{-1}\sum_sN(s)\), and put \(R(s)=N(s)/\barN\).  For \(r\ge2\), define \(\Gamma_r=L^{-1}\sum_sR(s)^r\).  Then
\[
        \max_s R(s)\le (L\Gamma_r)^{1/r}.
\]
Conversely, if \(\max_sR(s)\le M\), then \(\Gamma_r\le M^{r-1}\).
\end{lemma}

\begin{proof}
The first inequality follows from \(\max_sR(s)^r\le\sum_sR(s)^r=L\Gamma_r\).  For the converse, \(R(s)^r\le M^{r-1}R(s)\), and \(L^{-1}\sum_sR(s)=1\).
\end{proof}

\begin{corollary}[Finite moment floor]\label{cor:finite-moment-floor}
Suppose a finite row has remaining Q bit margin \(\Delta\), meaning that one needs \(\max_sR(s)\le2^\Delta\).  A moment proof of order \(r\) can certify this only if
\[
        \log_2\Gamma_r+\log_2 L\le r\Delta,
\]
where \(L\) is the number of realized targets in the selected
normalization.  In an ambient prefix proof this specializes to
\(\log_2\Gamma_r^{\rm amb}+R\log_2\abs\B\le r\Delta\); transferring an
image-normalized moment to that statement additionally costs
\((r-1)\log_2(\abs\B^R/L)\).
\end{corollary}

\begin{proof}
The moment sandwich gives \(\log_2\max_sR(s)\le(\log_2L+\log_2\Gamma_r)/r\).  To make this at most \(\Delta\) one needs the displayed inequality.
\end{proof}

The finite corollary explains why a small adjacent margin cannot be closed
by a fixed-order moment when \(R\log\abs\B\) is large.

\begin{definition}[Explanation image and retained-support occupancy]
\label{def:explanation-occupancy}
Fix \(a\ge k+1\), a received pair \(r=(r_0,r_1)\), and a cell
\(\Ccal\subseteq\Wcal_a(r)\).  Its \emph{explanation image} and
\emph{slope image} are
\[
 \Rcal(\Ccal)=\{(\gamma,h):\exists S\ (\gamma,S,h)\in\Ccal\},
 \qquad Z(\Ccal)=\{\gamma:\exists h\ (\gamma,h)\in\Rcal(\Ccal)\}.
\]
For \(\rho=(\gamma,h)\in\Rcal(\Ccal)\), its
\emph{retained-support occupancy} is
\[
 \nu_{\Ccal}(\rho)
   =\abs{\{S\in\tbinom Da:(\gamma,S,h)\in\Ccal\}}.
\]
Thus the witness projection factors canonically as
\[
 \Ccal\longrightarrow\Rcal(\Ccal)\longrightarrow Z(\Ccal),
 \qquad (\gamma,S,h)\longmapsto(\gamma,h)\longmapsto\gamma.
\]
The first map forgets an exact-agreement support; the second forgets the
explaining polynomial.  Distinct explanation states may have the same
slope, so neither map is silently treated as a bijection.  A
\emph{ray compiler} means a proved estimate for the final slope image,
obtained from these actual fibers or from a direct incidence bound.
\end{definition}

\begin{lemma}[Exact occupancy compiler]
\label{lem:exact-occupancy-compiler}
With the notation above,
\[
 \abs{\Ccal}=\sum_{\rho\in\Rcal(\Ccal)}\nu_{\Ccal}(\rho),
 \qquad
 \abs{\Rcal(\Ccal)}
   =\sum_{w\in\Ccal}
       \frac1{\nu_{\Ccal}(\pi_{\rm exp}(w))},
 \qquad
 \abs{Z(\Ccal)}\le\abs{\Rcal(\Ccal)}.             \tag{RC$_{\rm occ}$}
\]
Here \(\pi_{\rm exp}(\gamma,S,h)=(\gamma,h)\).  In particular, if every
retained explanation state has occupancy at least \(H\ge1\), then
\[
        \abs{Z(\Ccal)}\le
        \floor{\frac{\abs{\Ccal}}H}.                \tag{RC1}
\]
Also, because one noncommon exact-\(a\) support carries at most one slope
and the explaining polynomial on \(a\ge k+1\) points is unique,
\[
 \abs{Z(\Ccal)}\le
 \abs{\{S:\exists\gamma,h\ (\gamma,S,h)\in\Ccal\}}. \tag{RC2}
\]
All these statements are exact finite-row identities or inequalities.
\end{lemma}

\begin{proof}
Partition \(\Ccal\) by the fibers of \(\pi_{\rm exp}\).  This gives the
first identity; summing the reciprocal of the fiber size over each fiber
gives the second.  Projection \((\gamma,h)\mapsto\gamma\) can only
decrease cardinality.  If every first fiber has size at least \(H\),
\textup{(RC1)} follows.  Finally, if the same noncommon support carried two
distinct slopes, slope elimination would simultaneously explain the
received pair on that support.  For fixed support and slope, uniqueness of
interpolation gives at most one \(h\).  This proves \textup{(RC2)}.
\end{proof}

\begin{lemma}[Max-occupancy add-back]
\label{lem:max-occupancy-addback}
Let \(\Ccal\subseteq\Wcal_a(r)\) be partitioned into profile cells
\(\Ccal=\coprod_{\lambda\in\Lambda}\Ccal_\lambda\), and put
\(P=\abs\Lambda\).  For every explanation state
\(\rho=(\gamma,h)\in\Rcal(\Ccal)\), let
\[
 M(\rho)=\abs{\{S:(\gamma,S,h)\in\Ccal\}}.
\]
Assign \(\rho\) to a profile \(\lambda(\rho)\) on which
\(\nu_{\Ccal_\lambda}(\rho)\) is maximal, and let
\(\Rcal_\lambda^\star\) be the assigned states.  Then
\[
 \nu_{\Ccal_\lambda}(\rho)\ge\ceil{M(\rho)/P}
 \qquad(\rho\in\Rcal_\lambda^\star).               \tag{RC$_{\rm max}$1}
\]
Consequently, if \(M(\rho)\ge H\) for every retained state, then
\[
 \abs{Z(\Ccal)}
 \le\sum_{\lambda\in\Lambda}\abs{\Rcal_\lambda^\star}
 \le\sum_{\lambda\in\Lambda}
   \floor{\frac{\abs{\Ccal_\lambda^\star}}
                   {\ceil{H/P}}},                  \tag{RC$_{\rm max}$2}
\]
where \(\Ccal_\lambda^\star\) consists of the witnesses in
\(\Ccal_\lambda\) lying above assigned states.  Call \(\Ccal\)
\emph{explanation-state complete} if, whenever
\((\gamma,h)\in\Rcal(\Ccal)\), every exact-\(a\) noncommon witness
\((\gamma,S,h)\in\Wcal_a(r)\) also belongs to \(\Ccal\).  If \(\Ccal\)
is explanation-state complete, every retained state has exact agreement
\(b\ge a\), and the received pair is not simultaneously explained on any
\(a\)-subset of that agreement set, then one may take \(H=\binom ba\).
\end{lemma}

\begin{proof}
The \(P\) profile occupancies of a fixed state sum to \(M(\rho)\), so
their maximum is at least \(\ceil{M(\rho)/P}\).  Apply
\cref{lem:exact-occupancy-compiler} on each assigned cell and sum.  Under
the final hypotheses, explanation-state completeness retains every
\(a\)-subset of the \(b\)-point agreement set, and every such subset is
noncommon.  Hence \(M(\rho)=\binom ba\).
\end{proof}

\begin{definition}[Support interpolation functionals]
\label{def:pair-to-ray-map}
Let \(T\subseteq D\), \(\abs T=m\ge k\), put \(w=m-k\), and let
\(Q_T(X)=\prod_{x\in T}(X-x)\).  For a word \(u\), define
\[
 \mathcal A_T(u)=
 \left(
   \sum_{x\in T}\frac{u(x)x^r}{Q_T'(x)}
 \right)_{0\le r<w}\in\F^w.
 \tag{12.1}\label{eq:interpolation-functional}
\]
\end{definition}

\begin{lemma}[Slope elimination]\label{lem:saturation-quotient-rays}
A word \(u\) is explained by \(\RS_\F(D,k)\) on \(T\) if and only if
\(\mathcal A_T(u)=0\).  Hence \(u_0+\gamma u_1\) is explained on \(T\)
if and only if
\[
 \mathcal A_T(u_0)+\gamma\mathcal A_T(u_1)=0.
 \tag{12.2}\label{eq:slope-elimination}
\]
If \(T\) is not a common-support explanation, it supports at most one
finite slope.
\end{lemma}

\begin{proof}
The vectors \((Q_T'(x)^{-1}x^r)_{x\in T}\), \(0\le r<w\), form a
parity-check matrix for \(\RS_\F(T,k)\), proving the first assertion.  The
second is linearity.  If two distinct slopes solved
\eqref{eq:slope-elimination}, both functionals would vanish and the pair
would be explained on \(T\).
\end{proof}

\begin{lemma}[Second moment identity for primitive pairs]\label{lem:second-moment-identity}
For any prefix profile \(\lambda\), the number of ordered pairs of residual
supports with common prefix syndrome is exactly
\[
        \sum_s Q_\lambda(s)^2,
        \qquad
        Q_\lambda(s)=\abs{\Omc_\lambda\cap\Phi_\lambda^{-1}(s)}.
\]
and it is bounded by
\[
        \bigl(\max_sQ_\lambda(s)\bigr)\sum_sQ_\lambda(s).
\]
\end{lemma}

\begin{proof}
Partition the residual supports by their prefix value.  Ordered pairs in one
part give \(Q_\lambda(s)^2\), and summing proves the identity.  The final
inequality is elementary.  No assertion about MCA rays is made.
\end{proof}

Q therefore controls a support-pair statistic.  It does not justify dividing
that statistic by \(\abs{\Omega^0_\lambda}\): a ray can have a single
representative at exact agreement.

\begin{proposition}[Q and SP do not determine rays]
\label{prop:q-sp-no-ray}
Let \(\Phi:\Omega\to\Sigma\) be any support-profile map with
\(\Omega\ne\varnothing\), put
\(Q(s)=\abs{\Phi^{-1}(s)}\), and
\({\rm SP}=\sum_sQ(s)^2\).  Regard \(\Omega\) as an abstract set of
candidate noncommon supports, and suppose the only retained
slope-incidence axiom is that each support carries at most one slope.  For
every nonempty finite slope set \(\Gamma\) and every
\(1\le M\le\min\{\abs\Omega,\abs\Gamma\}\), within the class of abstract
incidence structures obeying these retained data there is a partial
support-to-slope map having exactly \(M\) distinct values.  There is also
a one-value map.  Hence no ray bound can be deduced from Q, SP, and that
axiom alone; an RS-specific bound must use additional algebraic incidence
information.
\end{proposition}

\begin{proof}
Choose \(M\) distinct supports and \(M\) distinct slopes and map them
bijectively, leaving all other supports unlabeled.  This changes neither
\(\Phi\), its fibers, nor SP, and it respects one-slope-per-support.
Mapping the same chosen supports to one slope gives identical Q and SP
data with a one-ray image.  The missing datum is the image incidence, not
another moment of the prefix fibers.
\end{proof}

\begin{proposition}[Exact pair-to-ray multiplicity requirement]
\label{prop:pair-ray-multiplicity}
Let
\(\mathcal P=\{(A,B)\in\Omega^2:\Phi(A)=\Phi(B)\}\), so
\(\abs{\mathcal P}={\rm SP}\).  Let \(Z\) be a bad-slope set and
\(I\subseteq Z\times\mathcal P\) an incidence relation, and let
\(H,J\ge1\).  If every
\(\gamma\in Z\) is incident with at least \(H\) pairs and every pair is
incident with at most \(J\) slopes, then
\[
 \abs Z\le\frac{J\,{\rm SP}}H
       \le\frac{J(\max_sQ(s))\abs\Omega}{H}.
 \tag{12.3}\label{eq:pair-ray-multiplicity}
\]
Consequently this pair-count argument yields a ray estimate at the same
profile scale provided \(J\abs\Omega/H=e^{o(n)}\); alternatively one may
prove a direct transverse-secant estimate.  Merely proving that every ray
has one representing pair is insufficient.
\end{proposition}

\begin{proof}
Double counting gives
\(H\abs Z\le\abs I\le J\abs{\mathcal P}=J\,{\rm SP}\).  The second
inequality follows from
\({\rm SP}\le(\max_sQ(s))\sum_sQ(s)=(\max_sQ(s))\abs\Omega\).
\end{proof}

\begin{corollary}[Exact finite pair-to-ray constant]
\label{cor:finite-pair-ray-compiler}
For a residual profile \(\lambda\), put
\(m_\lambda=\sum_sQ_\lambda(s)\),
\(L_\lambda=\abs{\Phi_\lambda(\Omega^0_\lambda)}\), and
\(\barN_\lambda=\abs{\Omega^0_\lambda}/L_\lambda\).  For an integer
\(q\ge2\), define
\[
 \Gamma_{q,\lambda}=L_\lambda^{-1}\sum_s
 \left(\frac{Q_\lambda(s)}{\barN_\lambda}\right)^q.
\]
If the pair-to-ray incidence has lower ray degree \(H_\lambda\) and upper
pair degree \(J_\lambda\), then
\[
 \abs{Z_\lambda^\circ}\le
 \floor{\frac{J_\lambda m_\lambda}{H_\lambda}\,
  \barN_\lambda
  \bigl(L_\lambda\Gamma_{q,\lambda}\bigr)^{1/q}}.   \tag{RC$_{\rm pair}$}
\]
Thus the finite compiler uses the actual residual mass and literal
incidence degrees; it contains no unnamed polynomial or asymptotic loss.
\end{corollary}

\begin{proof}
The moment sandwich gives
\(\max_sQ_\lambda(s)\le\barN_\lambda
(L_\lambda\Gamma_{q,\lambda})^{1/q}\).  Hence
\(\sum_sQ_\lambda(s)^2\le m_\lambda\max_sQ_\lambda(s)\).
Apply \cref{prop:pair-ray-multiplicity} and take the integer floor.
\end{proof}

\begin{proposition}[Curve-degree ray compiler]
\label{prop:curve-degree-ray-compiler}
Let \(V\) be a finite-dimensional space of locator polynomials and, for
\(x\in D\), let
\(H_x=\{[L]\in\Pj(V):L(x)=0\}\) be its \emph{evaluation hyperplane}.
Suppose a residual slope set \(Z\) admits an injective assignment into a
disjoint union of projective integral curves
\[
              Z\hookrightarrow\coprod_{i=1}^sX_i(\F),
\]
where ``integral curve'' means an irreducible reduced one-dimensional
projective variety, together with locator morphisms
\(\psi_i:X_i\to\Pj(V)\).  Put
\[
 \delta_i=\deg\psi_i^*\mathcal O_{\Pj(V)}(1),
\]
the degree on \(X_i\) of the pullback of a projective hyperplane.  For
every \(\gamma\in Z\), if its assigned point is \(y\in X_i(\F)\), require
\(\psi_i(y)=[L_{\gamma,y}]\), where \(L_{\gamma,y}\) is a locator from an
actual residual witness to \(\gamma\).  For each \(i\), put
\[
 G_i=\{x\in D:\psi_i(X_i)\subseteq H_x\},
 \qquad D_i^{\rm mov}=D\setminus G_i.
\]
Assume the locator attached to every chosen representative on \(X_i\)
has at least \(h_i\ge1\) roots in \(D_i^{\rm mov}\).  Then
\[
 \abs Z\le
 \floor{\sum_{i=1}^s
       \frac{\abs{D_i^{\rm mov}}\delta_i}{h_i}}
 \le\floor{n\sum_{i=1}^s\frac{\delta_i}{h_i}}.
 \tag{RC$_{\rm curve}$}
\]
In particular, a cover of subexponential total locator degree gives a
subexponential ray budget when the \(h_i\) have no exponential loss.  The
standard degree-one pencil \([s:t]\mapsto[sA+tB]\) is the special case
\(X=\Pj^1\) and \(\delta=1\).
\end{proposition}

\begin{proof}
Choose one representative point for each slope.  For fixed \(i\) and
\(x\in D_i^{\rm mov}\), the pullback of \(H_x\) is a nonzero effective
divisor of degree \(\delta_i\), and hence contains at most \(\delta_i\)
distinct rational points.  Each chosen point on \(X_i\) supplies at least
\(h_i\) moving-root incidences.  Thus the number assigned to \(X_i\) is
at most \(\abs{D_i^{\rm mov}}\delta_i/h_i\).  Sum over \(i\) and take
the integer floor.
\end{proof}

\begin{corollary}[Automatic exact-support locator-curve compiler]
\label{cor:automatic-locator-curve-ray}
Fix a received line and exact agreement \(a\ge k+1\).  For each slope in
a residual set \(Z\), choose one noncommon witness support
\(S_\gamma\in\binom Da\) and its locator \(Q_{S_\gamma}\).  Suppose the
locator points are covered by
\(\psi_i(X_i(\F))\), where
\(\psi_i:X_i\to\Pj(V)\) are nonconstant morphisms from projective
integral curves.  Put
\[
 \delta_i=\deg\psi_i^*\mathcal O(1),\qquad
 G_i=\{x\in D:\psi_i(X_i)\subseteq H_x\},
 \qquad g_i=\abs{G_i}.
\]
Then
\[
 \abs Z\le
 \#\{i:g_i=a\}+
 \floor{\sum_{i:g_i<a}\frac{(n-g_i)\delta_i}{a-g_i}}. \tag{RC$_{\rm loc}$}
\]
No separate slope-to-curve injectivity hypothesis is required.  A
component with \(g_i=a\) contains at most one selected degree-\(a\)
locator, because its \(a\) common roots determine the monic support
locator uniquely; this accounts for the first term.  Components with
\(g_i>a\) contain no selected degree-\(a\) locator and are discarded.
\end{corollary}

\begin{proof}
Two distinct slopes cannot use the same chosen noncommon support.  Since a
monic squarefree locator determines its root set, their projective locator
points are distinct as well.  Assign each locator to the first covering
component and choose one rational preimage; this is the injective
assignment required by \cref{prop:curve-degree-ray-compiler}.  On
component \(i\), every selected degree-\(a\) locator has the \(g_i\)
common roots and at least \(a-g_i\) roots outside them.  If \(g_i=a\),
there is only one such monic locator.  On the other components apply
\textup{(RC$_{\rm curve}$)} with \(h_i=a-g_i\).
\end{proof}

\begin{corollary}[Automatic exact-support one-pencil compiler]
\label{cor:automatic-one-pencil-ray}
In the setting above, suppose all selected locators lie in the projective
pencil spanned by \(A,B\), and let \(g<a\) be the number of common roots
of \(A,B\) in \(D\).  Then
\[
                         \abs Z\le\floor{(n-g)/(a-g)}. \tag{RC$_{\rm pen}$}
\]
The same argument applies to residual locators of degree \(e\) whose roots
lie in a fixed active domain \(T\subseteq D\), provided the fixed profile
data together with the residual locator determine the exact support.  If
\(N=\abs T\) and \(g_T\) is the number of common roots of \(A,B\) in
\(T\), then
\[
                         \abs Z\le\floor{(N-g_T)/(e-g_T)}.
\]
\end{corollary}

\begin{proof}
The exact-support locator assignment is injective by the preceding proof.
Apply \cref{lem:pencil-payment}, or the curve corollary with
\(X=\Pj^1\) and \(\delta=1\).
\end{proof}

By \cref{thm:syndrome-secant-exact}, the genuinely missing invariant can
be taken to be \(\Theta_t(y_0,y_1)\).  Equivalently one may construct a
higher-dimensional split-pencil \emph{eliminant}, meaning a polynomial in
the slope parameter obtained after eliminating locator and support
variables, whose degree bounds that
transverse secant incidence.  By
\cref{lem:independent-union-rays}, only the dependent-union, or
maximum-distance-separable (MDS) circuit,
strata can support more than \(t+1\) rays; this localizes the remaining
geometry.  One fixed circuit is bounded by
\cref{thm:single-mds-circuit-ray}, but an uncontrolled family of circuit
strata or a higher-dimensional balanced core is not thereby paid.

Here an \emph{MDS-circuit stratum} is a chart on which the relevant union
of generalized Vandermonde parity-check columns has a minimal linear
dependence.  Equivalently, deleting any one column restores independence.
The name identifies the only dependence pattern left by the MDS property;
it does not by itself place all witnesses in one fixed circuit or provide
a count for an unrestricted union of such strata.

The next hypothesis packages exactly this missing conversion.  Its number
\(H_\lambda\) is a lower bound on how many certified support pairs represent
one bad ray, while \(J_\lambda\) is an upper bound on how many bad rays one
pair can represent.  The ratio in \textup{(RC)} is what makes the resulting
double count stay at the profile scale.

\begin{hypothesis}[Ray compiler]\label{hyp:ray-compiler}
For every received line and primitive residual profile \(\lambda\), let
\(Z_\lambda^\circ\) be its actual first-match set of unpaid bad slopes and
put
\[
\Pcal_\lambda=\{(A,B)\in(\Omc_\lambda)^2:
                       \Phi_\lambda(A)=\Phi_\lambda(B)\}.
\]
Put \(m_\lambda=\abs{\Omc_\lambda}=\sum_sQ_\lambda(s)\).
The row satisfies \textup{(RC)} on this profile if either it has the direct
bound \(\abs{Z_\lambda^\circ}\le e^{o(n)}(1+\barN_\lambda)\), or there are
an incidence
\(I_\lambda\subseteq Z_\lambda^\circ\times\Pcal_\lambda\) and numbers
\(H_\lambda,J_\lambda\ge1\) such that
\[
 \deg_{I_\lambda}(\gamma)\ge H_\lambda,\qquad
 \deg_{I_\lambda}(A,B)\le J_\lambda,\qquad
 \frac{J_\lambda m_\lambda}{H_\lambda}=e^{o(n)}.
 \tag{RC}\label{eq:ray-compiler}
\]
All bounds are uniform in the received line and realized profile.  The
incidence must arise from the RS witness geometry; its existence is not a
consequence of Q or SP.  In particular, one representing pair per ray does
not supply the required lower degree.
\end{hypothesis}

\begin{proposition}[Q plus RC implies rays]\label{prop:q-implies-sp}
If primitive Q is discharged and \textup{(RC)} holds on every residual
profile, then all primitive ray ledgers are discharged
at total cost \(e^{o(n)}\mathfrak E_n\).
\end{proposition}

\begin{proof}
The boundary prefix fibers are bounded by Q.  For every remaining unpaid
slope choose one noncommon support, as permitted by
\cref{lem:saturation-quotient-rays}; \textup{(RC)} bounds their distinct slope set by
\(e^{o(n)}\barN_\lambda\).  Sum over the subexponential profile set and use
\eqref{eq:profile-envelope}.
\end{proof}

\begin{proposition}[Primitive residual numerator]\label{prop:primitive-residual-numerator}
For a ledger-admissible sequence, the total primitive contribution is at
most \(e^{o(n)}\mathfrak E_n\).
\end{proposition}

\begin{proof}
The certified Sidon payment and \cref{thm:primitive-q} give Q on the
primitive leaves.  The profile sum is part of \(\mathfrak E_n\), and
\cref{prop:q-implies-sp} applies using \textup{(RC)}.
\end{proof}

\begin{proposition}[Asymptotic MCA numerator bound]\label{prop:numerator-bound}
For a ledger-admissible sequence at agreement \(a_n\),
\[
        B_{C_n}^{\rm MCA}(a_n)
        \le e^{o(n)}\mathfrak E_n(a_n).
\]
\end{proposition}

\begin{proof}
Apply the first-match ledger.  The algebraic and direct-ray cells are paid by
\textup{(A2)}, the Fourier/Sidon cell is certified by \textup{(A4)}, and the remaining primitive
part is \cref{prop:primitive-residual-numerator}.  The first-match bound
adds the distinct-slope budgets.
\end{proof}

This is also a detailed proof of \cref{thm:main-smooth-circle}; a closed
ledger is the theorem-verified special case.  The threshold statement uses
the entropy calculation and the lower construction.

\section{Entropy bookkeeping and the lower side}\label{sec:frontier}

Let \(\rho_n=k_n/n\), \(K_n=k_n+1\),
\(\beta_n=\log_2\abs{\B_n}\), and write
\(a_n=K_n+g_nn+O(1)\).  Stirling's formula gives
\[
 \log_2\barN_1(a_n)
 =n\bigl(H_2(\rho_n+g_n)-\beta_ng_n\bigr)
  +O(\log n+\beta_n),
 \tag{13.1}\label{eq:target-entropy}
\]
where \(H_2\) is the bit-valued entropy fixed in the Introduction.  Let
\(\varnothing\ne\Gamma_n\subseteq\F_n\), let \(0\le\eps_n^*\le1\), and
put \(B_n^*=\floor{\eps_n^*\abs{\Gamma_n}}\).  Write \(\ell_n=o(n)\) for
the ledger overhead in bits, so that its certified bound is
\(2^{\ell_n}\mathfrak E_n\).  Agreement \(a\) is \emph{safe} if
\(B_{C_n,\Gamma_n}^{\rm MCA}(a)\le B_n^*\) and unsafe otherwise.  Since
the numerator is nonincreasing in \(a\), when a safe agreement exists its
integer threshold is
\[
 a_n^*=\min\{a\in\{k_n+1,\ldots,n\}:
 B_{C_n,\Gamma_n}^{\rm MCA}(a)\le B_n^*\}.
\]
A target crossing is a comparison with \(\log_2B_n^*\), not automatically
the zero of the entropy expression.  Safety at \(a_n\ge k_n+1\) is
impossible when
\(B_n^*<\min\{\abs{\Gamma_n},n-a_n+1\}\) by
\cref{prop:universal-tangent-floor}.

We call the target \emph{viable at agreement \(a\)} when it is not already
below this universal floor, namely when
\(B_n^*\ge\min\{\abs{\Gamma_n},n-a+1\}\).  Viability is only a necessary
condition for safety; it is not an upper bound.  A \emph{target reserve}
is the explicit bit slack in one of the two comparisons: on the safe side
it is the amount by which \(\log_2 B_n^*\) exceeds
\(\ell_n+\log_2\mathfrak E_n(a)\), and on the unsafe side it is the amount
by which the logarithm of a proved bad-slope construction exceeds
\(\log_2B_n^*\).  An \emph{agreement reserve} \(s_n\) is a number of
integer agreement steps taken past a candidate crossing; away from
tangency (that is, when the derivative at the crossing is nonzero), the
derivative of the entropy exponent converts those steps into
the corresponding bit reserve.

For a concave entropy level function, the superlevel set can have two
endpoints.  Its \emph{right crossing} is the larger, equivalently the
supremum of the superlevel set.  It is an \emph{interior right crossing}
when it lies strictly inside \([0,1-\rho_n]\) and satisfies the level
equation there.  If its derivative is nonzero, the crossing is transverse;
if the derivative vanishes, it is tangential and a linear agreement reserve
need not produce a linear bit gain.

\begin{lemma}[Target-aware safe side]\label{lem:safe-side}
Under a closed ledger, agreement \(a_n\) is safe at target \(B_n^*\) as
soon as
\[
        2^{\ell_n}\mathfrak E_n(a_n)\le B_n^*.
        \tag{13.2}\label{eq:exact-safe-budget}
\]
If \(\mathfrak E_n(a)=e^{o(n)}\barN_1(a)\) in a window and
\(\log(1+B_n^*)=o(n)\), the leading crossing in that window is the zero of
\(H_2(\rho_n+g)-\beta_ng\), with a reserve large enough to dominate
\(\ell_n+O(\log n+\beta_n)+\log_2(1+B_n^*)\).  The zero-crossing
description additionally requires target viability
\(B_n^*\ge\min\{\abs{\Gamma_n},n-a_n+1\}\), and in particular
\(B_n^*\ge1\).
\end{lemma}

\begin{proof}
The first assertion is \cref{prop:numerator-bound}.  For the second, use
\eqref{eq:target-entropy}.  Away from a tangential endpoint, an agreement
reserve \(s_n\) changes the entropy by
\(s_n\abs{H_2'(\rho_n+g)-\beta_n}+o(s_n)\).  Requiring this quantity to
dominate the written logarithmic costs gives \eqref{eq:exact-safe-budget}.
This formulation remains valid when \(\beta_n\) varies; the phrase
``\(o(n)\) corridor'' alone is not quantitative enough.
\end{proof}

\begin{proposition}[Exact unsafe test]\label{prop:simple-pole-lower}
Let
\[
 L_n=\left\lceil
 \binom n{a_n}\abs{\B_n}^{-(a_n-k_n-1)}
 \right\rceil,
 \qquad q_n=\abs{\F_n}>n.
\]
If
\[
 \left\lceil \frac{\abs{\Gamma_n}}{q_n}
 \left\lceil
 \frac{L_n(q_n-n)}{q_n-n+k_n(L_n-1)}
 \right\rceil\right\rceil>B_n^*,
 \tag{13.3}\label{eq:exact-unsafe-budget}
\]
then agreement \(a_n\) is unsafe at the target.  The same statement holds
with \(L_n\) replaced by any larger identity, quotient, Chebyshev, or
remainder-profile list proved for the dimension-\(k_n+1\) code.
\end{proposition}

\begin{proof}
The identity list is \cref{prop:exact-prefix-list}, and
\cref{thm:collision-aware-pole} produces a bad-slope set \(Z\subseteq\F_n\)
of the inner-ceiling size.  Replacing a line \((r_0,r_1)\) by
\((r_0+\delta r_1,r_1)\) translates \(Z\) by \(-\delta\).  Averaging over
\(\delta\in\F_n\) gives \(\abs Z\abs{\Gamma_n}/q_n\) challenge slopes, so
some translate has at least the outer-ceiling number.  The profile variants
use the same conversion after their list sizes have been established.
\end{proof}

\begin{theorem}[Unconditional exact support-envelope bracket]
\label{thm:unconditional-support-envelope-bracket}
Let \(C=\RS_\F(D,k)\), let \(D\subseteq\B\subseteq\F\), put
\(q=\abs\F>n\), fix \(\varnothing\ne\Gamma\subseteq\F\), and let
\(B^*=\floor{\eps\abs\Gamma}\).  For \(k+1\le a\le n\), define
\[
 \begin{split}
 L(a)&=\ceil{\binom na\abs\B^{-(a-k-1)}},\\
 P(a)&=\left\lceil\frac{\abs\Gamma}{q}
       \left\lceil\frac{L(a)(q-n)}
       {q-n+k(L(a)-1)}\right\rceil\right\rceil,\\
 U(a)&=\min\left\{\abs\Gamma,\binom na\right\}.
 \end{split}                                       \tag{SB1}
\]
If \(a_-<a_+\) satisfy
\[
                         P(a_-)>B^*,\qquad U(a_+)\le B^*, \tag{SB2}
\]
then the global target threshold exists and
\[
                         a_-<a^*\le a_+.            \tag{SB3}
\]
If \(a_+=a_-+1\), then \(a^*=a_+\) exactly.  In particular,
\[
                  \binom n{k+1}\le B^*\quad\Longrightarrow\quad
                  a^*=k+1.                         \tag{SB4}
\]
These are unconditional finite comparisons; \textup{(SB2)} is the
literal target reserve and contains no asymptotic placeholder.
\end{theorem}

\begin{proof}
The exact prefix list and pole-line construction give
\(B_{C,\Gamma}^{\rm MCA}(a)\ge P(a)\) by
\cref{prop:simple-pole-lower}.  The exact support atlas gives
\(B_{C,\Gamma}^{\rm MCA}(a)\le U(a)\) by
\cref{prop:exact-support-upper}.  Thus \(a_-\) is unsafe and \(a_+\) is
safe; monotonicity proves \textup{(SB3)} and the adjacent conclusion.
For \textup{(SB4)}, the first permitted agreement is already safe.
\end{proof}

\begin{proof}[Proof of
\cref{thm:intro-asymptotic-rs-mca,cor:intro-identity-frontier}]
At \(a_{+,n}\), the profile-envelope hypothesis is precisely
\eqref{eq:exact-safe-budget}, so
\(B_{C_n,\Gamma_n}^{\rm MCA}(a_{+,n})\le B_n^*\).  At
\(a_{-,n}\), the profile-list hypothesis is
\eqref{eq:exact-unsafe-budget}, with the proved list size for the selected
identity, quotient, Chebyshev, or remainder profile, so
\(B_{C_n,\Gamma_n}^{\rm MCA}(a_{-,n})>B_n^*\).  The numerator is
nonincreasing in the agreement threshold: every witness agreeing on at
least \(a+1\) coordinates is also a witness at threshold \(a\).
Consequently the safe set is nonempty, its least element satisfies
\(a_{-,n}<a_n^*\le a_{+,n}\), and substituting
\(\delta_n^*=1-a_n^*/n\) proves
\eqref{eq:intro-threshold-bracket}.

If \(a_{+,n}-a_{-,n}=o(n)\), any common crossing \(a_n\) between them
satisfies \(a_n^*=a_n+o(n)\).  Thus
\(a_n=k_n+1+g_nn+o(n)\) and \(k_n/n\to\rho\) give
\(\delta_n^*=1-\rho-g_n+o(1)\).  In the identity-dominant window,
Stirling's formula \eqref{eq:target-entropy} says
\[
 \frac1n\log_2\barN_1
 =H_2(\rho_n+g)-\beta_ng+o(1).
\]
Comparison with \(n^{-1}\log_2(1+B_n^*)=\tau_n\) yields
\eqref{eq:intro-target-crossing}; concavity selects its rightmost crossing
as in \eqref{eq:intro-target-crossing}.  The corollary assumes, rather than
deduces from exponent notation, that quantitative reserves furnish the
exact safe and unsafe tests within \(o(n)\) of that crossing.  The bracket
therefore gives its claimed threshold.  When
\(\log(1+B_n^*)=o(n)\), \(\tau_n=o(1)\), and continuity of the rightmost
level-set endpoint gives
\(g_{T,n}=g^*(\rho_n,\beta_n)+o(1)\), which proves the zero-exponent
specialization.
\end{proof}

\begin{proof}[Proof of \cref{thm:main-ledger}]
The numerator bound is \cref{prop:numerator-bound}.  If agreements
\(a_-<a_+\) satisfy \eqref{eq:exact-unsafe-budget} at \(a_-\) and
\eqref{eq:exact-safe-budget} at \(a_+\), monotonicity localizes the integer
threshold to \([a_-,a_+]\).  Under the additional hypotheses in
\cref{thm:intro-asymptotic-rs-mca,cor:intro-identity-frontier}, these
agreements can be chosen \(o(n)\) apart around the target-adjusted entropy
crossing, giving the stated radius corridor.
\end{proof}

\section{Finite-row audit and exact scope}\label{sec:finite-scope}

For a finite challenge set \(\Gamma\), the target is the integer
\(B^{\ast}=\floor{\eps^{\ast}\abs\Gamma}\).  The compiler itself is now
exact: \cref{thm:exact-finite-profile-compiler} expresses its upper budget
through the actual residual moment, realized image, residual mass,
pair-incidence degrees, and direct cell budgets.  The finite Q target is
row-sharp: if the entire remaining budget is assigned to Q, one must prove
\[
        \max_s\abs{\Omc\cap\Phi^{-1}(s)}\le B^{\ast}
\]
inside the actual first-match residual, not merely inside the raw null
fiber \(\Phi^{-1}(0)\).  A fixed-order moment may fail this target unless
its order is comparable with \(w\log\abs\B/\Delta\), where \(\Delta\) is
the remaining bit margin; \cref{cor:finite-moment-floor} gives the exact
inequality.

The explanation projection is likewise exact.  Raw witnesses first map to
states \((\gamma,h)\) and only then to slopes; the two fibers are not
conflated.  The one-parameter moving-root pencil and the curve-degree
compiler pay the charts to which they apply, but a finite proof must still
verify that each residual balanced-core chart has such a cover or has an
independent direct or pair-incidence payment.  The separating-pole example
\cref{cor:exact-prefix-ray-realization} shows that exact-agreement
occupancy can be one.

Partial-occupancy \emph{support} exhaustion no longer remains open:
\cref{thm:exact-partial-occupancy} gives it with no loss.  What remains
row-specific is the nontrivial-character sum \textup{(PO4)}, realized-image
coverage, overlap/payment factors in \cref{lem:exact-profile-addback}, and
the distinct-slope projection.  The unrestricted ambient
identity-normalized Sidon statement is already false by
\cref{thm:smooth-quotient-obstruction}, and the displayed
degree-uniform PF certificate does not reach the deep polynomial-field
prefix by \cref{prop:pf-deep-prefix-barrier}.  Thus a closed ledger remains
a substantive row theorem even though its support census and finite
constant bookkeeping are exact.

\section{Summary of the proof chain}\label{sec:summary}

The unconditional analytic chain is short.  Once a primitive Boolean leaf
has a logarithmic moment satisfying \(\log L/q=o(N)\) and a Sidon moment
payment, any
positive-rate max-fiber failure produces a large high-energy fiber.  BSG
then gives a large subset with \(e^{o(N)}\) difference doubling, whereas
Boolean-cube difference growth forces \(\abs{A-A}\ge\abs A^{3/2}\), a
contradiction.  Thus the Sidon hypothesis implies primitive Q.  To reach an
MCA numerator one must additionally supply the algebraic distinct-slope
budgets and \textup{(RC)}; to reach a threshold one must then compare the
full profile envelope with the target and use the exact collision-aware
lower test.  The quotient construction in \cref{sec:quotient-obstruction}
shows why none of these additional comparisons may be omitted.

The remaining bookkeeping is exact.  Partial occupancies partition the
whole support slice by \cref{thm:exact-partial-occupancy};
\cref{lem:exact-profile-addback} then shows precisely where image coverage,
payment constants, and overlap multiplicity enter.  At finite length,
\cref{thm:exact-finite-profile-compiler} replaces every asymptotic compiler
loss by a literal moment, residual-mass, incidence-degree, and cell-budget
factor.  These theorems close support add-back and constant tracking, but
they do not manufacture the effective Fourier or residual slope estimates.


\section{Ledger interfaces and notation}\label{sec:interfaces}

For clarity we collect the interfaces between the pieces of the proof.  The objects below are the quantities a row-specific verification must output.

\paragraph{MCA interface.}
Input: a received line \((r_1,r_2)\), an agreement \(a\), and the Reed--Solomon row \(\RS_\F(D,k)\).  Output: an ordered first-match list of slope supersets \(E_i\) covering all bad slopes.  The required bound is \(\sum_i\abs{E_i}\le e^{o(n)}\mathfrak E_n(a)\).  The individual witnesses may overlap; slopes do not after first match.

\paragraph{Algebraic-cell interface.}
Input: a constructible stratum \(\Ccal\subseteq\mathfrak W_a\).  Output: a slope bound for first-match witnesses in \(\Ccal\).  The permitted reasons for such a bound are degree of a finite projection, entropy loss in support choices, field descent, determinantal rank loss, or saturation of many raw witnesses to fewer explanation states or slopes.  A cell is not paid merely because it has a name; the bound is part of the data.

\paragraph{Primitive-prefix interface.}
Input: a primitive profile \(\lambda\).  Output: a fixed-density Boolean
slice, an explicitly verified prefix map, its realized image size
\(L_\lambda\), the image scale
\(\barN_\lambda=\abs{\Omega^0_\lambda}/L_\lambda\), and the
moment-accessibility conditions of \cref{thm:primitive-q}.  An ambient
scale may be substituted only after \textup{(FI)} or a direct ambient
theorem has been proved.  A Sidon moment payment then yields the stated Q
bound; merely assigning a name to a heavy leaf does not.

\paragraph{Sidon/Fourier interface.}
Input: the same primitive leaf.  Output: a proof that the low-energy heavy
moment \(\Gsid_{q,\sigma}\) is \(e^{o(Nq)}\).  One sufficient route is
effective-span normalization followed by \textup{(MI)} and
\textup{(MA)}; \textup{(PF)} is one sufficient certificate for the minor
part, while a dual major-arc label by itself is insufficient.  Once this
interface is supplied, the high-energy proof uses only BSG and
Boolean-cube difference growth.

\paragraph{Support-pair and ray interfaces.}
Q gives the exact support-pair bound \(\sum_sQ_\lambda(s)^2\).  It does
not saturate those pairs to slopes.  That final projection is the separate
ray interface \textup{(RC)}, or a direct distinct-slope theorem.
Exact alternatives are retained-support occupancy
\cref{lem:exact-occupancy-compiler}, pair incidence
\cref{cor:finite-pair-ray-compiler}, and the curve-degree compiler
\cref{prop:curve-degree-ray-compiler}.

\paragraph{Entropy-frontier interface.}
Input: the asymptotic numerator \(e^{o(n)}\mathfrak E_n(a)\), the target,
and a target-viable agreement.  Output: the safe side of the threshold.
When the identity profile dominates, the leading calculation is
\[
        \log_2\barN_n=n(H_2(\rho+g_n)-\beta_ng_n)+o(n).
\]
Crossing the relevant target-adjusted level by a quantitative reserve that
dominates every error proves safety; a bare \(o(n)\) displacement does not.

The conditional compiler can therefore be read as
\[
\begin{gathered}
\text{paid algebraic profiles}
+\text{Sidon moment payment}
\longrightarrow \text{primitive Q},\\
\text{primitive Q}+\textup{(RC)}
+\text{profile envelope}
\longrightarrow \text{MCA numerator},\\
\text{numerator}+\text{target reserve}+\text{exact lower test}
\longrightarrow \text{threshold corridor}.
\end{gathered}
\]
The analytic implication to primitive Q and the formal compiler implications
are proved here.  The premises \textup{(A2)}, \textup{(MI)},
\textup{(MA)}, and
\textup{(RC)} are row-specific enumerative inputs; smooth/circle geometry
organizes them but does not establish them unconditionally.

\begin{remark}[Why the theorem is not phrased as a black-box inverse theorem]
A single black-box theorem of the form ``\Renyi{} collision excess implies an approximate group'' would be false for the ordinary moment, because of Sidon-heavy fibers.  The correct inverse statement is ledger-valued: collision excess implies either a Sidon/Fourier cell or a high-energy fiber; the high-energy fiber is impossible in a large Boolean primitive leaf.  This ledger-valued form is the reason the paper is phrased in moduli terms.
\end{remark}



\section{Coordinate form of the witness atlas}\label{sec:coordinate-atlas}

This section expands the moduli language into coordinates.  It is included to make clear that the cells used above are not informal cases but constructible pieces of a single incidence problem.  Fix a row \(C=\RS_\F(D,k)\), \(D=\{t_1,\ldots,t_n\}\), and an agreement \(a\).  For an \(a\)-support \(S\) write
\[
        Q_S(X)=\prod_{t\in S}(X-t)
        =X^a+q_1(S)X^{a-1}+\cdots+q_a(S).
\]
The prefix map is the projection
\(S\mapsto(q_1(S),\ldots,q_w(S))\), where \(w=a-k-1\) when using the
dimension-\((k+1)\)-to-\(k\) simple-pole conversion.  The shift by one
comes from passing from \(k\) to \(K=k+1\): the first \(w\) coefficients
are those not absorbed by a degree-\(<K\) explanation.

\begin{definition}[Support-prefix incidence]\label{def:support-prefix-incidence}
Let \(\Gr_a(D)\) denote the set of \(a\)-subsets of \(D\).
The \emph{support-prefix incidence} is the graph
\[
 \mathcal I_w=\{(S,z)\in\Gr_a(D)\times\A^w:
       z=(q_1(S),\ldots,q_w(S))\}.
\]
Its projection to \(\A^w\) is the \emph{support-prefix image}.
For \(z\in\B^w\), write \(N_w(z)=\abs{\{S\in\Gr_a(D): (q_1(S),\ldots,q_w(S))=z\}}\).
\end{definition}

For the unprofiled locator-list problem put
\(L_w=\abs{\{\Phi_w(S):S\in\binom Da\}}\).  Primitive Q is exactly the
assertion
\(\max_zN_w(z)\le e^{o(n)}\binom na/L_w\).  Under the full-image
certificate \(L_w\ge e^{-o(n)}\abs\B^w\), this is equivalent at
exponential accuracy to the familiar ambient bound
\(\max_zN_w(z)\le e^{o(n)}\binom na\abs\B^{-w}\).  An arbitrary
received-line MCA witness instead satisfies the interpolation-functional
pencil \eqref{eq:slope-elimination}; identifying all of its boundary charts
with locator-prefix fibers is part of the atlas and ray hypotheses, not an
automatic elimination step.  The following dictionary applies to the
unprofiled prefix map in large characteristic.

\begin{lemma}[Newton dictionary]\label{lem:newton-dictionary-expanded}
Assume the characteristic is larger than \(w\).  For \(1\le j\le w\),
the first \(w\) elementary coefficients \((q_1,\ldots,q_w)\) and the first
\(w\) power sums \(p_j(S)=\sum_{t\in S}t^j\) determine each other by
triangular polynomial automorphisms.  Consequently the prefix fibers may
be counted using the linear map
\[
        x\longmapsto \biggl(\sum_{t\in D}x_t t,
        \sum_{t\in D}x_t t^2,
        \ldots,
        \sum_{t\in D}x_t t^w\biggr)
\]
\(on the fixed-density slice\) with only bounded-complexity coordinate changes.
\end{lemma}

\begin{proof}
Newton's identities give \(p_j+q_1p_{j-1}+\cdots+q_{j-1}p_1+jq_j=0\).
Thus \(q_j\) is determined recursively from \(p_1,\ldots,p_j\) when
\(j\) is invertible, and conversely.  No small-characteristic
linearization is claimed here; one must retain elementary coordinates or
prove a suitable divided-power replacement separately.
\end{proof}

\begin{definition}[MCA witness scheme]\label{def:mca-witness-scheme}
For a received line \(r=(r_0,r_1)\), the exact-\(a\) MCA witness scheme
means the fiber \(\Wcal_a(r)\) of the universal locally closed incidence
in \cref{def:universal-witness-incidence}.  Its support coordinate belongs
to the finite reduced support scheme
\(\mathfrak S_a(D)=\coprod_{S\in\binom Da}\Spec\F\), and its slope
projection is \((\gamma,S,h)\mapsto\gamma\).  By
\cref{lem:exact-agreement-reduction}, its slope image is the
threshold-\(a\) MCA-bad set whenever \(a\ge k+1\).
\end{definition}

For fixed received data and \(S\), the closed incidence equations in
\cref{def:mca-witness-scheme} are linear in the coefficients of \(h\) and
in \(\gamma\); support-wise nontriviality is the separate open condition
of \cref{def:explanation}.  Eliminating
\(h\) gives the exact support-dependent equations
\eqref{eq:interpolation-functional}--\eqref{eq:slope-elimination}; their
weights involve \(Q_S'\) and the received values.  A reduction of these
equations to fixed top locator coefficients is an additional chart theorem
and is required by \textup{(A2)} or \textup{(RC)}.

\begin{proposition}[Atlas criterion]\label{prop:first-match-atlas-finite}
Suppose a smooth or circle sequence admits a witness-exhaustive
first-match atlas
of the algebraic types in \cref{sec:cell-catalogue}, together with its
analytic Sidon component and primitive leaves, and suppose the total
description entropy is \(o(n)\).  Then the atlas has \(e^{o(n)}\) profiles
and satisfies the census part of \textup{(A2)}.
\end{proposition}

\begin{proof}
The assumed atlas is exhaustive.  Its algebraic pieces are constructible by
\cref{lem:constructible-cells}; its Sidon component is analytic rather than
constructible.  The description-entropy hypothesis and
\cref{lem:profile-multiplicity} give the subexponential count.  The
proposition deliberately does not derive exhaustiveness from smoothness.
\end{proof}

\begin{lemma}[Slope multiplicity after support fixation]\label{lem:slope-multiplicity-fixed-support}
Fix a support \(S\) with \(\abs S\ge k+1\).  For a received line \((r_1,r_2)\), either the pair is explained on \(S\), or the set of slopes \(\gamma\) for which \(r_1+\gamma r_2\) is explained on \(S\) has size at most one.
\end{lemma}

\begin{proof}
If two distinct slopes \(\gamma_1,\gamma_2\) are explained by polynomials \(h_1,h_2\) of degree less than \(k\) on \(S\), then \((h_1-h_2)/(\gamma_1-\gamma_2)\) explains \(r_2\) on \(S\), and \(h_1-\gamma_1 r_2\) explains \(r_1\) on \(S\).  Since \(\abs S\ge k+1\), the interpolants are unique, so the pair is explained on \(S\).  Thus in the MCA-bad locus, each support contributes at most one slope.
\end{proof}

This lemma says that one fixed noncommon support yields at most one slope.
It does not put all bad supports into one prefix fiber and does not control
how many supports represent different slopes.  Consequently Q remains a
support-fiber estimate, and \textup{(RC)} or a direct slope projection is
still required.

\section{The Li--Wan--Weil calculation and residual stability}\label{sec:li-wan-details}

This section spells out two interfaces that are easy to obscure in a
ledger proof.  First, the Li--Wan sieve is used through the
elementary-coefficient estimate already proved in
\cref{thm:prefix-flatness-power-sum}; the finite numerical certificate is
\textup{(PF)}, while the asymptotic compiler requires effective
\textup{(MI)} plus \textup{(MA)} or the direct Sidon premise of
\cref{thm:sidon-resolved-payment}.  Second, the primitive residual
\(\Omc\) is a first-match subset of a full fixed-density profile slice; it
need not be Zariski open.  All prefix-flat estimates are applied to the
full slice, so arbitrary first-match deletion cannot create a larger fiber
than the full slice already had.

\begin{lemma}[Cycle-index coefficient bound]\label{lem:cycle-index-bound}
Let \(a_t\) be complex numbers with \(\abs{a_t}\le1\), indexed by a finite set \(T\).  If \(\abs{\sum_t a_t^j}\le L_j\) for \(1\le j\le m\), then
\[
        \abs{e_m(a_t:t\in T)}
        \le [u^m]\exp\left(\sum_{j=1}^m\frac{L_j}{j}u^j\right).
\]
In particular, if \(L_j\le L\) for all \(j\le m\), then \(\abs{e_m}\le\binom{L+m-1}{m}\).
\end{lemma}

\begin{proof}
Use the identity \(\sum_{r\ge0}e_r u^r=\exp(\sum_{j\ge1}(-1)^{j-1}(\sum_ta_t^j)u^j/j)\), truncate at degree \(m\), and compare coefficients after taking absolute values.  The uniform bound gives \(\exp(L\sum_{j\ge1}u^j/j)=(1-u)^{-L}\).
\end{proof}

\begin{lemma}[Residual monotonicity]\label{lem:residual-monotonicity}
Let \(\Omega^0\) be a full profile slice, let \(\Omc\subseteq\Omega^0\),
and let \(\Phi:\Omega^0\to\B^R\) be a syndrome map.  If
\(\abs{\Omega^0\cap\Phi^{-1}(s)}\le U\) for every \(s\), then the same
holds on \(\Omc\).  Consequently every logarithmic moment of the residual,
normalized at the ambient scale
\(\barN_0^{\rm amb}=\abs{\Omega^0}\abs\B^{-R}\), is bounded by the
corresponding full-slice max-fiber bound.  If
\(U=e^{o(N)}\barN_0^{\rm amb}\), then \textup{(FI)} follows and the same
conclusion holds in image normalization.
\end{lemma}

\begin{proof}
The fiber inclusion gives the first assertion.  An ambient max-fiber bound
gives the ambient moment bound by the upper half of the moment sandwich.
It also proves \textup{(FI)} by \cref{rem:flatness-certifies-image}; now
apply \cref{lem:image-ambient-moment-conversion}.
\end{proof}

\begin{proposition}[Smooth and circle prefix-flatness criterion]\label{prop:smooth-circle-prefix-flatness-criterion}
Consider a primitive smooth or circle chart with full active set \(T\),
density \(m/\abs T\in[\alpha,1-\alpha]\), prefix depth \(R\), base field
\(\B\), and minor-arc degree bound \(D\).  Suppose
\(m<\operatorname{char}\B\), the corresponding \(\Lambda\) satisfies
\textup{(PF)}, and the major characters satisfy the ambient instance of
\textup{(MA)}.  Then the
full profile slice and every first-match residual satisfy the ambient
max-fiber estimate and hence image-normalized Q.
\end{proposition}

\begin{proof}
Weil and the cycle-index estimate bound the minor-character aggregate,
while \textup{(MA)} bounds the major-character aggregate.  Fourier
inversion gives the full-slice estimate, and residual monotonicity transfers
it to the residual.
\end{proof}

\begin{remark}[How to read the numerical condition]\label{rem:pf-numerical}
Condition \textup{(PF)} says exactly that the Li--Wan--Weil error term is subexponential relative to the random prefix fiber.  Written exponentially, it is
\[
        \abs\B^R\binom{\Lambda+m-1}{m}\le e^{o(N)}\binom Nm.
\]
For finite certificates this is the minor-character check.  It must be
paired with an exact major-character estimate.  If \textup{(PF)} or
\textup{(MA)} fails, the Fourier argument gives no flatness conclusion;
one must prove the Sidon moment bound by another method.
\end{remark}

\begin{proposition}[From prefix flatness to Sidon payment]\label{prop:equidistribution-to-sidon}
Under the hypotheses of
\cref{prop:smooth-circle-prefix-flatness-criterion}, put
\(L=\abs{\Phi(\Omega^0)}\) and
\(\barN_0^{\rm img}=\abs{\Omega^0}/L\).  For every fixed \(\sigma>0\)
and logarithmic \(q\),
\[
        L^{-1}\sum_{\Delta_s\le e^{-\sigma N}}
        \left(\frac{f_s}{\barN_0^{\rm img}}\right)^q=e^{o(Nq)}.
\]
\end{proposition}

\begin{proof}
This is \cref{cor:sidon-payment-from-prefix-flatness}.  The low-energy
restriction can only reduce the image-normalized full moment.
\end{proof}

\section{Threshold compiler details}\label{sec:threshold-details}

The asymptotic theorem uses two numerical facts repeatedly.  The first converts a max-fiber failure into a logarithmic moment failure; the second converts the Q max-fiber bound into the SP second-moment bound.  They are elementary, but they are the bookkeeping bridge between the moduli ledger and the MCA numerator.

\begin{lemma}[Largest fiber and logarithmic moments]\label{lem:largest-fiber-log-detail}
Let \(N(z)\) be nonnegative integers with \(\sum_zN(z)=M\), and let \(L\) be the number of targets.  Put \(\bar N=M/L\) and \(R=\max_z N(z)/\bar N\).  For every \(q\ge2\),
\[
        R^{q}/L\le \sum_z (N(z)/\bar N)^q/L\le R^{q-1}.
\]
Consequently a positive-rate max-fiber failure \(R\ge e^{\eta n}\) gives a positive-rate logarithmic moment failure for every \(q\to\infty\) with \(\log L/q=o(n)\).
\end{lemma}

\begin{proof}
The lower bound is the contribution of the largest fiber, divided by
\(L\).  For the upper bound,
\(N(z)^q\le(\max N)^{q-1}N(z)\), so the normalized sum is at most
\(R^{q-1}\).  Taking logarithms in the lower bound gives the stated
moment-accessibility condition.
\end{proof}

\begin{lemma}[Q-to-SP second moment]\label{lem:q-to-sp-detail}
If \(\max_z N(z)\le \kappa\bar N\), then \(M^{-1}\sum_zN(z)^2\le\kappa\bar N\).  In particular \(\kappa=e^{o(n)}\) gives the SP scale \(e^{o(n)}\bar N\).
\end{lemma}

\begin{proof}
One has \(\sum_zN(z)^2\le(\max_zN(z))\sum_zN(z)=\kappa\bar N M\).  Dividing by \(M\) proves the claim.
\end{proof}

\begin{definition}[Integer staircase]\label{def:integer-staircase-detail}
At agreement \(a=K+w\), the identity profile has scale
\(\barN_1=\binom na\abs\B^{-w}\).  Its entropy exponent is
\(n\Hb(a/n)-w\log_2\abs\B+o(n)\).  It is a safe numerator budget only
after the full profile envelope has been proved to be
\(e^{o(n)}\barN_1\) and the other compiler hypotheses hold.
\end{definition}

\begin{proposition}[Target-aware crossing criterion]\label{prop:entropy-crossing-detail}
Assume the full profile envelope is \(e^{o(n)}\barN_1\) in a window,
the conditional numerator hypothesis of
\cref{thm:main-smooth-circle} holds there, and the target
is viable.  An agreement is certified safe by
\eqref{eq:exact-safe-budget} and certified unsafe by
\eqref{eq:exact-unsafe-budget}.  The zero of
\(\Hb(\rho+g)-\beta g\) describes the leading location only when the
target has subexponential logarithm and a quantitative reserve dominates
all errors and pole collisions.
\end{proposition}

\begin{proof}
This is \cref{lem:safe-side,prop:simple-pole-lower}.  Concavity identifies
the leading zero, but an unspecified \(o(1)\) displacement need not
dominate the error terms.  The lower side uses the exact list size and the
collision-aware conversion, not merely the sign of the identity exponent.
\end{proof}

\begin{proposition}[First-match summation at the frontier]\label{prop:first-match-sum-detail}
Suppose a witness-exhaustive ordered ledger has \(e^{o(n)}\) profiles and their
distinct-slope budgets sum to \(e^{o(n)}\mathfrak E_n(a)\), including all
primitive ray budgets.  Then the total bad-slope numerator is
\(e^{o(n)}\mathfrak E_n(a)\).
\end{proposition}

\begin{proof}
This is the union bound of \cref{lem:first-match-bound}; the supremum in
\eqref{eq:profile-envelope} makes it uniform over received lines.
\end{proof}

\section{Finite-row interface and scope}\label{sec:finite-interface}

The theorem in this paper is asymptotic.  It deliberately separates exponent-level closure from finite adjacent-row certification.  The finite problem is not a different mathematical object, but it requires the same cells with constants, exact divisor counts, and exact margins.  This section records the interface so that the asymptotic proof can be specialized without changing its logic.

A finite row supplies a number \(B^*\), the desired bad-slope ceiling.  A
certificate consists of a witness-exhaustive ordered atlas
\(\Ccal_i\), exact supersets \(E_i\supseteq Z_i^\circ\) of its actual
first-match projections, and integer budgets \(\abs{E_i}\le U_i\) with
\(\sum_iU_i\le B^*\).  The proof of \cref{lem:first-match-bound} is already
finite; only the estimates behind the payments need sharpening.  Quotient
cells require exact divisor-lattice recursion.  Planted cells require exact
binomial ratios rather than entropy inequalities.  Tangent and rank-pivot
cells require explicit eliminants.  Fourier/Sidon cells require explicit
prefix-equidistribution error terms.  Split-pencil cells require exact
shift-pair enumeration or a certified upper envelope.

\begin{proposition}[Finite specialization principle]\label{prop:finite-specialization-principle}
Let a finite smooth or circle row admit a witness-exhaustive first-match atlas,
exact distinct-slope budgets for its algebraic profiles, an exact Sidon
moment estimate yielding Q, and an exact \textup{(RC)} or direct ray bound.
If the resulting first-match sum is at most \(B^*\), then
\(B_C^{\rm MCA}(a)\le B^*\).
\end{proposition}

\begin{proof}
This is \cref{lem:first-match-bound} with exact budgets.  The exhaustive
atlas, major-character estimate, and ray compiler are structural inputs,
not constant-tracking details.
\end{proof}

\begin{theorem}[Exact finite profile compiler]
\label{thm:exact-finite-profile-compiler}
Fix a finite row, an agreement \(a\), and a received line.  Let an ordered
witness-exhaustive first-match atlas have algebraic or direct cells with
exact integer budgets \(D_i\), and primitive cells indexed by \(\lambda\).
For a primitive cell put
\[
 M_\lambda=\abs{\Omega^0_\lambda},\quad
 W_\lambda=\abs{\Omc_\lambda},\quad
 L_\lambda=\abs{\Phi_\lambda(\Omega^0_\lambda)},\quad
 \barN_\lambda=M_\lambda/L_\lambda.
\]
For an integer \(q_\lambda\ge2\), define the exact residual moment
\[
 \Gamma_{q_\lambda,\lambda}
 =\frac1{L_\lambda}\sum_s
 \left(\frac{Q_\lambda(s)}{\barN_\lambda}\right)^{q_\lambda},
 \qquad
 Q_\lambda(s)=\abs{\Omc_\lambda\cap\Phi_\lambda^{-1}(s)}.
\]
Let
\[
 \Pcal_\lambda=
 \{(A,B)\in(\Omc_\lambda)^2:
              \Phi_\lambda(A)=\Phi_\lambda(B)\}.
\]
Whenever a primitive cell has a certified incidence
\(I_\lambda\subseteq Z_\lambda^\circ\times\Pcal_\lambda\) such that
every \(\gamma\in Z_\lambda^\circ\) has
\(\deg_{I_\lambda}(\gamma)\ge H_\lambda\ge1\), while every
\((A,B)\in\Pcal_\lambda\) has
\(\deg_{I_\lambda}(A,B)\le J_\lambda\), its first-match slope set
satisfies
\[
 \abs{Z_\lambda^\circ}
 \le U_\lambda:=
 \floor{\frac{J_\lambda W_\lambda}{H_\lambda}\,
 \barN_\lambda
 \bigl(L_\lambda\Gamma_{q_\lambda,\lambda}\bigr)^{1/q_\lambda}}.
 \tag{FC1}
\]
Alternatively it may carry a direct slope bound \(D_\lambda\).  Assume at
least one of \(U_\lambda,D_\lambda\) is available for every primitive
cell, and let \(\widetilde U_\lambda\) be the available bound or the
minimum when both are available.  Consequently
\[
 B_C^{\rm MCA}(a)
 \le \max_{\textup{received lines}}
 \left(\sum_iD_i+\sum_\lambda\widetilde U_\lambda\right).       \tag{FC2}
\]
For a challenge set with ceiling \(B^*\), agreement \(a\) is safe whenever
the right side of \textup{(FC2)}, after intersecting each slope set with
the challenge set, is at most \(B^*\).  If the exact lower construction is
unsafe at \(a\) and \textup{(FC2)} is safe at \(a+1\), then the finite
agreement threshold is exactly \(a+1\).
\end{theorem}

\begin{proof}
The moment sandwich gives
\[
 \max_sQ_\lambda(s)
 \le\barN_\lambda
      \bigl(L_\lambda\Gamma_{q_\lambda,\lambda}\bigr)^{1/q_\lambda}.
\]
The exact support-pair identity and
\(\sum_sQ_\lambda(s)=W_\lambda\) give
\(\sum_sQ_\lambda(s)^2\le
W_\lambda\max_sQ_\lambda(s)\).  Double counting the certified ray
incidence yields
\(H_\lambda\abs{Z_\lambda^\circ}
 \le J_\lambda\sum_sQ_\lambda(s)^2\), proving \textup{(FC1)}.
First-match slope sets are disjoint, so summing the exact cell budgets and
maximizing over received lines proves \textup{(FC2)}.  The final assertion
follows from monotonicity in the agreement threshold.
\end{proof}

\begin{lemma}[Exact adjacent identity ratio]
\label{lem:exact-adjacent-identity-ratio}
Put \(K=k+1\) and define
\(I(u)=\binom nu\abs\B^{-(u-K)}\) for \(K\le u\le n\).  If
\(K\le a<n\), then
\[
 \frac{I(a+1)}{I(a)}=
 \frac{n-a}{(a+1)\abs\B},
 \qquad
 \log_2\frac{I(a)}{I(a+1)}
 =\log_2\abs\B+\log_2\frac{a+1}{n-a}.              \tag{AR1}
\]
The exact prefix-list lower bound is \(\ceil{I(a)}\).  Thus an adjacent
certificate uses \textup{(AR1)} and the literal budgets in
\textup{(FC2)}, not a derivative of an entropy approximation.
\end{lemma}

\begin{proof}
Use \(\binom n{a+1}/\binom na=(n-a)/(a+1)\) and note that the prefix
depth increases from \(w\) to \(w+1\).  The list statement is
\cref{prop:exact-prefix-list}.
\end{proof}

\begin{lemma}[Finite Boolean high-energy constant]
\label{lem:finite-boolean-high-energy}
Assume a chosen quantitative BSG theorem has the finite form
\[
 E(F)\ge\abs F^3/K
 \Longrightarrow
 \exists F'\subseteq F:\quad
 \abs{F'}\ge c_0K^{-A}\abs F,
 \quad
 \abs{F'-F'}\le c_1K^B\abs{F'}.
\]
Then every \(F\subseteq\{0,1\}^T\) satisfying the energy hypothesis obeys
the explicit bound
\[
                 \abs F\le c_0^{-1}c_1^2K^{A+2B}.   \tag{FC3}
\]
Thus the finite inverse branch retains the actual constants of the selected
quantitative BSG theorem; a symbol \(K^{O(1)}\) is not an adjacent-row
certificate.
\end{lemma}

\begin{proof}
Boolean difference growth gives
\(\abs{F'}^{3/2}\le c_1K^B\abs{F'}\), hence
\(\abs{F'}\le c_1^2K^{2B}\).  Combine this with
\(\abs{F'}\ge c_0K^{-A}\abs F\).
\end{proof}

The largest practical finite cost is the moment-to-max conversion.  In asymptotics, a logarithmic moment order \(q\to\infty\) eliminates a fixed positive exponential max-fiber gap.  In a printed adjacent row with margin \(\Delta\) bits and prefix depth \(w\beta\), a direct moment hierarchy may require order comparable to \(w\beta/\Delta\).  This is why the finite program benefits from exact prefix equidistribution, exact split-pencil charts, and row-specific eliminants rather than a blind high-moment bound.

\begin{remark}[Why the asymptotic proof is still useful for finite rows]
The framework supplies a diagnostic list, but it does not prove that the
list is exhaustive for every row.  A finite certificate must verify the
atlas, Fourier aggregates, and ray projection explicitly.
\end{remark}

\section{Conditional proof chain}\label{sec:complete-chain}

We record the logical dependencies without turning geometric labels into
unproved enumerative bounds.

\begin{theorem}[Domain geometry supplies profile maps]\label{thm:domain-to-ledger-final}
Multiplicative cosets and circle twin cosets have the explicit power and
Chebyshev folding maps of \cref{sec:smooth-circle-domains}.  These maps
classify natural quotient profiles and their scaled quotient coefficient
fields.  They do
not, by themselves, prove an exhaustive first-match atlas, a major-arc
aggregate, or a distinct-slope budget.
\end{theorem}

\begin{proof}
This is the content of \cref{def:map-smooth,lem:cheb-smooth}; the limitations
are witnessed quantitatively by \cref{thm:smooth-quotient-obstruction}.
\end{proof}

\begin{theorem}[Primitive residue with Sidon payment implies Q]\label{thm:primitive-to-q-final}
A fixed-density primitive Boolean leaf satisfying the moment-accessibility,
image-scale, and Sidon moment hypotheses of \cref{thm:primitive-q}
satisfies \(\max_zN_w(z)\le e^{o(n)}\barN_\lambda\).
\end{theorem}

\begin{proof}
This is \cref{thm:primitive-q}; \textup{(PF)}, in the range where it
certifies \textup{(MI)}, together with \textup{(MA)} is one sufficient
way to verify its Sidon premise.
\end{proof}

\begin{theorem}[Paid profiles, Q, and RC imply the numerator]\label{thm:q-sp-to-mca-final}
For a ledger-admissible sequence, the paid algebraic profile budgets,
primitive Q bounds, and \textup{(RC)} imply
\(B_{C_n}^{\rm MCA}(a_n)\le e^{o(n)}\mathfrak E_n(a_n)\).
\end{theorem}

\begin{proof}
This is \cref{prop:q-implies-sp,prop:numerator-bound}.  The second moment
alone is not used as a substitute for \textup{(RC)}.
\end{proof}

\begin{corollary}[Conditional target-adjusted frontier]\label{cor:frontier-final}
Under the additional profile-envelope comparison (identity-dominance when
the identity formula is invoked), target-reserve, viability, and
exact lower-side hypotheses of
\cref{thm:intro-asymptotic-rs-mca,cor:intro-identity-frontier}, the threshold
lies in the corridor determined by
\eqref{eq:exact-safe-budget}--\eqref{eq:exact-unsafe-budget}.  It reduces to
the zero of \(\Hb(\rho+g)-\beta g\) only in the stated identity-dominant,
subexponential-target regime.
\end{corollary}

\begin{proof}
Apply \cref{thm:q-sp-to-mca-final,prop:entropy-crossing-detail}.
\end{proof}


\section{Algebraic-geometric payment criteria}\label{sec:expanded-ledger-proofs}

This section records what a row-specific payment proof must establish.
Constructibility, codimension in an ambient witness space, and support-pair
counts are not by themselves distinct-slope estimates.  Each criterion
below therefore separates the algebraic description from the missing
projection or profile comparison.

\subsection*{Quotient profiles}
For a multiplicative coset \(D=\theta H\), let \(H_c\le H\) have order
\(c\) and use the normalized quotient
\(\pi_c(x)=(x/\theta)^c\in H^c\).  If \(S\) is a union of \(H_c\)-orbits,
then \(S=\pi_c^{-1}(\bar S)\) and, with \(a=\abs S\),
\(Q_S(X)=\theta^aQ_{\bar S}((X/\theta)^c)\).  Thus its top coefficients
are lacunary.  A witness descends only when the received and explaining
data satisfy the corresponding fiber-invariance equations; otherwise the
noninvariant residues must be retained and counted, not automatically
routed to another paid cell.

The quotient budget is obtained from the entropy inequality
\[
        \binom{n/c}{a/c}\le \exp\{(n/c)\Hb(a/n)+O(\log n)\}.
\]
When \(c\ge2\) is fixed, the support count is exponentially smaller than
\(\binom na\), but the quotient also has fewer prefix equations and may be
defined over a smaller field.  Its correct budget is therefore its own
term \(\barN_\lambda\) in \(\mathfrak E_n\).  It need not be comparable
with the identity term; \cref{thm:smooth-quotient-obstruction} gives an
explicit exponential separation.

\begin{lemma}[Pure-power quotient lacunarity]\label{lem:quotient-lacunarity-expanded}
In the normalized multiplicative quotient above, the first \(w\) locator
coefficients satisfy exactly the zero relations in degrees not divisible
by \(c\); hence at least \(w-\ceil{w/c}\) independent linear relations
hold.  No analogous count is asserted for an arbitrary degree-\(c\) map
without an explicit triangular coefficient calculation.
\end{lemma}

\begin{proof}
Expand
\(Q_S(X)=\theta^aQ_{\bar S}((X/\theta)^c)\).  Only powers congruent to
\(a\) modulo \(c\) occur, and the surviving top coefficients depend
triangularly on those of \(Q_{\bar S}\).
\end{proof}

\subsection*{Tangent and common-line charts}
A tangent chart appears when the slope projection of the witness incidence
has \emph{ramification}, meaning that its differential has smaller than
the expected rank.  Write the interpolation equations on a support \(S\)
as \(A_S h+\gamma b_S=r_{1,S}\), where \(A_S\) is the Vandermonde
evaluation matrix for degree \(<k\) polynomials and \(b_S=r_{2,S}\).  If
\(\abs S=k+1\), the existence of \((h,\gamma)\) is the vanishing of one
determinant.  If the determinant vanishes to first order under a support
deformation, the derivative determinant vanishes as well.  These two
equations are the tangent cell.  At larger support size, the same condition
is expressed by the rank of the augmented matrix
\((A_S,b_S,r_{1,S})\).

The common-line locus contributes no bad slope on that support.  A
neighboring tangent locus, however, is not automatically paid: its defining
minors must have the claimed independence and its slope projection must be
controlled.

\begin{lemma}[Ramification payment criterion]\label{lem:ramification-payment-expanded}
Suppose a tangent chart is cut out by \(h\) independent bounded-degree
minors and the projection of that chart to slopes has fibers large enough
that the resulting codimension gives a factor \(\abs\B^{-h+o(n)}\) in the
distinct-slope image.  Then the chart is paid at its natural profile scale.
The existence of one nonzero determinant gives only codimension one and
does not imply this projection estimate.
\end{lemma}

\begin{proof}
Independent bounded-degree equations give the finite-field point loss by
the usual degree estimate.  The assumed projection-fiber bound converts
that point loss into a slope-image loss.  Both hypotheses are necessary;
dependence among minors need not have a quotient explanation.
\end{proof}

\subsection*{Extension and subfield profiles}
Let \(\F/\F_0\) be a finite extension.  A locator or slope profile is \(\F_0\)-defined when its coefficient vector is fixed by the Frobenius \(x\mapsto x^{\abs{
\F_0}}\).  Such loci are closed in affine coefficient space.  Galois
orbits organize their support profiles, but field of definition alone does
not guarantee a favorable numerator exponent; the resulting profile must
be counted at its own field and prefix depth.

\begin{lemma}[Weil restriction interface]\label{lem:weil-restriction-expanded}
A bounded-degree extension-line locus becomes a bounded-complexity
base-field incidence after Weil restriction.  This coordinate change does
not itself pay the locus; a base-field point count and a slope-projection
bound at the appropriate profile scale are still required.
\end{lemma}

\begin{proof}
Choose an \(\F_0\)-basis of \(\F\).  Coefficient variables and equations
expand into bounded tuples and systems over \(\F_0\), so complexity grows
by a bounded factor.  Nothing in this operation controls the number of
rational slopes; that is the additional hypothesis stated above.
\end{proof}

\subsection*{Effective-image rank collapse}
The prefix map on a verified primitive chart may have an expected image
dimension \(w\).  An \emph{effective-image rank-collapse locus} is
described by a rank drop of its Jacobian or incidence matrix.  It is
distinct from saturation, which records witnesses producing the same
explanation state or final slope under the two projections of
\cref{def:explanation-occupancy}.  To pay rank collapse one must exhibit independent
minors and control the slope projection as in
\cref{lem:rank-defect-payment}; a name such as quotient, planted, or
field descent is not a substitute for that verification.

\begin{lemma}[Vandermonde column rank]\label{lem:vandermonde-image-rank-expanded}
Let \(U\subseteq D\) be active coordinates in a verified weighted
Vandermonde leaf, and let \(v_t=\rho(t)(1,t,\ldots,t^{R-1})\) with
\(\rho(t)\ne0\).  If the parameters \(t\in U\) are distinct, every subset
of at most \(R\) columns is linearly independent.
\end{lemma}

\begin{proof}
For \(r\le R\), the determinant is
\((\prod_j\rho(t_j))\prod_{i<j}(t_j-t_i)\ne0\).  This proves only the
column-independence assertion.  Classification and payment of a
higher-dimensional image-rank locus remain separate atlas and projection
problems.
\end{proof}

\subsection*{Split pencils}
A split-pencil chart is most transparent in locator language.  Let \(A\) and \(B\) be monic split polynomials with disjoint root sets in a residual domain \(D'\), both of degree \(e\).  A depth-\(w\) shift pair satisfies \(\deg(A-B)\le e-w-1\), which is equivalent to equality of the first \(w\) nonleading coefficients.  Thus SP is the second moment of the prefix fibers.  The top stratum \(e=w+1\) has \(A-B=c\in\F^\times\).  The polynomial family \(A-c\) must remain split over \(D'\), and this is a one-dimensional pencil in the projective locator space.

\begin{lemma}[Split-pencil accounting]\label{lem:split-pencil-alternatives-expanded}
For an actual one-dimensional locator pencil, the moving-root bound
\eqref{eq:moving-root-bound} controls parameters.  For a general family,
Q controls only the number of support pairs with common prefix.  A
distinct-slope conclusion requires \textup{(RC)}, and no exhaustive
decomposition of higher-dimensional balanced cores is asserted.
\end{lemma}

\begin{proof}
The one-pencil assertion is \cref{prop:split-pencil-payment}.  The exact
common-prefix pair identity is \cref{lem:second-moment-identity}.  Neither
statement supplies the missing higher-dimensional exhaustion or the
pair-to-ray fiber bound.
\end{proof}

\section{Smooth and circle profile examples}\label{sec:examples}

The two domain families below supply folding geometry.  Without the
remaining checks in \cref{def:admissible-sequence}, they are not automatic
instances of the conditional compiler.

\paragraph{Multiplicative smooth rows.}
Let \(\B_n\) be a finite field and let \(D_n=\theta_nH_n\subseteq\B_n^\times\) be a coset of a cyclic subgroup of order \(n\).  The quotient scales are divisors \(a_n\mid n\) for which the quotient size \(N_n=n/a_n\) tends to infinity and the first-match Fourier alternative of \cref{thm:sidon-resolved-payment} is applied at the prefix depth used; the displayed Li--Wan/Weil error is the corresponding finite certificate when one wants constants.  Power maps \(x\mapsto x^{a_n}\) give exact uniform fibers.  The stabilizer lattice is the divisor lattice of \(H_n\), and quotient descent is exact because a polynomial of degree less than \(k\) that is constant on enough fibers must be a pullback from the quotient row.

\paragraph{Circle line rounds.}
Let \(p\equiv3\pmod4\), let \(\UU\le\F_{p^2}^\times\) be the norm-one
torus, and let \(\Dcal=gH\cup g^{-1}H\) be a twin coset with
\(g^{2a}\notin H^{(a)}\) at the scales used.  The \(x\)-coordinate domain
\(D=\chi(\Dcal)\) is Chebyshev smooth by \cref{lem:cheb-smooth}.
Multiplicative quotient profiles therefore have Chebyshev analogues with
\(X^a\) replaced by \(T_a\).  Minor-character sums pull back to the torus
with the factor two in \cref{thm:circle-prefix-equidistribution}; major
characters still require \textup{(MA)}.

\paragraph{Circle codes.}
The bivariate circle code of degree \(w\) on a circle-domain set is
diagonally equivalent to \(\RS(E',2w+1)\) on the torus domain.  Diagonal
equivalence preserves agreement supports and slopes, so all MCA numerators
are unchanged.  The only visible difference is parity: \(2w+1\) is odd,
while dyadic folding scales are even.  This is why
\cref{lem:map-smooth-fiber} is stated without assuming \(a\mid k\).

\begin{proposition}[Standard smooth/circle checklist]\label{prop:standard-admissible}
A multiplicative or circle row is ledger-admissible once it verifies all of
\textup{(A1)}--\textup{(A7)}.  Uniform power or Chebyshev fibers verify the
geometric part of \textup{(A1)}, and a Weil parameter satisfying
\textup{(PF)} verifies only the minor-character part of \textup{(A4)}.
The atlas payments, \textup{(MA)}, image-scale conditions,
\textup{(RC)}, and full profile envelope remain separate checks.
\end{proposition}

\begin{proof}
This is a restatement of \cref{def:admissible-sequence}, with the two
automatic geometric observations just proved.  The counterexample shows
that no stronger conclusion follows from smoothness alone.
\end{proof}

\section{Discussion of proof robustness}\label{sec:robustness}

The proof deliberately avoids a false inverse statement.  Ordinary
\Renyi{} excess splits into a Sidon-heavy part and a high-energy part.  The
former is an explicit analytic hypothesis, verified by Fourier flatness
only after effective \textup{(MI)} and \textup{(MA)} are proved;
\textup{(PF)} is a sufficient minor certificate in its numerical range.
The latter branch is killed by the Boolean support slice.  This analytic
implication is domain-independent.

The algebraic-geometric language also prevents double counting.  Quotient, tangent, and saturation loci overlap heavily.  A support may be periodic and planted; a tangent chart may also lie in an extension locus; a split pencil may be Chebyshev-invariant.  The first-match assignment converts an overlapping constructible cover into a disjoint counting cover without requiring the strata themselves to be disjoint.  This is essential for finite certification and harmless asymptotically.

Finally, finite certification replaces each asymptotic input by an exact
integer bound.  The paper closes the high-energy Boolean branch, but it
does not close the Sidon, algebraic-projection, or higher-dimensional ray
branches for every smooth/circle row.


\section{Smooth/circle ingredients as theorem packages}\label{sec:ingredient-packages}

For reference, and to make the integration with the smooth-domain and circle-domain work explicit, we isolate the exact theorem packages used by the proof.  Each package has already appeared in the body of the argument, but this section records them in the order in which a reader would verify an unfamiliar row.

\begin{theorem}[Stabilizer partition package]\label{thm:stabilizer-package}
Let \(D\) be either a multiplicative coset or a Chebyshev twin-coset \(x\)-domain.  Every support \(S\subseteq D\) has a maximal periodicity scale \(c(S)\).  If \(c(S)=c>1\), then \(S\) is the pullback of a unique support on the quotient domain of size \(n/c\).  If \(c(S)=1\), then no nontrivial quotient map in the domain tower preserves \(S\).
\end{theorem}

\begin{proof}
In the multiplicative case, the acting group is cyclic, so stabilizers are exactly the subgroups \(H_c\).  A set fixed by \(H_c\) is a union of \(H_c\)-orbits and hence a pullback under \(x\mapsto x^c\); maximality of \(c(S)\) is maximality of the stabilizer.  In the circle case, \cref{lem:cheb-smooth} identifies the Chebyshev fibers with the images of the same subgroup tower on the norm-one torus.  The \(x\)-coordinate map is two-to-one in a way compatible with inversion, so complete torus fibers descend to complete Chebyshev fibers and the same stabilizer argument applies.
\end{proof}

\begin{theorem}[Invariant quotient-descent package]\label{thm:exact-descent-package}
Let \(S\) be periodic of scale \(c\), and suppose the two received words
and the explaining polynomial are pullbacks through the same quotient map.
Then the witness descends exactly; agreement is divided by \(c\), and
pair-explanation commutes with pullback.
\end{theorem}

\begin{proof}
Write each datum as a composition with the quotient coordinate.  Equality
on a union of complete fibers is then exactly equality on the quotient
support, and pullback preserves simultaneous explanation.  Support
periodicity without data invariance does not imply this theorem.
\end{proof}

\begin{theorem}[Fully primitive residual criterion]\label{thm:aperiodic-package}
If an exhaustive atlas has separately identified all quotient, planted,
tangent, extension, rank-defect, saturation, and pencil loci, its
\emph{fully primitive residual} has the negated properties by definition.
Smoothness alone does not prove
that these loci form such an exhaustive subexponential atlas.
\end{theorem}

\begin{proof}
The conclusion is the definition of the complement.  Constructibility of
the algebraic loci follows from \cref{lem:constructible-cells};
exhaustiveness and description entropy are the hypotheses.
\end{proof}

\begin{theorem}[Effective prefix flatness package]\label{thm:prefix-flatness-package}
On a primitive smooth/circle leaf satisfying effective \textup{(MI)} and
\textup{(MA)}, every full-slice and residual prefix fiber is at most
\(e^{o(n)}\barN_\lambda\), and the Sidon moment is
\(e^{o(nq)}\).  The same conclusion follows from \textup{(PF)} plus
\textup{(MA)} whenever the pointwise certificate is numerically valid.
\end{theorem}

\begin{proof}
The effective max-fiber and image-coverage assertions are
\cref{prop:effective-mi-ma-flatness}.  If \(L\) is the realized
full-slice image size and \(M=\binom Nm\), then
\(L\ge e^{-o(N)}A_{\rm eff}\), while every residual fiber is at most
\(e^{o(N)}M/A_{\rm eff}\).  Relative to the image-normalized full-slice
mean \(M/L\), every residual fiber is therefore \(e^{o(N)}\)-bounded.
The upper moment sandwich gives \(e^{o(Nq)}\) for every logarithmic
\(q\), and hence for the Sidon submoment.  The pointwise sufficient route is
\cref{thm:bounded-prefix-equidistribution,thm:circle-prefix-equidistribution,prop:equidistribution-to-sidon}.
\end{proof}

\begin{theorem}[Primitive Boolean package]\label{thm:boolean-package}
Let \(\Omc\subseteq\{0,1\}^T\) be a fixed-density leaf satisfying the
moment-accessibility, Sidon moment, and image-scale hypotheses of
\cref{thm:primitive-q}.  Then
\(\max_s\abs{\Omc\cap\Phi^{-1}(s)}\le e^{o(n)}\barN\).
\end{theorem}

\begin{proof}
This is \cref{thm:primitive-q}.
\end{proof}

The packages isolate exactly what is proved and what remains an input.  The
stabilizer and invariant quotient-descent statements organize quotient profiles;
effective \textup{(MI)} plus \textup{(MA)} pays a verified prefix leaf,
with \textup{(PF)} available as a sufficient certificate; and the
Boolean theorem closes the high-energy branch.  Exhaustive algebraic slope
budgets and \textup{(RC)} remain conditions of the compiler.

\section{A row-by-row verification template}\label{sec:verification-template}

A new row is checked against \textup{(A1)}--\textup{(A7)}.  Besides the
folding tower and prefix depths, this requires an exhaustive
subexponential atlas with direct slope budgets, effective \textup{(MI)} and
\textup{(MA)} (with \textup{(PF)} as a sufficient minor certificate) or
a direct Sidon estimate, the primitive image-scale and
moment conditions, \textup{(RC)} or direct ray bounds, and the complete
profile envelope.

\begin{proposition}[Verification template]\label{prop:verification-template}
Suppose a row verifies \textup{(A1)}--\textup{(A7)} of
\cref{def:admissible-sequence}, with every asymptotic bound uniform over
received lines.  Then \cref{thm:main-smooth-circle} applies.
\end{proposition}

\begin{proof}
This is the conditional compiler theorem.
\end{proof}

The finite version replaces every input by an exact inequality.  In
particular, a second-moment identity is not a finite ray certificate without
an exact pair-to-ray fiber bound.

\section{Cell-budget audit}\label{sec:complete-cell-budgets}

This section lists the estimates a closure proof would have to certify.
Every profile is budgeted at its own \(1+\barN_\lambda\) and retained in
\(\mathfrak E_n\); the identity scale is not used as a universal budget.

\begin{definition}[Profile complexity]\label{def:profile-complexity}
A profile type and a realized profile are the objects of
\cref{def:profile-cell}.  Their description data include the type
\(\lambda\) and every realized parameter \(\xi\) needed to specify the
incidence, such as a quotient map, planted core, tangent multiplicity
vector, field of definition, rank-pivot shape, saturation image, or
split-pencil partition.  A profile class has subexponential complexity if
the total number of nonempty realized pairs \((\lambda,\xi)\) in a row is
\(e^{o(n)}\), uniformly in the received line, and the algebraic description
complexity is subexponential as in \cref{def:cell}.
\end{definition}

The following observation is used repeatedly.  If a cell with a fixed
profile has cost \(e^{-c n}\) relative to the uncut prefix scale, then any
subexponential family of such cells is paid.  If the saving is only
\(e^{-o(n)}\), it is still paid provided the exact comparison with
\(1+\barN\) is nonpositive at the frontier.  This is the source of the
integer-staircase checks in finite-row certificates.

\begin{lemma}[Exact weighted profile add-back]
\label{lem:exact-profile-addback}
Let \(\{\Omega_\lambda\}_{\lambda\in\Lambda}\) be a finite family of
nonempty full profile slices, put
\(\Omega=\bigcup_\lambda\Omega_\lambda\), and let
\[
 \mu=\max_{S\in\Omega}
       \abs{\{\lambda:S\in\Omega_\lambda\}}.
\]
For each \(\lambda\), let
\(\Phi_\lambda:\Omega_\lambda\to\Sigma_\lambda\), put
\(L_\lambda=\abs{\Phi_\lambda(\Omega_\lambda)}\), and suppose its actual
first-match slope budget satisfies
\[
 U_\lambda\le \kappa_\lambda
 \left(1+\frac{\abs{\Omega_\lambda}}{L_\lambda}\right).
\]
Then, with no asymptotic notation,
\[
 \sum_{\lambda\in\Lambda}U_\lambda
 \le 2\sum_{\lambda\in\Lambda}
       \kappa_\lambda\frac{\abs{\Omega_\lambda}}{L_\lambda}.       \tag{AB1}
\]
If all maps take values in one formal target of size \(A\) and
\(\eta_\lambda=A/L_\lambda\), then
\[
 \sum_\lambda U_\lambda
 \le \frac2A\sum_\lambda
       \kappa_\lambda\eta_\lambda\abs{\Omega_\lambda}.           \tag{AB2}
\]
In particular, if \(\kappa_\lambda\le\kappa\) and
\(L_\lambda\ge A/\eta\) for every \(\lambda\), then
\[
 \sum_\lambda U_\lambda\le
 2\kappa\eta\mu\frac{\abs\Omega}{A}.                            \tag{AB3}
\]
For a partition, \(\mu=1\).  Disjoint first-match slope projections do
not imply \(\mu=1\) for their full support slices.
\end{lemma}

\begin{proof}
For a nonempty slice, \(L_\lambda\le\abs{\Omega_\lambda}\), so
\(1+\abs{\Omega_\lambda}/L_\lambda
 \le2\abs{\Omega_\lambda}/L_\lambda\).  Summing proves
\textup{(AB1)}.  Substitute
\(L_\lambda=A/\eta_\lambda\) for \textup{(AB2)}.  The uniform bounds
and \(\sum_\lambda\abs{\Omega_\lambda}\le\mu\abs\Omega\) prove
\textup{(AB3)}.
\end{proof}

\begin{remark}[Image coverage is necessary]
There is no absolute constant \(C\) for which every partition and every
map satisfy
\[
 \sum_\lambda\left(1+
 \frac{\abs{\Omega_\lambda}}{\abs{\Phi(\Omega_\lambda)}}\right)
 \le C\left(1+
 \frac{\abs\Omega}{\abs{\Phi(\Omega)}}\right).
\]
Indeed, take \(\Phi\) to be a bijection on an \(M\)-element set and
partition its domain into singletons.  The two sides before multiplication
by \(C\) are \(2M\) and \(2\), respectively.  Thus a common coarse
boundary theorem or the explicit image-coverage factors in
\cref{lem:exact-profile-addback} are logically necessary; exact support
add-back alone is insufficient.
\end{remark}

\begin{lemma}[Subexponential profile summation]\label{lem:profile-summation}
Let \(\{E_\alpha\}_{\alpha\in I_n}\) be slope sets with
\(\abs{I_n}=e^{o(n)}\) and
\(\abs{E_\alpha}\le e^{o(n)}(1+\barN)\) uniformly.  Then
\(\abs{\bigcup_{\alpha}E_\alpha}\le e^{o(n)}(1+\barN)\).
\end{lemma}

\begin{proof}
The union is bounded by
\(\sum_\alpha\abs{E_\alpha}\le
e^{o(n)}e^{o(n)}(1+\barN)=e^{o(n)}(1+\barN)\).
\end{proof}

\subsection*{Quotient cells}
A quotient cell is attached to a finite map \(\pi:D\to D'\) with uniform fiber size \(d>1\).  In the multiplicative case \(\pi(x)=x^d\) on a coset of a cyclic group.  In the circle case \(\pi\) is the Chebyshev map \(x=z+z^{-1}\mapsto z^d+z^{-d}\) after passing to the norm-one torus.  A witness lies in the quotient cell if, after deleting a remainder support of size \(r\), its locator is constant on the fibers of \(\pi\).  Equivalently, its support is the full inverse image of a support \(S'\subseteq D'\), up to \(r\) marked exceptions.

Write \(N=n/d\).  For a fixed canonical remainder \(R\), put
\(r=\abs R\), \(p=\abs{\pi(R)}\), and \(m=(a-r)/d\).  The exact number
of compatible complete-fiber supports is
\[
                         \binom{N-p}{m}.              \tag{CB1}
\]
If the canonical remainder ranges over all partial occupancies of total
weight \(r\) meeting exactly \(p\) fibers, the exact coarse-slice count is
\[
 \binom Np\binom{N-p}m
 [x^r]\bigl((1+x)^d-1-x^d\bigr)^p.                  \tag{CB2}
\]
An exceptional set of size \(b\), with \(t\) selected exceptional points,
multiplies \textup{(CB2)} by \(\binom bt\).  These are the corresponding
specializations of \textup{(PO1)}.

For a fixed-remainder profile with proved boundary map \(\Phi_R\), the
image-normalized scale is
\[
 \bar N_R^{\rm img}
 =\frac{\binom{N-p}{m}}
        {\abs{\Phi_R(\Omega_R^0)}}.                  \tag{CB3}
\]
For a coarse \((t,m,p,r)\)-profile, its numerator in \textup{(CB3)} is
instead the full count in \textup{(PO1)}.  An ambient candidate divides
the same exact support numerator by the formal target size.  The row must
still prove the boundary depth and coefficient field, realized image,
lift multiplicity, and distinct-slope projection.  In particular, none of
these image-normalized terms is automatically bounded by the identity
scale, as \cref{thm:smooth-quotient-obstruction} shows.

\subsection*{Planted cells}
A planted cell must range over an algebraically controlled family
\(\Pcal_b\), not all \(\binom nb\) blocks.  For a fixed block its residual
image scale is computed from \(\binom{n-b}{a-b}\) divided by the realized
image size of the proved residual prefix map.  The ambient prefix depth may
be used only under \textup{(FI)}.  Payment is exactly the criterion of
\cref{prop:planted-payment-repaired}; neither sublinearity nor
constructibility alone supplies it.

\subsection*{Tangent and common-line cells}
A tangent cell is a determinantal locus in the witness incidence.  Rank
loss \(h\) need not have codimension \(h\); independent minors and a
slope-projection estimate must be exhibited.  The common-line locus itself
has no bad slopes on that support.  These are the hypotheses of
\cref{lem:ramification-payment-expanded}, not automatic consequences of
smoothness.

\subsection*{Extension and field-descent cells}
An extension profile may be defined over a proper intermediate field
\(\B'\), in which case its image-normalized profile scale can be larger
than the identity scale.  Frobenius coincidence equations pay it only if their
independence and their effect on the slope image are proved.  Otherwise the
profile is retained at its natural field in \(\mathfrak E_n\).

\subsection*{Differential-locator and effective-image rank-collapse cells}
Differential-locator and effective-image rank-collapse loci are rank-defect
loci.  A drop of
\(h\) does not imply codimension \(h\), and a smaller effective image can
make fibers larger.  Payment therefore requires the independent-minor and
slope-projection hypotheses of \cref{lem:rank-defect-payment}.  Ray
saturation is separate and is paid by the image--fiber count for the
witness-to-ray map.  Ordinary
Vandermonde column independence does not classify every persistent
higher-dimensional defect.

\subsection*{Split-pencil cells}
A one-dimensional split pencil is controlled by the moving-root incidence
bound after common roots are removed.  In higher dimension the exact
second moment counts support pairs but gives no inequality of the displayed
frontier size and no ray count.  Such profiles require a direct
balanced-core exhaustion or \textup{(RC)}.

\begin{proposition}[Conditional cell-budget theorem]\label{prop:cell-budget-theorem}
Suppose an exhaustive atlas has \(e^{o(n)}\) profiles and each profile
satisfies a direct distinct-slope estimate at most
\(e^{o(n)}\barN_\lambda\).  Then the total algebraic contribution is at
most \(e^{o(n)}\mathfrak E_n(a)\).
\end{proposition}

\begin{proof}
Sum the assumed slope estimates over the profile atlas and use
\eqref{eq:profile-envelope}.  The preceding paragraphs explain sufficient
criteria and the places where a direct proof is still needed.
\end{proof}

\section{Smooth-domain ledger closure in detail}\label{sec:smooth-domain-detail}

Let \(D=\theta H\), where \(H\subseteq\B^\times\) is cyclic of order \(n\).  A smooth-domain quotient is the map \(x\mapsto x^d\), \(d\mid n\).  The fiber size is exactly \(d\), and the quotient domain is \(\theta^dH^d\).  A support is quotient-saturated if it is a union of fibers.  A locator \(Q_S(X)\) is quotient-saturated precisely when it lies in \(\B[X^d]\) after the scalar change \(X\mapsto \theta X\), up to the planted remainder support.

\begin{lemma}[Exact invariant multiplicative quotient descent]\label{lem:exact-multiplicative-descent}
Let \(D=\theta H\), \(d\mid\abs H\), and \(\pi(x)=x^d\).  Suppose that
after deleting a marked set \(R\), the support is a union of complete
\(\pi\)-fibers and the restrictions of both received words and the
explaining polynomial are pullbacks through \(\pi\).  Then the witness
descends exactly to \(\pi(D)\), with agreement divided by \(d\).
Conversely, pullback gives lifts; a subexponential lift count requires a
separate \(o(n)\)-bit profile bound.
\end{lemma}

\begin{proof}
All three data are constant on the same complete fibers, so their equality
on the original support is exactly the quotient equality.  Pullback proves
the converse.  Without the invariance hypothesis, nontrivial kernel
characters remain and must be counted rather than discarded.
\end{proof}

\begin{lemma}[Smooth support stabilizers]\label{lem:smooth-stabilizer-exhaustion}
If a support in \(D=\theta H\) is invariant under a nontrivial subgroup of
\(H\), it is a union of quotient fibers.  This classifies the support; it
does not imply that the received or explaining data descend, nor does it
classify arbitrary positive-density partial fibers.
\end{lemma}

\begin{proof}
The stabilizer is a subgroup, and an invariant set is a union of its
orbits, which are precisely the quotient fibers.
\end{proof}

\begin{proposition}[Conditional smooth-domain algebraic budget]\label{prop:smooth-domain-algebraic-closure}
If a multiplicative smooth row verifies the exhaustive atlas, invariant
quotient-descent/lift counts, planted and determinantal projection criteria, and
one-pencil or direct balanced-core ray bounds, then its algebraic cells
contribute at most \(e^{o(n)}\mathfrak E_n\).
\end{proposition}

\begin{proof}
Apply the verified profile estimates and
\cref{prop:cell-budget-theorem}.  Smoothness supplies the quotient maps,
not the omitted projection estimates.
\end{proof}

For the analytic component, let \(\psi\) be a nontrivial additive character of \(\B\).  A primitive prefix Fourier coefficient has the form
\[
       \binom Na^{-1}e_a\bigl(\psi(P(t)):t\in T\bigr),
\]
where \(P\) is a nonzero polynomial of degree less than the prefix
equation count \(R\).
Li--Wan and Weil give the minor-character estimate when \textup{(PF)}
holds, thereby furnishing one sufficient route to \textup{(MI)}.  The
full conclusion requires effective \textup{(MI)} and \textup{(MA)}; a
major-phase label does not provide either aggregate.

\begin{corollary}[Conditional smooth-domain compiler]\label{cor:smooth-domain-full-ledger}
A multiplicative smooth row satisfying the algebraic hypotheses above,
effective \textup{(MI)} plus \textup{(MA)} or a direct Sidon payment, the primitive-Q
conditions, and \textup{(RC)} has a closed asymptotic ledger.
\end{corollary}

\begin{proof}
Combine the stated hypotheses with
\cref{thm:primitive-q,prop:q-implies-sp,prop:cell-budget-theorem}.
\end{proof}

\section{Circle-domain ledger closure in detail}\label{sec:circle-domain-detail}

Let \(\Tcal\) be the norm-one torus over \(\B\), and write \(x=z+z^{-1}\).  A twin coset is a union \(\xi H\cup \xi^{-1}H^{-1}\) stable under \(z\mapsto z^{-1}\), projected to the \(x\)-line.  The Chebyshev polynomial \(T_d\) satisfies \(T_d(z+z^{-1})=z^d+z^{-d}\).  Hence the quotient map on the \(x\)-line is the projection of the multiplicative quotient \(z\mapsto z^d\) on the torus.  The only additional issue is the two-to-one ramification at \(z=\pm1\), which contributes only bounded exceptional fibers in the asymptotic rows and is assigned to a planted cell.

\begin{lemma}[Invariant Chebyshev quotient descent]\label{lem:chebyshev-descent}
Let \(D\) be a circle twin-coset domain and let \(d\) divide the torus
coset order.  If the support and all received and explaining data are
pullbacks through \(T_d\), the witness descends exactly to the Chebyshev
quotient.  A lift census requires its own description-entropy bound.
\end{lemma}

\begin{proof}
Lift to the torus, apply invariant multiplicative quotient descent, and take
invariants under \(z\mapsto z^{-1}\).  The identity
\(T_d(z+z^{-1})=z^d+z^{-d}\) identifies the quotient coordinate.
\end{proof}

\begin{lemma}[Circle diagonal uniformization]\label{lem:circle-diagonal-uniformization}
For the standard circle code and evaluation set of
\cref{lem:circle-uniformization}, the diagonal multiplier gives an exact
equivalence to the torus Reed--Solomon row and preserves MCA numerators.
\end{lemma}

\begin{proof}
This is \cref{lem:circle-uniformization}.
\end{proof}

\begin{proposition}[Conditional circle-domain compiler]\label{prop:circle-domain-full-ledger}
A circle twin-coset row with a proved exhaustive subexponential atlas and
complete profile envelope, satisfying the circle analogues of the
algebraic profile criteria, effective \textup{(MI)} plus \textup{(MA)} or a direct
Sidon payment, the primitive image-scale conditions, and \textup{(RC)},
has a closed asymptotic ledger.  The same holds for a standard circle-code
row under the exact diagonal equivalence.
\end{proposition}

\begin{proof}
Pull the verified algebraic and analytic estimates to the torus, apply the
conditional smooth compiler equivariantly under inversion, and use
\cref{lem:circle-diagonal-uniformization}.  No unverified major character
or balanced-core chart is removed by this reduction.
\end{proof}

\section{Why ledger closure remains an input}\label{sec:no-closed-hypothesis}

Smooth and circle geometry supplies uniform folding maps, but it does not
prove the enumerative statements that drive the upper bound: exhaustion of
the corresponding boundary-map and algebraic slope projections, the
primitive effective minor/major aggregates, and the higher-dimensional
transverse-secant bound.  The underlying partial-occupancy \emph{support}
partition itself is exact by \cref{thm:exact-partial-occupancy}; what does
not follow from that census is image coverage, Fourier flatness, or ray
payment.  The following theorem records the exact conditional status.

\begin{theorem}[Smooth/circle compiler, expanded form]\label{thm:smooth-circle-ledger-closure-expanded}
Let \((C_n,a_n)\) be a ledger-admissible smooth/circle sequence.  Then
\[
 B_{C_n}^{\rm MCA}(a_n)\le e^{o(n)}\mathfrak E_n(a_n).
\]
\end{theorem}

\begin{proof}
This is \cref{thm:main-smooth-circle}.  The smooth and
circle propositions above explain how the required verified estimates
transport through power, Chebyshev, and diagonal maps.
\end{proof}

\begin{corollary}[Identity-dominant specialization]\label{cor:final-no-ledger-assumption}
If, in addition,
\[
 \mathfrak E_n(a_n)=e^{o(n)}\binom n{a_n}
 \abs{\B_n}^{-(a_n-k_n-1)},
\]
then the numerator has the corresponding
identity-scale upper bound.  The quotient counterexample proves that this
additional comparison cannot be deleted.
\end{corollary}

\begin{proof}
Substitute the profile comparison into the preceding theorem.
\end{proof}


\section{Algebraic ledger criteria: planted, determinantal, and circle edge cases}\label{sec:algebraic-repairs}

This section records the precise forms of three payments that are often misstated if one counts supports rather than moduli.  They are included to ensure that the first-match ledger is not hiding exponential candidate families.

\begin{definition}[Algebraically planted block]\label{def:algebraically-planted}
A planted block of size \(b\) is algebraically planted if it is the zero
set on \(D\) of a polynomial in a constructible family \(\Pcal_b\) of
support locators, common factors, ramification polynomials, polynomials
cutting out quotient fibers, or received-line resultants.  Here a
\emph{received-line resultant} is the polynomial obtained by eliminating
the explaining-polynomial variables from the received-line interpolation
equations; its zeros are the coordinates forced by that elimination.  We
require \(\abs{\Pcal_b(\B)}\le e^{o(n)}\) for every profile size occurring
in the profiled asymptotic row datum.  Arbitrary choices of \(b\)-subsets
of \(D\) are not planted cells.
\end{definition}

\begin{proposition}[Planted payment without arbitrary-block overcount]\label{prop:planted-payment-repaired}
Let \(P\) range over an algebraically planted family \(\Pcal_b\), and
suppose the residual chart after fixing \(P\) has image-normalized profile scale
\(\barN(P)\).  If the total description entropy is \(o(n)\),
\(\sum_P\barN(P)\le e^{o(n)}\mathfrak E_n\), and the projection from each
chart to distinct slopes has cost \(e^{o(n)}\barN(P)\), then the planted
family is paid.  These hypotheses must also be checked for positive-density
blocks; arbitrary blocks are excluded.
\end{proposition}

\begin{proof}
Sum the assumed distinct-slope budgets over \(P\).  The statement is
deliberately phrased in terms of description entropy and natural profile
scales, since \(b=o(n)\) alone does not imply
\(b\log\abs\B=o(n)\).
\end{proof}

\begin{proposition}[Determinantal payment with a nonzero pivot minor]\label{prop:determinantal-payment-repaired}
Let \(X\) be a bounded-complexity affine chart over \(\B\), and let
\(M:X\to\operatorname{Mat}_{u\times v}\) be polynomial of degree
\(\exp(o(n))\).  A nonzero \(r\times r\) minor makes
\(\{\rank M<r\}\) a proper constructible locus.  If \(h\) independent
pivot equations give codimension \(h\) and a separate slope-projection
estimate converts the point loss into image loss, then the cell gains
\(\abs\B^{-h+o(n)}\).  Dependence among the minors yields no automatic
quotient or tangent classification.
\end{proposition}

\begin{proof}
The minors define the rank locus.  Under the stated independence,
bounded-degree point counting gives
\(\abs\B^{\dim X-h}\exp(o(n))\) chart points.  The separately assumed
projection estimate gives the slope-image bound.  No conclusion is drawn
when either hypothesis fails.
\end{proof}

\begin{lemma}[Circle ramification and twin-coset collapse]\label{lem:circle-edge-cases-repaired}
For a circle twin coset \(gH\cup g^{-1}H\), the exceptional cases
\(u=\pm1\), \(gH=g^{-1}H\), and \(g^{2a}\in H^{(a)}\) are respectively a
bounded ramification profile, a dihedral collapse, and a Chebyshev quotient
profile.  They are constructible and must be budgeted at their natural
profile scales.
\end{lemma}

\begin{proof}
The points \(u=\pm1\) are fixed by the involution and contribute at most
two ramification points.  The equality \(gH=g^{-1}H\) is equivalent to
\(g^2\in H\), so the twin coset is dihedrally collapsed.  If
\(g^{2a}\in H^{(a)}\), the two image cosets under \(u\mapsto u^a\)
coincide.  These observations classify the profiles; their quotient terms
are retained in \(\mathfrak E_n\).
\end{proof}

\begin{proposition}[Conditional algebraic ledger from verified criteria]\label{prop:closed-algebraic-ledger-repaired}
For a smooth or circle row with an exhaustive subexponential atlas, suppose
every quotient, planted, tangent, extension, differential, saturation, and
one-pencil profile satisfies the corresponding direct slope criterion
above, while every higher-dimensional balanced core has a direct ray bound.
Then the algebraic contribution is \(e^{o(n)}\mathfrak E_n\).
\end{proposition}

\begin{proof}
Sum the verified profile budgets.  The one-pencil estimate is
\cref{prop:split-pencil-payment}; no subdivision theorem for
higher-dimensional pencils is used or claimed.
\end{proof}

\section{Finite source theorems and certificates}\label{sec:source-integration}

The finite constructions in \cite{Cho26CapV13,Cho26Grande} supply the
uniform-fiber viewpoint, locator-prefix coordinates, and finite ledger
interfaces.  The unconditional ingredients used here are the syndrome
description and deep theorem, the locator-prefix list, the collision-aware
pole line, the map-smooth list construction, the bounded-remainder normal
form, the polynomial-field counterexample, the exact support upper bound,
the adjacent-hyperplane construction, the support-pair identity, and the
one-pencil moving-root bound.  The exact partial-occupancy support
partition and its bounded-boundary closure are
\cref{thm:exact-partial-occupancy,thm:canonical-partial-occupancy-atlas}.
Bounded-kernel, single-circuit, exact-support curve/pencil, and deep
invariant quotient payments are unconditional in their displayed ranges;
the remaining determinantal and higher-dimensional balanced-core payments
occur only under the explicit conditional criteria above.

For circle rows, Chebyshev smoothness supplies the uniform-fiber maps.
The stereographic equivalence of
\cref{prop:stereographic-circle-equivalence} transfers the quotient
obstruction in its stated range, while the exact diagonal equivalence of
\cref{lem:circle-uniformization} transfers support-wise numerators for the
compatible torus presentation.

The primitive analytic residue splits Q into a Sidon-heavy branch and a
high-energy Boolean branch.  The high-energy branch is impossible by BSG
and Boolean-cube difference growth.  The Sidon moment remains an explicit
input outside the low-boundary and subsquare-root regimes, where
\cref{thm:small-effective-dual-closure,thm:unconditional-shallow-mi-ma}
verify it.  In general it is certified by effective \textup{(MI)} plus
\textup{(MA)}, with \textup{(PF)} as one
sufficient minor certificate when its numerical range holds.  The
polynomial-field construction proves that no such ambient
identity-normalized payment can hold before quotient, subfield, planted,
and remainder profiles are routed.  Thus the primitive theorem resolves
the high-energy branch, while the ambient identity-normalized claim is
false without profile routing.

A finite certificate tabulates the exact effective Fourier masses, direct
algebraic slope projections, residual moments, and ray incidence degrees
in \textup{(FC1)}--\textup{(FC2)}.  These quantities are not recovered by
inserting a universal constant into an asymptotic support count.

\section{A compact dependency graph}\label{sec:dependency-graph}

For clarity we record the logical dependencies in one place.  Three paths
are unconditional and bypass the general residual hypothesis: the deep
syndrome theorem gives
\cref{thm:intro-unconditional-asymptotic-deep}; the support upper bound
and pole-line lower bound give
\cref{thm:unconditional-asymptotic-rs-mca,thm:unconditional-support-envelope-bracket};
and the canonical partial-occupancy atlas under \textup{(PO8)} gives the
coarse upper bound \textup{(PO9)} with direct support-to-slope payment.
Under the active-size strengthening \textup{(PO10)}, its empty minor set
and all-major aggregate also verify the slice versions of MI and MA.  This upper path does not
locate a threshold without \textup{(SB2)}.  The remaining
domain-independent analytic step is
\[
\begin{aligned}
&\text{Sidon moment payment}+\Gord_q\text{ large}
 \Rightarrow E(F_s)\ge \abs{F_s}^3e^{-o(n)},\\
&E(F_s)\ge \abs{F_s}^3e^{-o(n)}
 \Rightarrow \text{BSG subset }A
 \Rightarrow \abs{A-A}\le e^{o(n)}\abs A.
\end{aligned}
\]
Since \(A\subseteq\{0,1\}^T\), Boolean-cube difference growth gives
\(\abs{A-A}\ge\abs A^{3/2}\), contradicting
\(\abs A=e^{\Omega(n)}\).  Hence primitive Q holds.  The compiler is
\[
 \text{paid algebraic profiles}+\text{primitive Q}+\textup{(RC)}
 \Rightarrow \text{MCA numerator},
\]
followed, under the required profile-envelope comparison and target/lower
hypotheses, by the
threshold comparison.  The implications are proved; their displayed
premises are not claimed for every smooth/circle row.  Indeed the
unrestricted ambient identity-normalized Sidon premise and
identity-profile dominance are disproved in
\cref{thm:smooth-quotient-obstruction}.  The abstract statement
\cref{prop:q-sp-no-ray} shows that Q, SP, and the one-slope-per-support
axiom alone do not logically imply a ray bound, while
\cref{thm:aperiodic-ray-obstruction} shows within Reed--Solomon geometry
that quotient-aperiodicity does not force retained-support multiplicity.
Neither statement disproves a correctly normalized RS-specific
\textup{(RC)}; the projection is proved here on the bounded-kernel,
single-circuit, curve, pencil, and deep-quotient strata, while its
unrestricted balanced-core form remains open.  Exact
partial-occupancy add-back and effective normalization
remove two bookkeeping gaps, but \cref{prop:pf-deep-prefix-barrier}
shows that the remaining deep minor estimate needs genuinely
phase-dependent or aggregate input.

\begin{thebibliography}{GGSW26}

\bibitem[ABF26]{ABF26}
G. Arnon, D. Boneh, and G. Fenzi,
\newblock Open problems in list decoding and correlated agreement,
\newblock IACR Cryptology ePrint Archive, Report 2026/680, 2026,
revised July 6, 2026.

\bibitem[BS94]{BalogSzemeredi}
A. Balog and E. Szemerédi,
\newblock A statistical theorem of set addition,
\newblock \emph{Combinatorica} 14 (1994), no. 3, 263--268.

\bibitem[BCHKS25]{BCHKS25}
E. Ben-Sasson, D. Carmon, U. Hab\"ock, S. Kopparty, and S. Saraf,
\newblock On proximity gaps for Reed--Solomon codes,
\newblock IACR Cryptology ePrint Archive, Report 2025/2055, 2025.

\bibitem[BCIKS20]{BCIKS20}
E. Ben-Sasson, D. Carmon, Y. Ishai, S. Kopparty, and S. Saraf,
\newblock Proximity gaps for Reed--Solomon codes,
\newblock \emph{J. ACM} 70 (2023), no.~5, 1--57,
\newblock doi:10.1145/3614423; conference version in
\emph{61st IEEE Symposium on Foundations of Computer Science
(FOCS 2020)} and extended version, IACR ePrint 2020/654.

\bibitem[BCGM25]{BordageChiesaGuanManzur25}
S. Bordage, A. Chiesa, Z. Guan, and I. Manzur,
\newblock All polynomial generators preserve distance with mutual
correlated agreement,
\newblock IACR Cryptology ePrint Archive, Report 2025/2051, 2025,
revised May 2026.

\bibitem[Cho26a]{Cho26CapV13}
P. Chojecki,
\newblock Universal field-size caps and a two-sided sandwich for mutual
correlated agreement on smooth Reed--Solomon domains,
\newblock manuscript, July 2026.

\bibitem[Cho26b]{Cho26Grande}
P. Chojecki,
\newblock Proof-audited and expanded bounds and reductions for RS--MCA,
\newblock manuscript, July 2026.

\bibitem[CS25]{CS25}
E. Crites and A. Stewart,
\newblock On Reed--Solomon proximity gaps conjectures,
\newblock IACR Cryptology ePrint Archive, Report 2025/2046, 2025.

\bibitem[CGHLL26]{CGHLL26}
D. Carmon, L. Goldberg, U. Hab\"ock, L. Lerer, and I. Lesokhin,
\newblock S-two whitepaper,
\newblock IACR Cryptology ePrint Archive, Report 2026/532, 2026.

\bibitem[GCXK25]{GaoCaiXuKan25}
Y. Gao, D. Cai, Y. Xu, and H. Kan,
\newblock From list-decodability to proximity gaps,
\newblock IACR Cryptology ePrint Archive, Report 2025/870, 2025.

\bibitem[Gow98]{GowersSzemeredi}
W. T. Gowers,
\newblock A new proof of Szemerédi's theorem for arithmetic progressions of length four,
\newblock \emph{Geom. Funct. Anal.} 8 (1998), no. 3, 529--551.

\bibitem[GG25]{GoyalGuruswami25}
R. Goyal and V. Guruswami,
\newblock Optimal proximity gaps for subspace-design codes and (random)
Reed--Solomon codes,
\newblock IACR Cryptology ePrint Archive, Report 2025/2054, 2025,
revised March 2026.

\bibitem[GGH26]{GoyalGuruswamiHsieh26}
R. Goyal, V. Guruswami, and J.-T. Hsieh,
\newblock Explicit constant-alphabet subspace design codes,
\newblock arXiv:2604.15218, 2026.

\bibitem[GGSW26]{GoyalGuruswamiSunWootters26}
R. Goyal, V. Guruswami, Y. Sun, and M. Wootters,
\newblock Locality of curve-decoding and improved proximity gaps,
\newblock arXiv:2607.08516, 2026.

\bibitem[Goh24]{GohEnergy}
M. K. Goh,
\newblock On an entropic analogue of additive energy,
\newblock arXiv:2406.18798.

\bibitem[GMT23]{GreenMannersTao}
B. Green, F. Manners, and T. Tao,
\newblock Sumsets and entropy revisited,
\newblock arXiv:2306.13403.

\bibitem[GMR+22]{GreenMatolcsiRuzsaShakanZhelezov}
B. Green, D. Matolcsi, I. Z. Ruzsa, G. Shakan, and D. Zhelezov,
\newblock A weighted Prékopa--Leindler inequality and sumsets with quasicubes,
\newblock in \emph{Analysis at Large}, Springer, 2022, 125--129; arXiv:2003.04077.

\bibitem[HLP24]{HLP24}
U. Hab\"ock, D. Levit, and S. Papini,
\newblock Circle STARKs,
\newblock IACR Cryptology ePrint Archive, Report 2024/278, 2024.

\bibitem[ACFY25]{WHIR}
G. Arnon, A. Chiesa, G. Fenzi, and E. Yogev,
\newblock WHIR: Reed--Solomon proximity testing with super-fast verification,
\newblock in \emph{EUROCRYPT 2025}, 214--243;
IACR ePrint 2024/1586.

\bibitem[Hab25]{Habock25}
U. Hab\"ock,
\newblock A note on mutual correlated agreement for Reed--Solomon codes,
\newblock IACR Cryptology ePrint Archive, Report 2025/2110, 2025.

\bibitem[LW10]{LiWanSieve}
J. Li and D. Wan,
\newblock A new sieve for distinct coordinate counting,
\newblock \emph{Sci. China Math.} 53 (2010), no. 9, 2351--2362.

\bibitem[LW12]{LiWanSubset}
J. Li and D. Wan,
\newblock Counting subset sums of finite abelian groups,
\newblock \emph{J. Combin. Theory Ser. A} 119 (2012), no. 1, 170--182.

\bibitem[MRSZ22]{MatolcsiRuzsaShakanZhelezov}
D. Matolcsi, I. Z. Ruzsa, G. Shakan, and D. Zhelezov,
\newblock An analytic approach to cardinalities of sumsets,
\newblock \emph{Combinatorica} 42 (2022), 203--236; arXiv:2003.04075.

\bibitem[Tao10]{TaoEntropy}
T. Tao,
\newblock Sumset and inverse sumset theorems for Shannon entropy,
\newblock \emph{Combin. Probab. Comput.} 19 (2010), no. 4, 603--639; arXiv:0906.4387.

\bibitem[TV06]{TaoVu}
T. Tao and V. Vu,
\newblock \emph{Additive Combinatorics},
\newblock Cambridge Studies in Advanced Mathematics 105, Cambridge University Press, 2006.

\bibitem[PP26]{ProximityPrize}
{Ethereum Foundation},
\newblock The Proximity Prize: \$1M Reed--Solomon Challenge,
\newblock preliminary version, 2026,
\newblock \url{https://proximityprize.org/}, accessed July 10, 2026.

\bibitem[Wei48]{Weil}
A. Weil,
\newblock On some exponential sums,
\newblock \emph{Proc. Natl. Acad. Sci. USA} 34 (1948), 204--207.

\bibitem[YZ26]{YuanZhu26}
C. Yuan and R. Zhu,
\newblock A syndrome-space approach to proximity gaps and correlated
agreement for random linear codes,
\newblock arXiv:2605.07595, 2026.

\end{thebibliography}

\end{document}
